% Generated mechanically from frozen input p05/ko/perfect.tex.
% Do not edit this generated file; edit the bound locale source instead.

% OK, start here.
%

\setcounter{chapter}{35}
\chapter{스킴의 유도 범주}



\phantomsection
\label{perfect-section-phantom}


\section{들어가며}
\label{perfect-section-introduction}

\noindent
이 장에서는 스킴 위 가군의 유도 범주를 다룬다.
여기서 논의하는 내용의 대부분은
\cite{TT}, \cite{Bokstedt-Neeman}, \cite{BvdB}, 그리고 \cite{LN}에서
찾을 수 있다. 물론 그 밖에도 많은 참고문헌이 있다.


\section{규약}
\label{perfect-section-conventions}

\noindent
$\mathcal{A}$가 아벨 범주이고 $M$이 $\mathcal{A}$의 대상이면,
차수 $0$에서 $M$이고 다른 모든 차수에서 영인 복합체에 대응하는
$K(\mathcal{A})$ 및/또는 $D(\mathcal{A})$의 대상도 $M$으로 표기한다.

\medskip\noindent
환 $A$가 주어지면 $K(A)$는 $A$-가군 복합체의 호모토피 범주를,
$D(A)$는 이에 연관된 유도 범주를 나타낸다. 마찬가지로 환 달린 공간
$(X, \mathcal{O}_X)$가 주어지면 $K(\mathcal{O}_X)$는
$\mathcal{O}_X$-가군 복합체의 호모토피 범주를,
$D(\mathcal{O}_X)$는 이에 연관된 유도 범주를 나타낸다.










\section{준연접 가군의 유도 범주}
\label{perfect-section-derived-quasi-coherent}

\noindent
이 절에서는 스킴 $X$ 위의 준연접 가군과 모든 가군 사이의 관계를
논의한다. 참고문헌으로 \cite[Appendix B]{TT}가 있다. Schemes, Section
\ref{schemes-section-quasi-coherent}의 논의에 따르면 매장
$\QCoh(\mathcal{O}_X) \subset \textit{Mod}(\mathcal{O}_X)$
은 $\QCoh(\mathcal{O}_X)$를 $\mathcal{O}_X$-가군 범주의 약한 Serre
부분범주로 나타낸다. 코호몰로지 층들이 준연접인 복합체의 부분범주를
$$
D_\QCoh(\mathcal{O}_X) \subset D(\mathcal{O}_X)
$$
로 표기한다. Derived Categories, Section
\ref{derived-section-triangulated-sub}을 보라. 따라서 표준 함자
\begin{equation}
\label{perfect-equation-compare}
D(\QCoh(\mathcal{O}_X))
\longrightarrow
D_\QCoh(\mathcal{O}_X)
\end{equation}
를 얻는다. Derived Categories, Equation
(\ref{derived-equation-compare})을 보라.

\begin{lemma}
\label{perfect-lemma-quasi-coherence-direct-sums}
$X$를 스킴이라 하자. 그러면 $D_\QCoh(\mathcal{O}_X)$에는 직합이 존재한다.
\end{lemma}

\begin{proof}
Injectives, Lemma \ref{injectives-lemma-derived-products}에 의해 유도 범주
$D(\mathcal{O}_X)$에는 직합이 존재하고, 이는 임의의 대표들에 대해 항별
직합을 취하여 계산된다. 모든 Grothendieck 아벨 범주에서 직합을 취하는
함자는 완전하므로, 직합의 코호몰로지 층은 코호몰로지 층들의 직합임이
분명하다. 준연접층의 직합은 준연접이므로 보조정리가 따른다. Schemes,
Section \ref{schemes-section-quasi-coherent}을 보라.
\end{proof}

\noindent
유도 극한에 관한 몇 가지 정보가 필요하다. 아래 보조정리에서 유도 극한은
일반적으로 $D_\QCoh$의 대상이 아니라는 점을 미리 밝혀 둔다.

\begin{lemma}
\label{perfect-lemma-Rlim-quasi-coherent}
$X$를 스킴이라 하자. $(K_n)$을 $D_\QCoh(\mathcal{O}_X)$의 역계라 하고,
$D(\mathcal{O}_X)$에서의 유도 극한을 $K = R\lim K_n$이라 하자. 모든
$q \in \mathbf{Z}$와 $n \geq 1$에 대해 $H^q(K_{n + 1}) \to H^q(K_n)$이
전사라고 가정하자. 그러면
\begin{enumerate}
\item $H^q(K) = \lim H^q(K_n)$,
\item $R\lim H^q(K_n) = \lim H^q(K_n)$이고,
\item 모든 아핀 열린집합 $U \subset X$와 $p > 0$에 대해
$H^p(U, \lim H^q(K_n)) = 0$이다.
\end{enumerate}
\end{lemma}

\begin{proof}
$\mathcal{B}$를 $X$의 아핀 열린집합들로 이루어진 집합이라 하자.
$H^q(K_n)$은 준연접이므로 Cohomology of Schemes, Lemma
\ref{coherent-lemma-quasi-coherent-affine-cohomology-zero}에 의해
$U \in \mathcal{B}$에 대해 $H^p(U, H^q(K_n)) = 0$이다. 또한 Schemes,
Lemma \ref{schemes-lemma-equivalence-quasi-coherent}에 의해
$U \in \mathcal{B}$에 대한 사상
$H^0(U, H^q(K_{n + 1})) \to H^0(U, H^q(K_n))$은 전사이다. 방금 확인한
조건들을 Cohomology, Lemma
\ref{cohomology-lemma-derived-limit-suitable-system}에 적용하면 (1)이
따른다. (2)와 (3)은 Cohomology, Lemma
\ref{cohomology-lemma-inverse-limit-is-derived-limit}에서 따른다.
\end{proof}

\noindent
다음 보조정리는 $D_\QCoh(\mathcal{O}_X)$의 대상 위에서 오른쪽 유도 함자를
``계산''하는 데 도움이 된다.

\begin{lemma}
\label{perfect-lemma-nice-K-injective}
$X$를 스킴이라 하고 $E$를 $D_\QCoh(\mathcal{O}_X)$의 대상이라 하자.
그러면 Derived Categories, Remark
\ref{derived-remark-map-into-derived-limit-truncations}의 사상
$E \to R\lim \tau_{\geq -n}E$는 동형사상이다\footnote{특히 Derived
Categories, Lemma \ref{derived-lemma-difficulty-K-injectives}에 의해
$E$에는 K-단사 대표가 존재한다.}.
\end{lemma}

\begin{proof}
$E$의 $i$번째 코호몰로지 층을 $\mathcal{H}^i = H^i(E)$로 표기하자.
$\mathcal{B}$를 $X$의 아핀 열린부분집합들로 이루어진 집합이라 하자.
그러면 모든 $p > 0$, $i \in \mathbf{Z}$, $U \in \mathcal{B}$에 대해
$H^p(U, \mathcal{H}^i) = 0$이다. Cohomology of Schemes, Lemma
\ref{coherent-lemma-quasi-coherent-affine-cohomology-zero}를 보라. 따라서
Cohomology, Lemma \ref{cohomology-lemma-is-limit-dimension}에서 보조정리가
따른다.
\end{proof}

\begin{lemma}
\label{perfect-lemma-application-nice-K-injective}
$X$를 스킴이라 하자. $F : \textit{Mod}(\mathcal{O}_X) \to \textit{Ab}$를
가법 함자라 하고 $N \geq 0$을 정수라 하자. 다음을 가정하자.
\begin{enumerate}
\item $F$는 가산 직접곱과 가환한다.
\item 모든 $p \geq N$과 준연접 $\mathcal{F}$에 대해
$R^pF(\mathcal{F}) = 0$이다.
\end{enumerate}
그러면 $E \in D_\QCoh(\mathcal{O}_X)$에 대해 다음이 성립한다.
\begin{enumerate}
\item $i \leq a$이면 $H^i(RF(\tau_{\leq a}E)) \to H^i(RF(E))$는
동형사상이다.
\item $i \geq b$이면 $H^i(RF(E)) \to H^i(RF(\tau_{\geq b - N + 1}E))$는
동형사상이다.
\item 어떤 $-\infty \leq a \leq b \leq \infty$에 대해
$i \not \in [a, b]$일 때 $H^i(E) = 0$이면,
$i \not \in [a, b + N - 1]$일 때 $H^i(RF(E)) = 0$이다.
\end{enumerate}
\end{lemma}

\begin{proof}
(1)은 Derived Categories, Lemma
\ref{derived-lemma-negative-vanishing}이다.

\medskip\noindent
(2)의 증명. $E_n = \tau_{\geq -n}E$로 쓰자. Lemma
\ref{perfect-lemma-nice-K-injective}에 의해 $E = R\lim E_n$이다. 따라서 Injectives,
Lemma \ref{injectives-lemma-RF-commutes-with-Rlim}에 의해
$D(\textit{Ab})$에서 $RF(E) = R\lim RF(E_n)$이다. 그러므로 모든
$i \in \mathbf{Z}$에 대해 짧은 완전열
$$
0 \to R^1\lim H^{i - 1}(RF(E_n)) \to H^i(RF(E)) \to \lim H^i(RF(E_n)) \to 0
$$
을 얻는다. More on Algebra, Remark
\ref{more-algebra-remark-compare-derived-limit}를 보라. (2)를 증명하기 위해
왼쪽 항이 영이고, $i \geq b$인 임의의 $b$에 대해 오른쪽 항이
$H^i(RF(E_{-b + N - 1}))$와 같음을 보이겠다.

\medskip\noindent
모든 $n$에 대해 완전 삼각형
$$
H^{-n}(E)[n] \to E_n \to E_{n - 1} \to H^{-n}(E)[n + 1]
$$
(Derived Categories, Remark
\ref{derived-remark-truncation-distinguished-triangle})이
$D(\mathcal{O}_X)$에 존재한다. $H^{-n}(E)$는 준연접이므로
$$
H^i(RF(H^{-n}(E)[n])) = R^{i + n}F(H^{-n}(E)) = 0
$$
$i + n \geq N$일 때 위 등식이 성립하고,
$$
H^i(RF(H^{-n}(E)[n + 1])) = R^{i + n + 1}F(H^{-n}(E)) = 0
$$
$i + n + 1 \geq N$일 때 이 등식이 성립한다. 따라서 사상
$$
H^i(RF(E_n)) \to H^i(RF(E_{n - 1}))
$$
은 $n \geq N - i$일 때 동형사상이다. 그러므로 모든 계
$H^i(RF(E_n))$은 ML 조건을 만족하고, 위 짧은 완전열의 $R^1\lim$ 항은
영이다(More on Algebra, Section \ref{more-algebra-section-Rlim}의 논의를
보라). 또한 계 $H^i(RF(E_n))$은 원하는 대로 $n = N - i - 1$부터
상수이다.

\medskip\noindent
(3)의 증명. $E$에 관한 가정 아래 $\tau_{\leq a - 1}E = 0$이므로 (1)에서
$i \leq a - 1$일 때 $H^i(RF(E))$가 소멸함을 얻는다. 마찬가지로
$\tau_{\geq b + 1}E = 0$이므로 (2)에서 $i \geq b + N$일 때
$H^i(RF(E))$가 소멸함을 얻는다.
\end{proof}

\noindent
다음 보조정리는 이 장의 많은 결과에서 핵심 재료가 된다.

\begin{lemma}
\label{perfect-lemma-affine-compare-bounded}
$X = \Spec(A)$를 아핀 스킴이라 하자. 도식
$$
\xymatrix{
D(\QCoh(\mathcal{O}_X)) \ar[rr]_{(\ref{perfect-equation-compare})}
& &
D_\QCoh(\mathcal{O}_X) \ar[ld]^{R\Gamma(X, -)} \\
& D(A) \ar[lu]^{\widetilde{\ \ }}
}
$$
의 모든 함자는 삼각 범주의 동치이다. 또한
$E$가 $D_\QCoh(\mathcal{O}_X)$의 대상이면
$H^0(X, E) = H^0(X, H^0(E))$이다.
\end{lemma}

\begin{proof}
함자 $R\Gamma(X, -)$는 함자 $D(\mathcal{O}_X) \to D(A)$를 주므로,
제한을 통해 함자
\begin{equation}
\label{perfect-equation-back}
R\Gamma(X, -) : D_\QCoh(\mathcal{O}_X) \longrightarrow D(A).
\end{equation}
를 준다. $X$ 위의 준연접 가군과 $A$-가군의 범주 사이의 동치를 통하여
이 함자가 (\ref{perfect-equation-compare})의 준역함자임을 보이겠다.

\medskip\noindent
설명. $A$가 주는 환의 층을 가진 한 점 공간을 $(Y, \mathcal{O}_Y)$로
표기하고, 환 달린 공간의 자명한 사상을
$\pi : (X, \mathcal{O}_X) \to (Y, \mathcal{O}_Y)$로 표기하자. 그러면
$R\Gamma(X, -)$를 $R\pi_*$와 동일시할 수 있고, 동치
$\textit{Mod}(\mathcal{O}_Y) = \text{Mod}_A = \QCoh(\mathcal{O}_X)$
를 통하여 함자 (\ref{perfect-equation-compare})를
$L\pi^* = \pi^* = \widetilde{\ }$와 동일시할 수 있다(Modules, Lemma
\ref{modules-lemma-construct-quasi-coherent-sheaves} 및 Schemes, Lemmas
\ref{schemes-lemma-compare-constructions}와
\ref{schemes-lemma-equivalence-quasi-coherent}를 보라). 따라서 함자들
$$
\xymatrix{
D(A) \ar@<1ex>[r] & D(\mathcal{O}_X) \ar@<1ex>[l]
}
$$
은 수반이다(Cohomology, Lemma \ref{cohomology-lemma-adjoint}). 특히 표준
수반 사상
$$
a : \widetilde{R\Gamma(X, E)} \longrightarrow E
$$
$E$가 $D(\mathcal{O}_X)$의 대상일 때의 위 사상과
$$
b : M^\bullet \longrightarrow R\Gamma(X, \widetilde{M^\bullet})
$$
$A$-가군 복합체 $M^\bullet$에 대한 이 사상을 얻는다.

\medskip\noindent
$E$를 $D_\QCoh(\mathcal{O}_X)$의 대상이라 하자. Cohomology of Schemes,
Lemma \ref{coherent-lemma-quasi-coherent-affine-cohomology-zero}에 의해
함자 $F(-) = \Gamma(X, -)$와 $N = 1$에 Lemma
\ref{perfect-lemma-application-nice-K-injective}를 적용할 수 있다. 따라서
$$
H^0(R\Gamma(X, E)) = H^0(R\Gamma(X, \tau_{\geq 0}E)) = \Gamma(X, H^0(E))
$$
이다(마지막 등식은 표준 절단의 정의에 따른다). 이를 사용하여 수반 사상
$a$와 $b$가 동형사상 $H^0(a)$와 $H^0(b)$를 유도함을 보이겠다. 그러면
이 명제는 이동 아래 불변이므로 $a$와 $b$는 준동형사상이고, 보조정리가
증명된다.

\medskip\noindent
두 경우 모두 $\widetilde{\ }$가 완전 함자라는 사실을 사용한다(Schemes,
Lemma \ref{schemes-lemma-spec-sheaves}). 즉, 이 사실로부터
$$
H^0\left(\widetilde{R\Gamma(X, E)}\right) =
\widetilde{H^0(R\Gamma(X, E))} =
\widetilde{\Gamma(X, H^0(E))}
$$
를 얻는데, $H^0(E)$가 준연접이므로 이는 $H^0(E)$와 같다. 따라서
$H^0(a)$는 동형사상이다. 다른 방향에서는
$$
H^0(R\Gamma(X, \widetilde{M^\bullet})) =
\Gamma(X, H^0(\widetilde{M^\bullet})) =
\Gamma(X, \widetilde{H^0(M^\bullet)}) =
H^0(M^\bullet)
$$
이므로 $H^0(b)$가 동형사상임이 증명된다.
\end{proof}

\begin{lemma}
\label{perfect-lemma-affine-K-flat}
$X = \Spec(A)$를 아핀 스킴이라 하자. $K^\bullet$이 $A$-가군의 K-평탄
복합체이면 $\widetilde{K^\bullet}$은 $\mathcal{O}_X$-가군의 K-평탄
복합체이다.
\end{lemma}

\begin{proof}
More on Algebra, Lemma \ref{more-algebra-lemma-base-change-K-flat}에 의해
모든 $\mathfrak p \in \Spec(A)$에 대해
$K^\bullet \otimes_A A_\mathfrak p$는 $A_\mathfrak p$-가군의 K-평탄
복합체이다. 그러므로 Cohomology, Lemma
\ref{cohomology-lemma-check-K-flat-stalks}(및 Schemes, Lemma
\ref{schemes-lemma-spec-sheaves})로부터 $\widetilde{K^\bullet}$이
K-평탄임을 얻는다.
\end{proof}

\begin{lemma}
\label{perfect-lemma-quasi-coherence-pushforward}
$f : X \to Y$가 환 사상 $A \to B$로 주어지는 아핀 스킴의 사상이면 도식
$$
\xymatrix{
D(B) \ar[d] \ar[r] & D_\QCoh(\mathcal{O}_X) \ar[d]^{Rf_*} \\
D(A) \ar[r] & D_\QCoh(\mathcal{O}_Y)
}
$$
은 가환이다.
\end{lemma}

\begin{proof}
Cohomology, Lemma \ref{cohomology-lemma-Leray-unbounded}에 의해
$R\Gamma(Y, Rf_*K) = R\Gamma(X, K)$임을 사용하면 Lemma
\ref{perfect-lemma-affine-compare-bounded}에서 따른다.
\end{proof}

\begin{lemma}
\label{perfect-lemma-quasi-coherence-pullback}
$f : Y \to X$를 스킴의 사상이라 하자.
\begin{enumerate}
\item 함자 $Lf^*$는 $D_\QCoh(\mathcal{O}_X)$를
$D_\QCoh(\mathcal{O}_Y)$ 안으로 보낸다.
\item $X$와 $Y$가 아핀이고 $f$가 환 사상 $A \to B$로 주어지면 도식
$$
\xymatrix{
D(B) \ar[r] & D_\QCoh(\mathcal{O}_Y) \\
D(A) \ar[r] \ar[u]^{- \otimes_A^\mathbf{L} B} &
D_\QCoh(\mathcal{O}_X) \ar[u]_{Lf^*}
}
$$
은 가환이다.
\end{enumerate}
\end{lemma}

\begin{proof}
먼저 도식
$$
\xymatrix{
D(B) \ar[r] & D(\mathcal{O}_Y) \\
D(A) \ar[r] \ar[u]^{- \otimes_A^\mathbf{L} B} &
D(\mathcal{O}_X) \ar[u]_{Lf^*}
}
$$
이 가환임을 증명하자. 이는 Lemma \ref{perfect-lemma-affine-K-flat}과 해당 함자들의
구성에서 분명하다. (1)을 보기 위해 $E$를
$D_\QCoh(\mathcal{O}_X)$의 대상이라 하자. $Lf^*E$의 코호몰로지 층들이
준연접임을 확인하는 일은 $X$ 위에서 국소적으로 할 수 있다. $Lf^*$는 열린
부분스킴으로의 제한과 양립함에 유의하라. 따라서 (2)에서처럼 $f$가 아핀
스킴의 사상이라고 가정할 수 있다. 그러면 Lemma
\ref{perfect-lemma-affine-compare-bounded}에 의해 $E$가 $A$-가군의 복합체에서
옴을 알 수 있다. 증명의 첫 도식이 가환하므로 $Lf^*E$도 마찬가지이며,
(1)이 성립한다.
\end{proof}

\begin{lemma}
\label{perfect-lemma-quasi-coherence-tensor-product}
$X$를 스킴이라 하자.
\begin{enumerate}
\item $D_\QCoh(\mathcal{O}_X)$의 대상 $K, L$에 대해 유도 텐서곱
$K \otimes^\mathbf{L}_{\mathcal{O}_X} L$은
$D_\QCoh(\mathcal{O}_X)$에 속한다.
\item $X = \Spec(A)$가 아핀이면
$$
\widetilde{M^\bullet} \otimes_{\mathcal{O}_X}^\mathbf{L} \widetilde{K^\bullet}
=
\widetilde{M^\bullet \otimes_A^\mathbf{L} K^\bullet}
$$
$A$-가군의 임의의 두 복합체 $K^\bullet$, $M^\bullet$에 대해 위 등식이
성립한다.
\end{enumerate}
\end{lemma}

\begin{proof}
(2)의 등식은 Lemma \ref{perfect-lemma-affine-K-flat}과 유도 텐서곱의 구성에서
즉시 따른다. (1)을 보기 위해 $K, L$을 $D_\QCoh(\mathcal{O}_X)$의
대상이라 하자. $K \otimes^\mathbf{L} L$이
$D_\QCoh(\mathcal{O}_X)$에 속하는지는 $X$ 위에서 국소적으로 확인할 수
있으므로, $X = \Spec(A)$가 아핀이라고 가정할 수 있다. Lemma
\ref{perfect-lemma-affine-compare-bounded}에 의해 $K$와 $L$을 $A$-가군의
복합체들로 나타낼 수 있다. 이제 (2)에서 결론이 따른다.
\end{proof}





\section{전체 직접상}
\label{perfect-section-total-direct-image}

\noindent
다음 보조정리는 Cohomology of Schemes, Lemma
\ref{coherent-lemma-quasi-coherence-higher-direct-images}에 대응한다.

\begin{lemma}
\label{perfect-lemma-quasi-coherence-direct-image}
$f : X \to S$를 스킴의 사상이라 하자. $f$가 준분리이고 준콤팩트라고
가정하자.
\begin{enumerate}
\item 함자 $Rf_*$는 $D_\QCoh(\mathcal{O}_X)$를
$D_\QCoh(\mathcal{O}_S)$ 안으로 보낸다.
\item $S$가 준콤팩트이면 다음 성질을 갖는 정수 $N = N(X, S, f)$가
존재한다. $D_\QCoh(\mathcal{O}_X)$의 대상 $E$가 $m > 0$에 대해
$H^m(E) = 0$을 만족하면, $m \geq N$에 대해 $H^m(Rf_*E) = 0$이다.
\item 실제로 $S$가 준콤팩트이면, 스킴의 모든 사상 $S' \to S$에 대해
함자 $R(f')_*$가 같은 결론을 만족하도록 하는 $N = N(X, S, f)$를 찾을
수 있다. 여기서 $f' : X' \to S'$는 $f$의 밑변환이다.
\end{enumerate}
\end{lemma}

\begin{proof}
$E$를 $D_\QCoh(\mathcal{O}_X)$의 대상이라 하자. (1)을 증명하려면
$Rf_*E$의 코호몰로지 층들이 준연접임을 보여야 한다. 이 문제는 $S$ 위에서
국소적이므로 $S$가 준콤팩트라고 가정할 수 있다. Cohomology of Schemes,
Lemma \ref{coherent-lemma-quasi-coherence-higher-direct-images}에서와 같이
$N = N(X, S, f)$를 택하자. 그러면 모든 준연접 $\mathcal{O}_X$-가군
$\mathcal{F}$와 모든 $p \geq N$에 대해 $R^pf_*\mathcal{F} = 0$이고,
이 사실은 밑변환 뒤에도 성립한다.

\medskip\noindent
먼저 $E$가 아래로 유계라고 가정하자. 이렇게 주어진 $E$와 위에서 택한
$N$에 대해 (1), (2), (3)이 성립함을 보이겠다. 이 경우에는 예를 들어
스펙트럼 열
$$
R^pf_*H^q(E) \Rightarrow R^{p + q}f_*E
$$
(Derived Categories, Lemma \ref{derived-lemma-two-ss-complex-functor}),
$R^pf_*H^q(E)$의 준연접성, 그리고 $p \geq N$일 때
$R^pf_*H^q(E)$가 소멸한다는 사실을 사용하여 (1), (2), (3)이 이 경우에
성립함을 알 수 있다.

\medskip\noindent
다음으로 (2)와 (3)을 증명하자. $m > 0$일 때 $H^m(E) = 0$이라고 하자.
$U \subset S$를 아핀 열린집합이라 하자. Cohomology of Schemes, Lemma
\ref{coherent-lemma-quasi-coherence-higher-direct-images-application}과
$N$의 선택에 의해, 모든 $p \geq N$과 준연접 $\mathcal{O}_X$-가군
$\mathcal{F}$에 대해 $H^p(f^{-1}(U), \mathcal{F}) = 0$이다. 따라서 함자
$\Gamma(f^{-1}(U), -)$에 Lemma
\ref{perfect-lemma-application-nice-K-injective}를 적용하면
$$
R\Gamma(U, Rf_*E) = R\Gamma(f^{-1}(U), E)
$$
의 코호몰로지가 차수 $\geq N$에서 소멸함을 알 수 있다. 이는 모든 아핀
열린집합 $U \subset S$에 대해 성립하므로, $m \geq N$일 때
$H^m(Rf_*E) = 0$이라는 결론을 얻는다.

\medskip\noindent
이제 일반적인 경우에 (1)을 증명하자. $D(\mathcal{O}_X)$에 완전 삼각형
$$
\tau_{\leq -n - 1}E \to E \to \tau_{\geq -n}E \to
(\tau_{\leq -n - 1}E)[1]
$$
이 있음을 상기하라. Derived Categories, Remark
\ref{derived-remark-truncation-distinguished-triangle}을 보라. (2)에 의해
$Rf_*\tau_{\leq -n - 1}E$의 코호몰로지 층들은 차수 $\geq -n + N$에서
소멸한다. 따라서 정수 $q$가 주어지면 어떤 $n$에 대해 $R^qf_*E$가
$R^qf_*\tau_{\geq -n}E$와 같음을 알 수 있고, 위의 결과를 적용할 수 있다.
\end{proof}

\begin{lemma}
\label{perfect-lemma-acyclicity-lemma}
$f : X \to S$를 준분리 준콤팩트 스킴 사상이라 하자.
$\mathcal{F}^\bullet$을 각 항이 $f_*$에 대해 오른쪽 비순환인 준연접
$\mathcal{O}_X$-가군의 복합체라 하자. 그러면
$f_*\mathcal{F}^\bullet$은 $D(\mathcal{O}_S)$에서
$Rf_*\mathcal{F}^\bullet$을 나타낸다.
\end{lemma}

\begin{proof}
표준 사상 $f_*\mathcal{F}^\bullet \to Rf_*\mathcal{F}^\bullet$은 항상
존재한다. 이 사상이 코호몰로지 층들 위에서 동형사상임을 보이면 된다.
명제는 이동 아래 불변이므로
$H^0(f_*(\mathcal{F}^\bullet)) \to R^0f_*\mathcal{F}^\bullet$
가 동형사상임을 보이면 충분하다. 이 명제는 $S$ 위에서 국소적이므로
$S$가 아핀이라고 가정할 수 있다. Lemma
\ref{perfect-lemma-quasi-coherence-direct-image}에 의해 충분히 큰 모든 $n$에 대해
$R^0f_*\mathcal{F}^\bullet = R^0f_*\tau_{\geq -n}\mathcal{F}^\bullet$이다.
따라서 $\mathcal{F}^\bullet$이 아래로 유계라고 가정할 수 있다. 가정에
의해 각 $\mathcal{F}^n$은 $f_*$에 대해 오른쪽 비순환이므로, Leray의
비순환성 보조정리(Derived Categories, Lemma
\ref{derived-lemma-leray-acyclicity})에 의해
$f_*\mathcal{F}^\bullet \to Rf_*\mathcal{F}^\bullet$은 준동형사상이다.
\end{proof}

\begin{lemma}
\label{perfect-lemma-acyclicity-lemma-global}
$X$를 준분리 준콤팩트 스킴이라 하자. $\mathcal{F}^\bullet$을 각 항이
$\Gamma(X, -)$에 대해 오른쪽 비순환인 준연접 $\mathcal{O}_X$-가군의
복합체라 하자. 그러면 $\Gamma(X, \mathcal{F}^\bullet)$은
$D(\Gamma(X, \mathcal{O}_X)$에서 $R\Gamma(X, \mathcal{F}^\bullet)$을
나타낸다.
\end{lemma}

\begin{proof}
표준 사상 $X \to \Spec(\Gamma(X, \mathcal{O}_X))$에 Lemma
\ref{perfect-lemma-acyclicity-lemma}를 적용한다. 몇 가지 세부사항은 생략한다.
\end{proof}

\begin{lemma}
\label{perfect-lemma-spectral-sequence}
$X$를 준분리 준콤팩트 스킴이라 하자. $D_\QCoh(\mathcal{O}_X)$의 임의의
대상 $K$에 대해 Cohomology, Example
\ref{cohomology-example-spectral-sequence}의 스펙트럼 열
$$
E_2^{i, j} = H^i(X, H^j(K)) \Rightarrow H^{i + j}(X, K)
$$
은 유계이고 수렴한다.
\end{lemma}

\begin{proof}
Cohomology, Lemma \ref{cohomology-lemma-spectral-sequence-filtered-object}에
따라 $\tau_{\leq -p}K$가 주는 여과를 사용하여 스펙트럼 열을 구성하면,
주어진 $n \in \mathbf{Z}$에 대해 다음 두 사실을 보이면 충분함을 알 수 있다.
$$
H^n(X, \tau_{\leq -p}K) = 0 \text{단 } p \gg 0
$$
그리고
$$
H^n(X, K) = H^n(X, \tau_{\leq -p}K) \text{단 } p \ll 0
$$
첫 번째는 $F = \Gamma(X, -)$와 Cohomology of Schemes, Lemma
\ref{coherent-lemma-quasi-coherence-higher-direct-images}의 경계에 Lemma
\ref{perfect-lemma-application-nice-K-injective}를 적용하면 따른다. 두 번째는
임의의 환 달린 공간 $(X, \mathcal{O}_X)$과 임의의
$K \in D(\mathcal{O}_X)$에 대해 $-p \leq n$이면 성립한다.
\end{proof}

\begin{lemma}
\label{perfect-lemma-quasi-coherence-pushforward-direct-sums}
$f : X \to S$를 스킴의 준분리 준콤팩트 사상이라 하자. 그러면
$Rf_* : D_\QCoh(\mathcal{O}_X) \to D_\QCoh(\mathcal{O}_S)$
는 직합과 교환한다.
\end{lemma}

\begin{proof}
$E_i$를 $D_\QCoh(\mathcal{O}_X)$의 대상들의 족이라 하고
$E = \bigoplus E_i$로 놓자. 사상
$$
\bigoplus Rf_*E_i \longrightarrow Rf_*E
$$
이 동형임을 보이고자 한다. 이 사상이 차수 $0$의 코호몰로지 층들에
동형을 유도함을 보이면 보조정리가 따른다. 이를 증명할 때 $S$ 위에서
국소적으로 작업해도 되므로, $S$가 준콤팩트라고 가정할 수 있고 그렇게 한다.
Lemma \ref{perfect-lemma-quasi-coherence-direct-image}에서와 같은 정수 $N$을 택하자.
인용한 보조정리에 의해 $R^0f_*E = R^0f_*\tau_{\geq -N}E$이고
$R^0f_*E_i = R^0f_*\tau_{\geq -N}E_i$이다. 또한
$\tau_{\geq -N}E = \bigoplus \tau_{\geq -N}E_i$임에 유의하자.
따라서 모든 $E_i$의 차수 $< -N$인 코호몰로지 층이 소멸한다고 가정할 수 있다.
이제 스펙트럼 열
$$
R^pf_*H^q(E) \Rightarrow R^{p + q}f_*E
\quad\text{그리고}\quad
R^pf_*H^q(E_i) \Rightarrow R^{p + q}f_*E_i
$$
(Derived Categories, Lemma \ref{derived-lemma-two-ss-complex-functor})을 사용하면
준연접 층들의 직합인 경우로 환원된다. 이 경우는
Cohomology of Schemes, Lemma \ref{coherent-lemma-colimit-cohomology}에서 다룬다.
\end{proof}









\section{아핀 사상}
\label{perfect-section-affine-morphisms}

\noindent
이 절에서는 스킴의 아핀 사상을 따른 직접상에 관한 몇 가지 정보를 모은다.

\begin{lemma}
\label{perfect-lemma-pushforward-affine-morphism}
$f : X \to S$를 스킴의 아핀 사상이라 하자. $\mathcal{F}^\bullet$을
준연접 $\mathcal{O}_X$-가군의 복합체라 하자. 그러면
$f_*\mathcal{F}^\bullet = Rf_*\mathcal{F}^\bullet$.
\end{lemma}

\begin{proof}
Lemma \ref{perfect-lemma-acyclicity-lemma}와 Cohomology of Schemes, Lemma
\ref{coherent-lemma-relative-affine-vanishing}를 결합한다.
다른 증명은 $S$ 위에서 아핀 국소적으로 작업하고
Lemma \ref{perfect-lemma-quasi-coherence-pushforward}를 사용하는 것이다.
\end{proof}

\begin{lemma}
\label{perfect-lemma-affine-morphism}
$f : X \to S$를 스킴의 아핀 사상이라 하자. 그러면
$Rf_* : D_\QCoh(\mathcal{O}_X) \to D_\QCoh(\mathcal{O}_S)$
는 동형을 반영한다.
\end{lemma}

\begin{proof}
이 명제는 $Rf_*\alpha$가 동형이면 $D_\QCoh(\mathcal{O}_X)$의 사상
$\alpha : E \to F$도 동형이라는 뜻이다. 이는 코호몰로지 층들에서 확인할 수 있다.
특히 이 문제는 $S$ 위에서 국소적이다. 따라서 $S$와 그에 따라 $X$가
아핀이라고 가정할 수 있다. 이 경우 명제는 Lemma
\ref{perfect-lemma-affine-compare-bounded}에 주어진 유도 범주
$D_\QCoh(\mathcal{O}_X)$와 $D_\QCoh(\mathcal{O}_S)$의 기술에서 명백하다.
세부사항 일부는 생략한다.
\end{proof}

\begin{lemma}
\label{perfect-lemma-affine-morphism-pull-push}
$f : X \to S$를 스킴의 아핀 사상이라 하자.
$D_\QCoh(\mathcal{O}_S)$의 $E$에 대해
$Rf_* Lf^* E = E \otimes^\mathbf{L}_{\mathcal{O}_S} f_*\mathcal{O}_X$.
\end{lemma}

\begin{proof}
$f$가 아핀이므로 사상 $f_*\mathcal{O}_X \to Rf_*\mathcal{O}_X$는 동형이다
(Cohomology of Schemes, Lemma \ref{coherent-lemma-relative-affine-vanishing}).
다음 사상에 수반하는 표준 사상이 있다.
$E \otimes^\mathbf{L} f_*\mathcal{O}_X =
E \otimes^\mathbf{L} Rf_*\mathcal{O}_X \to Rf_* Lf^* E$
$$
Lf^*(E \otimes^\mathbf{L} Rf_*\mathcal{O}_X) =
Lf^*E \otimes^\mathbf{L} Lf^*Rf_*\mathcal{O}_X \longrightarrow
Lf^* E \otimes^\mathbf{L} \mathcal{O}_X = Lf^* E
$$
이 사상은 $1 : Lf^*E \to Lf^*E$와 표준 사상
$Lf^*Rf_*\mathcal{O}_X \to \mathcal{O}_X$에서 나온다. 이렇게 구성한 사상이
동형인지 확인할 때 $S$ 위에서 국소적으로 작업해도 된다. 따라서 $S$와 그에 따라
$X$가 아핀이라고 가정할 수 있다. 이 경우 명제는 Lemmas
\ref{perfect-lemma-affine-compare-bounded}와 \ref{perfect-lemma-quasi-coherence-pullback}에
주어진 유도 범주 $D_\QCoh(\mathcal{O}_X)$와 $D_\QCoh(\mathcal{O}_S)$의
기술 및 함자 $Lf^*$의 기술에서 명백하다. 세부사항 일부는 생략한다.
\end{proof}

\noindent
$Y$를 스킴이라 하자. $\mathcal{A}$를 $\mathcal{O}_Y$-대수의 층이라 하자.
제한 함자 $D(\mathcal{A}) \to D(\mathcal{O}_X)$ 아래에서
$D_\QCoh(\mathcal{O}_X)$의 역상을 $D_\QCoh(\mathcal{A})$로 나타내겠다.
다시 말해, $K \in D(\mathcal{A})$가 $D_\QCoh(\mathcal{A})$에 속할 필요충분조건은
그 코호몰로지 층들이 $\mathcal{O}_X$-가군으로서 준연접인 것이다.
$\mathcal{A}$ 자체가 준연접이면, 이는 코호몰로지 층들이
$\mathcal{A}$-가군으로서 준연접일 것을 요구하는 것과 같다. Morphisms, Lemma
\ref{morphisms-lemma-affine-equivalence-modules}를 보라.

\begin{lemma}
\label{perfect-lemma-affine-morphism-equivalence}
$f : X \to Y$를 스킴의 아핀 사상이라 하자. 그러면 $f_*$는 동치
$$
\Phi : D_\QCoh(\mathcal{O}_X) \longrightarrow D_\QCoh(f_*\mathcal{O}_X)
$$
를 유도하며, 이를 $D_\QCoh(f_*\mathcal{O}_X) \to D_\QCoh(\mathcal{O}_Y)$와
합성하면 $Rf_* : D_\QCoh(\mathcal{O}_X) \to D_\QCoh(\mathcal{O}_Y)$가 된다.
\end{lemma}

\begin{proof}
$Rf_*$는 $K \in D_\QCoh(\mathcal{O}_X)$를 나타내는 $\mathcal{O}_X$-가군의
K-단사 복합체 $\mathcal{I}^\bullet$을 택하고 $f_*\mathcal{I}^\bullet$을 취하여
$K$에서 계산한다는 것을 상기하자. 따라서 $\Phi(K)$를
$f_*\mathcal{O}_X$-가군의 복합체로 본 복합체 $f_*\mathcal{I}^\bullet$로 둔다.
명백한 환 달린 공간의 사상을
$g : (X, \mathcal{O}_X) \to (Y, f_*\mathcal{O}_X)$로 나타내자.
그러면 $g$는 환 달린 공간의 평탄 사상이고(줄기에 관한 기술은 아래를 보라),
$\Phi$는 $Rg_*$를 $D_\QCoh(\mathcal{O}_X)$에 제한한 것이다.
$Lg^*$가 준역함자라고 주장한다. 먼저 $g^*$가 준연접 가군을 준연접 가군으로
보내므로(Modules, Lemma \ref{modules-lemma-pullback-quasi-coherent}),
$Lg^*$는 $D_\QCoh(f_*\mathcal{O}_X)$를 $D_\QCoh(\mathcal{O}_X)$ 안으로 보낸다.
증명을 마치려면 수반 사상
$$
Lg^*\Phi(K) = Lg^*Rg_*K \to K
\quad\text{그리고}\quad
M \to Rg_*Lg^*M = \Phi(Lg^*M)
$$
이 $K \in D_\QCoh(\mathcal{O}_X)$와 $M \in D_\QCoh(f_*\mathcal{O}_X)$에 대해
동형임을 보이면 충분하다. 이는 국소적 문제이므로 $Y$와 그에 따라 $X$가
아핀이라고 가정할 수 있다.

\medskip\noindent
$Y = \Spec(B)$이고 $X = \Spec(A)$라고 가정하자.
$\mathfrak p = x \in \Spec(A) = X$를
$\mathfrak q = y \in \Spec(B) = Y$로 가는 점이라 하자. 그러면
$(f_*\mathcal{O}_X)_y = A_\mathfrak q$이고 $\mathcal{O}_{X, x} = A_\mathfrak p$이므로
$g$는 평탄하다. 따라서 $g^*$는 완전하고, $D(f_*\mathcal{O}_X)$의 임의의
$M$에 대해 $H^i(Lg^*M) = g^*H^i(M)$이다.
$K \in D_\QCoh(\mathcal{O}_X)$에 대해서는
$$
H^i(\Phi(K)) = H^i(Rf_*K) = f_*H^i(K)
$$
임을 고차 직접상의 소멸
(Cohomology of Schemes, Lemma \ref{coherent-lemma-relative-affine-vanishing})과
Lemma \ref{perfect-lemma-application-nice-K-injective}에서 알 수 있다(작은 세부사항은 생략한다).
따라서 다음 사상들이 동형임을 보이면 충분하다.
$$
g^*g_*\mathcal{F} \to \mathcal{F}
\quad\text{그리고}\quad
\mathcal{G} \to g_*g^*\mathcal{G}
$$
여기서 $\mathcal{F}$는 준연접 $\mathcal{O}_X$-가군이고
$\mathcal{G}$는 준연접 $f_*\mathcal{O}_X$-가군이다. 이는 Morphisms, Lemma
\ref{morphisms-lemma-affine-equivalence-modules}에서 따른다.
\end{proof}





\section{닫힌 부분집합에 지지된 코호몰로지}
\label{perfect-section-cohomology-support}

\noindent
스킴과 준연접 가군에 대해 Cohomology, Sections
\ref{cohomology-section-cohomology-support}와
\ref{cohomology-section-cohomology-support-bis}의 내용을 자세히 다룬다.

\begin{definition}
\label{perfect-definition-supported-on}
$X$를 스킴이라 하고, $E$를 $D(\mathcal{O}_X)$의 대상이라 하며,
$T \subset X$를 닫힌 부분집합이라 하자. 코호몰로지 층 $H^i(E)$들이
$T$ 위에 지지되어 있으면 $E$가 {\it $T$ 위에 지지되어 있다}고 한다.
\end{definition}

\noindent
이 정의의 상황에서 Cohomology, Section
\ref{cohomology-section-cohomology-support-bis}의 논의 일부를 되풀이한다.
$X$를 스킴이라 하고 $T \subset X$를 닫힌 부분집합이라 하자.
지지가 $T$에 포함되는 $\mathcal{O}_X$-가군들의 범주는 모든
$\mathcal{O}_X$-가군의 범주의 Serre 부분범주이다. Homology, Definition
\ref{homology-definition-serre-subcategory}와 Modules, Lemma
\ref{modules-lemma-support-section-closed}를 보라. 따라서 Derived Categories,
Section \ref{derived-section-triangulated-sub}의 논의를 적용하면,
$T$ 위에 지지된 대상들로 이루어진 $D(\mathcal{O}_X)$의 엄밀히 충만하고
포화된 삼각 부분범주를 얻는다. 아래에서는 이 범주를
$D_T(\mathcal{O}_X)$로 나타내겠다.

\medskip\noindent
Definition \ref{perfect-definition-supported-on}의 상황에서 포함사상을
$i : T \to X$로 나타내자. 이 상황에서는 함자
$R\mathcal{H}_T : D(\mathcal{O}_X) \to D(i^{-1}\mathcal{O}_X)$
가 있고, 이는 $i_* : D(i^{-1}\mathcal{O}_X) \to D(\mathcal{O}_X)$의
오른쪽 수반임을 Cohomology, Section
\ref{cohomology-section-cohomology-support-bis}에서 상기하자.

\begin{lemma}
\label{perfect-lemma-support-quasi-coherent}
$X$를 스킴이라 하자. $X \setminus T$가 $X$의 역콤팩트 열린집합이 되는
닫힌 부분집합 $T \subset X$를 택하고, $i : T \to X$를 포함사상이라 하자.
\begin{enumerate}
\item $D_\QCoh(\mathcal{O}_X)$의 $E$에 대해
$i_*R\mathcal{H}_T(E)$는 $D_{\QCoh, T}(\mathcal{O}_X)$에 속한다.
\item 함자
$i_* \circ R\mathcal{H}_T : D_\QCoh(\mathcal{O}_X) \to
D_{\QCoh, T}(\mathcal{O}_X)$는 포함 함자
$D_{\QCoh, T}(\mathcal{O}_X) \to D_\QCoh(\mathcal{O}_X)$의 오른쪽 수반이다.
\end{enumerate}
\end{lemma}

\begin{proof}
$U = X \setminus T$로 놓고 포함사상을 $j : U \to X$로 나타내자.
Cohomology, Lemma \ref{cohomology-lemma-triangle-sections-with-support-sheaves}에
따라 $D(\mathcal{O}_X)$에 완전 삼각형
$$
i_*R\mathcal{H}_T(E) \to E \to Rj_*(E|_U) \to i_*R\mathcal{H}_Z(E)[1]
$$
이 있다. Lemma \ref{perfect-lemma-quasi-coherence-direct-image}에 의해 복합체
$Rj_*(E|_U)$의 코호몰로지 층들은 준연접이다(여기서 $U$가 $X$에서
역콤팩트라는 사실을 사용한다). 따라서 (1)이 참이다. (2)는 이 사실과
Cohomology, Lemma \ref{cohomology-lemma-complexes-with-support-on-closed}의
함자들의 수반성에서 따른다.
\end{proof}

\begin{lemma}
\label{perfect-lemma-support-direct-sums}
$X$를 스킴이라 하자. $X \setminus T$가 $X$의 역콤팩트 열린집합이 되는
닫힌 부분집합 $T \subset X$를 택하자. 그러면
$D_\QCoh(\mathcal{O}_X)$의 대상들의 족 $E_i$, $i \in I$에 대해
$R\mathcal{H}_T(\bigoplus E_i) = \bigoplus R\mathcal{H}_T(E_i)$.
\end{lemma}

\begin{proof}
$U = X \setminus T$로 놓고 포함사상을 $j : U \to X$로 나타내자.
Cohomology, Lemma \ref{cohomology-lemma-triangle-sections-with-support-sheaves}에
따라 임의의 $D(\mathcal{O}_X)$의 $E$에 대해 $D(\mathcal{O}_X)$에 완전 삼각형
$$
i_*R\mathcal{H}_T(E) \to E \to Rj_*(E|_U) \to i_*R\mathcal{H}_Z(E)[1]
$$
이 있다. Lemma \ref{perfect-lemma-quasi-coherence-pushforward-direct-sums}에 의해
함자 $E \mapsto Rj_*(E|_U)$는 $D_\QCoh(\mathcal{O}_X)$에서 직합과 교환한다.
따라서 함자 $i_* \circ R\mathcal{H}_T$도 마찬가지이다(세부사항은 생략한다).
$i_* : D(i^{-1}\mathcal{O}_X) \to D_T(\mathcal{O}_X)$가 동치이므로
(Cohomology, Lemma \ref{cohomology-lemma-complexes-with-support-on-closed})
결론이 따른다.
\end{proof}

\begin{remark}
\label{perfect-remark-support-c-equations}
$X$를 스킴이라 하자. $f_1, \ldots, f_c \in \Gamma(X, \mathcal{O}_X)$라 하자.
$f_1, \ldots, f_c$로 잘라낸 닫힌 부분스킴을 $Z \subset X$로 나타내자.
$0 \leq p < c$와 $1 \leq i_0 < \ldots < i_p \leq c$에 대해
$f_{i_0} \ldots f_{i_p}$가 가역인 열린 부분스킴을
$U_{i_0 \ldots i_p} \subset X$로 나타낸다. 임의의 $\mathcal{O}_X$-가군
$\mathcal{F}$에 대해 다음과 같이 놓는다.
$$
\mathcal{F}_{i_0 \ldots i_p} =
(U_{i_0 \ldots i_p} \to X)_*(\mathcal{F}|_{U_{i_0 \ldots i_p}})
$$
이 상황에서 {\it 확장 교대 {\v C}ech 복합체}는 다음
$\mathcal{O}_X$-가군의 복합체이다.
\begin{equation}
\label{perfect-equation-extended-alternating}
0 \to \mathcal{F} \to
\bigoplus\nolimits_{i_0} \mathcal{F}_{i_0} \to
\ldots \to
\bigoplus\nolimits_{i_0 < \ldots < i_p} \mathcal{F}_{i_0 \ldots i_p} \to
\ldots \to \mathcal{F}_{1 \ldots c} \to 0
\end{equation}
여기서 $\mathcal{F}$는 차수 $0$에 둔다. 사상들은 다음과 같이 구성한다.
$1 \leq i_0 < \ldots < i_{p + 1} \leq c$와 $0 \leq j \leq p + 1$이 주어지면
표준 사상
$$
\mathcal{F}_{i_0 \ldots \hat i_j \ldots i_{p + 1}} \to
\mathcal{F}_{i_0 \ldots i_p}
$$
이 있으며, 이는 포함사상
$U_{i_0 \ldots i_p} \subset U_{i_0 \ldots \hat i_j \ldots i_{p + 1}}$에서 나온다.
확장 교대 복합체의 미분은 이 표준 사상들에 부호 $(-1)^j$를 붙여 사용한다.
\end{remark}

\begin{lemma}
\label{perfect-lemma-extended-alternating-zero}
$X$, $f_1, \ldots, f_c \in \Gamma(X, \mathcal{O}_X)$ 및 $\mathcal{F}$가
Remark \ref{perfect-remark-support-c-equations}에서와 같다고 하자. 복합체
(\ref{perfect-equation-extended-alternating})를 $X \setminus Z$ 위로 제한하면
비순환 복합체가 된다.
\end{lemma}

\noindent
더 일반적으로, 유한 열린 덮개 $X \setminus Z = U_1 \cup \ldots \cup U_c$에서
출발하여 Remark \ref{perfect-remark-support-c-equations}에서처럼 정의한 임의의
확장 교대 {\v C}ech 복합체에 대해서도 이 보조정리가 성립한다.

\begin{proof}
$W \subset X \setminus Z$를 열린 부분집합이라 하자. 층의 복합체
(\ref{perfect-equation-extended-alternating})를 $W$에서 계산하면 복합체
$$
\mathcal{F}(W) \to \bigoplus\nolimits_{i_0} \mathcal{F}(U_{i_0} \cap W) \to
\bigoplus\nolimits_{i_0 < i_1} \mathcal{F}(U_{i_0i_1} \cap W) \to \ldots
$$
를 얻는다. 다시 말해, 덮개 $W = \bigcup U_i \cap W$와
$\{1, \ldots, c\}$의 표준 순서에 대한 확장 순서 {\v C}ech 복합체를 얻는다.
Cohomology, Section \ref{cohomology-section-alternating-cech}를 보라.
Cohomology, Lemma \ref{cohomology-lemma-alternating-cech-trivial}에 의해
$W$가 어떤 $1 \leq i \leq c$에 대해 $V(f_i)$에 포함되는 순간 이 복합체는
영과 호모토픽이다. 이것으로 증명이 끝난다.
\end{proof}

\begin{remark}
\label{perfect-remark-extended-alternating-map-to-support}
$X$, $f_1, \ldots, f_c \in \Gamma(X, \mathcal{O}_X)$ 및 $\mathcal{F}$가
Remark \ref{perfect-remark-support-c-equations}에서와 같다고 하자. 복합체
(\ref{perfect-equation-extended-alternating})를 $\mathcal{F}^\bullet$로 나타내자.
Lemma \ref{perfect-lemma-extended-alternating-zero}에 의해 $\mathcal{F}^\bullet$의
코호몰로지 층들은 $Z$ 위에 지지되므로 $\mathcal{F}^\bullet$은
$D_Z(\mathcal{O}_X)$의 대상이다. 한편 등식 $\mathcal{F}^0 = \mathcal{F}$는
$D(\mathcal{O}_X)$에서 표준 사상 $\mathcal{F}^\bullet \to \mathcal{F}$를 정한다.
$i_* \circ R\mathcal{H}_Z$는 포함 함자
$D_Z(\mathcal{O}_X) \to D(\mathcal{O}_X)$의 오른쪽 수반이므로(Cohomology,
Lemma \ref{cohomology-lemma-complexes-with-support-on-closed}를 보라),
표준 가환 그림
$$
\xymatrix{
\mathcal{F}^\bullet \ar[rd] \ar[rr] & & \mathcal{F} \\
& i_*R\mathcal{H}_Z(\mathcal{F}) \ar[ru]
}
$$
을 얻으며, 이는 $D(\mathcal{O}_X)$에서 $\mathcal{O}_X$-가군
$\mathcal{F}$에 대해 함자적이다.
\end{remark}

\begin{lemma}
\label{perfect-lemma-extended-alternating-represented}
$X$, $f_1, \ldots, f_c \in \Gamma(X, \mathcal{O}_X)$ 및 $\mathcal{F}$가
Remark \ref{perfect-remark-support-c-equations}에서와 같다고 하자.
$\mathcal{F}$가 준연접이면 복합체 (\ref{perfect-equation-extended-alternating})는
$D_Z(\mathcal{O}_X)$에서 $i_* R\mathcal{H}_Z(\mathcal{F})$를 나타낸다.
\end{lemma}

\begin{proof}
복합체 (\ref{perfect-equation-extended-alternating})를 $\mathcal{F}^\bullet$로 나타내자.
보조정리의 명제는 Remark \ref{perfect-remark-extended-alternating-map-to-support}의 사상
$\mathcal{F}^\bullet \to i_*R\mathcal{H}_Z(\mathcal{F})$가 동형이라는 뜻이다.
$\mathcal{F}^\bullet$은 $D_Z(\mathcal{O}_X)$에 속하므로(인용한 비고를 보라),
Cohomology, Lemma \ref{cohomology-lemma-complexes-with-support-on-closed}에 의해
$i_*R\mathcal{H}_Z(\mathcal{F}^\bullet) = \mathcal{F}^\bullet$이다.
사상 $U_{i_0 \ldots i_p} \to X$는 $X$의 아핀 열린집합 위에서 함수
$f_{i_0} \ldots f_{i_p}$를 가역화하여 주어지므로 아핀이다. 따라서
$$
\mathcal{F}_{i_0 \ldots i_p} =
(U_{i_0 \ldots i_p} \to X)_*\mathcal{F}|_{U_{i_0 \ldots i_p}} =
R(U_{i_0 \ldots i_p} \to X)_*\mathcal{F}|_{U_{i_0 \ldots i_p}}
$$
임을 Cohomology of Schemes, Lemma
\ref{coherent-lemma-relative-affine-vanishing}와 $\mathcal{F}$가 준연접이라는
가정에서 알 수 있다. Cohomology, Lemma
\ref{cohomology-lemma-sections-support-in-closed-disjoint-open}에 의해
$R\mathcal{H}_Z(\mathcal{F}_{i_0 \ldots i_p}) = 0$이라고 결론짓는다.
따라서 $p > 0$이면 $i_*R\mathcal{H}_Z(\mathcal{F}^p) = 0$이다.
모두 합치면 원하는 대로
$$
\mathcal{F}^\bullet = i_*R\mathcal{H}_Z(\mathcal{F}^\bullet) =
i_*R\mathcal{H}_Z(\mathcal{F})
$$
를 얻는다.
\end{proof}

\begin{lemma}
\label{perfect-lemma-supported-trivial-vanishing}
$X$를 스킴이라 하자. $T \subset X$를 국소적으로 구조층의 원소 최대 $c$개로
잘라낼 수 있는 닫힌 부분집합이라 하자. 그러면 임의의 준연접
$\mathcal{O}_X$-가군 $\mathcal{F}$와 $i > c$에 대해
$\mathcal{H}^i_Z(\mathcal{F}) = 0$이다.
\end{lemma}

\begin{proof}
이는 Lemma \ref{perfect-lemma-extended-alternating-represented}에 주어진
$R\mathcal{H}_T(\mathcal{F})$의 국소적 기술에서 바로 따른다.
\end{proof}

\begin{lemma}
\label{perfect-lemma-supported-vanishing}
$X$를 스킴이라 하자. $Z \subset X$를 $c$개의 원소를 갖는 코줄 정칙 수열로
국소적으로 잘라낼 수 있는 닫힌 부분집합이라 하자. 그러면 모든 평탄 준연접
$\mathcal{O}_X$-가군 $\mathcal{F}$에 대해 $i \not = c$이면
$\mathcal{H}^i_Z(\mathcal{F}) = 0$이다.
\end{lemma}

\begin{proof}
Lemma \ref{perfect-lemma-extended-alternating-represented}에 주어진
$R\mathcal{H}_Z(\mathcal{F})$의 기술에 의해 이는 다음 대수적 명제로
귀결된다. 환 $R$, 코줄 정칙 수열 $f_1, \ldots, f_c \in R$ 및 평탄
$R$-가군 $M$이 주어지면, 확장 교대 {\v C}ech 복합체
$M \to \bigoplus\nolimits_{i_0} M_{f_{i_0}} \to
\bigoplus\nolimits_{i_0 < i_1} M_{f_{i_0}f_{i_1}} \to
\ldots \to M_{f_1 \ldots f_c}$
More on Algebra, Section \ref{more-algebra-section-alternating-cech}에서
유래한 이 복합체는 차수 $c$에서만 코호몰로지를 갖는다. More on Algebra, Lemma
\ref{more-algebra-lemma-vanishing-extended-alternating-koszul}
에 의해 $R$의 확장 교대 {\v C}ech 복합체에 대해 원하는 소멸을 얻는다.
$M$에 대한 복합체는 이를 평탄 $R$-가군 $M$과 텐서하여 얻어진다
(More on Algebra, Lemma
\ref{more-algebra-lemma-extended-alternating-form-module})
이므로 결론이 따른다.
\end{proof}

\begin{remark}
\label{perfect-remark-supported-map-c-equations}
$X$, $f_1, \ldots, f_c \in \Gamma(X, \mathcal{O}_X)$ 및 $\mathcal{F}$가
Remark \ref{perfect-remark-support-c-equations}에서와 같다고 하자.
표준 $\mathcal{O}_X|_Z$-선형 사상
$$
c_{f_1, \ldots, f_c} :
i^*\mathcal{F}
\longrightarrow
\mathcal{H}^c_Z(\mathcal{F})
$$
이 $\mathcal{F}$에 대해 함자적으로 존재한다. 확장 교대 {\v C}ech 복합체
(\ref{perfect-equation-extended-alternating})를 $\mathcal{F}^\bullet$로 나타내면
다음 표준 사상이 있다.
$\mathcal{F}^\bullet \to i_*R\mathcal{H}_Z(\mathcal{F})$
이는 Remark \ref{perfect-remark-extended-alternating-map-to-support}의 사상이다.
이로부터 표준 사상
$$
\Coker\left(\bigoplus \mathcal{F}_{1 \ldots \hat i \ldots c} \to
\mathcal{F}_{1 \ldots c}\right)
\longrightarrow
i_*\mathcal{H}^c_Z(\mathcal{F})
$$
을 코호몰로지 층의 차수 $c$에서 얻는다. $\mathcal{F}$의 국소 단면 $s$가
주어지면 다음 국소 단면을 생각할 수 있다.
$$
\frac{s}{f_1 \ldots f_c}
$$
이는 $\mathcal{F}_{1 \ldots c}$의 단면이다. 위에 표시한 코커널에서 이 단면의
류는 다음 사상의 상을 법으로 한 $s$에만 의존한다.
$(f_1, \ldots, f_c) : \mathcal{F}^{\oplus c} \to \mathcal{F}$.
 $i_*i^*\mathcal{F}$는 다음 사상의 코커널과 같으므로,
$(f_1, \ldots, f_c) : \mathcal{F}^{\oplus c} \to \mathcal{F}$
 $\mathcal{O}_X$-가군 사상
$i_*i^*\mathcal{F} \to i_*\mathcal{H}_Z^c(\mathcal{F})$.
$i_*$가 충실충만하므로 사상 $c_{f_1, \ldots, f_c}$를 얻는다.
\end{remark}

\begin{example}
\label{perfect-example-affine-supported-map-c-equations}
$X = \Spec(A)$를 아핀이라 하고, $f_1, \ldots, f_c \in A$라 하자.
$A$-가군 $M$에 대해 $\mathcal{F} = \widetilde{M}$로 놓자. Remark
\ref{perfect-remark-supported-map-c-equations}의 $c_{f_1, \ldots, f_c}$는 다음 사상으로
기술할 수 있다.
$$
M/(f_1, \ldots, f_c)M
\longrightarrow
\Coker\left(
\bigoplus M_{f_1 \ldots \hat f_i \ldots f_c} \to M_{f_1 \ldots f_c}
\right)
$$
$s \in M$의 류를 코커널에서 $s/f_1 \ldots f_c$의 류로 보내는 사상이다.
\end{example}

\begin{lemma}
\label{perfect-lemma-supported-map-determinant}
$X$, $f_1, \ldots, f_c \in \Gamma(X, \mathcal{O}_X)$ 및 $\mathcal{F}$가
Remark \ref{perfect-remark-support-c-equations}에서와 같다고 하자.
$1 \leq i, j \leq c$에 대해 $a_{ji} \in \Gamma(X, \mathcal{O}_X)$라 하고
$g_j = \sum_{i = 1, \ldots, c} a_{ji}f_i$로 놓자. $g_1, \ldots, g_c$가
스킴으로 $Z$를 잘라낸다고 가정하자. $\mathcal{F}$가 준연접이면
$$
c_{f_1, \ldots, f_c} = \det(a_{ji}) c_{g_1, \ldots, g_c}
$$
$c_{f_1, \ldots, f_c}$와 $c_{g_1, \ldots, g_c}$는 Remark
\ref{perfect-remark-supported-map-c-equations}에서와 같다.
\end{lemma}

\begin{proof}
증명하고자 하는 것은 임의의 $s \in \mathcal{F}(X)$에 대해
$\mathcal{H}_Z(\mathcal{F})$의 전역 단면으로서
$c_{f_1, \ldots, f_c}(s) = \det(a_{ij}) c_{g_1, \ldots, g_c}(s)$라는 것이다.
그러면 $X$의 임의의 열린 부분스킴 위에서도 같은 결과를 얻으므로 충분하다.
이를 위해 $1 \leq i_0 < \ldots < i_p \leq c$ 및
$1 \leq j_0 < \ldots < j_q \leq c$에 대해 다음 열린 부분스킴들을 나타내자.
$U_{i_0 \ldots i_p} \subset X$, $V_{j_0 \ldots j_q} \subset X$,
$W_{i_0 \ldots i_p, j_0 \ldots j_q} \subset X$는 각각
$f_{i_0} \ldots f_{i_p}$가 가역인 곳, $g_{j_0} \ldots g_{j_q}$가 가역인 곳,
그리고 $f_{i_0} \ldots f_{i_p}g_{j_0} \ldots g_{j_q}$가 가역인 곳이다.
$\mathcal{F}_{i_0 \ldots i_p}$, 각각 $\mathcal{F}'_{j_0 \ldots j_q}$,
$\mathcal{F}''_{i_0 \ldots i_p, j_0 \ldots j_q}$는 각각
$\mathcal{F}$를 $U_{i_0 \ldots i_p}$, $V_{j_0 \ldots j_q}$,
$W_{i_0 \ldots i_p, j_0 \ldots j_q}$에 제한한 것의 $X$로의 직접상이다.
그러면 다음 세 개의 확장 교대 {\v C}ech 복합체를 얻는다.
$$
\mathcal{F}^\bullet :
\mathcal{F} \to \bigoplus\nolimits_{i_0} \mathcal{F}_{i_0} \to
\bigoplus\nolimits_{i_0 < i_1} \mathcal{F}_{i_0i_1} \to \ldots
$$
그리고
$$
(\mathcal{F}')^\bullet :
\mathcal{F} \to \bigoplus\nolimits_{j_0} \mathcal{F}'_{j_0} \to
\bigoplus\nolimits_{j_0 < j_1} \mathcal{F}'_{j_0j_1} \to \ldots
$$
그리고
$$
(\mathcal{F}'')^\bullet :
\mathcal{F} \to
\bigoplus\nolimits_{i_0} \mathcal{F}_{i_0} \oplus
\bigoplus\nolimits_{j_0} \mathcal{F}'_{j_0} \to
\bigoplus\nolimits_{i_0 < i_1} \mathcal{F}_{i_0i_1} \oplus
\bigoplus\nolimits_{i_0, j_0} \mathcal{F}''_{i_0, j_0} \oplus
\bigoplus\nolimits_{j_0 < j_1} \mathcal{F}'_{j_0j_1} \to
\ldots
$$
그 미분은 (\ref{perfect-equation-extended-alternating})을 정의할 때 사용한 것들이다.
복합체의 사상
$$
(\mathcal{F}'')^\bullet \to \mathcal{F}^\bullet
\quad\text{그리고}\quad
(\mathcal{F}'')^\bullet \to (\mathcal{F}')^\bullet
$$
은 각 항에서 사영사상으로 주어진다(따라서 차수 $0$에서 항등사상을 유도한다).
Lemma \ref{perfect-lemma-extended-alternating-represented}에 의해 이 복합체들은 각각
$i_*R\mathcal{H}_Z(\mathcal{F})$를 나타내며, 이 사상들은 이 대상 위의
항등사상을 나타낸다. 따라서 원소
$$
\sigma \in H^c((\mathcal{F}'')^\bullet(X))
$$
를 찾아 이 두 사상에 의해 각각 $c_{f_1, \ldots, f_c}(s)$와
$\det(a_{ji})c_{g_1, \ldots, g_c}(s)$로 보내면 충분하다.
$\sigma$에 대한 코사이클을 명시적으로 쓸 수 있다. 즉 다음과 같이 둔다.
$$
\sigma_{1 \ldots c} = \frac{s}{f_1 \ldots f_c} \in \mathcal{F}_{1 \ldots c}(X)
\quad\text{그리고}\quad
\sigma'_{1 \ldots c} = \frac{\det(a_{ji})s}{g_1 \ldots g_c} \in
\mathcal{F}'_{1 \ldots c}(X)
$$
그리고 다음과 같이 둔다.
$$
\sigma_{i_0 \ldots i_p, j_0 \ldots j_{c - p - 2}} =
\frac{\lambda(i_0 \ldots i_p, j_0 \ldots j_{c - p - 2})s}%
{f_{i_0} \ldots f_{i_p}g_{j_0} \ldots g_{j_{c - p - 2}}}
\in \mathcal{F}''_{i_0 \ldots i_p, j_0 \ldots j_{c - p - 2}}(X)
$$
여기서 $\lambda(i_0 \ldots i_p, j_0 \ldots j_{c - p - 2})$는 다음 형식적
표현에서 $e_1 \wedge \ldots \wedge e_c$의 계수이다.
$$
e_{i_0} \wedge \ldots \wedge e_{i_p} \wedge
(a_{j_01} e_1 + \ldots + a_{j_0c}e_c) \wedge \ldots \wedge
(a_{j_{c - p - 2}1} e_1 + \ldots + a_{j_{c - p - 2}c}e_c)
$$
 $\sigma$가 코사이클임을 확인하려면
$1 \leq i_0 < \ldots < i_p \leq c$ 및
$1 \leq j_0 < \ldots < j_{c - p - 1} \leq c$에 대해 다음을 보여야 한다.
\begin{align*}
0 & =
\sum\nolimits_{a = 0, \ldots, p} (-1)^a
f_{i_a} \lambda(i_0 \ldots \hat i_a \ldots i_p, j_0 \ldots j_{c - p - 1}) \\
& +
\sum\nolimits_{b = 0, \ldots, c - p - 1} (-1)^{p + b + 1}g_{j_b}
\lambda(i_0 \ldots i_p, j_0 \ldots \hat j_b \ldots j_{c - p - 1})
\end{align*}
이를 보는 가장 쉬운 방법은 아마 다음 형식적 표현이
$$
\xi = e_{i_0} \wedge \ldots \wedge e_{i_p} \wedge
(a_{j_01} e_1 + \ldots + a_{j_0c}e_c) \wedge \ldots \wedge
(a_{j_{c - p - 1}1} e_1 + \ldots + a_{j_{c - p - 1}c}e_c)
$$
 $e_1, \ldots, e_c$ 위 자유 가군의 $(c + 1)$차 외적 거듭제곱의 원소이므로
 $0$이라는 것과, 위 표현이 $e_i \to f_i$를 보내는 코줄 미분에 의한
$\xi$의 상이라는 사실로 논증하는 것이다. 세부사항은 생략한다.
\end{proof}

\begin{lemma}
\label{perfect-lemma-supported-map-global}
$X$를 스킴이라 하자. $Z \to X$를 유한 표시인 닫힌 매장이라 하고,
그 conormal 층 $\mathcal{C}_{Z/X}$가 계수 $c$인 국소 자유라고 하자.
그러면 표준 사상
$$
c :
\wedge^c(\mathcal{C}_{Z/X})^\vee \otimes_{\mathcal{O}_Z} i^*\mathcal{F}
\longrightarrow
\mathcal{H}_Z^c(\mathcal{F})
$$
이 준연접 가군 $\mathcal{F}$에 대해 함자적으로 존재한다.
\end{lemma}

\begin{proof}
이는 Remark \ref{perfect-remark-supported-map-c-equations}의 구성과
Lemma \ref{perfect-lemma-supported-map-determinant}에서 보인 아이디얼 층 생성원
선택의 독립성에서 따른다. 세부사항은 생략한다.
\end{proof}

\begin{remark}
\label{perfect-remark-supported-functorial}
$g : X' \to X$를 스킴의 사상이라 하자. 다음을 택하자.
$f_1, \ldots, f_c \in \Gamma(X, \mathcal{O}_X)$.
$f'_i = g^\sharp(f_i) \in \Gamma(X', \mathcal{O}_{X'})$로 놓자.
 $f_1, \ldots, f_c$ 및 $f'_1, \ldots, f'_c$로 잘라낸 닫힌 부분스킴을 각각
$Z \subset X$, $Z' \subset X'$로 나타내자.
그러면 $Z' = Z \times_X X'$이다. 유도된 사상 $h : Z' \to Z$를 스킴의
사상으로 나타내자. $\mathcal{F}$를 $\mathcal{O}_X$-가군이라 하고
$\mathcal{F}' = g^*\mathcal{F}$로 놓자. 이 상황에서 $\mathcal{F}$가
준연접이면 다음 그림
$$
\xymatrix{
(i')^{-1}\mathcal{O}_{X'} \otimes_{h^{-1}i^{-1}\mathcal{O}_X}
h^{-1}\mathcal{H}^c_Z(\mathcal{F}) \ar[r] &
\mathcal{H}_{Z'}^c(\mathcal{F}') \\
h^*i^*\mathcal{F} \ar[r] \ar[u]_-{c_{f_1, \ldots, f_c}} &
(i')^*\mathcal{F}' \ar[u]^-{c_{f'_1, \ldots, f'_c}}
}
$$
은 가환하며, 위쪽 수평 사상은 Cohomology, Remark
\ref{cohomology-remark-support-functorial}에서 코호몰로지 층 차수 $c$에
대해 주어진 사상이다. 구체적으로 Remark \ref{perfect-remark-support-c-equations}에서
$\mathcal{F}, f_1, \ldots, f_c$ 및 $\mathcal{F}', f'_1, \ldots, f'_c$를
사용하여 구성한 확장 교대 {\v C}ech 복합체를 각각
$\mathcal{F}^\bullet$, $(\mathcal{F}')^\bullet$로 나타내자.
그러면 $(\mathcal{F}')^\bullet = g^*\mathcal{F}^\bullet$이다.
$\mathcal{F}$가 준연접이라는 가정 없이도 다음 그림
$$
\xymatrix{
i'_* L(g|_{Z'})^* R\mathcal{H}_Z(\mathcal{F}) \ar[r] \ar@{=}[d] &
i'_*R\mathcal{H}_{Z'}(Lg^*\mathcal{F}) \ar[d] \\
Lg^*i_*R\mathcal{H}_Z(\mathcal{F}) &
i'_*R\mathcal{H}_{Z'}(\mathcal{F}') \\
Lg^*(\mathcal{F}^\bullet) \ar[u] \ar[r] &
(\mathcal{F}')^\bullet \ar[u]
}
$$
은 가환한다. 여기서 $g|_{Z'} : (Z', (i')^{-1}\mathcal{O}_{X'}) \to
(Z, i^{-1}\mathcal{O}_X)$는 환 달린 공간의 유도된 사상이다.
위쪽 수평 사상과 등호에 대한 설명은 Cohomology, Remark
\ref{cohomology-remark-support-functorial}에서 주어진다.
위쪽을 향한 화살표들은 Remark \ref{perfect-remark-extended-alternating-map-to-support}에서
온다. 아래쪽 수평 사상은 $Lg^*\mathcal{F}^\bullet \to
g^*\mathcal{F}^\bullet = (\mathcal{F}')^\bullet$이고 아래쪽을 향한 화살표는
$Lg^*\mathcal{F} \to g^*\mathcal{F} = \mathcal{F}'$에서 유도된다.
그림이 가환하는 이유는 그림을 어느 방향으로 돌아도 두 사상
$Lg^*\mathcal{F}^\bullet \to
i'_*R\mathcal{H}_{Z'}(\mathcal{F}')$와
$i'_*R\mathcal{H}_{Z'}(\mathcal{F}') \to \mathcal{F}'$의 합성이
표준 사상 $Lg^*\mathcal{F}^\bullet \to \mathcal{F}'$가 되기 때문이다.
세부사항은 생략한다. 이제 첫 번째 그림의 가환성은 이 그림을 코호몰로지 층의
차수 $c$에서 살펴보고 다음 사상의 구성이
$i^*\mathcal{F} \to
\Coker(\bigoplus \mathcal{F}_{1 \ldots \hat i \ldots c} \to
\mathcal{F}_{1 \ldots c})$
Remark \ref{perfect-remark-supported-map-c-equations}에서 사용한 구성과
밑당김에 양립한다는 사실을 이용하면 따른다.
\end{remark}








\section{준연접화자}
\label{perfect-section-coherator}

\noindent
$X$를 스킴이라 하자. {\it 준연접화자}는 다음 함자이다.
$$
Q_X :
\textit{Mod}(\mathcal{O}_X)
\longrightarrow
\QCoh(\mathcal{O}_X)
$$
이는 포함 함자
$\QCoh(\mathcal{O}_X) \to \textit{Mod}(\mathcal{O}_X)$.
이는 모든 스킴 $X$에 대해 존재하며, 더 나아가 모든 준연접 가군
$\mathcal{F}$에 대해 수반 사상 $Q_X(\mathcal{F}) \to \mathcal{F}$가
동형이다. Properties, Proposition \ref{properties-proposition-coherator}를 보라.
$Q_X$는 (오른쪽 수반이므로) 왼쪽 완전이므로 그 오른쪽 유도 확장
$$
RQ_X :
D(\mathcal{O}_X)
\longrightarrow
D(\QCoh(\mathcal{O}_X)).
$$
 $Q_X$가 포함 함자
$\QCoh(\mathcal{O}_X) \to \textit{Mod}(\mathcal{O}_X)$
의 오른쪽 수반이라는 사실에서, Derived Categories, Lemma
\ref{derived-lemma-derived-adjoint-functors}에 의해 $RQ_X$는 표준 함자
$D(\QCoh(\mathcal{O}_X)) \to D(\mathcal{O}_X)$의 오른쪽 수반이다.

\medskip\noindent
이 절에서는 함자 $RQ_X$를 연구한다. Section
\ref{perfect-section-better-coherator}에서는 (밀접하게 관련된) 포함 함자
$D_\QCoh(\mathcal{O}_X) \to D(\mathcal{O}_X)$의 오른쪽 수반(존재하는 경우)을
연구한다.

\begin{lemma}
\label{perfect-lemma-affine-pushforward}
$f : X \to Y$를 스킴의 아핀 사상이라 하자. 그러면 $f_*$는 유도 함자
$f_* : D(\QCoh(\mathcal{O}_X)) \to D(\QCoh(\mathcal{O}_Y))$.
를 정의하며, 이 함자는 다음 성질을 갖는다.
$$
\xymatrix{
D(\QCoh(\mathcal{O}_X)) \ar[d]_{f_*} \ar[r] &
D_\QCoh(\mathcal{O}_X) \ar[d]^{Rf_*} \\
D(\QCoh(\mathcal{O}_Y)) \ar[r] &
D_\QCoh(\mathcal{O}_Y)
}
$$
은 가환한다.
\end{lemma}

\begin{proof}
함자
$f_* : \QCoh(\mathcal{O}_X) \to \QCoh(\mathcal{O}_Y)$
는 완전하다. Cohomology of Schemes, Lemma
\ref{coherent-lemma-relative-affine-vanishing}을 보라. 따라서 $f_*$는
유도 함자
$f_* : D(\QCoh(\mathcal{O}_X)) \to D(\QCoh(\mathcal{O}_Y))$
를 임의의 대표 복합체에 단순히 $f_*$를 적용하여 정의한다. Derived Categories,
Lemma \ref{derived-lemma-right-derived-exact-functor}를 보라.
그림의 가환성은 Lemma \ref{perfect-lemma-pushforward-affine-morphism}에서 따른다.
\end{proof}

\begin{lemma}
\label{perfect-lemma-flat-pushforward-coherator}
$f : X \to Y$를 스킴의 사상이라 하자. $f$가 준콤팩트, 준분리, 평탄하다고
가정하자. 다음을 나타내자.
$$
\Phi : D(\QCoh(\mathcal{O}_X)) \to D(\QCoh(\mathcal{O}_Y))
$$
 $f_* : \QCoh(\mathcal{O}_X) \to \QCoh(\mathcal{O}_Y)$의 오른쪽 유도 함자를
나타내면 $RQ_Y \circ Rf_* = \Phi \circ RQ_X$이다.
\end{lemma}

\begin{proof}
이는 $RQ_Y \circ Rf_*$와 $\Phi \circ RQ_X$가 같은 함자
$D(\QCoh(\mathcal{O}_Y)) \to D(\mathcal{O}_X)$의 오른쪽 수반임을 보여
증명한다.

\medskip\noindent
 $f$가 준콤팩트이고 준분리이므로 $f_*$는 준연접성을 보존한다. Schemes,
Lemma \ref{schemes-lemma-push-forward-quasi-coherent}를 보라.
 $\QCoh(\mathcal{O}_X)$가 Grothendieck 아벨 범주임을 상기하자
(Properties, Proposition \ref{properties-proposition-coherator}). 따라서
$D(\QCoh(\mathcal{O}_X))$의 임의의 $K$는 $\QCoh(\mathcal{O}_X)$의
K-단사 복합체 $\mathcal{I}^\bullet$로 나타낼 수 있다. 다음을 보라.
Injectives, Theorem
\ref{injectives-theorem-K-injective-embedding-grothendieck}.
그러면 $\Phi(K) = f_*\mathcal{I}^\bullet$로 정의할 수 있다.

\medskip\noindent
 $f$가 평탄하므로 함자 $f^*$는 완전하다. 따라서 $f^*$는
$f^* : D(\mathcal{O}_Y) \to D(\mathcal{O}_X)$ 및
$f^* : D(\QCoh(\mathcal{O}_Y)) \to D(\QCoh(\mathcal{O}_X))$를 정의한다.
함자 $f^* = Lf^* : D(\mathcal{O}_Y) \to D(\mathcal{O}_X)$는
$Rf_* : D(\mathcal{O}_X) \to D(\mathcal{O}_Y)$의 왼쪽 수반이다.
Cohomology, Lemma \ref{cohomology-lemma-adjoint}를 보라.
마찬가지로 함자
$f^* : D(\QCoh(\mathcal{O}_Y)) \to D(\QCoh(\mathcal{O}_X))$
는 Derived Categories, Lemma \ref{derived-lemma-derived-adjoint-functors}에 의해
$\Phi : D(\QCoh(\mathcal{O}_X)) \to D(\QCoh(\mathcal{O}_Y))$의
왼쪽 수반이다.

\medskip\noindent
$A$를 $D(\QCoh(\mathcal{O}_Y))$의 대상이라 하고
$E$를 $D(\mathcal{O}_X)$의 대상이라 하자. 그러면
\begin{align*}
\Hom_{D(\QCoh(\mathcal{O}_Y))}(A, RQ_Y(Rf_*E))
& =
\Hom_{D(\mathcal{O}_Y)}(A, Rf_*E) \\
& =
\Hom_{D(\mathcal{O}_X)}(f^*A, E) \\
& =
\Hom_{D(\QCoh(\mathcal{O}_X))}(f^*A, RQ_X(E)) \\
& =
\Hom_{D(\QCoh(\mathcal{O}_Y))}(A, \Phi(RQ_X(E)))
\end{align*}
이로부터 원하는 결론이 따른다.
\end{proof}

\begin{lemma}
\label{perfect-lemma-affine-coherator}
$X = \Spec(A)$를 아핀 스킴이라 하자. 그러면
\begin{enumerate}
\item $Q_X : \textit{Mod}(\mathcal{O}_X) \to \QCoh(\mathcal{O}_X)$
는 다음 함자이다.
이는 $\mathcal{F}$를 $A$-가군 $\Gamma(X, \mathcal{F})$에 대응하는
준연접 $\mathcal{O}_X$-가군으로 보낸다.
\item $RQ_X : D(\mathcal{O}_X) \to D(\QCoh(\mathcal{O}_X))$
는 $E$를 $D(A)$의 대상 $R\Gamma(X, E)$에 대응하는 준연접
$\mathcal{O}_X$-가군의 복합체로 보내는 함자이고,
\item $D_\QCoh(\mathcal{O}_X)$에 제한하면 함자 $RQ_X$는
(\ref{perfect-equation-compare})의 준역함자를 정의한다.
\end{enumerate}
\end{lemma}

\begin{proof}
$Q_X$는 다음 함자이다.
$$
\mathcal{F} \mapsto \widetilde{\Gamma(X, \mathcal{F})}
$$
Schemes, Lemma \ref{schemes-lemma-compare-constructions}에 의해서이다.
이는 즉시 (1)과 (2)를 뜻한다. 세 번째 주장은 Lemma
\ref{perfect-lemma-affine-compare-bounded}의 (증명에서) 따른다.
\end{proof}

\noindent
이제 함자 $D(\QCoh(\mathcal{O}_X)) \to D_\QCoh(\mathcal{O}_X)$가
동치가 되는 조건을 증명할 준비가 되었다.

\begin{lemma}
\label{perfect-lemma-argument-proves}
 $X$를 준콤팩트 준분리 스킴이라 하자. 다음을 가정하자.
 $X$의 모든 아핀 열린집합 $U \subset X$에 대해 오른쪽 유도 함자
$$
\Phi : D(\QCoh(\mathcal{O}_U)) \to D(\QCoh(\mathcal{O}_X))
$$
왼쪽 완전 함자 $j_* : \QCoh(\mathcal{O}_U) \to \QCoh(\mathcal{O}_X)$가
다음 가환 그림에 들어간다고 가정하자.
$$
\xymatrix{
D(\QCoh(\mathcal{O}_U)) \ar[d]_\Phi \ar[r]_{i_U} &
D_\QCoh(\mathcal{O}_U) \ar[d]^{Rj_*} \\
D(\QCoh(\mathcal{O}_X)) \ar[r]^{i_X} &
D_\QCoh(\mathcal{O}_X)
}
$$
그러면 함자 (\ref{perfect-equation-compare})
$$
D(\QCoh(\mathcal{O}_X))
\longrightarrow
D_\QCoh(\mathcal{O}_X)
$$
는 동치이고 준역함자는 $RQ_X$이다.
\end{lemma}

\begin{proof}
$E$를 $D_\QCoh(\mathcal{O}_X)$의 대상이라 하고 $A$를
$D(\QCoh(\mathcal{O}_X))$의 대상이라 하자. 수반 사상
$$
A \to RQ_X(i_X(A))
\quad\text{그리고}\quad
i_X(RQ_X(E)) \to E
$$
이 동형임을 보여야 한다. 다음 가정
$H_n$을 생각하자. 즉, $E$와 $i_X(A)$가 (Definition
\ref{perfect-definition-supported-on})의 의미에서 $X$의 $n$개 아핀 열린집합의
합집합에 포함되는 $X$의 닫힌 부분집합 위에 지지되면 위 수반 사상들이
동형이라는 가정이다. $H_n$을 $n$에 대한 귀납법으로 증명한다.

\medskip\noindent
기초 단계: $n = 0$. 이 경우 $E = 0$이므로 사상
$i_X(RQ_X(E)) \to E$는 동형이다. 마찬가지로 $i_X(A) = 0$이다.
따라서 $i_X(A)$의 코호몰로지 층들은 영이다. 포함 함자
$\QCoh(\mathcal{O}_X) \to \textit{Mod}(\mathcal{O}_X)$가 충실충만하고
완전하므로 $A$의 코호몰로지 대상들도 영이다. 즉 $A = 0$이고
$A \to RQ_X(i_X(A))$도 동형이다.

\medskip\noindent
귀납 단계. $E$와 $i_X(A)$가 $U_1 \cup \ldots \cup U_n$에 포함되는
닫힌 부분집합 $T$ 위에 지지되어 있다고 하자. 여기서 $U_i \subset X$는
아핀 열린집합이다. $U = U_n$으로 놓자. 다음 완전 삼각형을 생각하자.
$$
A \to \Phi(A|_U) \to A' \to A[1]
\quad\text{그리고}\quad
E \to Rj_*(E|_U) \to E' \to E[1]
$$
여기서 $\Phi$는 보조정리의 명제에서와 같다. $E \to Rj_*(E|_U)$가
$U = U_n$ 위에서 준동형임에 유의하자. 가정에 의해
$i_X \circ \Phi = Rj_* \circ i_U$이고 $i_X(A)|_U = i_U(A|_U)$이므로,
$i_X(A) \to i_X(\Phi(A|_U))$는 $U$ 위에서 준동형이다.
따라서 $i_X(A')$와 $E'$는 $X$의 $T \setminus U$라는 닫힌 부분집합 위에
지지되며, 이는 $U_1 \cup \ldots \cup U_{n - 1}$에 포함된다. 귀납 가정에 의해
$A'$와 $E'$에 대해서는 명제가 참이다. Derived Categories, Lemma
\ref{derived-lemma-third-isomorphism-triangle}에 의해 다음 사상들이
동형임을 보이는 것으로 충분하다.
$$
\Phi(A|_U) \to RQ_X(i_X(\Phi(A|_U)))
\quad\text{그리고}\quad
i_X(RQ_X(Rj_*E|_U)) \to Rj_*(E|_U)
$$
은 동형이다. 가정과 Lemma \ref{perfect-lemma-flat-pushforward-coherator}에 의해
(포함 사상 $j : U \to X$는 평탄, 준콤팩트, 준분리이다) 다음을 얻는다.
$$
RQ_X(i_X(\Phi(A|_U))) = RQ_X(Rj_*(i_U(A|_U))) = \Phi(RQ_U(i_U(A|_U)))
$$
그리고
$$
i_X(RQ_X(Rj_*(E|_U))) = i_X(\Phi(RQ_U(E|_U))) = Rj_*(i_U(RQ_U(E|_U)))
$$
마지막으로 사상
$$
A|_U \to RQ_U(i_U(A|_U))
\quad\text{그리고}\quad
i_U(RQ_U(E|_U)) \to E|_U
$$
은 Lemma \ref{perfect-lemma-affine-coherator}에 의해 동형이다. 결과가 따른다.
\end{proof}

\begin{proposition}
\label{perfect-proposition-quasi-compact-affine-diagonal}
$X$를 아핀 대각선을 갖는 준콤팩트 스킴이라 하자. 그러면 함자
(\ref{perfect-equation-compare})
$$
D(\QCoh(\mathcal{O}_X))
\longrightarrow
D_\QCoh(\mathcal{O}_X)
$$
는 동치이고 준역함자는 $RQ_X$이다.
\end{proposition}

\begin{proof}
$U \subset X$를 아핀 열린집합이라 하자. 그러면 $U \to X$는
Morphisms, Lemma \ref{morphisms-lemma-affine-permanence}에 의해 아핀이다.
따라서 Lemma \ref{perfect-lemma-affine-pushforward}에 의해 Lemma
\ref{perfect-lemma-argument-proves}의 가정이 성립하므로 결론을 얻는다.
\end{proof}

\begin{lemma}
\label{perfect-lemma-direct-image-coherator}
$f : X \to Y$를 스킴의 사상이라 하자. $X$와 $Y$가 준콤팩트이고
아핀 대각선을 갖는다고 가정하자. 다음을 나타내자.
$$
\Phi : D(\QCoh(\mathcal{O}_X)) \to D(\QCoh(\mathcal{O}_Y))
$$
 $f_* : \QCoh(\mathcal{O}_X) \to \QCoh(\mathcal{O}_Y)$의 오른쪽 유도 함자이다.
그러면 다음 그림
$$
\xymatrix{
D(\QCoh(\mathcal{O}_X)) \ar[d]_\Phi \ar[r] &
D_\QCoh(\mathcal{O}_X) \ar[d]^{Rf_*} \\
D(\QCoh(\mathcal{O}_Y)) \ar[r] &
D_\QCoh(\mathcal{O}_Y)
}
$$
은 가환한다.
\end{lemma}

\begin{proof}
그림의 수평 화살표들은 Proposition
\ref{perfect-proposition-quasi-compact-affine-diagonal}에 의해 범주의 동치이다.
따라서 이 범주들을 동일시할 수 있다(아핀 대각선을 갖는 다른 준콤팩트
스킴들에 대해서도 마찬가지이다). 보조정리의 명제는
$D(\QCoh(\mathcal{O}_X))$의 모든 $K$에 대해 표준 사상
$\Phi(K) \to Rf_*(K)$가 동형이라는 것이다. $K_1 \to K_2 \to K_3 \to K_1[1]$가
$D(\QCoh(\mathcal{O}_X))$의 완전 삼각형이고 두 대상에 대해 명제가
성립하면 나머지 대상에 대해서도 성립한다는 점에 유의하자.

\medskip\noindent
 $U \subset X$를 아핀 열린집합이라 하자. $X$의 대각선이 아핀이므로
포함 사상 $j : U \to X$는 아핀이다(Morphisms, Lemma
\ref{morphisms-lemma-affine-permanence}). 마찬가지로 합성
$g = f \circ j : U \to Y$도 아핀이다. $\mathcal{I}^\bullet$를
$\QCoh(\mathcal{O}_U)$의 K-단사 복합체라 하자. $j_* :
\QCoh(\mathcal{O}_U) \to \QCoh(\mathcal{O}_X)$가 완전한 왼쪽 수반
$j^* : \QCoh(\mathcal{O}_X) \to \QCoh(\mathcal{O}_U)$를 가지므로,
$j_*\mathcal{I}^\bullet$는 $\QCoh(\mathcal{O}_X)$의 K-단사 복합체이다. 다음을 보라.
Derived Categories, Lemma \ref{derived-lemma-adjoint-preserve-K-injectives}.
따라서
$$
\Phi(j_*\mathcal{I}^\bullet) =
f_*j_*\mathcal{I}^\bullet =
g_*\mathcal{I}^\bullet
$$
Lemma \ref{perfect-lemma-affine-pushforward}에 의해 $j_*\mathcal{I}^\bullet$는
$Rj_*\mathcal{I}^\bullet$를 나타내고 $g_*\mathcal{I}^\bullet$는
$Rg_*\mathcal{I}^\bullet$를 나타낸다. 한편 $Rf_* \circ Rj_* = Rg_*$이다.
따라서 $f_*j_*\mathcal{I}^\bullet$는 $Rf_*(j_*\mathcal{I}^\bullet)$를 나타낸다.
그러므로 $U$ 위 준연접 가군들의 복합체 $\mathcal{G}^\bullet$에 대해
$j_*\mathcal{G}^\bullet$ 꼴인 임의의 복합체에서 보조정리가 성립한다.
( $\mathcal{G}^\bullet \to \mathcal{I}^\bullet$가 준동형이면 $j_*$가
준연접 가군에서 완전 함자이므로 $j_*\mathcal{G}^\bullet \to
j_*\mathcal{I}^\bullet$도 준동형이라는 점에 유의하자.)

\medskip\noindent
$\mathcal{F}^\bullet$를 준연접 $\mathcal{O}_X$-가군들의 복합체라 하자.
$T \subset X$를 모든 $p$에 대해 $\mathcal{F}^p$의 지지가 $T$에 포함되는
닫힌 부분집합이라 하자. $T \subset U_1 \cup \ldots \cup U_n$이 되게 하는
아핀 열린집합 $U_1, \ldots, U_n$의 최소 개수 $n$에 대한 귀납법을 사용한다.
기초 단계 $n = 0$은 자명하다. $n \geq 1$이면 $U = U_1$로 놓고 위와 같이
열린 매장 $j : U \to X$를 나타내자. 다음 복합체의 사상을 생각하자.
$c : \mathcal{F}^\bullet \to j_*j^*\mathcal{F}^\bullet$.
그러면 다음 두 개의 짧은 완전열을 얻는다.
$$
0 \to \Ker(c) \to \mathcal{F}^\bullet \to \Im(c) \to 0
$$
그리고
$$
0 \to \Im(c) \to j_*j^*\mathcal{F}^\bullet \to \Coker(c) \to 0
$$
$\Ker(c)$와 $\Coker(c)$는 $T \setminus U \subset U_2 \cup \ldots \cup U_n$인
닫힌 부분집합 위에 지지되므로 귀납법에 의해 이들에 대해 결과가 성립한다.
앞 단락의 논의에 의해 $j_*j^*\mathcal{F}^\bullet$에 대해서도 결과가 성립한다.
두 짧은 완전열에 대응하는 완전 삼각형을 살펴보고 증명의 처음 비고를
사용하면 결론을 얻는다.
\end{proof}

\begin{remark}[주의]
\label{perfect-remark-warning-coherator}
 $X$를 아핀 대각선을 갖는 준콤팩트 스킴이라 하자. Proposition
\ref{perfect-proposition-quasi-compact-affine-diagonal}에 의해
$D(\QCoh(\mathcal{O}_X)) = D_\QCoh(\mathcal{O}_X)$임을 알고 있지만,
이 자료에서는 이상한 일이 일어날 수 있고 실수하기 쉽다. 이 등식이
유도 함자들도 같다는 뜻이라고 부주의하게 가정하는 것이 한 가지 함정이다.
예를 들어 준콤팩트 열린집합 $U \subset X$가 있다고 하자. 다음 오른쪽
유도 함자들을 생각할 수 있다.
$$
R^i(\QCoh)\Gamma(U, -) : \QCoh(\mathcal{O}_X) \to \textit{Ab}
$$
왼쪽 완전 함자 $\Gamma(U, -)$에서 유래한다. 이는 보편적인
$\delta$-함자이고, $H^i(U, -)$ 함자들($X$ 위의 모든 아벨 층에 대해
정의됨)을 $\QCoh(\mathcal{O}_X)$에 제한하면 $\delta$-함자를 이루므로,
표준 변환
$$
t^i : R^i(\QCoh)\Gamma(U, -) \to H^i(U, -).
$$
를 얻는다. $X = \Spec(A)$가 아핀인 경우에도 이 변환들은 일반적으로
동형이 아니다! 실제로 오른쪽 유도 함자의 구성과
$\QCoh(\mathcal{O}_X)$와 $\text{Mod}_A$의 동치에 의해 $I$가 단사
$A$-가군이면 $R^1(\QCoh)\Gamma(U, \widetilde{I}) = 0$이다.
그러나 Examples, Lemma \ref{examples-lemma-nonvanishing}는 어떤 $A$, $I$,
$U$에 대해 $H^1(U, \widetilde{I}) \not = 0$임을 보인다.
\end{remark}




\section{Noether 스킴의 준연접화자}
\label{perfect-section-coherator-Noetherian}

\noindent
Noether 스킴의 경우 다음 보조정리를 사용할 수 있다.

\begin{lemma}
\label{perfect-lemma-injective-quasi-coherent-sheaf-Noetherian}
$X$를 Noether 스킴이라 하자. $\mathcal{J}$를
$\QCoh(\mathcal{O}_X)$의 단사 대상이라 하자. 그러면 $\mathcal{J}$는
$\mathcal{O}_X$-가군의 플라스크 층이다.
\end{lemma}

\begin{proof}
$U \subset X$를 열린 부분집합이라 하고 $s \in \mathcal{J}(U)$를 단면이라 하자.
$\mathcal{I} \subset \mathcal{O}_X$를 $X \setminus U$의 축약된 유도 스킴 구조를
정의하는 아이디얼의 준연접 층이라 하자(Schemes, Definition
\ref{schemes-definition-reduced-induced-scheme}). Cohomology of Schemes, Lemma
\ref{coherent-lemma-homs-over-open}에 의해 단면 $s$는 어떤 $n$에 대한 사상
$\sigma : \mathcal{I}^n \to \mathcal{J}$에 대응한다. $\mathcal{J}$가
$\QCoh(\mathcal{O}_X)$의 단사 대상이므로 $\sigma$를 사상
$\tilde s : \mathcal{O}_X \to \mathcal{J}$로 확장할 수 있다. 그러면 $\tilde s$는
$s$로 제한되는 $\mathcal{J}$의 전역 단면에 대응한다.
\end{proof}

\begin{lemma}
\label{perfect-lemma-Noetherian-pushforward}
$f : X \to Y$를 Noether 스킴의 사상이라 하자. 준연접 층에서의 $f_*$는
다음 오른쪽 유도 확장을 갖는다.
$\Phi : D(\QCoh(\mathcal{O}_X)) \to D(\QCoh(\mathcal{O}_Y))$
그리고 다음 그림은
$$
\xymatrix{
D(\QCoh(\mathcal{O}_X)) \ar[d]_{\Phi} \ar[r] &
D_\QCoh(\mathcal{O}_X) \ar[d]^{Rf_*} \\
D(\QCoh(\mathcal{O}_Y)) \ar[r] &
D_\QCoh(\mathcal{O}_Y)
}
$$
가환한다.
\end{lemma}

\begin{proof}
 $X$와 $Y$가 Noether 스킴이므로 이 사상은 준콤팩트이고 준분리이다
(Properties, Lemma \ref{properties-lemma-locally-Noetherian-quasi-separated} 및
Schemes, Remark \ref{schemes-remark-quasi-compact-and-quasi-separated}).
따라서 $f_*$는 준연접성을 보존한다. Schemes, Lemma
\ref{schemes-lemma-push-forward-quasi-coherent}를 보라. 이제
$K$를 $D(\QCoh(\mathcal{O}_X))$의 대상으로 하자. $\QCoh(\mathcal{O}_X)$가
Grothendieck 아벨 범주이므로(Properties, Proposition
\ref{properties-proposition-coherator}), $K$를 각 $\mathcal{I}^n$이
$\QCoh(\mathcal{O}_X)$의 단사 대상인 K-단사 복합체 $\mathcal{I}^\bullet$로
나타낼 수 있다. 다음을 보라.
Injectives, Theorem
\ref{injectives-theorem-K-injective-embedding-grothendieck}.
따라서 함자 $\Phi$는 다음과 같이 정의된다.
$$
\Phi(K) = f_*\mathcal{I}^\bullet
$$
여기서 우변은 $D(\QCoh(\mathcal{O}_Y))$의 대상으로 본다. 보조정리의 증명을
끝내려면 표준 사상
$$
f_*\mathcal{I}^\bullet \longrightarrow Rf_*\mathcal{I}^\bullet
$$
가 $D(\mathcal{O}_Y)$에서 동형임을 보이면 충분하다. Lemma
\ref{perfect-lemma-acyclicity-lemma}에 의해 이를 위해 모든 $n \in \mathbf{Z}$에 대해
$\mathcal{I}^n$이 $f_*$-오른쪽 비순환임을 보이면 충분하다.
이는 Lemma \ref{perfect-lemma-injective-quasi-coherent-sheaf-Noetherian}에 의해
$\mathcal{I}^n$이 플라스크이고, Cohomology, Lemma
\ref{cohomology-lemma-flasque-acyclic-pushforward}에 의해 플라스크 가군이
$f_*$-오른쪽 비순환이기 때문이다.
\end{proof}

\begin{proposition}
\label{perfect-proposition-Noetherian}
$X$를 Noether 스킴이라 하자. 그러면 함자 (\ref{perfect-equation-compare})
$$
D(\QCoh(\mathcal{O}_X))
\longrightarrow
D_\QCoh(\mathcal{O}_X)
$$
는 동치이고 준역함자는 $RQ_X$이다.
\end{proposition}

\begin{proof}
이는 Lemma \ref{perfect-lemma-argument-proves}와 Lemma
\ref{perfect-lemma-Noetherian-pushforward}에서 따른다.
\end{proof}






\section{코줄 복합체}
\label{perfect-section-koszul}

\noindent
$A$를 환이라 하고 $f_1, \ldots, f_r$를 $A$의 원소들의 수열이라 하자.
More on Algebra, Definition \ref{more-algebra-definition-koszul-complex}에서
코줄 복합체 $K_\bullet(f_1, \ldots, f_r)$를 정의했다. 이는 차수
$r, \ldots, 0$에 놓인 사슬 복합체이다. $K^{-n}(f_1, \ldots, f_r) =
K_n(f_1, \ldots, f_r)$로 놓고 같은 미분을 사용하여 이를 코사슬 복합체
$K^\bullet(f_1, \ldots, f_r)$로 바꾼다. 이 절의 나머지에서는 모든 복합체를
코사슬 복합체로 한다.

\medskip\noindent
$I^\bullet(f_1, \ldots, f_r)$라는 복합체를 정의하여 다음 완전 삼각형을 갖게 하자.
$$
I^\bullet(f_1, \ldots, f_r) \to
A \to
K^\bullet(f_1, \ldots, f_r) \to
I^\bullet(f_1, \ldots, f_r)[1]
$$
$K(A)$에서. 다시 말해 다음과 같이 놓는다.
$$
I^i(f_1, \ldots, f_r) =
\left\{
\begin{matrix}
K^{i - 1}(f_1, \ldots, f_r) & i \leq 0\text{일 때} \\
0 & \text{그 밖의 경우}
\end{matrix}
\right.
$$
그리고 $K^\bullet(f_1, \ldots, f_r)$의 미분의 음수를 사용한다.
완전 삼각형의 사상들은 자명한 것들이다. $I^0(f_1, \ldots, f_r) =
A^{\oplus r} \to A$가 $i$번째 인자에 $f_i$를 곱하는 사상임에 유의하자.
따라서 $I^\bullet(f_1, \ldots, f_r) \to A$는 다음과 같이 인수분해된다.
$$
I^\bullet(f_1, \ldots, f_r) \to I \to A
$$
$I = (f_1, \ldots, f_r)$이다. 실제로 다음 짧은 완전열이 있다.
$$
0 \to H^{-1}(K^\bullet(f_1, \ldots, f_s)) \to
H^0(I^\bullet(f_1, \ldots, f_s)) \to I \to 0
$$
그리고 모든 $i < 0$에 대해
$H^i(I^\bullet(f_1, \ldots, f_r)) = H^{i - 1}(K^\bullet(f_1, \ldots, f_r))$이다.
$A$의 원소들의 두 번째 수열 $g_1, \ldots, g_r$가 주어지면 다음 표준
사상들이 있다.
$$
I^\bullet(f_1g_1, \ldots, f_rg_r) \to I^\bullet(f_1, \ldots, f_r)
\quad\text{그리고}\quad
K^\bullet(f_1g_1, \ldots, f_rg_r) \to K^\bullet(f_1, \ldots, f_r)
$$
이 사상들은 위에서 기술한 사상들과 양립한다. 첫 번째 사상은
$I^0(f_1g_1, \ldots, f_rg_r) = A^{\oplus r}$의 $i$번째 직합 성분에서 $g_i$를
곱하는 사상이다. 특히 $f_1, \ldots, f_r$가 주어지면 다음 복합체의 역계를 얻는다.
\begin{equation}
\label{perfect-equation-system}
I^\bullet(f_1, \ldots, f_r) \leftarrow
I^\bullet(f_1^2, \ldots, f_r^2) \leftarrow
I^\bullet(f_1^3, \ldots, f_r^3) \leftarrow \ldots
\end{equation}
이는 다음 내용에서 중요한 역할을 한다. 다음 보조정리들을 쉽게 서술하기 위해
표기를 고정하자.

\begin{situation}
\label{perfect-situation-complex}
여기서 $A$는 환이고 $f_1, \ldots, f_r$는 $A$의 원소들의 수열이다.
$X = \Spec(A)$ 및 $U = D(f_1) \cup \ldots \cup D(f_r) \subset X$로 놓는다.
주어진 $U$의 열린 덮개 $\mathcal{U} : U = \bigcup_{i = 1, \ldots, r} D(f_i)$를
나타낸다.
\end{situation}

\noindent
첫 번째 보조정리는 위 복합체들로 $U$ 위 준연접 층의 코호몰로지를 계산할 수
있다는 것이다.

\begin{lemma}
\label{perfect-lemma-alternating-cech-complex}
Situation \ref{perfect-situation-complex}에서 $M$을 $A$-가군이라 하고 이에 대응하는
$\mathcal{O}_X$-가군을 $\mathcal{F}$로 나타내자. 그러면 다음 복합체의 표준
동형이 있다.
$$
\Psi : \colim_e \Hom_A(I^\bullet(f_1^e, \ldots, f_r^e), M)
\longrightarrow
\check{\mathcal{C}}_{alt}^\bullet(\mathcal{U}, \mathcal{F})
$$
$M$에 대해 함자적이며, $\Hom$-복합체의 미분은
$I^\bullet(f_1^e, \ldots, f_r^e)$의 미분의 반대 쌍대이다.
\end{lemma}

\begin{proof}
교대 {\v C}ech 복합체는 통상적인 {\v C}ech 복합체의 교대 코체인으로 이루어진
부분복합체이다. Cohomology, Section \ref{cohomology-section-alternating-cech}를
보라. 평소처럼 $\check{\mathcal{C}}_{alt}^\bullet(\mathcal{U}, \mathcal{F})$의
$p$-코체인을 $\{1, \ldots, r\}^{p + 1}$ 위의 교대 함수 $s$로 본다.
$(i_0, \ldots, i_p)$에서의 값 $s_{i_0\ldots i_p}$는
$M_{f_{i_0}\ldots f_{i_p}} = \mathcal{F}(U_{i_0\ldots i_p})$에 속한다.
한편 $\Hom^\bullet(I^\bullet(f_1^e, \ldots, f_r^e), M)$의
$p$-코체인 $t$는 사상 $t : \wedge^{p + 1}(A^{\oplus r}) \to M$이다.
$[i] \in A^{\oplus r}$를 $i$번째 기저 원소로 쓰고 다음과 같이 쓰자.
$$
[i_0, \ldots, i_p] = [i_0] \wedge \ldots \wedge [i_p]
\in \wedge^{p + 1}(A^{\oplus r})
$$
위와 같은 $t$에 대해 다음과 같이 놓는다.
$$
\Psi(t)_{i_0 \ldots i_p} =
(-1)^p \frac{t([i_0, \ldots, i_p])}{f_{i_0}^e\ldots f_{i_p}^e}
$$
 $\Psi(t)$가 교대 코체인임은 명백하다. 위 규칙은 다음 역계의 전이 사상
$$
I^\bullet(f_1^e, \ldots, f_r^e) \leftarrow
I^\bullet(f_1^{e + 1}, \ldots, f_r^{e + 1}),
$$
 (\ref{perfect-equation-system})가 $[i_0, \ldots, i_p]$를
$f_{i_0}\ldots f_{i_p}[i_0, \ldots, i_p]$로 보내는 것과 양립한다.
Algebra, Lemma \ref{algebra-lemma-localization-colimit}의
$M_{f_{i_0} \ldots f_{i_p}}$ 국소화 기술에 의해, $\Psi$ 규칙은 코체인
가군의 귀납극한에서 차수 $p$의 동형을 정의한다. 증명을 끝내려면 이 사상이
미분과 양립함을 보여야 한다. 이를 위해 위와 같은 $t$에 대해 계산하면
\begin{align*}
d(\Psi(t))_{i_0 \ldots i_{p + 1}}
& =
\sum\nolimits_{j = 0}^{p + 1} (-1)^j
\Psi(t)_{i_0\ldots \hat i_j \ldots i_{p + 1}} \\
& =
(-1)^p
\sum\nolimits_{j = 0}^{p + 1} (-1)^j t([i_0 \ldots \hat i_j \ldots i_{p + 1}])
(f_{i_0} \ldots \hat f_{i_j} \ldots f_{i_p})^{-e}
\end{align*}
 $I^\bullet(f_1^e, \ldots, f_r^e)$의 미분은
$K^\bullet(f_1, \ldots, f_r)$의 미분의 음수임을 상기하자. 따라서
\begin{align*}
\Psi(d(t))_{i_0 \ldots i_{p + 1}} & =
(-1)^{p + 1}
d(t)([i_0, \ldots, i_{p + 1}]) (f_{i_0} \ldots f_{i_{p + 1}})^{-e} \\
& =
(-1)^{p + 1}
t(d([i_0, \ldots, i_{p + 1}])) (f_{i_0} \ldots f_{i_{p + 1}})^{-e} \\
& =
(-1)^{p + 1}
t(-\sum\nolimits_{j = 0}^{p + 1}
(-1)^j f_{i_j}^e [i_0, \ldots, \hat i_j, \ldots i_{p + 1}])
(f_{i_0} \ldots f_{i_{p + 1}})^{-e} \\
& =
-(-1)^{p + 1}
\sum\nolimits_{j = 0}^{p + 1}
(-1)^j t([i_0, \ldots, \hat i_j, \ldots i_{p + 1}])
(f_{i_0} \ldots \hat f_{i_j} \ldots f_{i_p})^{-e}
\end{align*}
두 공식이 일치하므로 증명이 끝난다.
\end{proof}

\noindent
 $A$-가군의 유한 복합체 $I^\bullet$와 $A$-가군의 복합체 $M^\bullet$가
주어졌다고 하자. 이에 대응하는 Hom 복합체
$\Hom^\bullet(I^\bullet, M^\bullet)$가 있다. 이는 차수 $n$ 항이
다음으로 주어지는 복합체이다.
$$
\bigoplus\nolimits_{p + q = n} \Hom_A(I^{-q}, M^p)
$$
미분은 More on Algebra, Section \ref{more-algebra-section-hom-complexes}에서
기술한 대로이다. 복합체 $I^\bullet$에는 영이 아닌 항이 유한 개뿐이므로
위에 표시한 직합은 유한하다. 복합체의 {\v C}ech 복합체에 대응하는 전체
복합체를 취하는 규약은 Cohomology, Section
\ref{cohomology-section-cech-cohomology-of-complexes}에서와 같다.

\begin{lemma}
\label{perfect-lemma-alternating-cech-complex-complex}
Situation \ref{perfect-situation-complex}에서 $M^\bullet$를 $A$-가군의 복합체라 하고
이에 대응하는 $\mathcal{O}_X$-가군 복합체를 $\mathcal{F}^\bullet$로 나타내자.
그러면 다음 복합체의 표준 동형이 있다.
$$
\colim_e \Hom^\bullet(I^\bullet(f_1^e, \ldots, f_r^e), M^\bullet)
\longrightarrow
\text{Tot}(\check{\mathcal{C}}_{alt}^\bullet(\mathcal{U}, \mathcal{F}^\bullet))
$$
$M^\bullet$에 대해 함자적이다.
\end{lemma}

\begin{proof}
 $F^{\bullet, \bullet}$를 항 $F^{p, q} =
\mathcal{C}_{alt}^p(\mathcal{U}, \mathcal{F}^q)$로 갖는 이중 복합체라 하자.
이는 Cohomology, Section \ref{cohomology-section-cech-cohomology-of-complexes}에서
논의했다. $G^{\bullet, \bullet}$를 항
$G^{p, q} = \colim_e \Hom_A(I^{-p}(f_1^e, \ldots, f_r^e), M^q)$로 갖고
부호를 개입하지 않는 함자성으로 미분을 주는 이중 복합체라 하자.
Lemma \ref{perfect-lemma-alternating-cech-complex}의 증명에서 구성한 사상
$\psi^{p, q} : G^{p, q} \to F^{p, q}$는 동형이고 미분 $d_1$(보조정리에 의해)과
$d_2$(명백함)와 양립한다. 그러나 보조정리의 화살표 양쪽 복합체에서 미분
$d$의 부호는 다르다. 구체적으로 $g \in G^{p, q}$에 대해
$$
d(g) =  d_2(g) - (-1)^{p + q} d_1(g) 
$$
 (More on Algebra, Section \ref{more-algebra-section-hom-complexes}를 보라)이고,
$f \in F^{p, q}$에 대한 미분은
$$
d(f) = d_1(f) + (-1)^p d_2(f)
$$
따라서 $\psi^{p, q}$에 $(-1)^{pq + p(p - 1)/2}$를 곱하면 부호를 맞출 수 있다.
\end{proof}

\begin{lemma}
\label{perfect-lemma-alternating-cech-complex-complex-computes-cohomology}
Situation \ref{perfect-situation-complex}에서 $\mathcal{F}^\bullet$를 준연접
$\mathcal{O}_X$-가군들의 복합체라 하자. 그러면 다음 표준 동형이 있다.
$$
\text{Tot}(\check{\mathcal{C}}_{alt}^\bullet(\mathcal{U}, \mathcal{F}^\bullet))
\longrightarrow
R\Gamma(U, \mathcal{F}^\bullet)
$$
$D(A)$에서 $\mathcal{F}^\bullet$에 대해 함자적이다.
\end{lemma}

\begin{proof}
$\mathcal{B}$를 $U$의 아핀 열린집합들의 집합이라 하자. 아핀 스킴 위
준연접 가군의 고차 코호몰로지 군은 영이므로(Cohomology of Schemes, Lemma
\ref{coherent-lemma-quasi-coherent-affine-cohomology-zero}), 이는
Cohomology, Lemma \ref{cohomology-lemma-alternating-cech-complex-complex-ss}의
특수한 경우이다.
\end{proof}

\noindent
Situation \ref{perfect-situation-complex}에서 Lemma
\ref{perfect-lemma-affine-compare-bounded}의 동치를 통해 $A$-가군 복합체
$I^\bullet(f_1^e, \ldots, f_r^e)$에 대응하는 $D(\mathcal{O}_X)$의 대상을
$I_e$로 나타내자. 사상 (\ref{perfect-equation-system})은 다음 계를 준다.
$$
I_1 \leftarrow
I_2 \leftarrow
I_3 \leftarrow \ldots
$$
더 나아가 $U$에 제한하면 동형이 되는 호환되는 사상 계
$I_e \to \mathcal{O}_X$가 있다. 따라서 $D(\mathcal{O}_X)$의 모든 대상
$E$에 대해 다음 표준 사상이 있다.
\begin{equation}
\label{perfect-equation-comparison}
\colim_e \Hom_{D(\mathcal{O}_X)}(I_e, E) \longrightarrow H^0(U, E)
\end{equation}
이는 사상 $I_e \to E$를 $U$로 제한하여 보내고
$\Hom_{D(\mathcal{O}_U)}(\mathcal{O}_U, E|_U) = H^0(U, E)$임을
사용하여 구성한다.

\begin{proposition}
\label{perfect-proposition-represent-cohomology-class-on-open}
Situation \ref{perfect-situation-complex}에서 $D_\QCoh(\mathcal{O}_X)$의 모든 대상
$E$에 대해 사상 (\ref{perfect-equation-comparison})은 동형이다.
\end{proposition}

\begin{proof}
Lemma \ref{perfect-lemma-affine-compare-bounded}에 의해 $E$가 준연접 층들의
복합체 $\mathcal{F}^\bullet$로 주어진다고 가정할 수 있다.
$M^\bullet = \Gamma(X, \mathcal{F}^\bullet)$를 이에 대응하는
$A$-가군 복합체라 하자. 다음 보조정리들에 의해
보조정리 \ref{perfect-lemma-alternating-cech-complex-complex} 및
\ref{perfect-lemma-alternating-cech-complex-complex-computes-cohomology}에 의해 다음
준동형들이 있다.
$$
\colim_e \Hom^\bullet(I^\bullet(f_1^e, \ldots, f_r^e), M^\bullet)
\longrightarrow
\text{Tot}(\check{\mathcal{C}}_{alt}^\bullet(\mathcal{U}, \mathcal{F}^\bullet))
\longrightarrow
R\Gamma(U, \mathcal{F}^\bullet)
$$
More on Algebra, Lemma \ref{more-algebra-lemma-RHom-out-of-projective}와
Equation (\ref{more-algebra-equation-h0-RHom})에 의해 복합체
$\Hom^\bullet(I^\bullet(f_1^e, \ldots, f_r^e), M^\bullet)$의 $H^0$를 취하면
$D(A)$에서의 $\Hom$을 계산한다. 따라서 양변에서 $H^0$를 취하면
$$
\colim_e \Hom_{D(A)}(I^\bullet(f_1^e, \ldots, f_r^e), M^\bullet)
=
H^0(U, E)
$$
Lemma \ref{perfect-lemma-affine-compare-bounded}에 의해
$\Hom_{D(A)}(I^\bullet(f_1^e, \ldots, f_r^e), M^\bullet) =
\Hom_{D(\mathcal{O}_X)}(I_e, E)$이므로 보조정리가 따른다.
\end{proof}

\noindent
Situation \ref{perfect-situation-complex}에서 Lemma
\ref{perfect-lemma-affine-compare-bounded}의 동치를 통해 $A$-가군 복합체
$K^\bullet(f_1^e, \ldots, f_r^e)$에 대응하는 $D(\mathcal{O}_X)$의 대상을
$K_e$로 나타내자. 그러면 다음 완전 삼각형들과
$$
I_e \to \mathcal{O}_X \to K_e \to I_e[1]
$$
다음 계를 갖는다.
$$
K_1 \leftarrow
K_2 \leftarrow
K_3 \leftarrow \ldots
$$
이는 계 $(I_e)$와 호환된다. 더 나아가 호환되는 사상 계
$$
K_e \to H^0(K_e) = \mathcal{O}_X/(f_1^e, \ldots, f_r^e)
$$

\begin{lemma}
\label{perfect-lemma-represent-cohomology-class-on-closed}
Situation \ref{perfect-situation-complex}에서 $E$를
$D_\QCoh(\mathcal{O}_X)$의 대상이라 하자. $i = -r + 1, \ldots, 0$에 대해
$H^i(E)|_U = 0$이라고 가정하자. 그러면 $s \in H^0(X, E)$가 주어질 때
어떤 $e \geq 0$와 사상 $K_e \to E$가 존재하여
$s$가 $H^0(X, K_e) \to H^0(X, E)$의 상에 속한다.
\end{lemma}

\begin{proof}
 $U$가 $r$개의 아핀 열린집합으로 덮이므로 임의의 준연접 가군에 대해
$j \geq r$이면 $H^j(U, \mathcal{F}) = 0$이다(Cohomology of Schemes,
Lemma \ref{coherent-lemma-vanishing-nr-affines}). Lemma
\ref{perfect-lemma-application-nice-K-injective}에 의해 $H^0(U, E)$가
$H^0(U, \tau_{\geq -r + 1}E)$와 같음을 알 수 있다. 다음 스펙트럼 열이 있다.
$$
H^j(U, H^i(\tau_{\geq -r + 1}E)) \Rightarrow H^{i + j}(U, \tau_{\geq -N}E)
$$
Derived Categories, Lemma \ref{derived-lemma-two-ss-complex-functor}를 보라.
따라서 $E$의 코호몰로지 층이 소멸한다는 가정에 의해 $H^0(U, E) = 0$이다.
그러므로 $s|_U = 0$이다. $s$를 $D(\mathcal{O}_X)$에서의 사상
$\mathcal{O}_X \to E$로 생각하자. Proposition
\ref{perfect-proposition-represent-cohomology-class-on-open}에 의해 어떤 $e$에 대해
합성 $I_e \to \mathcal{O}_X \to E$가 영이다. 완전 삼각형
$I_e \to \mathcal{O}_X \to K_e \to I_e[1]$에 의해 사상 $K_e \to E$를 얻으며,
$s$는 합성 $\mathcal{O}_X \to K_e \to E$가 된다.
\end{proof}


\section{유사연접 및 완전 복합체}
\label{perfect-section-spell-out}

\noindent
이 절에서는
Cohomology의 절들
\ref{cohomology-section-strictly-perfect},
\ref{cohomology-section-pseudo-coherent},
\ref{cohomology-section-tor},
\ref{cohomology-section-perfect}에서 정의한 일반 개념들과
More on Algebra의 절들
\ref{more-algebra-section-pseudo-coherent},
\ref{more-algebra-section-tor},
\ref{more-algebra-section-perfect}에서 가군 복합체에 대해 대응하는 개념들 사이의
관계를 정리한다.

\begin{lemma}
\label{perfect-lemma-pseudo-coherent}
$X$를 스킴이라 하자. $E$가 $D(\mathcal{O}_X)$의
$m$-유사연접 대상이면, $i > m$에 대해 $H^i(E)$는 준연접
$\mathcal{O}_X$-가군이고 $H^m(E)$는 준연접
$\mathcal{O}_X$-가군의 몫이다.
$E$가 유사연접이면 $E$는
$D_\QCoh(\mathcal{O}_X)$의 대상이다.
\end{lemma}

\begin{proof}
$X$ 위에서 국소적으로 엄밀히 완전한 복합체 $\mathcal{E}^\bullet$가 존재하여
$i > m$일 때 $H^i(E)$가 $H^i(\mathcal{E}^\bullet)$와 동형이고,
$H^m(E)$가 $H^m(\mathcal{E}^\bullet)$의 몫이 된다. 층
$\mathcal{E}^i$는 유한 자유 가군의 직합 인자이므로 준연접이다.
따라서 보조정리가 따른다.
\end{proof}

\begin{lemma}
\label{perfect-lemma-pseudo-coherent-affine}
$X = \Spec(A)$를 아핀 스킴이라 하자. $M^\bullet$를
$A$-가군의 복합체라 하고 $E$를 $D(\mathcal{O}_X)$의 대응하는 대상으로
두자. 그러면 $E$가 $D(\mathcal{O}_X)$의 대상으로서
$m$-유사연접(각각 유사연접)인 것은 $M^\bullet$가 $A$-가군의
복합체로서 $m$-유사연접(각각 유사연접)인 것과 동치이다.
\end{lemma}

\begin{proof}
정의에서 $M^\bullet$가 $m$-유사연접이면 $E$도 그러함은 자명하다.
역을 보이기 위해 $E$가 $m$-유사연접이라고 하자. $X = \Spec(A)$는
준콤팩트이고 표준 열린집합들이 위상의 기저를 이루므로, 표준 열린 덮개
$X = D(f_1) \cup \ldots \cup D(f_n)$와 각 $D(f_i)$ 위의 엄밀히 완전한
복합체 $\mathcal{E}_i^\bullet$, 그리고
$H^j$에서 $j > m$이면 동형을, $H^m$에서는 전사를 유도하는 사상
$\alpha_i : \mathcal{E}_i^\bullet \to E|_{U_i}$를 잡을 수 있다.
열린 덮개를 세분한 뒤 Cohomology, Lemma
\ref{cohomology-lemma-local-actual}에 의해 각 $i$에 대해 $\alpha_i$가
복합체의 사상
$\mathcal{E}_i^\bullet \to \widetilde{M^\bullet}|_{U_i}$로 주어진다고
가정할 수 있다. Modules, Lemma
\ref{modules-lemma-direct-summand-of-locally-free-is-locally-free}에 의해
항들 $\mathcal{E}_i^n$은 유한 국소 자유 가군이다.
따라서 열린 덮개를 다시 세분하여 각 $\mathcal{E}_i^n$이 유한 자유
$\mathcal{O}_{U_i}$-가군이라고 가정할 수 있다.
정의로부터 $M^\bullet_{f_i}$는 $A_{f_i}$-가군의
$m$-유사연접 복합체이다. 이제
More on Algebra, Lemma \ref{more-algebra-lemma-glue-pseudo-coherent}를
적용하면 결론을 얻는다.

\medskip\noindent
``유사연접''인 경우는, 정의에 의해 $E$가 유사연접일 필요충분조건이
모든 $m$에 대해 $E$가 $m$-유사연접인 것이라는 사실과,
More on Algebra, Lemma \ref{more-algebra-lemma-pseudo-coherent}에 의해
$M^\bullet$에도 같은 사실이 성립한다는 데서 따른다.
\end{proof}

\begin{lemma}
\label{perfect-lemma-identify-pseudo-coherent-noetherian}
 $X$를 Noether 스킴이라 하자. $E$를
$D_\QCoh(\mathcal{O}_X)$의 대상으로 하자. $m \in \mathbf{Z}$에 대해
다음 조건들은 서로 동치이다.
\begin{enumerate}
\item $i \geq m$일 때 $H^i(E)$가 연접이고 $i \gg 0$일 때 영이며,
\item $E$가 $m$-유사연접이다.
\end{enumerate}
특히 $E$가 유사연접일 필요충분조건은 $E$가
$D^-_{\textit{Coh}}(\mathcal{O}_X)$의 대상인 것이다.
\end{lemma}

\begin{proof}
 $X$가 준콤팩트이므로 (1), (2) 모두에서 대상 $E$는 위로 유계이다.
따라서 문제는 $X$ 위에서 국소적이므로 $X$가 아핀이라고 가정할 수 있다.
어떤 Noether 환 $A$에 대해 $X = \Spec(A)$라 하자.
이 경우 Lemma \ref{perfect-lemma-affine-compare-bounded}에 의해 $E$는
$A$-가군 복합체 $M^\bullet$에 대응한다. Lemma
\ref{perfect-lemma-pseudo-coherent-affine}에 의해 $E$가 $m$-유사연접일
필요충분조건은 $M^\bullet$가 $m$-유사연접인 것이다. 한편
$H^i(E)$가 연접일 필요충분조건은 $H^i(M^\bullet)$가 유한
$A$-가군인 것이다(Properties, Lemma
\ref{properties-lemma-finite-type-module}). 따라서 결과는 More on
Algebra, Lemma \ref{more-algebra-lemma-Noetherian-pseudo-coherent}에서
따른다.
\end{proof}

\begin{lemma}
\label{perfect-lemma-tor-dimension-affine}
 $X = \Spec(A)$를 아핀 스킴이라 하자. $M^\bullet$를
$A$-가군의 복합체라 하고 $E$를 $D(\mathcal{O}_X)$의 대응하는
대상이라 하자. 그러면
\begin{enumerate}
\item $E$의 Tor 진폭이 $[a, b]$일 필요충분조건은 $M^\bullet$의
Tor 진폭이 $[a, b]$인 것이다.
\item $E$의 Tor 차원이 유한할 필요충분조건은 $M^\bullet$의
Tor 차원이 유한한 것이다.
\end{enumerate}
\end{lemma}

\begin{proof}
 (2)는 (1)에서 자명하게 따른다. (1)의 증명에서는 더 언급하지 않고
Lemma \ref{perfect-lemma-affine-compare-bounded}의 동치
$D(A) = D_\QCoh(X)$를 사용한다.
$M^\bullet$의 Tor 진폭이 $[a, b]$라고 하자. 그러면 More on Algebra,
Lemma \ref{more-algebra-lemma-tor-amplitude}에 의해 $K^\bullet$는 $D(A)$에서
$K^i = 0$ ($i \not \in [a, b]$)인 평탄 $A$-가군 복합체
$K^\bullet$와 동형이다. 따라서 $E$는 $\widetilde{K^\bullet}$와 동형이다.
각 $\widetilde{K^i}$가 평탄한 $\mathcal{O}_X$-가군이므로
Cohomology, Lemma \ref{cohomology-lemma-tor-amplitude}에 의해
$E$의 Tor 진폭은 $[a, b]$이다.

\medskip\noindent
 $E$의 Tor 진폭이 $[a, b]$라고 하자. 그러면 $E$는 유계이고 따라서
$M^\bullet$는 $K^-(A)$에 속한다. 그러므로 $M^\bullet$를 위로 유계인
$A$-가군 복합체로 바꾸어도 된다. 더 나아가 사영 분해를 택하여
$M^\bullet$가 위로 유계인 자유 $A$-가군 복합체라고 가정할 수 있다.
임의의 $A$-가군 $N$에 대해 다음이 성립한다.
$$
E \otimes_{\mathcal{O}_X}^\mathbf{L} \widetilde{N}
\cong
\widetilde{M^\bullet} \otimes_{\mathcal{O}_X}^\mathbf{L} \widetilde{N}
\cong
\widetilde{M^\bullet \otimes_A N}
$$
이는 $D(\mathcal{O}_X)$에서의 등식이다. 따라서 왼쪽 항의 코호몰로지
층들이 소멸한다는 사실로부터 $M^\bullet$의 Tor 진폭이 $[a, b]$임을
얻는다.
\end{proof}

\begin{lemma}
\label{perfect-lemma-tor-dimension-rel-affine}
 $f : X \to S$를 환 사상 $R \to A$에 대응하는 아핀 스킴 사상이라 하자.
$M^\bullet$를 $A$-가군의 복합체라 하고 $E$를 $D(\mathcal{O}_X)$의
대응하는 대상으로 하자. 그러면
\begin{enumerate}
\item $D(f^{-1}\mathcal{O}_S)$의 대상으로서 $E$의 Tor 진폭이 $[a, b]$일
필요충분조건은 $D(R)$의 대상으로서 $M^\bullet$의 Tor 진폭이 $[a, b]$인
것이다.
\item $D(f^{-1}\mathcal{O}_S)$의 대상으로서 $E$가 국소적으로 유한 Tor
차원을 가질 필요충분조건은 $D(R)$의 대상으로서 $M^\bullet$가 유한 Tor
차원을 갖는 것이다.
\end{enumerate}
\end{lemma}

\begin{proof}
 $\mathfrak p \subset R$ 위에 놓이는 소아이디얼 $\mathfrak q \subset A$를
잡자. 대응하는 점을 $x \in X$, $s = f(x) \in S$라 하자. 그러면
$(f^{-1}\mathcal{O}_S)_x = \mathcal{O}_{S, s} = R_\mathfrak p$이고
$E_x = M^\bullet_\mathfrak q$이다. 이를 염두에 두면 동치는 다음과 같이
보인다.

\medskip\noindent
 $M^\bullet$가 $R$-가군 복합체로서 Tor 진폭 $[a, b]$를 가지면,
$A$의 임의의 소아이디얼에서 $M^\bullet$를 국소화한 것도 같은 진폭을
갖는다. 따라서
Cohomology, Lemma \ref{cohomology-lemma-tor-amplitude-stalk}에 의해 $E$는
$f^{-1}\mathcal{O}_S$-가군 층의 복합체로서 Tor 진폭 $[a, b]$를 갖는다.
역으로 $D(f^{-1}\mathcal{O}_S)$의 대상으로 $E$의 Tor 진폭이 $[a, b]$라고
하자. 마지막으로 인용한 보조정리를 사용하면
$\mathfrak p \subset R$ 위의 모든 소아이디얼 $\mathfrak q \subset A$에 대해
$M^\bullet_\mathfrak q$가 $R_\mathfrak p$-가군 복합체로서 Tor 진폭
$[a, b]$를 가짐을 얻는다. More on Algebra, Lemma
\ref{more-algebra-lemma-tor-amplitude-localization}에 의해
$M^\bullet$는 $R$-가군 복합체로서 Tor 진폭 $[a, b]$를 갖는다.
이로써 (1)의 증명이 끝난다.

\medskip\noindent
 $X$가 준콤팩트이므로 $E$가 $f^{-1}\mathcal{O}_S$-가군 복합체로서
국소적으로 유한 Tor 차원을 가지면 실제로 $E$는 어떤 $a, b$에 대해
$f^{-1}\mathcal{O}_S$-가군 복합체로서 Tor 진폭 $[a, b]$를 갖는다.
따라서 (2)는 (1)에서 따른다.
\end{proof}

\begin{lemma}
\label{perfect-lemma-tor-qc-qs}
 $X$를 준분리 스킴이라 하자. $E$를
$D_\QCoh(\mathcal{O}_X)$의 대상이라 하고 $a \leq b$라 하자.
다음 조건들은 서로 동치이다.
\begin{enumerate}
\item $E$의 Tor 진폭은 $[a, b]$이다.
\item $\QCoh(\mathcal{O}_X)$의 모든 $\mathcal{F}$에 대해
$i \not \in [a, b]$이면
$H^i(E \otimes_{\mathcal{O}_X}^\mathbf{L} \mathcal{F}) = 0$이다.
\end{enumerate}
\end{lemma}

\begin{proof}
 (1)이 (2)를 함의함은 분명하다. (2)를 가정하자. $U \subset X$를
아핀 열린집합이라 하자. $X$가 준분리이므로 사상 $j : U \to X$는
준콤팩트이고 분리되어 있다. 따라서 $j_*$는 준연접 가군을 준연접
가군으로 보낸다(Schemes, Lemma
\ref{schemes-lemma-push-forward-quasi-coherent}). 그러므로 함자
$\QCoh(\mathcal{O}_X) \to \QCoh(\mathcal{O}_U)$는 본질적으로
전사이다. 이로부터 모든 준연접 $\mathcal{O}_U$-가군 $\mathcal{G}$에 대해
$i \not \in [a, b]$이면
$H^i(E|_U \otimes_{\mathcal{O}_U}^\mathbf{L} \mathcal{G})$가
소멸함을 얻는다. $U = \Spec(A)$라 하고 Lemma
\ref{perfect-lemma-affine-compare-bounded}에 의해 $E|_U$에 대응하는
$A$-가군 복합체를 $M^\bullet$라 하자.
위에서 $M^\bullet \otimes_A^\mathbf{L} N$의 코호몰로지 군이
$[a, b]$ 밖에서 소멸함을 보였으므로, 이는 $M^\bullet$의 Tor 진폭이
$[a, b]$라는 뜻이다. Lemma \ref{perfect-lemma-tor-dimension-affine}에 의해
$E|_U$의 Tor 진폭도 $[a, b]$이다. 보조정리가 증명되었다.
\end{proof}

\begin{lemma}
\label{perfect-lemma-perfect-affine}
 $X = \Spec(A)$를 아핀 스킴이라 하자. $M^\bullet$를 $A$-가군의
복합체라 하고 $E$를 $D(\mathcal{O}_X)$의 대응하는 대상으로 하자.
그러면 $E$가 $D(\mathcal{O}_X)$의 완전 대상일 필요충분조건은
$M^\bullet$가 $D(A)$의 대상으로서 완전한 것이다.
\end{lemma}

\begin{proof}
이는 Lemma \ref{perfect-lemma-pseudo-coherent-affine},
\ref{perfect-lemma-tor-dimension-affine}, Cohomology, Lemma
\ref{cohomology-lemma-perfect}, More on Algebra, Lemma
\ref{more-algebra-lemma-perfect}의 논리적 귀결이다.
\end{proof}

\noindent
스킴 위의 유사연접 복합체에 대한 위의 기술로부터 특정 내부 Hom들이
준연접임을 증명할 수 있다.

\begin{lemma}
\label{perfect-lemma-quasi-coherence-internal-hom}
 $X$를 스킴이라 하자.
\begin{enumerate}
\item $L$이 $D^+_\QCoh(\mathcal{O}_X)$에 속하고
$D(\mathcal{O}_X)$의 $K$가 유사연접이면
$R\SheafHom(K, L)$은 $D_\QCoh(\mathcal{O}_X)$에 속하고
국소적으로 아래로 유계이다.
\item $L$이 $D_\QCoh(\mathcal{O}_X)$에 속하고
$D(\mathcal{O}_X)$의 $K$가 완전하면
$R\SheafHom(K, L)$은 $D_\QCoh(\mathcal{O}_X)$에 속한다.
\item $X = \Spec(A)$가 아핀이고 $K, L \in D(A)$이면
$$
R\SheafHom(\widetilde{K}, \widetilde{L}) = \widetilde{R\Hom_A(K, L)}
$$
다음 두 경우에 위 등식이 성립한다.
\begin{enumerate}
\item $K$가 유사연접이고 $L$이 아래로 유계인 경우,
\item $K$가 완전하고 $L$이 임의인 경우.
\end{enumerate}
\item $X = \Spec(A)$이고 $K, L$이 $D(A)$에 속하면,
$R\SheafHom(\widetilde{K}, \widetilde{L})$의 $n$번째 코호몰로지 층은
다음 준층에 대응하는 층이다.
$$
X \supset D(f) \longmapsto \Ext^n_{A_f}(K \otimes_A A_f, L \otimes_A A_f)
$$
여기서 $f \in A$이다.
\end{enumerate}
\end{lemma}

\begin{proof}
 $\mathcal{O}_X$의 유도 범주에서 내부 Hom의 구성은 국소화와 가환한다
(Cohomology, Section \ref{cohomology-section-internal-hom}을 보라).
따라서 (1), (2)를 증명할 때 $X$를 아핀 열린집합으로 바꾸어도 된다.
Lemmas \ref{perfect-lemma-affine-compare-bounded},
\ref{perfect-lemma-pseudo-coherent-affine},
\ref{perfect-lemma-perfect-affine}에 의해 (1), (2)를 증명하려면 (3)을 증명하는
것으로 충분하다.

\medskip\noindent
 (3)은 $K$를 유한 사영 $A$-가군의 위로 유계인(각각 유한인) 복합체로,
$L$을 $A$-가군의 아래로 유계인(각각 임의인) 복합체로 나타내고,
Cohomology, Lemma
\ref{cohomology-lemma-Rhom-complex-of-direct-summands-finite-free}의
내부 Hom 계산을 적용하면 따른다.

\medskip\noindent
 (4)를 증명하기 위해, 임의의 환 달린 공간에서
$R\SheafHom(A, B)$의 $n$번째 코호몰로지 층이 다음 준층에 대응하는 층임을
상기하자.
$$
U \mapsto \Hom_{D(U)}(A|_U, B|_U[n]) =
\Ext^n_{D(\mathcal{O}_U)}(A|_U, B|_U)
$$
Cohomology, Section \ref{cohomology-section-internal-hom}을 보라.
한편 $\widetilde{K}$를 주 열린집합 $D(f)$에 제한한 것은
$K \otimes_A A_f$의 상이며, $L$에 대해서도 마찬가지이다.
따라서 (4)는 Lemma \ref{perfect-lemma-affine-compare-bounded}의 범주 동치에서
따른다.
\end{proof}

\begin{lemma}
\label{perfect-lemma-internal-hom-evaluate-tensor-isomorphism}
 $X$를 스킴이라 하자. $K, L, M$을 $D_\QCoh(\mathcal{O}_X)$의
대상이라 하자. 다음 사상은
$$
K \otimes_{\mathcal{O}_X}^\mathbf{L} R\SheafHom(M, L)
\longrightarrow
R\SheafHom(M, K \otimes_{\mathcal{O}_X}^\mathbf{L} L)
$$
Cohomology, Lemma \ref{cohomology-lemma-internal-hom-diagonal-better}의
사상이며, 다음 경우에 동형사상이다.
\begin{enumerate}
\item $M$이 완전하거나,
\item $K$가 완전하거나,
\item $M$이 유사연접이고 $L \in D^+(\mathcal{O}_X)$이며 $K$의 Tor
차원이 유한한 경우.
\end{enumerate}
\end{lemma}

\begin{proof}
Lemma \ref{perfect-lemma-quasi-coherence-internal-hom}은 (1), (3)을 아핀인 경우로
환원하며, 이 경우는 More on Algebra, Lemma
\ref{more-algebra-lemma-internal-hom-evaluate-tensor-isomorphism}에서
다룬다. (대수적 명제로 옮기기 위해 Lemmas
\ref{perfect-lemma-pseudo-coherent-affine}, \ref{perfect-lemma-perfect-affine},
\ref{perfect-lemma-tor-dimension-affine}도 사용해야 한다.)
$K$가 완전하지만 다른 가정을 하지 않으면 화살표의 어느 쪽이든
$D_\QCoh(\mathcal{O}_X)$에 속한다는 것을 알 수는 없다. 그러나 국소적으로
작업하여 $K$가 유한 자유 가군들의 유한 복합체인 경우로 환원할 수 있고,
그 경우에는 명제가 분명하므로 결과는 여전히 참이다.
\end{proof}





\section{연접 가군의 유도 범주}
\label{perfect-section-derived-coherent}

\noindent
$X$를 국소 Noether 스킴이라 하자. 이 경우 연접 $\mathcal{O}_X$-가군의 범주
$\textit{Coh}(\mathcal{O}_X) \subset \textit{Mod}(\mathcal{O}_X)$는
약한 Serre 부분범주이다. Homology, Section
\ref{homology-section-serre-subcategories} 및 Cohomology of Schemes, Lemma
\ref{coherent-lemma-coherent-abelian-Noetherian}를 보라. 코호몰로지 층들이
연접인 복합체들의 부분범주를
$$
D_{\textit{Coh}}(\mathcal{O}_X) \subset D(\mathcal{O}_X)
$$
로 나타내자. Derived Categories, Section
\ref{derived-section-triangulated-sub}을 보라. 그러면 표준 함자
\begin{equation}
\label{perfect-equation-compare-coherent}
D(\textit{Coh}(\mathcal{O}_X))
\longrightarrow
D_{\textit{Coh}}(\mathcal{O}_X)
\end{equation}
를 얻는다. Derived Categories, Equation (\ref{derived-equation-compare})을 보라.

\begin{lemma}
\label{perfect-lemma-coh-to-qcoh}
$X$를 Noether 스킴이라 하자. 그러면 함자
$$
D^-(\textit{Coh}(\mathcal{O}_X))
\longrightarrow
D^-_{\textit{Coh}(\mathcal{O}_X)}(\QCoh(\mathcal{O}_X))
$$
는 동치이다.
\end{lemma}

\begin{proof}
$\textit{Coh}(\mathcal{O}_X) \subset \QCoh(\mathcal{O}_X)$가 Serre
부분범주임에 유의하자. Homology, Definition
\ref{homology-definition-serre-subcategory}, Lemma
\ref{homology-lemma-characterize-serre-subcategory}, 그리고 Cohomology of
Schemes, Lemmas \ref{coherent-lemma-coherent-abelian-Noetherian} 및
\ref{coherent-lemma-coherent-Noetherian-quasi-coherent-sub-quotient}를 보라.
한편 $\mathcal{G} \to \mathcal{F}$가 준연접 $\mathcal{O}_X$-가군에서
연접 $\mathcal{O}_X$-가군으로의 전사이면, $\mathcal{F}$로 전사하는 연접
부분가군 $\mathcal{G}' \subset \mathcal{G}$가 존재한다. 실제로 Properties,
Lemma \ref{properties-lemma-quasi-coherent-colimit-finite-type}에 의해
$\mathcal{G}$를 그 연접 부분가군들의 여과 합집합으로 쓸 수 있고, 그중
하나가 요구를 만족한다. 따라서 이 보조정리는 Derived Categories, Lemma
\ref{derived-lemma-fully-faithful-embedding}에서 따른다.
\end{proof}

\begin{proposition}
\label{perfect-proposition-DCoh}
$X$를 Noether 스킴이라 하자. 그러면 함자들
$$
D^-(\textit{Coh}(\mathcal{O}_X))
\longrightarrow
D^-_{\textit{Coh}}(\mathcal{O}_X)
\quad\text{그리고}\quad
D^b(\textit{Coh}(\mathcal{O}_X))
\longrightarrow
D^b_{\textit{Coh}}(\mathcal{O}_X)
$$
은 동치이다.
\end{proposition}

\begin{proof}
다음 가환 그림을 생각하자.
$$
\xymatrix{
D^-(\textit{Coh}(\mathcal{O}_X)) \ar[r] \ar[d] &
D^-_{\textit{Coh}}(\mathcal{O}_X) \ar[d] \\
D^-(\QCoh(\mathcal{O}_X)) \ar[r] &
D^-_\QCoh(\mathcal{O}_X)
}
$$
Lemma \ref{perfect-lemma-coh-to-qcoh}에 의해 왼쪽 수직 화살표는 충실충만하다.
Proposition \ref{perfect-proposition-Noetherian}에 의해 아래쪽 화살표는 동치이다.
구성에 의해 오른쪽 수직 화살표는 충실충만하다. 따라서 위쪽 수평 화살표는
충실충만하다. $K$가 $D^-_{\textit{Coh}}(\mathcal{O}_X)$의 대상이면,
Proposition \ref{perfect-proposition-Noetherian}에 의해 이에 대응하는
$D^-(\QCoh(\mathcal{O}_X))$의 대상 $K'$는 연접 코호몰로지 층들을 갖는다.
그러므로 Lemma \ref{perfect-lemma-coh-to-qcoh}에 의해 $K'$는 왼쪽 수직 화살표의
본질적 상에 속하며, 따라서 위쪽 수평 화살표는 본질적으로 전사이다.
이로써 위로 유계인 경우의 증명이 끝난다. 유계인 경우는 위로 유계인 경우에서
즉시 따른다.
\end{proof}

\begin{lemma}
\label{perfect-lemma-direct-image-coherent}
$S$를 Noether 스킴이라 하고, $f : X \to S$를 국소 유한형인 스킴의
사상이라 하자. $E$를 $D^b_{\textit{Coh}}(\mathcal{O}_X)$의 대상이라 하며,
모든 $i$에 대해 $H^i(E)$의 지지가 $S$ 위에서 고유라고 하자. 그러면
$Rf_*E$는 $D^b_{\textit{Coh}}(\mathcal{O}_S)$의 대상이다.
\end{lemma}

\begin{proof}
다음 스펙트럼 열을 생각하자.
$$
R^pf_*H^q(E) \Rightarrow R^{p + q}f_*E
$$
Derived Categories, Lemma \ref{derived-lemma-two-ss-complex-functor}를 보라.
가정과 Cohomology of Schemes, Lemma
\ref{coherent-lemma-support-proper-over-base-pushforward}에 의해 층들
$R^pf_*H^q(E)$는 연접이다. 따라서 $R^{p + q}f_*E$는 연접이며, 즉
$Rf_*E \in D_{\textit{Coh}}(\mathcal{O}_S)$이다. 아래로 유계인 것은
자명하다. 위로 유계인 것은 Cohomology of Schemes, Lemma
\ref{coherent-lemma-quasi-coherence-higher-direct-images} 또는 Lemma
\ref{perfect-lemma-quasi-coherence-direct-image}에서 따른다.
\end{proof}

\begin{lemma}
\label{perfect-lemma-direct-image-coherent-bdd-below}
$S$를 Noether 스킴이라 하고, $f : X \to S$를 국소 유한형인 스킴의
사상이라 하자. $E$를 $D^+_{\textit{Coh}}(\mathcal{O}_X)$의 대상이라 하며,
모든 $i$에 대해 $H^i(E)$의 지지가 $S$ 위에서 고유라고 하자. 그러면
$Rf_*E$는 $D^+_{\textit{Coh}}(\mathcal{O}_S)$의 대상이다.
\end{lemma}

\begin{proof}
증명은 Lemma \ref{perfect-lemma-direct-image-coherent}의 증명과 같다. 또한 완전 함자
$Rf_*$가 완전 삼각형
$\tau_{\leq a}E \to E \to \tau_{\geq a + 1}E \to \tau_{\leq a}E[1]$
에 어떻게 작용하는지를 살펴봄으로써 Lemma
\ref{perfect-lemma-direct-image-coherent}에서 이 결론을 얻을 수도 있다.
\end{proof}

\begin{lemma}
\label{perfect-lemma-coherent-internal-hom}
$X$를 국소 Noether 스킴이라 하자. $L$이
$D^+_{\textit{Coh}}(\mathcal{O}_X)$에 속하고 $K$가
$D^-_{\textit{Coh}}(\mathcal{O}_X)$에 속하면, $R\SheafHom(K, L)$은
$D^+_{\textit{Coh}}(\mathcal{O}_X)$에 속한다.
\end{lemma}

\begin{proof}
$X$가 Noether 환 $A$의 스펙트럼인 경우를 증명하면 충분하다. Lemma
\ref{perfect-lemma-identify-pseudo-coherent-noetherian}에 의해 $K$는 유사연접이다.
그러면 Lemma \ref{perfect-lemma-quasi-coherence-internal-hom}을 사용하여 문제를
다음 대수적 문제로 옮길 수 있다. $L \in D^+_{\textit{Coh}}(A)$이고 $K$가
$D^-_{\textit{Coh}}(A)$에 속하면, $R\Hom_A(K, L)$은
$D^+_{\textit{Coh}}(A)$에 속한다. $L$은 아래로 유계이고 $K$는 위로
유계이므로 다음과 같은 수렴하는 스펙트럼 열이 있다.
$$
\Ext^p_A(K, H^q(L)) \Rightarrow \text{Ext}^{p + q}_A(K, L)
$$
또한 다음과 같은 수렴하는 스펙트럼 열들이 있다.
$$
\Ext^i_A(H^{-j}(K), H^q(L)) \Rightarrow \text{Ext}^{i + j}_A(K, H^q(L))
$$
Injectives, Remarks \ref{injectives-remark-spectral-sequences-ext} 및
\ref{injectives-remark-spectral-sequences-ext-variant}를 보라. Algebra, Lemma
\ref{algebra-lemma-ext-noetherian}에 의해 유한 $A$-가군 $M$, $N$에 대해
가군 $\Ext^p_A(M, N)$이 유한이므로 증명이 끝난다.
\end{proof}

\begin{lemma}
\label{perfect-lemma-perfect-on-noetherian}
$X$를 Noether 스킴이라 하자. $D(\mathcal{O}_X)$의 대상 $E$가 완전이라고
하자. 그러면 다음이 성립한다.
\begin{enumerate}
\item $E$는 $D^b_{\textit{Coh}}(\mathcal{O}_X)$에 속한다.
\item $L$이 $D_{\textit{Coh}}(\mathcal{O}_X)$에 속하면,
$E \otimes_{\mathcal{O}_X}^\mathbf{L} L$와
$R\SheafHom_{\mathcal{O}_X}(E, L)$은 $D_{\textit{Coh}}(\mathcal{O}_X)$에
속한다.
\item $L$이 $D^b_{\textit{Coh}}(\mathcal{O}_X)$에 속하면,
$E \otimes_{\mathcal{O}_X}^\mathbf{L} L$와
$R\SheafHom_{\mathcal{O}_X}(E, L)$은 $D^b_{\textit{Coh}}(\mathcal{O}_X)$에
속한다.
\item $L$이 $D^+_{\textit{Coh}}(\mathcal{O}_X)$에 속하면,
$E \otimes_{\mathcal{O}_X}^\mathbf{L} L$와
$R\SheafHom_{\mathcal{O}_X}(E, L)$은 $D^+_{\textit{Coh}}(\mathcal{O}_X)$에
속한다.
\item $L$이 $D^-_{\textit{Coh}}(\mathcal{O}_X)$에 속하면,
$E \otimes_{\mathcal{O}_X}^\mathbf{L} L$와
$R\SheafHom_{\mathcal{O}_X}(E, L)$은 $D^-_{\textit{Coh}}(\mathcal{O}_X)$에
속한다.
\end{enumerate}
\end{lemma}

\begin{proof}
$X$가 준콤팩트이므로 각 명제는 $X$의 임의의 열린 덮개의 각 원소 위에서
확인할 수 있다. 따라서 $E$가 유한 자유 가군들의 유계 복합체
$\mathcal{E}^\bullet$로 나타난다고 가정할 수 있다. Cohomology, Lemma
\ref{cohomology-lemma-perfect-on-locally-ringed}를 보라. 이 경우
$R\SheafHom_{\mathcal{O}_X}(E, L)$과
$E \otimes_{\mathcal{O}_X}^\mathbf{L} L$을 모두 $\mathcal{E}^\bullet$을
사용하여 복합체의 수준에서 계산할 수 있으므로 각 명제는 자명하다.
Cohomology, Lemmas \ref{cohomology-lemma-Rhom-strictly-perfect} 및
\ref{cohomology-lemma-bounded-flat-K-flat}를 보라. 몇 가지 세부사항은 생략한다.
\end{proof}

\begin{lemma}
\label{perfect-lemma-ext-finite}
$A$를 Noether 환이라 하고 $X$를 $A$ 위의 고유 스킴이라 하자.
$D^+_{\textit{Coh}}(\mathcal{O}_X)$의 $L$과
$D^-_{\textit{Coh}}(\mathcal{O}_X)$의 $K$에 대해 $A$-가군
$\Ext_{\mathcal{O}_X}^n(K, L)$은 유한이다.
\end{lemma}

\begin{proof}
다음을 상기하자.
$$
\Ext_{\mathcal{O}_X}^n(K, L) =
H^n(X, R\SheafHom_{\mathcal{O}_X}(K, L)) =
H^n(\Spec(A), Rf_*R\SheafHom_{\mathcal{O}_X}(K, L))
$$
Cohomology, Lemma \ref{cohomology-lemma-section-RHom-over-U} 및 Cohomology,
Section \ref{cohomology-section-Leray}를 보라. 따라서 결과는 Lemmas
\ref{perfect-lemma-coherent-internal-hom} 및
\ref{perfect-lemma-direct-image-coherent-bdd-below}에서 따른다.
\end{proof}

\begin{lemma}
\label{perfect-lemma-perfect-on-regular}
$X$를 정칙 스킴이라 하자. 그러면
$D^b_{\textit{Coh}}(\mathcal{O}_X)$의 모든 대상은 완전이다. $X$가
준콤팩트, 즉 Noether이고 정칙이면, 반대로 $D(\mathcal{O}_X)$의 모든 완전
대상은 $D^b_{\textit{Coh}}(\mathcal{O}_X)$에 속한다.
\end{lemma}

\begin{proof}
$K$를 $D^b_{\textit{Coh}}(\mathcal{O}_X)$의 대상이라 하자. $K$가 완전임을
확인하기 위해 $X$ 위에서 아핀 국소적으로 작업할 수 있다(Cohomology,
Section \ref{cohomology-section-perfect}을 보라). 그러면 Lemma
\ref{perfect-lemma-perfect-affine}과 More on Algebra, Lemma
\ref{more-algebra-lemma-regular-perfect}에 의해 $K$는 완전이다. 역은 Lemma
\ref{perfect-lemma-perfect-on-noetherian}이다.
\end{proof}




\section{복합체의 유한성 성질의 강하}
\label{perfect-section-descent-finiteness}

\noindent
이 절은 스킴의 유도 범주의 대상에 대한 Descent, Section
\ref{descent-section-descent-finiteness}의 유사물이다. 이러한 결과 가운데
가장 쉬운 것은 아마 다음일 것이다.

\begin{lemma}
\label{perfect-lemma-tor-amplitude-descends}
$f : X \to Y$를 스킴(또는 더 일반적으로 국소환 달린 공간)의 전사 평탄
사상이라 하자. $E \in D(\mathcal{O}_Y)$라 하고 $a, b \in \mathbf{Z}$라 하자.
그러면 $E$가 $[a, b]$에서 Tor 진폭을 가질 필요충분조건은 $Lf^*E$가
$[a, b]$에서 Tor 진폭을 갖는 것이다.
\end{lemma}

\begin{proof}
당김은 항상 Tor 진폭을 보존한다. Cohomology, Lemma
\ref{cohomology-lemma-tor-amplitude-pullback}을 보라. $[a, b]$에서의 Tor
진폭은 줄기에서 확인할 수 있다. Cohomology, Lemma
\ref{cohomology-lemma-tor-amplitude-stalk}을 보라. 평탄한 국소환 준동형은
Algebra, Lemma \ref{algebra-lemma-local-flat-ff}에 의해 충실히 평탄하다.
따라서 결과는 More on Algebra, Lemma
\ref{more-algebra-lemma-flat-descent-tor-amplitude}에서 따른다.
\end{proof}

\begin{lemma}
\label{perfect-lemma-pseudo-coherent-descends-fpqc}
$\{f_i : X_i \to X\}$를 스킴의 fpqc 덮개라 하자.
$E \in D_\QCoh(\mathcal{O}_X)$라 하고 $m \in \mathbf{Z}$라 하자. 그러면
$E$가 $m$-유사연접일 필요충분조건은 각 $Lf_i^*E$가 $m$-유사연접인 것이다.
\end{lemma}

\begin{proof}
당김은 항상 $m$-유사연접성을 보존한다. Cohomology, Lemma
\ref{cohomology-lemma-pseudo-coherent-pullback}을 보라. 반대로 모든 $i$에
대해 $Lf_i^*E$가 $m$-유사연접이라고 가정하자. $U \subset X$를 아핀
열린집합이라 하자. $E|_U$가 $m$-유사연접임을 보이면 충분하다.
$\{f_i : X_i \to X\}$가 fpqc 덮개이므로 $f_{a(j)}(V_j) \subset U$이고
$U = \bigcup f_{a(j)}(V_j)$인 유한 개의 아핀 열린집합
$V_j \subset X_{a(j)}$를 찾을 수 있다. $V = \coprod V_i$로 두자. 따라서
$X$를 $U$로, $\{f_i : X_i \to X\}$를 $\{V \to U\}$로 바꾸어 $X$가
아핀이고 덮개가 아핀 스킴의 하나의 전사 평탄 사상
$\{f : Y \to X\}$로 주어진다고 가정할 수 있다. 이 경우 결과는 Lemmas
\ref{perfect-lemma-affine-compare-bounded}와 \ref{perfect-lemma-pseudo-coherent-affine}를
통하여 More on Algebra, Lemma
\ref{more-algebra-lemma-flat-descent-pseudo-coherent}에서 따른다.
\end{proof}

\begin{lemma}
\label{perfect-lemma-pseudo-coherent-descends-fppf}
$\{f_i : X_i \to X\}$를 스킴의 fppf 덮개라 하자.
$E \in D(\mathcal{O}_X)$라 하고 $m \in \mathbf{Z}$라 하자. 그러면 $E$가
$m$-유사연접일 필요충분조건은 각 $Lf_i^*E$가 $m$-유사연접인 것이다.
\end{lemma}

\begin{proof}
당김은 항상 $m$-유사연접성을 보존한다. Cohomology, Lemma
\ref{cohomology-lemma-pseudo-coherent-pullback}을 보라. 반대로 모든 $i$에
대해 $Lf_i^*E$가 $m$-유사연접이라고 가정하자. $U \subset X$를 아핀
열린집합이라 하자. $E|_U$가 $m$-유사연접임을 보이면 충분하다.
$\{f_i : X_i \to X\}$가 fppf 덮개이므로 $f_{a(j)}(V_j) \subset U$이고
$U = \bigcup f_{a(j)}(V_j)$인 유한 개의 아핀 열린집합
$V_j \subset X_{a(j)}$를 찾을 수 있다. $V = \coprod V_i$로 두자. 따라서
$X$를 $U$로, $\{f_i : X_i \to X\}$를 $\{V \to U\}$로 바꾸어 $X$가
아핀이고 덮개가 하나의 유한 표시 전사 평탄 사상 $\{f : Y \to X\}$로
주어진다고 가정할 수 있다.

\medskip\noindent
$f$가 평탄하므로 유도 함자 $Lf^*$는 단순히 $f^*$로 주어지고 $f^*$는
완전하다. 따라서 $H^i(Lf^*E) = f^*H^i(E)$이다. $Lf^*E$가
$m$-유사연접이므로 $Lf^*E \in D^-(\mathcal{O}_Y)$이다. $f$가 전사이고
평탄하므로 $E \in D^-(\mathcal{O}_X)$이다. $H^i(E)$가 영이 아닌 가장 큰
정수 $i \in \mathbf{Z}$를 택하자. $i < m$이면 끝났다. 그렇지 않으면
Cohomology, Lemma \ref{cohomology-lemma-finite-cohomology}에 의해
$f^*H^i(E)$는 유한형 $\mathcal{O}_Y$-가군이다. 그러면 Descent, Lemma
\ref{descent-lemma-finite-type-descends-fppf}에 의해 $\mathcal{O}_X$-가군
$H^i(E)$는 유한형이다. 따라서 $X$를 유한 아핀 열린 덮개의 원소들로
바꾼 뒤, 다음 사상이 존재한다고 가정할 수 있다.
$$
\alpha : \mathcal{O}_X^{\oplus n}[-i] \longrightarrow E
$$
여기서 $H^i(\alpha)$는 전사이다. $C$를 $D(\mathcal{O}_X)$에서
$\alpha$의 사상뿔이라 하자. $Y$로 당기고 Cohomology, Lemma
\ref{cohomology-lemma-cone-pseudo-coherent}를 사용하면 $Lf^*C$가
$m$-유사연접임을 얻는다. 또한 $j \geq i$에 대해 $H^j(C) = 0$이다.
따라서 $i$에 대한 귀납법으로 $C$가 $m$-유사연접임을 알 수 있다. 다시
Cohomology, Lemma \ref{cohomology-lemma-cone-pseudo-coherent}를 사용하면
결론을 얻는다.
\end{proof}

\begin{lemma}
\label{perfect-lemma-perfect-descends-fpqc}
$\{f_i : X_i \to X\}$를 스킴의 fpqc 덮개라 하고
$E \in D(\mathcal{O}_X)$라 하자. 그러면 $E$가 완전일 필요충분조건은 각
$Lf_i^*E$가 완전인 것이다.
\end{lemma}

\begin{proof}
당김은 항상 완전 복합체를 보존한다. Cohomology, Lemma
\ref{cohomology-lemma-perfect-pullback}을 보라. 반대로 모든 $i$에 대해
$Lf_i^*E$가 완전이라고 가정하자. 그러면 각 $Lf_i^*E$의 코호몰로지 층들은
준연접이다. Lemma \ref{perfect-lemma-pseudo-coherent}와 Cohomology, Lemma
\ref{cohomology-lemma-perfect}를 보라. 사상들 $f_i$가 평탄하므로
$H^p(Lf_i^*E) = f_i^*H^p(E)$이다. 따라서 Descent, Proposition
\ref{descent-proposition-fpqc-descent-quasi-coherent}에 의해 $E$의
코호몰로지 층들은 준연접이다. 이제 이 보조정리는 Cohomology, Lemma
\ref{cohomology-lemma-perfect}, 그리고 Lemmas
\ref{perfect-lemma-tor-amplitude-descends} 및
\ref{perfect-lemma-pseudo-coherent-descends-fpqc}에서 형식적으로 따른다.
\end{proof}

\begin{lemma}
\label{perfect-lemma-closed-push-pseudo-coherent}
$i : Z \to X$를 환 달린 공간의 사상이라 하자. 바탕 위상공간에서 $i$가
닫힌 매장이고, $i_*\mathcal{O}_Z$가 $\mathcal{O}_X$-가군으로서
유사연접이라고 하자. $E \in D(\mathcal{O}_Z)$라 하자. 그러면 $E$가
$m$-유사연접일 필요충분조건은 $Ri_*E$가 $m$-유사연접인 것이다.
\end{lemma}

\begin{proof}
이 증명 전체에서 $i_*$가 완전 함자이고 따라서 $Ri_* = i_*$임을 사용한다.
Modules, Lemma \ref{modules-lemma-i-star-exact}를 보라.

\medskip\noindent
$E$가 $m$-유사연접이라고 가정하자. $x \in X$라 하자. 그 위에서 $i_*E$가
$m$-유사연접인 $x$의 이웃을 찾겠다. $x \not \in Z$이면 이는 자명하다.
따라서 $x \in Z$라고 가정할 수 있다. 열린집합 $x \in U \subset X$에 대한
$U \cap Z$들이 $Z$에서 $x$의 이웃들의 기본계를 이룬다는 사실을 사용한다.
$X$를 축소한 뒤 $E$가 위로 유계라고 가정할 수 있다. $H^p(E)$가 영이 아닌
가장 큰 정수 $p$에 대한 귀납법을 사용한다. $p < m$이면 보일 것이 없다.
$p \geq m$이면 Cohomology, Lemma
\ref{cohomology-lemma-finite-cohomology}에 의해 $H^p(E)$는 유한형
$\mathcal{O}_Z$-가군이다. 따라서 $X$를 축소한 뒤 전사
$\mathcal{O}_Z^{\oplus n} \to H^p(E)$를 유도하는 사상
$\mathcal{O}_Z^{\oplus n}[-p] \to E$를 택할 수 있다. 다음 완전 삼각형을
택하자.
$$
\mathcal{O}_Z^{\oplus n}[-p] \to E \to C \to \mathcal{O}_Z^{\oplus n}[-p + 1]
$$
Cohomology, Lemma \ref{cohomology-lemma-cone-pseudo-coherent}에 의해
$j \geq p$에 대해 $H^j(C) = 0$이고 $C$는 $m$-유사연접이다. 귀납법으로
$i_*C$가 $X$ 위에서 $m$-유사연접임을 알 수 있다.
$i_*\mathcal{O}_Z$도 $X$ 위에서 $m$-유사연접이므로, 완전 삼각형
$$
i_*\mathcal{O}_Z^{\oplus n}[-p] \to i_*E \to i_*C \to
i_*\mathcal{O}_Z^{\oplus n}[-p + 1]
$$
과 Cohomology, Lemma \ref{cohomology-lemma-cone-pseudo-coherent}로부터
$i_*E$가 $m$-유사연접임을 얻는다.

\medskip\noindent
$i_*E$가 $m$-유사연접이라고 가정하자. $z \in Z$라 하자. 그 위에서 $E$가
$m$-유사연접인 $z$의 이웃을 찾겠다. 열린집합 $z \in U \subset X$에 대한
$U \cap Z$들이 $Z$에서 $z$의 이웃들의 기본계를 이룬다는 사실을 사용한다.
$X$를 축소한 뒤 $i_*E$, 따라서 $E$가 위로 유계라고 가정할 수 있다.
$H^p(E)$가 영이 아닌 가장 큰 정수 $p$에 대한 귀납법을 사용한다. $p < m$이면
보일 것이 없다. $p \geq m$이면 Cohomology, Lemma
\ref{cohomology-lemma-finite-cohomology}에 의해
$H^p(i_*E) = i_*H^p(E)$는 유한형 $\mathcal{O}_X$-가군이다. $E$를 나타내는
$\mathcal{O}_Z$-가군의 복합체 $\mathcal{E}^\bullet$을 택하자. $X$를 축소한 뒤
전사 $\mathcal{O}_X^{\oplus n} \to i_*H^p(\mathcal{E}^\bullet)$을 유도하는
사상 $\alpha : \mathcal{O}_X^{\oplus n}[-p] \to i_*\mathcal{E}^\bullet$을
택할 수 있다. 수반성에 의해 전사
$\mathcal{O}_Z^{\oplus n} \to H^p(\mathcal{E}^\bullet)$을 유도하는 사상
$\alpha : \mathcal{O}_Z^{\oplus n}[-p] \to \mathcal{E}^\bullet$을 얻는다.
다음 완전 삼각형을 택하자.
$$
\mathcal{O}_Z^{\oplus n}[-p] \to E \to C \to \mathcal{O}_Z^{\oplus n}[-p + 1]
$$
$j \geq p$에 대해 $H^j(C) = 0$이다. 완전 삼각형
$$
i_*\mathcal{O}_Z^{\oplus n}[-p] \to i_*E \to i_*C \to
i_*\mathcal{O}_Z^{\oplus n}[-p + 1]
$$
과 $i_*\mathcal{O}_Z$가 유사연접이라는 사실 및 Cohomology, Lemma
\ref{cohomology-lemma-cone-pseudo-coherent}로부터 $i_*C$가
$m$-유사연접임을 얻는다. 귀납법으로 $C$가 $m$-유사연접임을 얻는다. 다시
Cohomology, Lemma \ref{cohomology-lemma-cone-pseudo-coherent}에 의해 $E$가
$m$-유사연접임을 얻는다.
\end{proof}

\begin{lemma}
\label{perfect-lemma-finite-push-pseudo-coherent}
$f : X \to Y$를 스킴의 유한 사상이라 하고, $f_*\mathcal{O}_X$가
$\mathcal{O}_Y$-가군으로서 유사연접이라고 하자\footnote{이는 $f$가
유사연접이라는 뜻이다. More on Morphisms, Lemma
\ref{more-morphisms-lemma-finite-pseudo-coherent}를 보라.}.
$E \in D_\QCoh(\mathcal{O}_X)$라 하자. 그러면 $E$가 $m$-유사연접일
필요충분조건은 $Rf_*E$가 $m$-유사연접인 것이다.
\end{lemma}

\begin{proof}
이는 More on Algebra, Lemma
\ref{more-algebra-lemma-finite-push-pseudo-coherent}를 스킴의 언어로 옮긴
것이다. 이 번역에는 Lemmas \ref{perfect-lemma-affine-compare-bounded} 및
\ref{perfect-lemma-pseudo-coherent-affine}를 사용한다.
\end{proof}


\section{복합체의 올림}
\label{perfect-section-lift}

\noindent
$U \subset X$를 환 달린 공간의 열린 부분공간이라 하고 포함 사상을
$j : U \to X$로 나타내자. $Rj_*$가 제한의 오른쪽 역이므로 함자
$D(\mathcal{O}_X) \to D(\mathcal{O}_U)$는 본질적으로 전사이다. 이 절에서는
이를 준연접 코호몰로지 층들을 갖는 복합체 등으로 확장한다.

\begin{lemma}
\label{perfect-lemma-lift-quasi-coherent}
$X$를 스킴이라 하고 $j : U \to X$를 준콤팩트 열린 매장이라 하자. 함자들
$$
D_\QCoh(\mathcal{O}_X) \to D_\QCoh(\mathcal{O}_U)
\quad\text{그리고}\quad
D^+_\QCoh(\mathcal{O}_X) \to D^+_\QCoh(\mathcal{O}_U)
$$
은 본질적으로 전사이다. $X$가 준콤팩트이면 함자들
$$
D^-_\QCoh(\mathcal{O}_X) \to D^-_\QCoh(\mathcal{O}_U)
\quad\text{그리고}\quad
D^b_\QCoh(\mathcal{O}_X) \to D^b_\QCoh(\mathcal{O}_U)
$$
도 본질적으로 전사이다.
\end{lemma}

\begin{proof}
Lemma \ref{perfect-lemma-quasi-coherence-direct-image}에 의해 $Rj_*$가
$D_\QCoh(\mathcal{O}_U)$를 $D_\QCoh(\mathcal{O}_X)$로 보내므로 보조정리
앞의 논증이 첫 번째 경우에 적용된다. $Rj_*$가
$D^+_\QCoh(\mathcal{O}_U)$를 $D^+_\QCoh(\mathcal{O}_X)$로 보내는 것은
자명하며, 이는 아래로 유계인 복합체들에 대한 명제를 준다. 마지막으로 $X$가
준콤팩트이면 Lemma \ref{perfect-lemma-quasi-coherence-direct-image}에 의해 $Rj_*$는
$D^-_\QCoh(\mathcal{O}_U)$를 $D^-_\QCoh(\mathcal{O}_X)$로 보낸다. 이 두
사실을 결합하면 마지막 명제를 얻는다.
\end{proof}

\begin{lemma}
\label{perfect-lemma-lift-coherent}
$X$를 Noether 스킴이라 하고 $j : U \to X$를 열린 매장이라 하자. 함자
$D^b_{\textit{Coh}}(\mathcal{O}_X) \to D^b_{\textit{Coh}}(\mathcal{O}_U)$
는 본질적으로 전사이다.
\end{lemma}

\begin{proof}
$K$를 $D^b_{\textit{Coh}}(\mathcal{O}_U)$의 대상이라 하자. Proposition
\ref{perfect-proposition-DCoh}에 의해 $K$를 연접 $\mathcal{O}_U$-가군들의 유계
복합체 $\mathcal{F}^\bullet$로 나타낼 수 있다. 어떤 $a \leq b$에 대해
$i \not \in [a, b]$이면 $\mathcal{F}^i = 0$이라고 하자. $j$가 준콤팩트이고
분리되어 있으므로 유계 복합체 $j_*\mathcal{F}^\bullet$의 항들은 $X$ 위의
준연접 가군이다. Schemes, Lemma
\ref{schemes-lemma-push-forward-quasi-coherent}를 보라. 다음과 같이 연접
부분가군 $\mathcal{G}^i \subset j_*\mathcal{F}^i$를 귀납적으로 택한다.
$i = a$일 때에는 $U$로의 제한이 $\mathcal{F}^a$인 임의의 연접 부분가군
$\mathcal{G}^a \subset j_*\mathcal{F}^a$를 택한다. 이는 Properties, Lemma
\ref{properties-lemma-extend}에 의해 가능하다. $i > a$일 때에는 먼저
$U$로의 제한이 $\mathcal{F}^i$인 임의의 연접 부분가군
$\mathcal{H}^i \subset j_*\mathcal{F}^i$를 택한 다음
$\mathcal{G}^i = \Im(\mathcal{H}^i \oplus \mathcal{G}^{i - 1}
\to j_*\mathcal{F}^i)$로 둔다. 그러면
$\mathcal{G}^\bullet \subset j_*\mathcal{F}^\bullet$은 연접
$\mathcal{O}_X$-가군들의 유계 복합체이고 그 $U$로의 제한은 원하는 대로
$\mathcal{F}^\bullet$이다.
\end{proof}

\begin{lemma}
\label{perfect-lemma-lift-pseudo-coherent}
$X$를 아핀 스킴이라 하고 $U \subset X$를 준콤팩트 열린 부분스킴이라 하자.
$D(\mathcal{O}_U)$의 임의의 유사연접 대상 $E$에 대해, $U$로의 제한이
$E$와 동형인 유한 자유 $\mathcal{O}_X$-가군들의 위로 유계인 복합체가
존재한다.
\end{lemma}

\begin{proof}
Lemma \ref{perfect-lemma-pseudo-coherent}에 의해 $E$는
$D_\QCoh(\mathcal{O}_U)$의 대상이다. Lemma
\ref{perfect-lemma-lift-quasi-coherent}에 의해 어떤 $D_\QCoh(\mathcal{O}_X)$의 대상
$E'$에 대해 $E = E'|U$라고 가정할 수 있다. $X = \Spec(A)$로 쓰자. Lemma
\ref{perfect-lemma-affine-compare-bounded}에 의해, 연관된 $\mathcal{O}_X$-가군의
복합체가 $E'$의 대표가 되는 $A$-가군 복합체 $M^\bullet$을 찾을 수 있다.

\medskip\noindent
$U = D(f_1) \cup \ldots \cup D(f_r)$인 $f_1, \ldots, f_r \in A$를 택하자.
Lemma \ref{perfect-lemma-pseudo-coherent-affine}에 의해 복합체 $M^\bullet_{f_j}$는
$A_{f_j}$-가군의 유사연접 복합체이다. $n$을 정수라 하자. 복합체의 사상
$\alpha : F^\bullet \to M^\bullet$이 주어졌다고 하자. 여기서
$F^\bullet$은 위로 유계이고, $i < n$이면 $F^i = 0$이며, 각 $F^i$는
유한 자유 $R$-가군이고, 다음 사상
$$
H^i(\alpha_{f_j}) : H^i(F^\bullet_{f_j}) \to H^i(M^\bullet_{f_j})
$$
은 $i > n$이면 동형이고 $i = n$이면 전사라고 하자. 그림으로 쓰면 다음과 같다.
$$
\xymatrix{
& F^n \ar[r] \ar[d]^\alpha & F^{n + 1} \ar[d]^\alpha \ar[r] & \ldots \\
M^{n-1} \ar[r] & M^n \ar[r] & M^{n + 1} \ar[r] & \ldots
}
$$
각 $M^\bullet_{f_j}$의 코호몰로지는 큰 차수에서 사라지므로 $n \gg 0$에
대해 이러한 사상을 찾을 수 있다. $n$에 대한 귀납법으로 이를 모든 $i$에
대해 $H^i(\alpha_{f_j})$가 동형인 복합체 사상
$F^\bullet \to M^\bullet$로 확장하겠다. 그러면
$\widetilde{F^\bullet}$을 택하여 보조정리를 얻는다.

\medskip\noindent
귀납 단계에서는 $F^{n - 1}$을 추가하여 위 그림을 확장한다. $C^\bullet$을
$\alpha$의 사상뿔이라 하자(Derived Categories, Definition
\ref{derived-definition-cone}). 코호몰로지의 긴 완전열에 의해 $i \geq n$이면
$H^i(C^\bullet_{f_j}) = 0$이다. More on Algebra, Lemma
\ref{more-algebra-lemma-cone-pseudo-coherent}에 의해 $C^\bullet_{f_j}$는
$(n - 1)$-유사연접이다. More on Algebra, Lemma
\ref{more-algebra-lemma-finite-cohomology}에 의해
$H^{n - 1}(C^\bullet_{f_j})$는 유한 $A_{f_j}$-가군이다. 합성
$F^{n - 1} \to C^{n - 1} \to C^n$이 영이고 $F^{n - 1}_{f_j}$가
$H^{n - 1}(C^\bullet_{f_j})$로 전사하도록 유한 자유 $A$-가군
$F^{n - 1}$과 $A$-가군 $\beta : F^{n - 1} \to C^{n - 1}$을 택하자.
(몇 가지 세부사항은 생략한다. 힌트: 분모를 없애라.)
$C^{n - 1} = M^{n - 1} \oplus F^n$이므로
$\beta = (\alpha^{n - 1}, -d^{n - 1})$로 쓸 수 있다. 합성
$F^{n - 1} \to C^{n - 1} \to C^n$이 영이라는 사실은 이 사상들이 다음
복합체의 사상을 이룬다는 것을 뜻한다.
$$
\xymatrix{
& F^{n - 1} \ar[d]^{\alpha^{n - 1}} \ar[r]_{d^{n - 1}} &
F^n \ar[r] \ar[d]^\alpha &
F^{n + 1} \ar[d]^\alpha \ar[r] & \ldots \\
\ldots \ar[r] &
M^{n - 1} \ar[r] & M^n \ar[r] & M^{n + 1} \ar[r] & \ldots
}
$$
더 나아가 이 사상들은 다음 완전 삼각형들의 사상을 정의한다.
$$
\xymatrix{
(F^n \to \ldots) \ar[r] \ar[d] &
(F^{n-1} \to \ldots) \ar[r] \ar[d] &
F^{n-1} \ar[r] \ar[d]_\beta &
(F^n \to \ldots)[1] \ar[d] \\
(F^n \to \ldots) \ar[r] &
M^\bullet \ar[r] &
C^\bullet \ar[r] &
(F^n \to \ldots)[1]
}
$$
따라서 $\beta$의 선택에 의해 복합체의 사상
$(F^{-1} \to \ldots) \to M^\bullet$은 $f_j$에서 국소화한 코호몰로지에
차수 $\geq n$에서 동형을, 차수 $n - 1$에서 전사를 유도한다. 이로써
보조정리의 증명이 끝난다.
\end{proof}

\noindent
다음 두 보조정리는 아마 다른 곳에 두어야 할 것이다.

\begin{lemma}
\label{perfect-lemma-vanishing-ext}
$X$를 준콤팩트 준분리 스킴이라 하고 $E \in D^b_\QCoh(\mathcal{O}_X)$라
하자. 모든 유한 국소 자유 $\mathcal{O}_X$-가군 $\mathcal{E}$와 모든
$n \geq n_0$에 대해 $\Ext^n_{D(\mathcal{O}_X)}(\mathcal{E}, E) = 0$이
되게 하는 정수 $n_0 > 0$이 존재한다.
\end{lemma}

\begin{proof}
$\Ext^n_{D(\mathcal{O}_X)}(\mathcal{E}, E) =
\Hom_{D(\mathcal{O}_X)}(\mathcal{E}, E[n])$임을 상기하자. 유도 범주의
사상에 대한 Mayer--Vietoris를 사용할 수 있다. Cohomology, Lemma
\ref{cohomology-lemma-mayer-vietoris-hom}을 보라. 따라서 $X = U \cup V$이고
어떤 경계 $n_0$에 대해 $E|_U$, $E|_V$, $E|_{U \cap V}$에 보조정리의
결론이 성립하면, $E$에 대해서는 경계 $n_0 + 1$로 결론이 성립한다. 그러므로
$X$가 아핀인 경우를 증명하면 충분하다. Cohomology of Schemes, Lemma
\ref{coherent-lemma-induction-principle}을 보라.

\medskip\noindent
$X = \Spec(A)$가 아핀이라고 가정하자. 연관된 준연접 가군의 복합체가
$E$를 나타내는 $A$-가군의 복합체 $M^\bullet$을 택하자. Lemma
\ref{perfect-lemma-affine-compare-bounded}를 보라. 어떤 $A$-가군 $P$에 대해
$\mathcal{E} = \widetilde{P}$로 쓰자. $\mathcal{E}$가 유한 국소 자유이므로
$P$는 유한 사영 $A$-가군이다. 다음 등식들이 성립한다.
\begin{align*}
\Hom_{D(\mathcal{O}_X)}(\mathcal{E}, E[n])
& = 
\Hom_{D(A)}(P, M^\bullet[n]) \\
& =
\Hom_{K(A)}(P, M^\bullet[n]) \\
& =
\Hom_A(P, H^n(M^\bullet))
\end{align*}
첫 번째 등식은 Lemma \ref{perfect-lemma-affine-compare-bounded}에 의해, 두 번째
등식은 Derived Categories, Lemma
\ref{derived-lemma-morphisms-from-projective-complex}에 의해, 마지막 등식은
$\Hom_A(P, -)$가 완전 함자라는 사실에 의해 성립한다. $E$와 따라서
$M^\bullet$이 유계이므로 충분히 큰 모든 $n$에 대해 영을 얻는다.
\end{proof}

\noindent
다음 보조정리는 강화할 수 있다(코호몰로지 층이 어떤 고정된 범위에서만
영이 아닌 모든 $L$에 대해 이 소멸은 균일하다).

\begin{lemma}
\label{perfect-lemma-ext-from-perfect-into-bounded-QCoh}
$X$를 준콤팩트 준분리 스킴이라 하고 $K$를 $D(\mathcal{O}_X)$의 완전
대상이라 하자. 그러면 다음이 성립한다.
\begin{enumerate}
\item $i \in [a, b]$에 대해 $H^i(L) = 0$인 임의의
$L \in D_\QCoh(\mathcal{O}_X)$에 대해
$\Hom_{D(\mathcal{O}_X)}(K, L) = 0$이 되게 하는 정수 $a \leq b$가
존재한다.
\item $L$이 유계이면 유한 개를 제외한 모든 $n$에 대해
$\Ext^n_{D(\mathcal{O}_X)}(K, L)$은 영이다.
\end{enumerate}
\end{lemma}

\begin{proof}
$\Ext^n_{D(\mathcal{O}_X)}(K, L) =
\Hom_{D(\mathcal{O}_X)}(K, L[n])$이므로 (2)는 (1)에서 따른다. (1)을
증명하자. $K$가 완전이므로 다음을 얻는다.
$$
\Hom_{D(\mathcal{O}_X)}(K, L) =
H^0(X, K^\vee \otimes_{\mathcal{O}_X}^\mathbf{L} L)
$$
여기서 $K^\vee$는 $K$의 ``쌍대'' 완전 복합체이다. Cohomology, Lemma
\ref{cohomology-lemma-dual-perfect-complex}를 보라. Lemmas
\ref{perfect-lemma-quasi-coherence-tensor-product} 및 \ref{perfect-lemma-pseudo-coherent}에
의해 $K^\vee \otimes_{\mathcal{O}_X}^\mathbf{L} L$은 $D_\QCoh(X)$에 속한다
(완전 복합체가 준연접 코호몰로지 층들을 가짐을 보기 위해 두 번째
보조정리를 사용한다). $K^\vee$가 $[a, b]$에서 Tor 진폭을 갖는다고 하자.
그러면 스펙트럼 열
$$
E_1^{p, q} = H^p(K^\vee \otimes_{\mathcal{O}_X}^\mathbf{L} H^q(L))
\Rightarrow
H^{p + q}(K^\vee \otimes_{\mathcal{O}_X}^\mathbf{L} L)
$$
에 의해 $q \in [j - b, j - a]$에 대해 $H^q(L) = 0$이면
$H^j(K^\vee \otimes_{\mathcal{O}_X}^\mathbf{L} L)$은 영이다. $N$을
Cohomology of Schemes, Lemma
\ref{coherent-lemma-vanishing-nr-affines-quasi-separated}의 정수 $d$라 하자.
코호몰로지 층들
$$
H^{-N}(K^\vee \otimes_{\mathcal{O}_X}^\mathbf{L} L),
\ H^{-N + 1}(K^\vee \otimes_{\mathcal{O}_X}^\mathbf{L} L),
\ \ldots,
\ H^0(K^\vee \otimes_{\mathcal{O}_X}^\mathbf{L} L)
$$
이 영이면 $H^0(X, K^\vee \otimes_{\mathcal{O}_X}^\mathbf{L} L)$은 영이다.
실제로 위에서 인용한 보조정리와 Lemma
\ref{perfect-lemma-application-nice-K-injective}에 의해
$$
H^0(X, K^\vee \otimes_{\mathcal{O}_X}^\mathbf{L} L) =
H^0(X, \tau_{\geq -N}(K^\vee \otimes_{\mathcal{O}_X}^\mathbf{L} L))
$$
이고, 코호몰로지 층들의 소멸에 의해 이는
$H^0(X, \tau_{\geq 1}(K^\vee \otimes_{\mathcal{O}_X}^\mathbf{L} L))$과 같다.
이는 Derived Categories, Lemma \ref{derived-lemma-negative-vanishing}에 의해
영이다. 따라서 $i \in [-b - N, -a]$에 대해 $H^i(L) = 0$이면
$\Hom_{D(\mathcal{O}_X)}(K, L)$은 영이다.
\end{proof}

\begin{lemma}
\label{perfect-lemma-lift-perfect-complex-plus-locally-free}
$X$를 아핀 스킴이라 하고 $U \subset X$를 준콤팩트 열린집합이라 하자.
$D(\mathcal{O}_U)$의 모든 완전 대상 $E$에 대해, $\mathcal{F}[-r] \oplus E$가
$D(\mathcal{O}_X)$의 완전 대상의 제한이 되게 하는 정수 $r$과 $U$ 위의
유한 국소 자유 층 $\mathcal{F}$가 존재한다.
\end{lemma}

\begin{proof}
$X = \Spec(A)$라고 하자. 완전 복합체는 유사연접임을 상기하자. Cohomology,
Lemma \ref{cohomology-lemma-perfect}를 보라. Lemma
\ref{perfect-lemma-lift-pseudo-coherent}에 의해, $D(\mathcal{O}_U)$에서 $E$가
$\mathcal{F}^\bullet|_U$와 동형이 되게 하는 유한 자유 $A$-가군들의 위로
유계인 복합체 $\mathcal{F}^\bullet$을 찾을 수 있다. Cohomology, Lemma
\ref{cohomology-lemma-perfect}와 $U$가 준콤팩트라는 사실에 의해 $E$는 유한
Tor 차원을 갖는다. $E$가 $[a, b]$에서 Tor 진폭을 갖는다고 하자. $r < a$를
택하고 다음과 같이 두자.
$$
\mathcal{K} = \Ker(\mathcal{F}^{r} \to \mathcal{F}^{r + 1})
= \Im(\mathcal{F}^{r - 1} \to \mathcal{F}^r).
$$
$E$가 $[a, b]$에서 Tor 진폭을 가지므로
$\mathcal{F} = \mathcal{K}|_U$는 평탄하다(Cohomology, Lemma
\ref{cohomology-lemma-last-one-flat}). 따라서 $\mathcal{F}$는 평탄하고 유한
표시이며, 그러므로 유한 국소 자유이다(Properties, Lemma
\ref{properties-lemma-finite-locally-free}). 따라서
$$
\mathcal{F} \to \mathcal{F}^r|_U \to \mathcal{F}^{r + 1}|_U \to \ldots
$$
은 $E$를 나타내는 $U$ 위의 엄밀 완전 복합체이다. 한편 복합체
$P = (\mathcal{F}^r \to \mathcal{F}^{r + 1} \to \ldots )$는 $X$ 위의
완전 복합체이다. 단순 절단을 사용하여 다음 완전 삼각형을 얻는다.
$$
P|_U \to E \to \mathcal{F}[-r - 1] \to (P|_U)[1]
$$
사상 $E \to \mathcal{F}[-r - 1]$이 $D(\mathcal{O}_U)$에서 영이면
$P|_U = \mathcal{F}[-r - 2] \oplus E$이다. Derived Categories, Lemma
\ref{derived-lemma-split}을 보라. 예를 들어 Lemma
\ref{perfect-lemma-ext-from-perfect-into-bounded-QCoh}에 의해 $r \ll 0$이면 이것이
성립한다.
\end{proof}

\begin{lemma}
\label{perfect-lemma-lift-map}
$X$를 아핀 스킴이라 하고 $U \subset X$를 준콤팩트 열린집합이라 하자.
$E, E'$을 $D_\QCoh(\mathcal{O}_X)$의 대상이라 하고 $E$가 완전이라고 하자.
모든 사상 $\alpha : E|_U \to E'|_U$에 대해 다음 사상들이 존재한다.
$$
E \xleftarrow{\beta} E_1 \xrightarrow{\gamma} E'
$$
여기서 이들은 $X$ 위 복합체들의 사상이고, $E_1$은 완전이며,
$\beta : E_1 \to E$의 $U$로의 제한은 동형이고,
$\alpha = \gamma|_U \circ \beta|_U^{-1}$이다. 더 나아가 $X$ 위의 어떤
완전 복합체 $I$에 대해 $E_1 = E \otimes_{\mathcal{O}_X}^\mathbf{L} I$라고
가정할 수 있다.
\end{lemma}

\begin{proof}
$X = \Spec(A)$ 및 $U = D(f_1) \cup \ldots \cup D(f_r)$로 쓰자. $E$를
나타내는 유한 사영 $A$-가군들의 유한 복합체 $M^\bullet$을 택하자(Lemma
\ref{perfect-lemma-perfect-affine}). $E'$을 나타내는 $A$-가군의 복합체
$(M')^\bullet$을 택하자(Lemma \ref{perfect-lemma-affine-compare-bounded}). 이 경우
복합체 $H^\bullet = \Hom_A(M^\bullet, (M')^\bullet)$은 $A$-가군의
복합체이고, 이에 연관된 준연접 $\mathcal{O}_X$-가군의 복합체는
$R\SheafHom(E, E')$을 나타낸다. Cohomology, Lemma
\ref{cohomology-lemma-Rhom-strictly-perfect}를 보라. 그러면 $\alpha$는
$H^0(U, R\SheafHom(E, E'))$의 원소 $s$를 정한다. Cohomology, Lemma
\ref{cohomology-lemma-section-RHom-over-U}를 보라. $s$에 대응하는 어떤
$e$와 사상
$$
\xi : I^\bullet(f_1^e, \ldots, f_r^e) \to \Hom_A(M^\bullet, (M')^\bullet)
$$
이 존재한다. Proposition \ref{perfect-proposition-represent-cohomology-class-on-open}을
보라. 다음 복합체에 연관된 준연접 $\mathcal{O}_X$-가군의 복합체에 대응하는
대상을 $E_1$이라 하자.
$$
\text{Tot}(I^\bullet(f_1^e, \ldots, f_r^e) \otimes_A M^\bullet)
$$
표준 사상 $I^\bullet(f_1^e, \ldots, f_r^e) \to A$를 사용하여
$E_1 \to E$를 얻고, $\xi$와 Cohomology, Lemma
\ref{cohomology-lemma-section-RHom-over-U}를 사용하여 $E_1 \to E'$을 얻는다.
\end{proof}

\begin{lemma}
\label{perfect-lemma-lift-perfect-complex-plus-shift}
$X$를 아핀 스킴이라 하고 $U \subset X$를 준콤팩트 열린집합이라 하자.
$D(\mathcal{O}_U)$의 모든 완전 대상 $F$에 대해 대상 $F \oplus F[1]$은
$D(\mathcal{O}_X)$의 완전 대상의 제한이다.
\end{lemma}

\begin{proof}
Lemma \ref{perfect-lemma-lift-perfect-complex-plus-locally-free}에 의해 어떤 유한
국소 자유 $\mathcal{O}_U$-가군 $\mathcal{F}$에 대해
$E|_U = \mathcal{F}[r] \oplus F$가 되게 하는 $D(\mathcal{O}_X)$의 완전
대상 $E$를 찾을 수 있다. Lemma \ref{perfect-lemma-lift-map}에 의해
$(E_1)|_U \cong E|_U$이고 $\alpha|_U$가 다음 사상이 되게 하는 완전
복합체의 사상 $\alpha : E_1 \to E$를 찾을 수 있다.
$$
\left(
\begin{matrix}
\text{id}_{\mathcal{F}[r]} & 0 \\
0 & 0
\end{matrix}
\right)
:
\mathcal{F}[r] \oplus F \to \mathcal{F}[r] \oplus F
$$
그러면 $\alpha$의 사상뿔이 원하는 대상이다.
\end{proof}

\begin{lemma}
\label{perfect-lemma-perfect-into-support-on-T}
$X$를 준콤팩트 준분리 스킴이라 하고 $f \in \Gamma(X, \mathcal{O}_X)$라
하자. 다음을 만족하는 $D_\QCoh(\mathcal{O}_X)$의 임의의 사상
$\alpha : E \to E'$에 대해
\begin{enumerate}
\item $E$는 완전이고,
\item $E'$은 $T = V(f)$ 위에 지지된다.
\end{enumerate}
$f^n \alpha  = 0$이 되게 하는 $n \geq 0$이 존재한다.
\end{lemma}

\begin{proof}
유도 범주의 사상에 대한 Mayer--Vietoris를 사용할 수 있다. Cohomology,
Lemma \ref{cohomology-lemma-mayer-vietoris-hom}을 보라. 따라서
$X = U \cup V$이고 $f|_U$, $f|_V$, $f|_{U \cap V}$에 대해 보조정리의
결과가 성립하면 $f$에 대해서도 성립한다. 그러므로 $X$가 아핀인 경우를
증명하면 충분하다. Cohomology of Schemes, Lemma
\ref{coherent-lemma-induction-principle}을 보라.

\medskip\noindent
$X = \Spec(A)$라 하자. 그러면 $f \in A$이다. 앞으로 별도의 언급 없이
Lemma \ref{perfect-lemma-affine-compare-bounded}의 동치 $D(A) = D_\QCoh(X)$를
사용한다. $E$를 유한 사영 $A$-가군들의 유한 복합체 $P^\bullet$로
나타내자. 이는 Lemma \ref{perfect-lemma-perfect-affine}에 의해 가능하다. $P^t$가
영이 아닌 가장 큰 정수 $t$를 택하자. 완전 삼각형
$$
P^t[-t] \to P^\bullet \to \sigma_{\leq t - 1}P^\bullet \to P^t[-t + 1]
$$
에 의해 $P^\bullet$의 길이에 대한 귀납법으로 $P^\bullet$이 하나의 영이
아닌 항만 갖는 경우로 환원할 수 있다. 이 사실과 이동 함자를 사용하면
$P^\bullet$이 차수 $0$에 놓인 하나의 유한 사영 $A$-가군 $P$로 이루어진
경우로 환원된다. $E'$을 $A$-가군의 복합체 $M^\bullet$로 나타내자. 그러면
$\alpha$는 사상 $P \to H^0(M^\bullet)$에 대응한다. 가정 (2)에 의해 가군
$H^0(M^\bullet)$은 $V(f)$ 위에 지지되므로, $H^0(M^\bullet)$의 모든 원소는 $f$의 어떤
거듭제곱에 의해 소멸된다. $P$가 유한 $A$-가군이므로 원하는 대로 어떤
$n$에 대해 사상 $f^n\alpha : P \to H^0(M^\bullet)$은 영이다.
\end{proof}

\begin{lemma}
\label{perfect-lemma-lift-perfect-complex-plus-shift-support}
$X$를 아핀 스킴이라 하고 $T \subset X$를 $X \setminus T$가 준콤팩트인
닫힌 부분집합이라 하자. $U \subset X$를 준콤팩트 열린집합이라 하자.
$T \cap U$ 위에 지지되는 $D(\mathcal{O}_U)$의 모든 완전 대상 $F$에 대해,
대상 $F \oplus F[1]$은 $T$ 위에 지지되는 $D(\mathcal{O}_X)$의 완전 대상
$E$의 제한이다.
\end{lemma}

\begin{proof}
$T = V(g_1, \ldots, g_s)$라고 하자. $g_j$를 어떤 거듭제곱으로 바꾼 뒤
$F$ 위에서 $g_j$의 곱셈이 영이라고 가정할 수 있다. Lemma
\ref{perfect-lemma-perfect-into-support-on-T}를 보라. Lemma
\ref{perfect-lemma-lift-perfect-complex-plus-shift}에서와 같은 $E$를 택하자. 사상
$g_j : E \to E$의 $U$로의 제한은 영임에 유의하자. 다음 완전 삼각형을
택하자.
$$
E \xrightarrow{g_1} E \to C_1 \to E[1]
$$
Derived Categories, Lemma \ref{derived-lemma-split}에 의해 대상 $C_1$의
$U$로의 제한은 $F \oplus F[1] \oplus F[1] \oplus F[2]$이다. 더 나아가
Derived Categories, Lemma \ref{derived-lemma-third-map-square-zero}에 의해
$g_1 : C_1 \to C_1$의 제곱은 영이다. 즉, 그림
$$
\xymatrix{
E \ar[r] \ar[d]_0 & C_1 \ar[d]_{g_1} \ar[r] & E[1] \ar[d]_0 \\
E \ar[r] & C_1 \ar[r] & E[1]
}
$$
은 합성 $E \xrightarrow{g_1} E \to C_1$과
$C_1 \to E[1] \xrightarrow{g_1} E[1]$이 영이므로 가환한다. 계속하여
$C_{i + 1}$을 사상 $g_i : C_i \to C_i$의 사상뿔로 두면, $U$로의 제한이
다음과 같은 $T$ 위에 지지되는 $X$ 위의 완전 복합체 $C_s$를 얻는다.
$$
F \oplus F[1]^{\oplus s} \oplus F[2]^{\oplus {s \choose 2}}
\oplus \ldots \oplus F[s]
$$
Lemma \ref{perfect-lemma-lift-map}에서와 같이 $\beta|_U$가 동형이고
$\gamma|_U \circ \beta|_U^{-1}$가 다음 사상이 되게 하는 완전 복합체의
사상들 $\beta : C' \to C_s$와 $\gamma : C' \to C_s$를 택하자.
$$
F \oplus F[1]^{\oplus s} \oplus F[2]^{\oplus {s \choose 2}}
\oplus \ldots \oplus F[s]
\to
F \oplus F[1]^{\oplus s} \oplus F[2]^{\oplus {s \choose 2}}
\oplus \ldots \oplus F[s]
$$
이 사상은 $F$ 위에서는 영이고 다른 모든 직합 인자 위에서는 항등이다.
Lemma \ref{perfect-lemma-lift-map}에 의해 $X$ 위의 어떤 완전 복합체 $I$에 대해
$C' = C_s \otimes^\mathbf{L} I$이기도 하다. 따라서
$g_j^2\text{id}_{C_s}$가 영이라는 사실은 $C'$에 대해서도 같은 사실을
함의한다. 그러므로 $C'$도 $T$ 위에 지지된다. 이제
$\text{Cone}(\gamma)$가 원하는 대상이다.
\end{proof}

\noindent
다음 보조정리의 특수한 경우는 \cite{Neeman-Grothendieck}에서 찾을 수 있다.

\begin{lemma}
\label{perfect-lemma-lift-map-from-perfect-complex-with-support}
$X$를 준콤팩트 준분리 스킴이라 하자. $U \subset X$를 준콤팩트
열린집합이라 하고, $T \subset X$를 $X \setminus T$가 $X$에서 역콤팩트인
닫힌 부분집합이라 하자. $E$를 $D_\QCoh(\mathcal{O}_X)$의 대상이라 하자.
$\alpha : P \to E|_U$를 사상이라 하며, 여기서 $P$는 $T \cap U$ 위에
지지되는 $D(\mathcal{O}_U)$의 완전 대상이라고 하자. 그러면 $T$ 위에
지지되는 $D(\mathcal{O}_X)$의 완전 대상 $R$과 사상 $\beta : R \to E$가
존재하여, $D(\mathcal{O}_U)$에서 $P$는 $R|_U$의 직합 인자이고 이는
$\alpha$와 $\beta|_U$에 양립한다.
\end{lemma}

\begin{proof}
$X$가 준콤팩트이므로 어떤 정수 $m$과 $X$의 아핀 열린집합들 $V_j$가 있어
$X = U \cup V_1 \cup \ldots \cup V_m$이다. $m$에 대한 귀납법으로
$m = 1$이라고 가정할 수 있다. 다시 말해 $V$가 아핀이고
$X = U \cup V$라고 가정할 수 있다. Lemma
\ref{perfect-lemma-lift-perfect-complex-plus-shift-support}에 의해 $T \cap V$ 위에
지지되는 $D(\mathcal{O}_V)$의 완전 대상 $Q$와 동형
$Q|_{U \cap V} \to (P \oplus P[1])|_{U \cap V}$를 택할 수 있다. Lemma
\ref{perfect-lemma-lift-map}에 의해 $Q$를 $Q \otimes^\mathbf{L} I$로 바꾸고(여전히
$T \cap V$ 위에 지지된다) 다음 사상이
$$
Q|_{U \cap V} \to (P \oplus P[1])|_{U \cap V}
\longrightarrow P|_{U \cap V}
\longrightarrow
E|_{U \cap V}
$$
$Q \to E|_V$로 올라간다고 가정할 수 있다. Cohomology, Lemma
\ref{cohomology-lemma-glue}에 의해 $a|_U$가
$P \oplus P[1] \to E|_U$와 동형이고 $a|_V$가 $Q \to E|_V$와 동형인
$D(\mathcal{O}_X)$의 사상 $a : R \to E$를 얻는다. 따라서 원하는 대로
$R$은 완전이고 $T$ 위에 지지된다.
\end{proof}

\begin{remark}
\label{perfect-remark-addendum}
Lemma \ref{perfect-lemma-lift-map-from-perfect-complex-with-support}의 증명에 의해
$$
R|_U = P \oplus P^{\oplus n_1}[1] \oplus \ldots \oplus P^{\oplus n_m}[m]
$$
인 어떤 $m \geq 0$과 $n_j \geq 0$이 존재한다. 따라서 $R|_U$의 가장 높은
차수 코호몰로지 층은 $P$의 그것과 같다. 사상
$P^{\oplus n_1}[1] \oplus \ldots \oplus P^{\oplus n_m}[m] \to R|_U$에 대해
이 구성을 반복하고 사상뿔을 취하며 귀납법을 사용하면, 임의로 주어진 차수
위에서 $R|_U$와 $P$의 코호몰로지 층들이 같게 할 수 있다.
\end{remark}



\section{완전 복합체에 의한 근사}
\label{perfect-section-approximation}

\noindent
이 절에서는 Neeman과 Lipman에 의한 관찰, 즉 유사연접 복합체를 완전 복합체로
``근사''할 수 있다는 것을 논의한다.

\begin{definition}
\label{perfect-definition-approximation-holds}
$X$를 스킴이라 하자. 다음 조건을 만족하는 삼중항 $(T, E, m)$을 생각하자.
\begin{enumerate}
\item $T \subset X$는 닫힌 부분집합이다.
\item $E$는 $D_\QCoh(\mathcal{O}_X)$의 대상이다.
\item $m \in \mathbf{Z}$이다.
\end{enumerate}
$T$에 지지된 $D(\mathcal{O}_X)$의 완전 대상 $P$와, $i > m$에 대해 동형
$H^i(P) \to H^i(E)$ 및 전사 $H^m(P) \to H^m(E)$를 유도하는 사상
$\alpha : P \to E$가 존재하면, 삼중항 $(T, E, m)$에 대해
{\it 근사가 성립한다}고 한다.
\end{definition}

\noindent
모든 삼중항에 대해 근사가 성립할 수는 없다. 실제로 삼중항 $(T, E, m)$에
대해 근사가 성립하면 다음이 자명하다.
\begin{enumerate}
\item $E$는 $m$-유사연접이다. Cohomology, Definition
\ref{cohomology-definition-pseudo-coherent}을 보라.
\item $i \geq m$에 대해 코호몰로지 층 $H^i(E)$는 $T$에 지지된다.
\end{enumerate}
더욱이 완전 복합체의 ``지지''는 그 여집합이 $X$에서 역콤팩트인 닫힌
부분스킴이다(세부사항은 생략한다). 따라서 $T$에 대한 이 가정 없이는 근사가
성립하리라고 기대할 수 없다. 이는 다음 정의의 조건들을 부분적으로 설명한다.

\begin{definition}
\label{perfect-definition-approximation}
$X$를 스킴이라 하자. $X \setminus T$가 $X$에서 역콤팩트인 임의의 닫힌
부분집합 $T \subset X$에 대해 다음 성질을 갖는 정수 $r$이 존재하면,
$X$에서 {\it 완전 복합체에 의한 근사가 성립한다}고 한다. 즉 Definition
\ref{perfect-definition-approximation-holds}에서와 같은 삼중항 $(T, E, m)$이
다음 조건을 만족할 때마다
\begin{enumerate}
\item $E$는 $(m - r)$-유사연접이고,
\item $i \geq m - r$에 대해 $H^i(E)$는 $T$에 지지되면,
\end{enumerate}
근사가 성립한다.
\end{definition}

\noindent
준콤팩트이고 준분리인 스킴에서 완전 복합체에 의한 근사가 성립함을 증명할
것이다. 두 번째 조건은 우리의 증명 방법에 필요한 것으로 보인다. 아래의
논증을 주의 깊게 분석하면 첫 번째 조건을 ``$E$가 $m$-유사연접이다''로
약화할 수 있을지도 모른다.

\begin{lemma}
\label{perfect-lemma-open}
$X$를 스킴이라 하고 $U \subset X$를 열린 부분스킴이라 하자.
$(T, E, m)$을 Definition \ref{perfect-definition-approximation-holds}에서와 같은
삼중항이라 하자. 다음을 가정하자.
\begin{enumerate}
\item $T \subset U$이다.
\item $(T, E|_U, m)$에 대해 근사가 성립한다.
\item $i \geq m$에 대한 층 $H^i(E)$는 $T$에 지지된다.
\end{enumerate}
그러면 $(T, E, m)$에 대해 근사가 성립한다.
\end{lemma}

\begin{proof}
$j : U \to X$를 포함 사상이라 하자. $P \to E|_U$가 $U$ 위의 삼중항
$(T, E|_U, m)$의 근사이면,
$j_!P = Rj_*P \to j_!(E|_U) \to E$는 $X$ 위의 $(T, E, m)$의 근사이다.
Cohomology, Lemmas \ref{cohomology-lemma-pushforward-restriction} 및
\ref{cohomology-lemma-pushforward-perfect}를 보라.
\end{proof}

\begin{lemma}
\label{perfect-lemma-approximation-affine}
$X$를 아핀 스킴이라 하자. 다음 조건을 갖는 정수 $r \geq 0$이 존재하는,
Definition \ref{perfect-definition-approximation-holds}에서와 같은 모든 삼중항
$(T, E, m)$에 대해 근사가 성립한다.
\begin{enumerate}
\item $E$는 $m$-유사연접이다.
\item $i \geq m - r + 1$에 대해 $H^i(E)$는 $T$에 지지된다.
\item $X \setminus T$는 $r$개의 아핀 열린집합의 합집합이다.
\end{enumerate}
특히 아핀 스킴에서는 완전 복합체에 의한 근사가 성립한다.
\end{lemma}

\begin{proof}
$X = \Spec(A)$라고 하자. $T = V(f_1, \ldots, f_r)$로 쓰자.
($r = 0$, 즉 $T = X$인 경우는 Lemma
\ref{perfect-lemma-pseudo-coherent-affine}와 정의들에서 곧바로 따른다.)
$(T, E, m)$을 보조정리에서와 같은 삼중항이라 하자. $H^t(E)$가 0이 아닌
가장 큰 정수를 $t$라 하자. $t$에 대한 귀납법을 사용한다. 기초 단계는
$t < m$이며, 이 경우 결과는 자명하다. 이제 $t \geq m$이라고 하자.
Cohomology, Lemma \ref{cohomology-lemma-finite-cohomology}에 의해 층
$H^t(E)$는 유한형이다. 이는 준연접이므로 유한개의 단면으로 생성된다
(Properties, Lemma \ref{properties-lemma-finite-type-module}).
모든 $s \in \Gamma(X, H^t(E)) = H^t(X, E)$에 대해(Lemma
\ref{perfect-lemma-affine-compare-bounded}의 증명을 보라), $s$가
$H^0(K_e) = H^t(K_e[-t]) \to H^t(E)$의 상에 놓이도록 하는 $e > 0$과
사상 $K_e[-t] \to E$를 찾을 수 있다. Lemma
\ref{perfect-lemma-represent-cohomology-class-on-closed}를 보라. 이 사상들의 유한
직합을 취하면 사상 $P \to E$를 얻는다. 여기서 $P$는 $T$에 지지된 완전
복합체이고, $i > t$에 대해 $H^i(P) = 0$이며, $H^t(P) \to E$는 전사이다.
완전 삼각형을 하나 선택하자.
$$
P \to E \to E' \to P[1]
$$
그러면 $E'$는 $m$-유사연접이고(Cohomology, Lemma
\ref{cohomology-lemma-cone-pseudo-coherent}), $i \geq t$에 대해
$H^i(E') = 0$이며, $i \geq m - r + 1$에 대해 $H^i(E')$는 $T$에
지지된다. 귀납법으로 $(T, E', m)$의 근사 $P' \to E'$를 찾는다.
합성 $P' \to E' \to P[1]$을 완전 삼각형
$P \to P'' \to P' \to P[1]$에 넣고, 사상 $P' \to E'$와
$P[1] \to P[1]$을 다음 완전 삼각형의 사상으로 확장한다.
$$
\xymatrix{
P \ar[r] \ar[d] & P'' \ar[d] \ar[r] & P' \ar[d] \ar[r] & P[1] \ar[d] \\
P \ar[r] &  E \ar[r] & E' \ar[r] & P[1]
}
$$
여기서 TR3을 사용하였다. 그러면 $P''$는 $T$에 지지된 완전 복합체이다
(Cohomology, Lemma \ref{cohomology-lemma-two-out-of-three-perfect}).
간단한 그림 추적으로 $P'' \to E$가 원하는 근사임을 알 수 있다.
\end{proof}

\begin{lemma}
\label{perfect-lemma-induction-step}
$X$를 스킴이라 하자. $X = U \cup V$를 열린 덮개라 하자. 여기서 $U$는
준콤팩트이고, $V$는 아핀이며, $U \cap V$는 준콤팩트이다. $U$에서 완전
복합체에 의한 근사가 성립하면 $X$에서도 근사가 성립한다.
\end{lemma}

\begin{proof}
$T \subset X$를 $X \setminus T$가 $X$에서 역콤팩트인 닫힌 부분집합이라
하자. $r_U$를 쌍 $(U, T \cap U)$에 맞춘 Definition
\ref{perfect-definition-approximation}의 정수라 하자. $T' = T \setminus U$로 놓자.
$T' \subset V$이고, $T$에 대한 가정에 의해
$V \setminus T' = (X \setminus T) \cap U \cap V$는 준콤팩트임에 유의하라.
$V \setminus T'$를 덮는 데 필요한 아핀의 수를 $r'$라 하자. 쌍 $(X, T)$에
대해 $r = \max(r_U, r')$가 작동한다고 주장한다.

\medskip\noindent
이를 보기 위해 $E$가 $(m - r)$-유사연접이고 $i \geq m - r$에 대해
$H^i(E)$가 $T$에 지지되는 삼중항 $(T, E, m)$을 선택하자.
$H^t(E)|_U$가 0이 아닌 가장 큰 정수를 $t$라 하자. ($U$가 준콤팩트이고
$E|_U$가 $(m - r)$-유사연접이므로 이러한 정수가 존재한다.) $t$에 대한
귀납법으로 $E$를 근사할 수 있음을 증명하겠다.

\medskip\noindent
기초 단계: $t \leq m - r'$이라 하자. 이는 $i \geq m - r'$에 대해
$H^i(E)$가 $T'$에 지지됨을 뜻한다. 따라서 Lemma
\ref{perfect-lemma-approximation-affine}는 $V$ 위의 $(T', E|_V, m)$의 근사
$P \to E|_V$가 존재함을 보장한다. Lemma \ref{perfect-lemma-open}을 적용하면
$(T', E, m)$을 근사할 수 있음을 알 수 있다. 이러한 근사는
$(T, E, m)$의 근사이기도 하다.

\medskip\noindent
귀납 단계. $(T \cap U, E|_U, m)$의 근사 $P \to E|_U$를 선택하자.
특히 이는 전사 $H^t(P) \to H^t(E|_U)$를 준다. Lemma
\ref{perfect-lemma-lift-perfect-complex-plus-shift-support}에 의해 $T \cap V$에
지지된 $D(\mathcal{O}_V)$의 완전 대상 $Q$와 동형
$Q|_{U \cap V} \to (P \oplus P[1])|_{U \cap V}$를 선택할 수 있다.
Lemma \ref{perfect-lemma-lift-map}에 의해 $Q$를 $Q \otimes^\mathbf{L} I$로
바꾸고, 다음 사상이
$$
Q|_{U \cap V} \to (P \oplus P[1])|_{U \cap V}
\longrightarrow P|_{U \cap V}
\longrightarrow
E|_{U \cap V}
$$
$Q \to E|_V$로 올라간다고 가정할 수 있다. Cohomology, Lemma
\ref{cohomology-lemma-glue}에 의해 $D(\mathcal{O}_X)$의 사상
$a : R \to E$를 찾을 수 있다. 여기서 $a|_U$는
$P \oplus P[1] \to E|_U$와 동형이고 $a|_V$는 $Q \to E|_V$와 동형이다.
따라서 $R$은 완전이고 $T$에 지지되며, 사상 $H^t(R) \to H^t(E)$는
$U$로 제한했을 때 전사이다. 완전 삼각형을 하나 선택하자.
$$
R \to E \to E' \to R[1]
$$
그러면 $E'$는 $(m - r)$-유사연접이고(Cohomology, Lemma
\ref{cohomology-lemma-cone-pseudo-coherent}), $i \geq t$에 대해
$H^i(E')|_U = 0$이며, $i \geq m - r$에 대해 $H^i(E')$는 $T$에
지지된다. 귀납법으로 $(T, E', m)$의 근사 $R' \to E'$를 찾는다.
합성 $R' \to E' \to R[1]$을 완전 삼각형
$R \to R'' \to R' \to R[1]$에 넣고, 사상 $R' \to E'$와
$R[1] \to R[1]$을 다음 완전 삼각형의 사상으로 확장한다.
$$
\xymatrix{
R \ar[r] \ar[d] & R'' \ar[d] \ar[r] & R' \ar[d] \ar[r] & R[1] \ar[d] \\
R \ar[r] &  E \ar[r] & E' \ar[r] & R[1]
}
$$
여기서 TR3을 사용하였다. 그러면 $R''$는 $T$에 지지된 완전 복합체이다
(Cohomology, Lemma \ref{cohomology-lemma-two-out-of-three-perfect}).
간단한 그림 추적으로 $R'' \to E$가 원하는 근사임을 알 수 있다.
\end{proof}

\begin{theorem}
\label{perfect-theorem-approximation}
$X$를 준콤팩트이고 준분리인 스킴이라 하자. 그러면 $X$에서 완전 복합체에
의한 근사가 성립한다.
\end{theorem}

\begin{proof}
Cohomology of Schemes, Lemma \ref{coherent-lemma-induction-principle}의
귀납 원리와 Lemmas \ref{perfect-lemma-induction-step} 및
\ref{perfect-lemma-approximation-affine}에서 따른다.
\end{proof}






\section{유도 범주의 생성}
\label{perfect-section-generating}

\noindent
이 절에서는 준콤팩트이고 준분리인 스킴의 유도 범주
$D_\QCoh(\mathcal{O}_X)$가 하나의 완전 대상으로 생성될 수 있음을 증명한다.
독자에게 Bondal과 van den Bergh의 훌륭한 논문에 있는 이 결과의 증명을
읽기를 권한다. \cite{BvdB}를 보라.

\begin{lemma}
\label{perfect-lemma-direct-summand-of-a-restriction}
$X$를 준콤팩트이고 준분리인 스킴이라 하고, $U$를 준콤팩트 열린
부분스킴이라 하자. $P$를 $D(\mathcal{O}_U)$의 완전 대상이라 하자.
그러면 $P$는 $D(\mathcal{O}_X)$의 어떤 완전 대상을 제한한 것의 직합인자이다.
\end{lemma}

\begin{proof}
Lemma \ref{perfect-lemma-lift-map-from-perfect-complex-with-support}의 특수한 경우이다.
\end{proof}

\begin{lemma}
\label{perfect-lemma-orthogonal-koszul-complex}
\begin{reference}
\cite[Proposition 6.1]{Bokstedt-Neeman}
\end{reference}
Situation \ref{perfect-situation-complex}에서 열린 매장을 $j : U \to X$로 나타내고,
$K$를 $A$ 위의 $f_1, \ldots, f_r$에 대한 코줄 복합체에 대응하는
$D(\mathcal{O}_X)$의 완전 대상이라 하자. $E \in D_\QCoh(\mathcal{O}_X)$에
대해 다음은 동치이다.
\begin{enumerate}
\item $E = Rj_*(E|_U)$이다.
\item 모든 $n \in \mathbf{Z}$에 대해
$\Hom_{D(\mathcal{O}_X)}(K[n], E) = 0$이다.
\end{enumerate}
\end{lemma}

\begin{proof}
완전 삼각형 $E \to Rj_*(E|_U) \to N \to E[1]$을 선택하자. 모든 $n$에 대해
$$
\Hom_{D(\mathcal{O}_X)}(K[n], Rj_*(E|_U)) =
\Hom_{D(\mathcal{O}_U)}(K|_U[n], E) = 0
$$
임에 유의하라. 실제로 $K|_U = 0$이다. 따라서 $N$에 대해 결과를 증명하면
충분하다. 다시 말해 $E$의 $U$로의 제한이 0이라고 가정해도 된다. 코줄
복합체들의 다음 완전 삼각형들이 존재함에 유의하라.
$$
K^\bullet(f_1^{e_1}, \ldots, f_i^{e'_i}, \ldots, f_r^{e_r}) \to
K^\bullet(f_1^{e_1}, \ldots, f_i^{e'_i + e''_i}, \ldots, f_r^{e_r}) \to
K^\bullet(f_1^{e_1}, \ldots, f_i^{e''_i}, \ldots, f_r^{e_r}) \to \ldots
$$
More on Algebra, Lemma \ref{more-algebra-lemma-koszul-mult}를 보라. 따라서 모든
$n \in \mathbf{Z}$에 대해 $\Hom_{D(\mathcal{O}_X)}(K[n], E) = 0$이면,
Lemma \ref{perfect-lemma-represent-cohomology-class-on-closed}에서와 같이 $K$를
$K_e$로 바꾸어도 같은 명제가 참이다. 그러므로 이 보조정리와, $E$가
$A$-가군의 복합체 $R\Gamma(X, E)$에 의해 결정된다는 사실에서 우리의
보조정리가 곧바로 따른다. Lemma \ref{perfect-lemma-affine-compare-bounded}를 보라.
\end{proof}

\begin{theorem}
\label{perfect-theorem-bondal-van-den-Bergh}
$X$를 준콤팩트이고 준분리인 스킴이라 하자. 범주
$D_\QCoh(\mathcal{O}_X)$는 하나의 완전 대상으로 생성될 수 있다. 더
정확히 말해, $D(\mathcal{O}_X)$의 완전 대상 $P$가 존재하여
$E \in D_\QCoh(\mathcal{O}_X)$에 대해 다음은 동치이다.
\begin{enumerate}
\item $E = 0$이다.
\item 모든 $n \in \mathbf{Z}$에 대해
$\Hom_{D(\mathcal{O}_X)}(P[n], E) = 0$이다.
\end{enumerate}
\end{theorem}

\begin{proof}
Cohomology of Schemes, Lemma \ref{coherent-lemma-induction-principle}의
귀납 원리를 사용하여 이를 증명하겠다.

\medskip\noindent
$X$가 아핀이면 $\mathcal{O}_X$는 완전 생성자이다. 이는 Lemma
\ref{perfect-lemma-affine-compare-bounded}에서 따른다.

\medskip\noindent
$X = U \cup V$가 열린 덮개이고, $U$는 준콤팩트이며, $U$에 대해 정리가
성립하고, $V$는 아핀 열린집합이라고 가정하자. $P$를
$D_\QCoh(\mathcal{O}_U)$의 생성자인 $D(\mathcal{O}_U)$의 완전 대상이라
하자. Lemma \ref{perfect-lemma-direct-summand-of-a-restriction}을 사용하면,
$U$로의 제한이 $P$를 한 직합인자로 갖는 직합인 $D(\mathcal{O}_X)$의 완전
대상 $Q$를 선택할 수 있다. $V = \Spec(A)$라고 하고 $Z = X \setminus U$라
하자. 이는 $V \setminus Z$가 준콤팩트인 $V$의 닫힌 부분집합이다.
$Z = V(f_1, \ldots, f_r)$가 되도록 $f_1, \ldots, f_r \in A$를 선택하자.
$K \in D(\mathcal{O}_V)$를 $A$ 위의 $f_1, \ldots, f_r$에 대한 코줄
복합체에 대응하는 완전 대상이라 하자. $K$는 닫힌집합 $Z \subset V$에
지지되므로, 직상 $K' = R(V \to X)_*K$는 그 $V$로의 제한이 $K$인
$D(\mathcal{O}_X)$의 완전 대상임에 유의하라(Cohomology, Lemma
\ref{cohomology-lemma-pushforward-perfect}). $Q \oplus K'$가
$D_\QCoh(\mathcal{O}_X)$의 생성자라고 주장한다.

\medskip\noindent
$Q \oplus K'$의 어떤 이동에서도 $E$로 가는 영이 아닌 사상이 없도록
$E$를 $D_\QCoh(\mathcal{O}_X)$의 대상이라 하자. Cohomology, Lemma
\ref{cohomology-lemma-pushforward-restriction}에 의해
$K' =  R(V \to X)_! K$이고, 따라서
$$
\Hom_{D(\mathcal{O}_X)}(K'[n], E) = \Hom_{D(\mathcal{O}_V)}(K[n], E|_V)
$$
이다. 그러므로 Lemma \ref{perfect-lemma-orthogonal-koszul-complex}에 의해 이 군들이
소멸하면 $E|_V$는 $R(U \cap V \to V)_*E|_{U \cap V}$와 동형이다.
이는 $E = R(U \to X)_*E|_U$를 함의한다(작은 세부사항은 생략한다).
이 경우
$$
\Hom_{D(\mathcal{O}_X)}(Q[n], E) = \Hom_{D(\mathcal{O}_U)}(Q|_U[n], E|_U)
$$
는 $\Hom_{D(\mathcal{O}_U)}(P[n], E|_U)$를 직합인자로 포함한다. 따라서
$P$의 선택에 의해 이 군들이 소멸하면 $E|_U$는 0이다. 그러므로 $E$는 0이다.
\end{proof}

\noindent
다음 결과는 정확히 같은 방법으로 증명되는 Theorem
\ref{perfect-theorem-bondal-van-den-Bergh}의 강화이다. 스킴 $X$의 닫힌 부분집합
$T$에 대해, $T$에 지지된 대상들로 이루어진 $D(\mathcal{O}_X)$의 엄밀
충만 포화 삼각 부분범주를 $D_T(\mathcal{O}_X)$로 나타냄을 상기하자
(Definition \ref{perfect-definition-supported-on}). 마찬가지로, 코호몰로지 층들이
준연접이고 $T$에 지지되는 복합체들로 이루어진 $D(\mathcal{O}_X)$의 엄밀
충만 포화 삼각 부분범주를 $D_{\QCoh, T}(\mathcal{O}_X)$로 나타낸다.

\begin{lemma}
\label{perfect-lemma-generator-with-support}
\begin{reference}
\cite[Theorem 6.8]{Rouquier-dimensions}
\end{reference}
$X$를 준콤팩트이고 준분리인 스킴이라 하자. $T \subset X$를
$X \setminus T$가 준콤팩트인 닫힌 부분집합이라 하자. 위의 표기법에서 범주
$D_{\QCoh, T}(\mathcal{O}_X)$는 하나의 완전 대상으로 생성된다.
\end{lemma}

\begin{proof}
Cohomology of Schemes, Lemma \ref{coherent-lemma-induction-principle}의
귀납 원리를 사용하여 이를 증명하겠다.

\medskip\noindent
$X = \Spec(A)$가 아핀이라고 가정하자. 이 경우
$T = V(f_1, \ldots, f_r)$가 되도록 하는 $f_1, \ldots, f_r \in A$가
존재한다. $K$를 Lemma \ref{perfect-lemma-orthogonal-koszul-complex}에서와 같은
$f_1, \ldots, f_r$에 대한 코줄 복합체라 하자. 그러면 $K$는 코호몰로지가
$T$에 지지되는 완전 대상이며, 따라서 $D_{\QCoh, T}(\mathcal{O}_X)$의
완전 대상이다. 한편 $E \in D_{\QCoh, T}(\mathcal{O}_X)$이고 모든 $n$에
대해 $\Hom(K, E[n]) = 0$이면, Lemma
\ref{perfect-lemma-orthogonal-koszul-complex}에 의해
$E = Rj_*(E|_{X \setminus T}) = 0$이다. 따라서 $K$는
$D_{\QCoh, T}(\mathcal{O}_X)$를 생성한다(Derived Categories, Definition
\ref{derived-definition-generators}의 삼각 범주의 생성자 정의에 의한다).

\medskip\noindent
$X = U \cup V$가 열린 덮개이고, $V$는 아핀이며, $U$는 준콤팩트이고,
$U$에 대해 보조정리가 성립한다고 가정하자. $P$를 $T \cap U$에 지지되고
$D_{\QCoh, T \cap U}(\mathcal{O}_U)$의 생성자인 $D(\mathcal{O}_U)$의 완전
대상이라 하자. Lemma \ref{perfect-lemma-lift-map-from-perfect-complex-with-support}을
사용하면, $T$에 지지되고 그 $U$로의 제한이 $P$를 한 직합인자로 갖는 직합인
$D(\mathcal{O}_X)$의 완전 대상 $Q$를 선택할 수 있다. $V = \Spec(B)$로 쓰고
$Z = X \setminus U$라 하자. 그러면 $Z$는 $V \setminus Z$가 준콤팩트인
$V$의 닫힌 부분집합이다. $X$가 준분리이므로 $Z \cap T$는
$W = V \setminus (Z \cap T)$가 준콤팩트인 $V$의 닫힌 부분집합이다. 따라서
$Z \cap T = V(g_1, \ldots, g_r)$가 되도록 $g_1, \ldots, g_s \in B$를
선택할 수 있다. $K \in D(\mathcal{O}_V)$를 $B$ 위의
$g_1, \ldots, g_s$에 대한 코줄 복합체에 대응하는 완전 대상이라 하자.
$K$가 닫힌집합 $(Z \cap T) \subset V$에 지지되므로, 직상
$K' = R(V \to X)_*K$는 그 $V$로의 제한이 $K$인 $D(\mathcal{O}_X)$의
완전 대상임에 유의하라(Cohomology, Lemma
\ref{cohomology-lemma-pushforward-perfect}). $Q \oplus K'$가
$D_{\QCoh, T}(\mathcal{O}_X)$의 생성자라고 주장한다.

\medskip\noindent
$Q \oplus K'$의 어떤 이동에서도 $E$로 가는 영이 아닌 사상이 없도록
$E$를 $D_{\QCoh, T}(\mathcal{O}_X)$의 대상이라 하자. Cohomology, Lemma
\ref{cohomology-lemma-pushforward-restriction}에 의해
$K' =  R(V \to X)_! K$이고, 따라서
$$
\Hom_{D(\mathcal{O}_X)}(K'[n], E) = \Hom_{D(\mathcal{O}_V)}(K[n], E|_V)
$$
이다. 그러므로 Lemma \ref{perfect-lemma-orthogonal-koszul-complex}에 의해
$E|_V = Rj_*E|_W$이다. 여기서 $j : W \to V$는 포함 사상이다. 그림은
$$
\xymatrix{
W \ar[r]_j & V & Z \cap T \ar[l] \ar[d] \\
U \cap V \ar[u]^{j'} \ar[ru]_{j''} & & Z \ar[lu]
}
$$
이다. $E$가 $T$에 지지되므로 $E|_W$는 $W$에서 닫힌
$T \cap W = T \cap U \cap V$에 지지됨을 알 수 있다. 따라서
$$
E|_V = Rj_*(E|_W) = Rj_*(Rj'_*(E|_{U \cap V})) = Rj''_*(E|_{U \cap V})
$$
이다. 여기서 두 번째 등식은 Cohomology, Lemma
\ref{cohomology-lemma-pushforward-restriction}의 (1)이다. 이는
$E = R(U \to X)_*E|_U$를 함의한다(작은 세부사항은 생략한다). 이 경우
$$
\Hom_{D(\mathcal{O}_X)}(Q[n], E) = \Hom_{D(\mathcal{O}_U)}(Q|_U[n], E|_U)
$$
는 $\Hom_{D(\mathcal{O}_U)}(P[n], E|_U)$를 직합인자로 포함한다. 따라서
$P$의 선택에 의해 이 군들이 소멸하면 $E|_U$는 0이다. 그러므로 $E$는 0이다.
\end{proof}



\section{생성자의 한 예}
\label{perfect-section-example-generator}

\noindent
이 절에서는 환 위의 사영공간의 유도 범주가 벡터 다발, 실제로는 구조층의
이동들의 직합으로 생성됨을 증명한다.

\medskip\noindent
다음 보조정리는 $\mathcal{L}$이 풍부하면
$\bigoplus_{n \geq 0} \mathcal{L}^{\otimes -n}$이 생성자임을 말한다.

\begin{lemma}
\label{perfect-lemma-nonzero-some-cohomology}
$X$를 스킴이라 하고 $\mathcal{L}$을 풍부한 가역
$\mathcal{O}_X$-가군이라 하자. $K$가 $D_\QCoh(\mathcal{O}_X)$의 0이 아닌
대상이면, 어떤 $n \geq 0$과 $p \in \mathbf{Z}$에 대해 코호몰로지 군
$H^p(X, K \otimes_{\mathcal{O}_X}^\mathbf{L} \mathcal{L}^{\otimes n})$
은 0이 아니다.
\end{lemma}

\begin{proof}
$X$는 풍부한 가역층을 가지므로 준콤팩트이고 분리임을 상기하자(Properties,
Definition \ref{properties-definition-ample} 및 Lemma
\ref{properties-lemma-affine-s-opens-cover-quasi-separated}). 따라서
Proposition \ref{perfect-proposition-quasi-compact-affine-diagonal}을 적용하여 $K$를
준연접 가군들의 복합체 $\mathcal{F}^\bullet$로 나타낼 수 있다.
다음이 0이 아니도록 하는 임의의 $p$를 선택하자.
$\mathcal{H}^p = \Ker(\mathcal{F}^p \to \mathcal{F}^{p + 1})/
\Im(\mathcal{F}^{p - 1} \to \mathcal{F}^p)$
줄기 $\mathcal{H}^p_x$가 0이 아닌 점 $x \in X$를 선택하자. $X_s$가 $x$의
아핀 열린 이웃이 되도록 $n \geq 0$과
$s \in \Gamma(X, \mathcal{L}^{\otimes n})$을 선택하자. 줄기
$\mathcal{H}^p_x$의 0이 아닌 원소로 가는 $\tau \in \mathcal{H}^p(X_s)$를
선택하자. $\mathcal{H}^p$가 준연접이고 $X_s$가 아핀이므로 이는 가능하다.
$X_s$ 위의 단면을 취하는 것은 준연접 가군에 대한 완전 함자이므로,
$\mathcal{F}^{p + 1}(X_s)$에서 0으로 가고 $\mathcal{H}^p(X_s)$에서 $\tau$로
가는 단면 $\tau' \in \mathcal{F}^p(X_s)$를 찾을 수 있다. Properties,
Lemma \ref{properties-lemma-invert-s-sections}에 의해
$\tau' \otimes s^{\otimes m}$이 단면
$\tau'' \in \Gamma(X, \mathcal{F}^p \otimes \mathcal{L}^{\otimes mn})$의
상이 되도록 하는 $m$이 존재한다. 같은 보조정리를 한 번 더 적용하면,
$\tau'' \otimes s^{\otimes l}$이
$\mathcal{F}^{p + 1} \otimes \mathcal{L}^{\otimes (m + l)n}$에서 0으로
가도록 하는 $l \geq 0$을 찾는다. 그러면 $\tau''$는 원하는 대로
$H^p(X, K \otimes^\mathbf{L}_{\mathcal{O}_X} \mathcal{L}^{(m + l)n})$
의 0이 아닌 류를 준다.
\end{proof}

\begin{lemma}
\label{perfect-lemma-construct-the-next-one}
$A$를 환이라 하고 $X = \mathbf{P}^n_A$라 하자. 모든 $a \in \mathbf{Z}$에
대해 $X$ 위의 벡터 다발들의 완전 복합체
$$
0 \to \mathcal{O}_X(a) \to \ldots
\to \mathcal{O}_X(a + i)^{\oplus {n + 1 \choose i}} \to
\ldots \to \mathcal{O}_X(a + n + 1) \to 0
$$
가 존재한다.
\end{lemma}

\begin{proof}
$\mathbf{P}^n_A$는 $\text{Proj}(A[X_0, \ldots, X_n])$임을 상기하자.
Constructions, Definition \ref{constructions-definition-projective-space}를
보라. $S = A[X_0, \ldots, X_n]$ 위의 $X_0, \ldots, X_n$에 대한 코줄 복합체
$$
K_\bullet = K_\bullet(A[X_0, \ldots, X_n], X_0, \ldots, X_n)
$$
를 생각하자. $X_0, \ldots, X_n$은 다항식환 $S$의 정칙열임이 자명하므로,
코줄 복합체 $K_\bullet$는 차수 $0$을 제외하고 완전이며, 차수 0에서의
코호몰로지는 $S/(X_0, \ldots, X_n)$이다(More on Algebra, Lemma
\ref{more-algebra-lemma-regular-koszul-regular}). $K_i$의 생성자들을 차수
$+i$에 놓으면 $K_\bullet$가 등급 가군의 복합체가 됨에 유의하라. 다시 말해
완전 복합체
$$
0 \to S(-n - 1) \to \ldots \to S(-n - 1 + i)^{\oplus {n \choose i}} \to \ldots
\to S \to S/(X_0, \ldots, X_n) \to 0
$$
를 얻는다. Constructions, Lemma \ref{constructions-lemma-proj-sheaves}의 완전
함자 $\tilde{\ }$를 적용하고 마지막 항이 이 함자의 핵에 있음을 사용하면,
완전 복합체
$$
0 \to \mathcal{O}_X(-n - 1) \to \ldots
\to \mathcal{O}_X(-n - 1 + i)^{\oplus {n + 1 \choose i}} \to
\ldots \to \mathcal{O}_X \to 0
$$
를 얻는다. 가역층 $\mathcal{O}_X(n + a + 1)$로 비틀면 보조정리의 완전
복합체들을 얻는다.
\end{proof}

\begin{lemma}
\label{perfect-lemma-generator-P1}
$A$를 환이라 하고 $X = \mathbf{P}^n_A$라 하자. 그러면
$$
E =
\mathcal{O}_X \oplus \mathcal{O}_X(-1) \oplus \ldots \oplus \mathcal{O}_X(-n)
$$
는 $D_\QCoh(X)$의 생성자이다(Derived Categories, Definition
\ref{derived-definition-generators}).
\end{lemma}

\begin{proof}
$K \in D_\QCoh(\mathcal{O}_X)$라 하자. 모든 $p \in \mathbf{Z}$에 대해
$\Hom(E, K[p]) = 0$이라고 가정하자. $K = 0$임을 보여야 한다. Derived
Categories, Lemma \ref{derived-lemma-right-orthogonal}에 의해 모든
$E' \in \langle E \rangle$와 $p \in \mathbf{Z}$에 대해
$\Hom(E', K[p])$가 0임을 알 수 있다. Lemma
\ref{perfect-lemma-construct-the-next-one}을 $a = -n - 1$로 적용하면
$\mathcal{O}_X(-n - 1) \in \langle E \rangle$임을 알 수 있다. 실제로 이는
각 항이 $E$의 직합인자들의 유한 직합인 유한 복합체와 준동형이다.
$a = -n - 2$로 논증을 반복하면
$\mathcal{O}_X(-n - 2) \in \langle E \rangle$임을 알 수 있다. 귀납법으로
모든 $m \geq 0$에 대해 $\mathcal{O}_X(-m) \in \langle E \rangle$임을 얻는다.
이제
$$
\Hom(\mathcal{O}_X(-m), K[p]) =
H^p(X, K \otimes_{\mathcal{O}_X}^\mathbf{L} \mathcal{O}_X(m)) =
H^p(X, K \otimes_{\mathcal{O}_X}^\mathbf{L} \mathcal{O}_X(1)^{\otimes m})
$$
이므로 Lemma \ref{perfect-lemma-nonzero-some-cohomology}에 의해 $K = 0$임을
결론짓는다. (여기에는 $\mathcal{O}_X(1)$이 $X$ 위의 풍부한 가역층이라는
사실도 사용하였다. 이는 Properties, Lemma
\ref{properties-lemma-open-in-proj-ample}에서 따른다.)
\end{proof}

\begin{remark}
\label{perfect-remark-pullback-generator}
$f : X \to Y$를 준콤팩트이고 준분리인 스킴들의 사상이라 하자.
$E \in D_\QCoh(\mathcal{O}_Y)$를 생성자라 하자(Theorem
\ref{perfect-theorem-bondal-van-den-Bergh}를 보라). 그러면 다음은 동치이다.
\begin{enumerate}
\item $K \in D_\QCoh(\mathcal{O}_X)$에 대해 $Rf_*K = 0$일 필요충분조건은
$K = 0$인 것이다.
\item $Rf_* : D_\QCoh(\mathcal{O}_X) \to D_\QCoh(\mathcal{O}_Y)$는
동형사상을 반영한다.
\item $Lf^*E$는 $D_\QCoh(\mathcal{O}_X)$의 생성자이다.
\end{enumerate}
(1)과 (2)의 동치는
$Rf_* : D_\QCoh(\mathcal{O}_X) \to D_\QCoh(\mathcal{O}_Y)$가 삼각 범주의
완전 함자라는 사실의 형식적 귀결이다. 마찬가지로 (1)과 (3)의 동치는
$Lf^*$가 $Rf_*$의 왼쪽 수반이라는 사실에서 형식적으로 따른다. $f$가
아핀이거나(Lemma \ref{perfect-lemma-affine-morphism}), $f$가 열린 매장이거나,
$f$가 그러한 사상들의 합성이면 이 조건들이 성립한다. 따라서 다음을 얻는다.
\begin{enumerate}
\item $X$가 준아핀 스킴이면 $\mathcal{O}_X$는
$D_\QCoh(\mathcal{O}_X)$의 생성자이다.
\item $X \subset \mathbf{P}^n_A$가 준콤팩트 국소 닫힌 부분스킴이면,
$\mathcal{O}_X \oplus \mathcal{O}_X(-1) \oplus \ldots \oplus \mathcal{O}_X(-n)$
는 Lemma \ref{perfect-lemma-generator-P1}에 의해 $D_\QCoh(\mathcal{O}_X)$의
생성자이다.
\end{enumerate}
\end{remark}




\section{콤팩트 대상과 완전 대상}
\label{perfect-section-compact}

\noindent
$X$를 유한 차원의 Noether 스킴이라 하자. Cohomology, Proposition
\ref{cohomology-proposition-vanishing-Noetherian}과 Cohomology on Sites,
Lemma \ref{sites-cohomology-lemma-when-jshriek-compact}에 의해, 모든 열린집합
$U \subset X$에 대해 가군의 층 $j_!\mathcal{O}_U$는
$D(\mathcal{O}_X)$의 콤팩트 대상이다. 이 층들은 보통 준연접이 아니므로 유도
범주 $D(\mathcal{O}_X)$의 완전 대상을 주지 않는다. 그러나 준연접 코호몰로지
층들을 갖는 복합체로 제한하면 이러한 일은 일어나지 않는다. 정확한 명제는
다음과 같다.

\begin{proposition}
\label{perfect-proposition-compact-is-perfect}
$X$를 준콤팩트이고 준분리인 스킴이라 하자. $D_\QCoh(\mathcal{O}_X)$의
대상이 콤팩트일 필요충분조건은 그것이 완전인 것이다.
\end{proposition}

\begin{proof}
$K$가 쌍대 $K^\vee$를 갖는 $D(\mathcal{O}_X)$의 완전 대상이면
(Cohomology, Lemma \ref{cohomology-lemma-dual-perfect-complex}), $M$에 대해
함자적인 등식
$$
\Hom_{D(\mathcal{O}_X)}(K, M) =
H^0(X, K^\vee \otimes_{\mathcal{O}_X}^\mathbf{L} M)
$$
을 갖는다. $K^\vee \otimes_{\mathcal{O}_X}^\mathbf{L} -$는 직합과 가환하고,
Lemma \ref{perfect-lemma-quasi-coherence-pushforward-direct-sums}에 의해 $H^0(X, -)$는
$D_\QCoh(\mathcal{O}_X)$에서 직합과 가환하므로, $K$는
$D_\QCoh(\mathcal{O}_X)$에서 콤팩트이다.

\medskip\noindent
역으로 $K$를 $D_\QCoh(\mathcal{O}_X)$의 콤팩트 대상이라 하자. $K$가
완전임을 보이려면 모든 아핀 열린집합 $U \subset X$에 대해 $K|_U$가
완전임을 보이면 충분하다. Cohomology, Lemma
\ref{cohomology-lemma-perfect-independent-representative}를 보라.
$j : U \to X$가 준콤팩트 분리 사상임에 유의하라. 따라서
$Rj_* : D_\QCoh(\mathcal{O}_U) \to D_\QCoh(\mathcal{O}_X)$
는 직합과 가환한다. Lemma \ref{perfect-lemma-quasi-coherence-pushforward-direct-sums}를
보라. 그러므로 $U$로의 제한과 $Rj_*$의 수반성으로부터 $K|_U$가
$D_\QCoh(\mathcal{O}_U)$의 콤팩트 대상임을 얻는다. 따라서 $X$가 아핀인
경우로 환원된다.

\medskip\noindent
$X = \Spec(A)$가 아핀이라고 가정하자. Lemma
\ref{perfect-lemma-affine-compare-bounded}에 의해 문제는 $D(A)$에 대한 같은 문제로
옮겨진다. $D(A)$에 대해서는 More on Algebra, Proposition
\ref{more-algebra-proposition-perfect-is-compact}이 그 결과이다.
\end{proof}

\begin{remark}
\label{perfect-remark-classical-generator}
$X$를 준콤팩트이고 준분리인 스킴이라 하자. $G$를
$D_\QCoh(\mathcal{O}_X)$의 생성자인 $D(\mathcal{O}_X)$의 완전 대상이라
하자. Theorem \ref{perfect-theorem-bondal-van-den-Bergh}에 의해 그러한 대상은 적어도
하나 존재한다. Lemma \ref{perfect-lemma-quasi-coherence-direct-sums}, Proposition
\ref{perfect-proposition-compact-is-perfect}, 그리고 Derived Categories,
Proposition \ref{derived-proposition-generator-versus-classical-generator}을
결합하면 $G$가 $D_{perf}(\mathcal{O}_X)$의 고전적 생성자임을 알 수 있다.
\end{remark}

\noindent
다음 결과는 Proposition \ref{perfect-proposition-compact-is-perfect}의 강화이다.
$T \subset X$를 스킴 $X$의 닫힌 부분집합이라 하자. 이전과 같이
$D_T(\mathcal{O}_X)$는 $T$에 지지된 대상들로 이루어진
$D(\mathcal{O}_X)$의 엄밀 충만 포화 삼각 부분범주를 나타낸다(Definition
\ref{perfect-definition-supported-on}). 직합을 취하는 것은 코호몰로지 층을 취하는
것과 가환하므로, $D_T(\mathcal{O}_X)$에는 직합이 있으며 이는
$D(\mathcal{O}_X)$에서의 직합과 같다.

\begin{lemma}
\label{perfect-lemma-compact-is-perfect-with-support}
$X$를 준콤팩트이고 준분리인 스킴이라 하자. $T \subset X$를
$X \setminus T$가 준콤팩트인 닫힌 부분집합이라 하자.
$D_{\QCoh, T}(\mathcal{O}_X)$의 대상이 콤팩트일 필요충분조건은 그것이
$D(\mathcal{O}_X)$의 대상으로서 완전인 것이다.
\end{lemma}

\begin{proof}
보조정리 앞의 설명에 의해 $D_{\QCoh, T}(\mathcal{O}_X)$는 직합을 갖는 삼각
범주이다. Proposition \ref{perfect-proposition-compact-is-perfect}에 의해 완전
대상들은 $D(\mathcal{O}_X)$의 콤팩트 대상들을 정의하고, 따라서 더구나
직합을 취하는 연산으로 보존되는 임의의 부분범주의 콤팩트 대상들을 정의한다.
역을 위해 $\mathcal{O}_X$-가군의 완전 복합체인 생성자
$E \in D_{\QCoh, T}(\mathcal{O}_X)$가 존재한다는 것을 사용하겠다. Lemma
\ref{perfect-lemma-generator-with-support}를 보라. 따라서 위에 의해 $E$는
콤팩트이다. 그러면 Derived Categories, Proposition
\ref{derived-proposition-generator-versus-classical-generator}에 의해 $E$는
$D_{\QCoh, T}(\mathcal{O}_X)$의 콤팩트 대상들로 이루어진 충만 부분범주의
고전적 생성자이다. 따라서 임의의 콤팩트 대상은 (a) 이동을 취하기, (b) 유한
직합을 취하기, (c) 사상뿔을 취하기, (d) 직합인자를 취하기로 이루어진 유한한
연산열을 통해 $E$로부터 구성할 수 있다. 이 연산들은 각각 완전성이라는
성질을 보존하므로 결과가 따른다.
\end{proof}

\begin{remark}
\label{perfect-remark-classical-generator-with-support}
$X$를 준콤팩트이고 준분리인 스킴이라 하자. $T \subset X$를
$X \setminus T$가 준콤팩트인 닫힌 부분집합이라 하자. $G$를
$D_{\QCoh, T}(\mathcal{O}_X)$의 완전 대상이며
$D_{\QCoh, T}(\mathcal{O}_X)$의 생성자인 대상이라 하자.
Lemma \ref{perfect-lemma-generator-with-support}에 의해 그러한 대상은 적어도 하나
존재한다. $D_{\QCoh, T}(\mathcal{O}_X)$가 직합을 갖는다는 사실을 Lemma
\ref{perfect-lemma-compact-is-perfect-with-support} 및 Derived Categories,
Proposition \ref{derived-proposition-generator-versus-classical-generator}과
결합하면 $G$가 $D_{perf, T}(\mathcal{O}_X)$의 고전적 생성자임을 알 수 있다.
\end{remark}

\noindent
다음 보조정리는 앞의 보조정리의 증명에 들어간 아이디어의 응용이다.

\begin{lemma}
\label{perfect-lemma-map-from-pseudo-coherent-to-complex-with-support}
$X$를 준콤팩트이고 준분리인 스킴이라 하자. $T \subset X$를
$U = X \setminus T$가 준콤팩트인 닫힌 부분집합이라 하자.
$\alpha : P \to E$를 다음 중 하나를 만족하는 $D_\QCoh(\mathcal{O}_X)$의
사상이라 하자.
\begin{enumerate}
\item $P$는 완전이고 $E$는 $T$에 지지된다.
\item $P$는 유사연접이고, $E$는 $T$에 지지되며, $E$는 아래로 유계이다.
\end{enumerate}
그러면 $\mathcal{O}_X$-가군의 완전 복합체 $I$와 사상
$I \to \mathcal{O}_X[0]$이 존재하여 $I \otimes^\mathbf{L} P \to E$는
0이고 $I|_U \to \mathcal{O}_U[0]$은 동형이다.
\end{lemma}

\begin{proof}
$\mathcal{D} = D_{\QCoh, T}(\mathcal{O}_X)$로 놓자. 두 경우 모두 복합체
$K = R\SheafHom(P, E)$는 $\mathcal{D}$의 대상이다. 준연접성에 대해서는
Lemma \ref{perfect-lemma-quasi-coherence-internal-hom}을 보라. $R\SheafHom$의 구성은
열린집합으로의 제한과 가환하므로 $K$가 $T$에 지지됨은 자명하다. 사상
$\alpha$는 $H^0(K) = \Hom_{D(\mathcal{O}_X)}(\mathcal{O}_X[0], K)$의 원소를
정의한다. 따라서 사상 $\alpha : \mathcal{O}_X[0] \to K$에 대해 결과를
증명하면 충분하다.

\medskip\noindent
$E \in \mathcal{D}$를 완전 생성자라 하자. Lemma
\ref{perfect-lemma-generator-with-support}를 보라. Derived Categories, Lemma
\ref{derived-lemma-write-as-colimit}에서처럼 생성자 $E$를 사용하여
$$
K = \text{hocolim} K_n
$$
로 쓰자. 함자 $\mathcal{D} \to D(\mathcal{O}_X)$는 직합과 가환하므로
$K = \text{hocolim} K_n$은 $D(\mathcal{O}_X)$에서도 성립한다.
$\mathcal{O}_X$는 $D(\mathcal{O}_X)$의 콤팩트 대상이므로, $\alpha$를
유도하는 $n$과 사상 $\alpha_n : \mathcal{O}_X \to K_n$을 찾는다. Derived
Categories, Lemma \ref{derived-lemma-commutes-with-countable-sums}를 보라.
주변 범주 $D(\mathcal{O}_X)$의 사상 $\mathcal{O}_X[0] \to K_n$에 Derived
Categories, Lemma \ref{derived-lemma-factor-through}을 적용하면,
$\alpha_n$은 $\mathcal{O}_X[0] \to Q \to K_n$으로 분해된다. 여기서 $Q$는
$\langle E \rangle$의 대상이다. 따라서 $Q$는 $T$에 지지된 완전 복합체이다.

\medskip\noindent
완전 삼각형을 하나 선택하자.
$$
I \to \mathcal{O}_X[0] \to Q \to I[1]
$$
구성에 의해 $I$는 완전이고, 사상 $I \to \mathcal{O}_X[0]$은 $U$ 위로
제한하면 동형이며, $\alpha$가 $Q$를 통해 분해되므로 합성 $I \to K$는
0이다. 이것으로 보조정리가 증명되었다.
\end{proof}






\section{가군 범주로서의 유도 범주}
\label{perfect-section-derived-is-dga}

\noindent
이 절에서는 앞서 얻은 결과들의 몇 가지 귀결을 도출한다. 그러기 전에 보조정리
두 개가 더 필요하다.

\begin{lemma}
\label{perfect-lemma-tensor-with-QCoh-complex}
$X$를 스킴이라 하자. $K^\bullet$를 코호몰로지 층들이 준연접인
$\mathcal{O}_X$-가군의 복합체라 하자. 자기준동형 미분 등급 대수를
$(E, d) = \Hom_{\text{Comp}^{dg}(\mathcal{O}_X)}(K^\bullet, K^\bullet)$
로 놓자. 그러면 Differential Graded Algebra, Lemma
\ref{dga-lemma-tensor-with-complex-derived}의 함자
$$
- \otimes_E^\mathbf{L} K^\bullet :
D(E, \text{d}) \longrightarrow D(\mathcal{O}_X)
$$
의 상은 $D_\QCoh(\mathcal{O}_X)$에 포함된다.
\end{lemma}

\begin{proof}
$P$를 성질 (P)를 갖는 미분 등급 $E$-가군이라 하고, $F_\bullet$를
Differential Graded Algebra, Section \ref{dga-section-P-resolutions}에서와
같은 $P$의 여과라 하자. 그러면
$$
P \otimes_E K^\bullet = \text{hocolim}\ F_iP \otimes_E K^\bullet
$$
이다. 각 $F_iP$는 등급 조각들이 $E[k]$의 직합인 유한 여과를 갖는다.
결과는 쉽게 따른다.
\end{proof}

\noindent
다음 결과는 \cite{BvdB}에서 가져온 것이다.

\begin{theorem}
\label{perfect-theorem-DQCoh-is-Ddga}
$X$를 준콤팩트이고 준분리인 스킴이라 하자. 0이 아닌 코호몰로지 군
$H^i(E)$가 유한개뿐인 미분 등급 대수 $(E, \text{d})$가 존재하여
$D_\QCoh(\mathcal{O}_X)$는 $D(E, \text{d})$와 동치이다.
\end{theorem}

\begin{proof}
$K^\bullet$를 완전이고 $D_\QCoh(\mathcal{O}_X)$를 생성하는
$\mathcal{O}$-가군의 K-단사 복합체라 하자. Theorem
\ref{perfect-theorem-bondal-van-den-Bergh}과 K-단사 분해의 존재에 의해 그러한 대상이
존재한다. 다음으로 놓으면 정리가 성립함을 보이겠다.
$$
(E, \text{d}) = \Hom_{\text{Comp}^{dg}(\mathcal{O}_X)}(K^\bullet, K^\bullet)
$$
여기서 $\text{Comp}^{dg}(\mathcal{O}_X)$는 $\mathcal{O}$-가군의 복합체들의
미분 등급 범주이다. Differential Graded Algebra, Section
\ref{dga-section-variant-base-change}를 보라. $K^\bullet$가 K-단사이므로 모든
$n \in \mathbf{Z}$에 대해
\begin{equation}
\label{perfect-equation-E-is-OK}
H^n(E) = \Ext^n_{D(\mathcal{O}_X)}(K^\bullet, K^\bullet)
\end{equation}
이다. Lemma \ref{perfect-lemma-ext-from-perfect-into-bounded-QCoh}에 의해 이 Ext들 중
유한개만 0이 아니다. Differential Graded Algebra, Lemma
\ref{dga-lemma-tensor-with-complex-derived}의 함자
$$
- \otimes_E^\mathbf{L} K^\bullet :
D(E, \text{d}) \longrightarrow D(\mathcal{O}_X)
$$
를 생각하자. $K^\bullet$가 완전이므로 이는 $D(\mathcal{O}_X)$의 콤팩트
대상을 정의한다. Proposition \ref{perfect-proposition-compact-is-perfect}를 보라.
(\ref{perfect-equation-E-is-OK})와 결합하면, Differential Graded Algebra, Lemmas
\ref{dga-lemma-fully-faithful-in-compact-case}에서 따르듯이 위 함자는
충실충만하다. Differential Graded Algebra, Lemmas
\ref{dga-lemma-tensor-with-complex-hom-adjoint}에 의해 이는 오른쪽 수반
$$
R\Hom(K^\bullet, - ) : D(\mathcal{O}_X) \longrightarrow D(E, \text{d})
$$
을 갖는데, 수반 함자에 대한 일반론에 의해 이는 왼쪽 준역함자이다. 한편
Lemma \ref{perfect-lemma-tensor-with-QCoh-complex}에서 다음을 얻는다.
$$
- \otimes_E^\mathbf{L} K^\bullet :
D(E, \text{d}) \longrightarrow D_\QCoh(\mathcal{O}_X)
$$
또한 $K^\bullet$를 $D_\QCoh(\mathcal{O}_X)$의 생성자로 선택했으므로, 수반을
$D_\QCoh(\mathcal{O}_X)$로 제한한 것의 핵은 0이다. 형식적 논증으로 원하는
동치를 얻는다. Derived Categories, Lemma
\ref{derived-lemma-fully-faithful-adjoint-kernel-zero}를 보라.
\end{proof}

\begin{remark}[지지를 갖는 변형]
\label{perfect-remark-DQCoh-is-Ddga-with-support}
$X$를 준콤팩트이고 준분리인 스킴이라 하자. $T \subset X$를
$X \setminus T$가 준콤팩트인 닫힌 부분집합이라 하자. Theorem
\ref{perfect-theorem-DQCoh-is-Ddga}의 유사 명제가 $D_{\QCoh, T}(\mathcal{O}_X)$에
대해 성립한다. 이는 Lemmas \ref{perfect-lemma-generator-with-support} 및
\ref{perfect-lemma-compact-is-perfect-with-support}, 그리고 지지를 갖는 Lemma
\ref{perfect-lemma-tensor-with-QCoh-complex}의 변형을 사용하여 정리의 증명과 정확히
같은 논증으로부터 따른다. 이것이 필요해지면 여기에 결과를 정확히 진술하고
상세한 증명을 제시할 것이다.
\end{remark}

\begin{remark}[dga의 유일성]
\label{perfect-remark-independence-choice}
$X$를 환 $R$ 위의 준콤팩트이고 준분리인 스킴이라 하자. Theorem
\ref{perfect-theorem-DQCoh-is-Ddga}의 증명에 있는 구성에 의해, 삼각 범주로서
$D_\QCoh(X)$가 $D(A, \text{d})$와 $R$-선형 동치가 되도록 하는 $R$ 위의
미분 등급 대수 $(A, \text{d})$가 존재한다. $(A, \text{d})$는 얼마나
유일한가라는 질문을 할 수 있다. 답은 $(A, \text{d})$가 유도 동치까지
잘 정의된다고만 말하는 것보다 (겨우) 조금 더 좋다. 실제로
$(B, \text{d})$가 두 번째 그러한 쌍이라고 하자. 그러면
$$
(A, \text{d}) = \Hom_{\text{Comp}^{dg}(\mathcal{O}_X)}(K^\bullet, K^\bullet)
$$
이고
$$
(B, \text{d}) = \Hom_{\text{Comp}^{dg}(\mathcal{O}_X)}(L^\bullet, L^\bullet)
$$
이다. 여기서 $K^\bullet$와 $L^\bullet$는 $D_\QCoh(\mathcal{O}_X)$의 완전
생성자에 대응하는 $\mathcal{O}_X$-가군의 어떤 K-단사 복합체들이다. 다음으로
놓자.
$$
\Omega = \Hom_{\text{Comp}^{dg}(\mathcal{O}_X)}(K^\bullet, L^\bullet)
\quad
\Omega' = \Hom_{\text{Comp}^{dg}(\mathcal{O}_X)}(L^\bullet, K^\bullet)
$$
그러면 $\Omega$는 미분 등급 $B^{opp} \otimes_R A$-가군이고, $\Omega'$은
미분 등급 $A^{opp} \otimes_R B$-가군이다. 더욱이 동치
$$
D(A, \text{d}) \to D_\QCoh(\mathcal{O}_X) \to
D(B, \text{d})
$$
는 함자 $- \otimes_A^\mathbf{L} \Omega'$으로 주어지고 준역함자도 마찬가지이다.
따라서 우리는 Differential Graded Algebra, Remark
\ref{dga-remark-hochschild-cohomology}의 상황에 놓인다. 이 주석이
필요해지면 여기에 정확한 명제와 상세한 증명을 제시할 것이다.
\end{remark}






\section{유사연접 복합체의 특징화 I}
\label{perfect-section-pseudo-coherent-hocolim}

\noindent
위의 방법들을 사용하면 유사연접 대상을 완전 대상에 의한 근사들의 유도
호모토피 극한으로 특징지을 수 있다.

\begin{lemma}
\label{perfect-lemma-pseudo-coherent-hocolim}
$X$를 준콤팩트이고 준분리인 스킴이라 하자.
$K \in D(\mathcal{O}_X)$라 하자. 다음은 동치이다.
\begin{enumerate}
\item $K$는 유사연접이다.
\item $K = \text{hocolim} K_n$이고, 여기서 $K_n$은 완전이며 모든 $n$에
대해 $\tau_{\geq -n}K_n \to \tau_{\geq -n}K$는 동형이다.
\end{enumerate}
\end{lemma}

\begin{proof}
함의 (2) $\Rightarrow$ (1)은 임의의 환 달린 공간에서 참이다. 실제로 (2)가
성립한다고 하자. 유도 범주의 완전 대상은 유사연접임을 상기하자. Cohomology,
Lemma \ref{cohomology-lemma-perfect}를 보라. 그러면 정의들로부터
$\tau_{\geq -n}K_n$은 $(-n + 1)$-유사연접이고, 따라서
$\tau_{\geq -n}K$는 $(-n + 1)$-유사연접이며, 따라서 $K$도
$(-n + 1)$-유사연접이다. 이는 모든 $n$에 대해 참이므로 $K$는 유사연접이다.
Cohomology, Definition \ref{cohomology-definition-pseudo-coherent}을 보라.

\medskip\noindent
(1)을 가정하자. 먼저 완전 복합체 $K_1$에 의한 $(X, K, -2)$의 근사
$K_1 \to K$를 선택한다. Definitions \ref{perfect-definition-approximation-holds} 및
\ref{perfect-definition-approximation}, 그리고 Theorem \ref{perfect-theorem-approximation}을
보라. 귀납적으로
$$
K_1 \to K_2 \to \ldots \to K_n \to K
$$
을 얻었다고 하자. 여기서 $K_i$는 완전이고 모든 $1 \leq i \leq n$에 대해
$\tau_{\geq -i}K_i \to \tau_{\geq -i}K$는 동형이다. 완전 대상 $K_n$에
대해 Lemma \ref{perfect-lemma-ext-from-perfect-into-bounded-QCoh}에서와 같은
$a \leq b$를 선택하자. $(X, K, \min(a - 1, -n - 1))$의 근사
$K_{n + 1} \to K$를 선택하고, 완전 삼각형을 하나 선택하자.
$$
K_{n + 1} \to K \to C \to K_{n + 1}[1]
$$
그러면 $C \in D_\QCoh(\mathcal{O}_X)$는 $i \geq a$에 대해
$H^i(C) = 0$이다. 따라서 $a, b$의 선택에 의해
$\Hom_{D(\mathcal{O}_X)}(K_n, C) = 0$임을 알 수 있다. 그러므로 합성
$K_n \to K \to C$는 0이다. 따라서 Derived Categories, Lemma
\ref{derived-lemma-representable-homological}에 의해 $K_n \to K$를
$K_{n + 1}$을 통해 분해할 수 있다. 이로써 귀납 단계가 증명된다.

\medskip\noindent
이제 $K = \text{hocolim} K_n$임을 증명해야 한다. 이는 함자
$H^i( - ) : D(\mathcal{O}_X) \to \textit{Mod}(\mathcal{O}_X)$와 우리의
$K_n$ 선택에 Derived Categories, Lemma
\ref{derived-lemma-cohomology-of-hocolim}을 적용하면 따른다.
\end{proof}

\begin{lemma}
\label{perfect-lemma-pseudo-coherent-hocolim-with-support}
$X$를 준콤팩트이고 준분리인 스킴이라 하자. $T \subset X$를
$X \setminus T$가 준콤팩트인 닫힌 부분집합이라 하자. $K \in D(\mathcal{O}_X)$를
$T$에 지지되는 대상이라 하자. 다음은 동치이다.
\begin{enumerate}
\item $K$는 유사연접이다.
\item $K = \text{hocolim} K_n$이고, 여기서 $K_n$은 완전이고 $T$에 지지되며,
모든 $n$에 대해 $\tau_{\geq -n}K_n \to \tau_{\geq -n}K$는 동형이다.
\end{enumerate}
\end{lemma}

\begin{proof}
근사의 선택에서 삼중항 $(T, K, m)$을 사용한다는 점을 제외하면 이 보조정리의
증명은 Lemma \ref{perfect-lemma-pseudo-coherent-hocolim}의 증명과 정확히 같다.
\end{proof}









\section{동치의 한 예}
\label{perfect-section-example-equivalence}

\noindent
Section \ref{perfect-section-example-generator}에서 환 $A$ 위의 사영공간
$\mathbf{P}^n_A$의 유도 범주가 벡터 다발, 실제로는 구조층의 이동들의
직합으로 생성됨을 증명하였다. 이 절에서는 이것이
$D_\QCoh(\mathcal{O}_{\mathbf{P}^n_A})$와 어떤 $A$-대수의 유도 범주 사이의
동치를 결정함을 증명한다.

\medskip\noindent
결과를 진술하기 전에 몇 가지 표기법이 필요하다. $A$를 환이라 하고
$X = \mathbf{P}^n_A = \text{Proj}(S)$라 하자. 여기서
$S = A[X_0, \ldots, X_n]$이다. Lemma \ref{perfect-lemma-generator-P1}에 의해
\begin{equation}
\label{perfect-equation-generator-Pn}
P =
\mathcal{O}_X \oplus \mathcal{O}_X(-1) \oplus \ldots \oplus \mathcal{O}_X(-n)
\end{equation}
가 $D_\QCoh(\mathcal{O}_X)$의 완전 생성자임을 안다. 자명한 곱셈과 덧셈을
갖는 (비가환) $A$-대수
\begin{equation}
\label{perfect-equation-algebra-for-Pn}
R = \Hom_{\mathcal{O}_X}(P, P) =
\left(
\begin{matrix}
S_0 & S_1 & S_2 & \ldots & \ldots \\
0 & S_0 & S_1 & \ldots & \ldots\\
0 & 0 & S_0 & \ldots & \ldots \\
\ldots & \ldots & \ldots & \ldots & \ldots \\
0 & \ldots & \ldots & \ldots & S_0
\end{matrix}
\right)
\end{equation}
를 생각하자. $P$를 보통 방식으로(즉 $P$를 차수 $0$에 놓고 다른 모든
차수에는 0을 놓아) $\mathcal{O}_X$-가군의 복합체로 보면
$$
R = \Hom_{\text{Comp}^{dg}(\mathcal{O}_X)}(P, P)
$$
이다. 여기서 오른쪽에서는 $R$을 영 미분을 갖는 $A$ 위의 미분 등급
대수로(즉 $R$을 차수 $0$에 놓고 다른 모든 차수에는 0을 놓아) 본다.
Differential Graded Algebra, Section \ref{dga-section-variant-base-change}의
논의에 따라 유도 함자
$$
- \otimes_R^\mathbf{L} P : D(R) \longrightarrow D(\mathcal{O}_X),
$$
를 얻는다. 특히 Differential Graded Algebra, Lemma
\ref{dga-lemma-tensor-with-complex-derived}를 보라. Lemma
\ref{perfect-lemma-tensor-with-QCoh-complex}에 의해 이 함자의 본질적 상은
$D_\QCoh(\mathcal{O}_X)$에 포함된다.

\begin{lemma}
\label{perfect-lemma-Pn-module-category}
\begin{reference}
\cite{Beilinson}
\end{reference}
$A$를 환이라 하고 $X = \mathbf{P}^n_A = \text{Proj}(S)$라 하자. 여기서
$S = A[X_0, \ldots, X_n]$이다. $P$를 (\ref{perfect-equation-generator-Pn})에서처럼,
$R$을 (\ref{perfect-equation-algebra-for-Pn})에서처럼 두면, 함자
$$
- \otimes_R^\mathbf{L} P : D(R) \longrightarrow D_\QCoh(\mathcal{O}_X)
$$
는 $R$을 $P$로 보내는 삼각 범주들의 $A$-선형 동치이다.
\end{lemma}

\noindent
말로 표현하면, 사영공간 위의 준연접 가군의 유도 범주는 어떤 (비가환)
대수 위의 가군의 유도 범주와 동치이다. 사영공간의 이 성질은 $A$ 위의 모든
사영 스킴 가운데 상당히 이례적인 것으로 보인다.

\begin{proof}
우리의 함자가 충실충만함을 증명하려면 $i \not = 0$에 대해
$\Ext^i_X(P, P)$가 0이고 $i = 0$에 대해 $R$과 같음을 증명하면 충분하다.
Differential Graded Algebra, Lemma
\ref{dga-lemma-fully-faithful-in-compact-case}를 보라. Lemma
\ref{perfect-lemma-ext-from-perfect-into-bounded-QCoh}의 증명에서와 같이
$$
\Ext^i_X(P, P) = H^i(X, P^\wedge \otimes P) =
\bigoplus\nolimits_{0 \leq a, b \leq n} H^i(X, \mathcal{O}_X(a - b))
$$
임을 알 수 있다. 사영공간의 코호몰로지 계산(Cohomology of Schemes, Lemma
\ref{coherent-lemma-cohomology-projective-space-over-ring})에 의해 이
$\Ext$-군들은 $i = 0$이 아니면 0이다. $i = 0$일 때는
(\ref{perfect-equation-algebra-for-Pn})에서 $R$을 바로 이렇게 정의했으므로 $R$을
다시 얻는다. Differential Graded Algebra, Lemma
\ref{dga-lemma-tensor-with-complex-hom-adjoint}에 의해 우리의 함자는 오른쪽
수반 $R\Hom(P, -) : D_\QCoh(\mathcal{O}_X) \to D(R)$을 갖는다. Lemma
\ref{perfect-lemma-generator-P1}에 의해 $P$가 $D_\QCoh(\mathcal{O}_X)$의
생성자이므로 $R\Hom(P, -)$의 핵은 0이다. 따라서 Derived Categories,
Lemma \ref{derived-lemma-fully-faithful-adjoint-kernel-zero}에 의해 우리의
함자는 삼각 범주들의 동치이다.
\end{proof}






\section{준연접화자의 재검토}
\label{perfect-section-better-coherator}

\noindent
Section \ref{perfect-section-coherator}에서 표준 함자
$D(\QCoh(\mathcal{O}_X)) \to D(\mathcal{O}_X)$의 오른쪽 수반 $RQ_X$를
구성하고 연구하였다. 이는 준연접화자
$Q_X : \textit{Mod}(\mathcal{O}_X) \to \QCoh(\mathcal{O}_X)$의 오른쪽 유도
확장으로 구성되었다. 이 절에서는 포함 함자
$$
D_\QCoh(\mathcal{O}_X) \longrightarrow D(\mathcal{O}_X)
$$
가 언제 오른쪽 수반을 갖는지 연구한다. 이 오른쪽 수반이 존재하면 이를
다음으로 나타내겠다\footnote{이는 아마 비표준 표기법일 것이다. 그러나 이미
준연접화자에 $Q_X$를, 그 유도 확장에 $RQ_X$를 사용하였다.}.
$$
DQ_X :
D(\mathcal{O}_X) \longrightarrow D_\QCoh(\mathcal{O}_X)
$$
준콤팩트이고 준분리인 스킴은 이러한 오른쪽 수반을 가짐을 알게 될 것이다.

\begin{lemma}
\label{perfect-lemma-better-coherator}
$X$를 준콤팩트이고 준분리인 스킴이라 하자. 포함 함자
$D_\QCoh(\mathcal{O}_X) \to D(\mathcal{O}_X)$는 오른쪽 수반 $DQ_X$를 갖는다.
\end{lemma}

\begin{proof}[첫 번째 증명]
Cohomology of Schemes, Lemma \ref{coherent-lemma-induction-principle}에서와
같은 귀납 원리를 사용하여 이를 증명하겠다.
$D(\QCoh(\mathcal{O}_X)) \to D_\QCoh(\mathcal{O}_X)$가 동치이면 보조정리는
참이다. 실제로 Section \ref{perfect-section-coherator}의 함자 $RQ_X$는 함자
$D(\QCoh(\mathcal{O}_X)) \to D(\mathcal{O}_X)$의 오른쪽 수반이다. 특히
우리의 보조정리는 아핀 스킴에 대해 참이다. Lemma
\ref{perfect-lemma-affine-coherator}를 보라. 따라서 다음을 보이면 충분하다.
$X = U \cup V$가 두 준콤팩트 열린집합의 합집합이고 보조정리가 $U$, $V$,
$U \cap V$에 대해 성립하면, 보조정리는 $X$에 대해서도 성립한다.

\medskip\noindent
수반이 존재할 필요충분조건은 $D(\mathcal{O}_X)$의 모든 대상 $K$에 대해
완전 삼각형
$$
E' \to E \to K \to E'[1]
$$
을 $D(\mathcal{O}_X)$에서 찾을 수 있고, 여기서 $E'$은
$D_\QCoh(\mathcal{O}_X)$에 속하며 모든 $D_\QCoh(\mathcal{O}_X)$의 $M$에
대해 $\Hom(M, K) = 0$인 것이다. Derived Categories, Lemma
\ref{derived-lemma-right-adjoint}를 보라. Cohomology, Lemma
\ref{cohomology-lemma-exact-sequence-j-star}의 $D(\mathcal{O}_X)$ 안의 완전
삼각형
$$
E \to Rj_{U, *}E|_U \oplus Rj_{V, *}E|_V \to
Rj_{U \cap V, *}E|_{U \cap V} \to E[1]
$$
을 생각하자. Derived Categories, Lemma \ref{derived-lemma-prepare-adjoint}에
의해 $Rj_{U, *}E|_U$, $Rj_{V, *}E|_V$, 그리고
$Rj_{U \cap V, *}E|_{U \cap V}$에 대해 원하는 완전 삼각형을 구성하면
충분하다. 이는 다음 문단에서 논의하는 명제로 환원한다.

\medskip\noindent
$j : U \to X$를 열린 매장이라 하자. 여기서 $U$는 보조정리가 성립하는
준콤팩트 열린집합이다. $L$을 $D(\mathcal{O}_U)$의 대상이라 하자. 그러면
$D(\mathcal{O}_X)$ 안의 완전 삼각형
$$
E' \to Rj_*L \to K \to E'[1]
$$
이 존재하고, 여기서 $E'$은 $D_\QCoh(\mathcal{O}_X)$에 속하며 모든
$D_\QCoh(\mathcal{O}_X)$의 $M$에 대해 $\Hom(M, K) = 0$이다. 이를 보기 위해
$D(\mathcal{O}_U)$ 안의 완전 삼각형
$$
L' \to L \to Q \to L'[1]
$$
을 선택하자. 여기서 $L'$은 $D_\QCoh(\mathcal{O}_U)$에 속하고, 모든
$D_\QCoh(\mathcal{O}_U)$의 $N$에 대해 $\Hom(N, Q) = 0$이다. Derived
Categories, Lemma \ref{derived-lemma-right-adjoint}의 명제가 필요충분조건이므로
이는 가능하다. $D(\mathcal{O}_X)$ 안의 완전 삼각형
$$
Rj_*L' \to Rj_*L \to Rj_*Q \to Rj_*L'[1]
$$
을 얻는다. Lemma \ref{perfect-lemma-quasi-coherence-direct-image}에 의해 $Rj_*L'$은
$D_\QCoh(\mathcal{O}_X)$에 속함에 유의하라. 한편 $M$이
$D_\QCoh(\mathcal{O}_X)$에 속하면
$$
\Hom(M, Rj_*Q) = \Hom(Lj^*M, Q) = 0
$$
이다. 실제로 Lemma \ref{perfect-lemma-quasi-coherence-pullback}에 의해 $Lj^*M$은
$D_\QCoh(\mathcal{O}_U)$에 속한다. 이것으로 증명이 끝난다.
\end{proof}

\begin{proof}[두 번째 증명]
Derived Categories, Proposition \ref{derived-proposition-brown}에 의해 수반이
존재한다. 가정들은 만족된다. 먼저 $D_\QCoh(\mathcal{O}_X)$는 직합을 갖고
직합은 포함 함자와 가환함에 유의하라(Lemma
\ref{perfect-lemma-quasi-coherence-direct-sums}). 한편
$D_\QCoh(\mathcal{O}_X)$는 완전 생성자를 갖고 Theorem
\ref{perfect-theorem-bondal-van-den-Bergh}, Proposition
\ref{perfect-proposition-compact-is-perfect}에 의해 완전 대상들이 콤팩트이므로
콤팩트하게 생성된다.
\end{proof}

\begin{lemma}
\label{perfect-lemma-pushforward-better-coherator}
$f : X \to Y$를 스킴들의 준콤팩트이고 준분리인 사상이라 하자. $X$와
$Y$에 대해 포함 함자 $D_\QCoh \to D$의 오른쪽 수반 $DQ_X$와 $DQ_Y$가
존재하면
$$
Rf_* \circ DQ_X = DQ_Y \circ Rf_*
$$
이다.
\end{lemma}

\begin{proof}
Lemma \ref{perfect-lemma-quasi-coherence-direct-image}에 의해 $Rf_*$가
$D_\QCoh(\mathcal{O}_X)$를 $D_\QCoh(\mathcal{O}_Y)$ 안으로 보내므로 명제는
의미가 있다. 또한 $Lf^*$도 마찬가지로 $D_\QCoh(\mathcal{O}_Y)$를
$D_\QCoh(\mathcal{O}_X)$ 안으로 보내고(Lemma
\ref{perfect-lemma-quasi-coherence-pullback}), 따라서 $Rf_* \circ DQ_X$와
$DQ_Y \circ Rf_*$가 모두
$Lf^* : D_\QCoh(\mathcal{O}_Y) \to D(\mathcal{O}_X)$의 오른쪽 수반이므로
명제는 참이다.
\end{proof}

\begin{remark}
\label{perfect-remark-explain-consequence}
$X$를 준콤팩트이고 준분리인 스킴이라 하자. $X = U \cup V$라 하고 $U$와
$V$를 준콤팩트 열린집합이라 하자. Lemma \ref{perfect-lemma-better-coherator}에 의해
함자 $DQ_X$, $DQ_U$, $DQ_V$, $DQ_{U \cap V}$가 존재한다. 더욱이 임의의
$K \in D(\mathcal{O}_X)$에 대해 표준 완전 삼각형
$$
DQ_X(K) \to Rj_{U, *}DQ_U(K|_U) \oplus Rj_{V, *}DQ_V(K|_V)
\to Rj_{U \cap V, *}DQ_{U \cap V}(K|_{U \cap V}) \to
$$
이 존재한다. 이는 Cohomology, Lemma
\ref{cohomology-lemma-exact-sequence-j-star}의 완전 삼각형에 완전 함자
$DQ_X$를 적용하고 Lemma \ref{perfect-lemma-pushforward-better-coherator}를 세 번
사용하면 따른다.
\end{remark}

\begin{lemma}
\label{perfect-lemma-boundedness-better-coherator}
$X$를 준콤팩트이고 준분리인 스킴이라 하자. Lemma
\ref{perfect-lemma-better-coherator}의 함자 $DQ_X$는 다음 유계성 성질을 갖는다.
정수 $N = N(X)$가 존재하여, $K$가 $D(\mathcal{O}_X)$에 속하고 $X$의 아핀
열린집합 $U$ 및 $i \not \in [a, b]$에 대해 $H^i(U, K) = 0$이면,
$i \not \in [a, b + N]$에 대해 코호몰로지 층 $H^i(DQ_X(K))$는 0이다.
\end{lemma}

\begin{proof}
Cohomology of Schemes, Lemma \ref{coherent-lemma-induction-principle}의
귀납 원리를 사용하여 이를 증명하겠다.

\medskip\noindent
$X$가 아핀이면 $N = 0$으로 보조정리가 참이다. 실제로 이때 $RQ_X = DQ_X$는
$R\Gamma(X, K)$에 연관된 준연접 층의 복합체를 취함으로써 주어진다. Lemmas
\ref{perfect-lemma-affine-compare-bounded} 및 \ref{perfect-lemma-affine-coherator}를 보라.

\medskip\noindent
$U, V$가 $X$의 준콤팩트 열린집합이고 보조정리가 $U$, $V$, $U \cap V$에
대해 성립한다고 하자. 각각의 정수를 $N(U)$, $N(V)$, $N(U \cap V)$라 하자.
이제 $K$가 $D(\mathcal{O}_X)$에 속하고 모든 아핀 열린집합 $W \subset X$와
모든 $i \not \in [a, b]$에 대해 $H^i(W, K) = 0$이라고 하자. 그러면
$K|_U$, $K|_V$, $K|_{U \cap V}$도 같은 성질을 갖는다. 따라서
$RQ_U(K|_U)$, $RQ_V(K|_V)$, $RQ_{U \cap V}(K|_{U \cap V})$의 코호몰로지
층들은 구간 $[a, b + \max(N(U), N(V), N(U \cap V))$ 밖에서 소멸한다.
Lemma \ref{perfect-lemma-quasi-coherence-direct-image}에 의해 함자 $Rj_{U, *}$,
$Rj_{V, *}$, $Rj_{U \cap V, *}$는 $D_\QCoh$에서 유한 코호몰로지 차원을
가지므로, $Rj_{U, *}DQ_U(K|_U)$, $Rj_{V, *}DQ_V(K|_V)$, 그리고
$Rj_{U \cap V, *}DQ_{U \cap V}(K|_{U \cap V})$의 코호몰로지 층들이 구간
$[a, b + N]$ 밖에서 소멸하도록 하는 $N$이 존재한다. 마지막으로 Remark
\ref{perfect-remark-explain-consequence}의 완전 삼각형으로부터 결론을 얻는다.
\end{proof}

\begin{example}
\label{perfect-example-inverse-limit-quasi-coherent}
$X$를 준콤팩트이고 준분리인 스킴이라 하자. $(\mathcal{F}_n)$을 준연접 층의
역계라 하자. $DQ_X$는 오른쪽 수반이므로 곱과 가환하고, 따라서 유도 극한과
가환한다. 그러므로
$$
DQ_X(R\lim \mathcal{F}_n) =
(R\lim\text{ (범주 }D_\QCoh(\mathcal{O}_X)\text{에서)})(\mathcal{F}_n)
$$
임을 알 수 있다. 여기서 첫 번째 $R\lim$은 $D(\mathcal{O}_X)$에서 취한다.
실제로 이를 $K = R\lim \mathcal{F}_n$으로 쓰자. 임의의 아핀 열린집합
$U \subset X$에 대해
$$
H^i(U, K) =
H^i(R\Gamma(U, R\lim \mathcal{F}_n)) =
H^i(R\lim R\Gamma(U, \mathcal{F}_n)) =
H^i(R\lim \Gamma(U, \mathcal{F}_n))
$$
이다. 실제로 코호몰로지는 유도 극한과 가환하고, 준연접 층
$\mathcal{F}_n$은 아핀 위에서 고차 코호몰로지가 없다. 아벨 군의 범주에서
$R\lim$을 계산하면 $i \in [0, 1]$이 아닌 한 $H^i(U, K) = 0$임을 알 수
있다. 마지막으로 Lemma \ref{perfect-lemma-boundedness-better-coherator}에 의해
$D_\QCoh(\mathcal{O}_X)$에서의 $R\lim$, 즉 위에 의한 $DQ_X(K)$가
$D^b_\QCoh(\mathcal{O}_X)$에 속한다고 결론짓는다.
\end{example}













\section{코호몰로지와 기저변환 IV}
\label{perfect-section-cohomology-and-base-change-perfect}

\noindent
이 절은 Cohomology of Schemes, Section
\ref{coherent-section-cohomology-and-base-change-perfect}의 논의를 계속한다.
먼저 스킴들의 준콤팩트이고 준분리인 사상과 준연접 코호몰로지 층들을 갖는
복합체에 대한 매우 일반적인 사영 공식이 있다.

\begin{lemma}
\label{perfect-lemma-cohomology-base-change}
$f : X \to Y$를 스킴들의 준콤팩트이고 준분리인 사상이라 하자.
$D_\QCoh(\mathcal{O}_X)$의 $E$와 $D_\QCoh(\mathcal{O}_Y)$의 $K$에 대해,
Cohomology, Equation (\ref{cohomology-equation-projection-formula-map})에서
정의된 사상
$$
Rf_*(E) \otimes_{\mathcal{O}_Y}^\mathbf{L} K
\longrightarrow
Rf_*(E \otimes_{\mathcal{O}_X}^\mathbf{L} Lf^*K)
$$
은 동형이다.
\end{lemma}

\begin{proof}
사상이 동형임을 확인하기 위해 $Y$ 위에서 국소적으로 작업해도 된다. 따라서
$Y$가 아핀인 경우로 환원한다.

\medskip\noindent
$K = \bigoplus K_i$가 어떤 복합체 $K_i \in D_\QCoh(\mathcal{O}_Y)$들의
직합이라고 하자. 명제가 각 $K_i$에 대해 성립하면 $K$에 대해서도 성립한다.
실제로 함자 $Lf^*$와 $\otimes^\mathbf{L}$는 구성에 의해 직합을 보존하고,
Lemma \ref{perfect-lemma-quasi-coherence-pushforward-direct-sums}에 의해 $Rf_*$는
(준연접 코호몰로지 층들을 갖는 복합체에 대해) 직합과 가환한다. 더욱이
$K \to L \to M \to K[1]$이 $D_\QCoh(Y)$의 완전 삼각형이라고 하자. 보조정리의
명제가 $K, L, M$ 중 두 대상에 대해 성립하면 세 번째 대상에 대해서도
성립한다(관련 함자들이 삼각 범주의 완전 함자이기 때문이다).

\medskip\noindent
$Y$가 아핀이라고 하고 $Y = \Spec(A)$로 쓰자. 함자
$\widetilde{\ } : D(A) \to D_\QCoh(\mathcal{O}_Y)$는 동치이다(Lemma
\ref{perfect-lemma-affine-compare-bounded}). $K \in D(A)$에 대해 보조정리의 명제가
$\widetilde{K}$에 대해 성립한다는 성질을 $T$라 하자. 위의 논의와 More on
Algebra, Remark \ref{more-algebra-remark-P-resolution}에 의해 $T$가 $A[k]$에
대해 성립함을 증명하면 충분하다. 구조층의 이동에 대해서는 보조정리의 명제가
자명하므로 이것으로 증명이 끝난다.
\end{proof}

\begin{definition}
\label{perfect-definition-tor-independent}
$S$를 스킴이라 하고 $X$, $Y$를 $S$ 위의 스킴이라 하자. 같은 점
$s \in S$로 가는 모든 $x \in X$와 $y \in Y$에 대해 환
$\mathcal{O}_{X, x}$와 $\mathcal{O}_{Y, y}$가 $\mathcal{O}_{S, s}$ 위에서
Tor 독립이면, $X$와 $Y$가 {\it $S$ 위에서 Tor 독립}이라고 한다(More on
Algebra, Definition \ref{more-algebra-definition-tor-independent}을 보라).
\end{definition}

\begin{lemma}
\label{perfect-lemma-tor-independent}
$f : X \to S$와 $g : Y \to S$를 스킴들의 사상이라 하자. 다음은 동치이다.
\begin{enumerate}
\item $X$와 $Y$는 $S$ 위에서 Tor 독립이다.
\item $f(U) \subset W$와 $g(V) \subset W$를 만족하는 모든 아핀 열린집합
$U \subset X$, $V \subset Y$, $W \subset S$에 대해 환 $\mathcal{O}_X(U)$와
$\mathcal{O}_Y(V)$는 $\mathcal{O}_S(W)$ 위에서 Tor 독립이다.
\item 아핀 열린 덮개 $S = \bigcup W_i$와, 각 $i$에 대해 아핀 열린 덮개
$f^{-1}(W_i) = \bigcup U_{ij}$ 및 $g^{-1}(W_i) = \bigcup V_{ik}$가 존재하여,
모든 $i, j, k$에 대해 환 $\mathcal{O}_X(U_{ij})$와
$\mathcal{O}_Y(V_{ik})$는 $\mathcal{O}_S(W_i)$ 위에서 Tor 독립이다.
\end{enumerate}
\end{lemma}

\begin{proof}
생략한다. 힌트: More on Algebra, Lemma
\ref{more-algebra-lemma-tor-independent}을 사용하라.
\end{proof}

\begin{lemma}
\label{perfect-lemma-flat-base-change-tor-independent}
$X \to S$와 $Y \to S$를 스킴들의 사상이라 하자. $S' \to S$를 스킴들의
사상이라 하고 $X' = X \times_S S'$ 및 $Y' = Y \times_S S'$로 나타내자.
$X$와 $Y$가 $S$ 위에서 Tor 독립이고 $S' \to S$가 평탄이면, $X'$와 $Y'$는
$S'$ 위에서 Tor 독립이다.
\end{lemma}

\begin{proof}
생략한다. 힌트: Lemma \ref{perfect-lemma-tor-independent}을 사용하고, 아핀
열린집합에서는 More on Algebra, Lemma
\ref{more-algebra-lemma-flat-base-change-tor-independent}을 사용하라.
\end{proof}

\begin{lemma}
\label{perfect-lemma-compare-base-change}
$g : S' \to S$를 스킴들의 사상이라 하자. $f : X \to S$를 준콤팩트이고
준분리인 사상이라 하자. 기저변환 그림
$$
\xymatrix{
X' \ar[r]_{g'} \ar[d]_{f'} &
X \ar[d]^f \\
S' \ar[r]^g &
S
}
$$
을 생각하자. $X$와 $S'$가 $S$ 위에서 Tor 독립이면, 모든
$E \in D_\QCoh(\mathcal{O}_X)$에 대해 표준 화살표
$Lg^*Rf_*E \to Rf'_*L(g')^*E$는 동형이다.
\end{lemma}

\begin{proof}
$D(\mathcal{O}_X)$의 임의의 대상 $E$에 대해 Cohomology, Remark
\ref{cohomology-remark-base-change}를 사용하여 표준 기저변환 사상
$Lg^*Rf_*E \to Rf'_*L(g')^*E$를 얻을 수 있다. 이것이 동형임을 확인하기
위해 $S'$ 위에서 국소적으로 작업해도 된다. 따라서 $g : S' \to S$가 아핀
스킴들의 사상이라고 가정해도 된다. 특히 $g$는 아핀이므로
$$
Rg_*Lg^*Rf_*E \to Rg_*Rf'_*L(g')^*E = Rf_*(Rg'_* L(g')^* E)
$$
이 동형임을 보이면 충분하다. Lemma \ref{perfect-lemma-affine-morphism}을 보라
($Rf'_*L(g')^*E$와 $Lg^*Rf_*E$가 준연접 코호몰로지 층들을 가짐을 보기
위해 Lemmas \ref{perfect-lemma-quasi-coherence-pullback},
\ref{perfect-lemma-quasi-coherence-tensor-product}, 그리고
\ref{perfect-lemma-quasi-coherence-direct-image}를 사용하라). $g'$도 아핀임에 유의하라
(Morphisms, Lemma \ref{morphisms-lemma-base-change-affine}). Lemma
\ref{perfect-lemma-affine-morphism-pull-push}에 의해 이 사상은
$$
Rf_*E \otimes_{\mathcal{O}_S}^\mathbf{L} g_*\mathcal{O}_{S'}
\longrightarrow
Rf_*(E \otimes_{\mathcal{O}_X}^\mathbf{L} g'_*\mathcal{O}_{X'})
$$
이 된다. $g'_*\mathcal{O}_{X'} = f^*g_*\mathcal{O}_{S'}$임에 유의하라(아핀
기저변환에 의한다. Cohomology of Schemes, Lemma
\ref{coherent-lemma-affine-base-change}를 보라). 따라서 Lemma
\ref{perfect-lemma-cohomology-base-change}에 의해
$Lf^*g_*\mathcal{O}_{S'} = f^*g_*\mathcal{O}_{S'}$임을 증명하면 충분하다.
이는 $X$와 $S'$가 $S$ 위에서 Tor 독립이라는 가정에서 따른다. 실제로 이를
확인하기 위해 $X$ 위에서 국소적으로 작업해도 되므로 $X$도 아핀이라고 가정할
수 있다. $X = \Spec(A)$, $S = \Spec(R)$, $S' = \Spec(R')$로 쓰자. 우리의
가정은 $A$와 $R'$가 $R$ 위에서 Tor 독립임을 함의한다(More on Algebra,
Lemma \ref{more-algebra-lemma-tor-independent}). 즉 $i > 0$에 대해
$\text{Tor}_i^R(A, R') = 0$이다. 다시 말해
$A \otimes_R^\mathbf{L} R' = A \otimes_R R'$이고, 이는 정확히
$Lf^*g_*\mathcal{O}_{S'} = f^*g_*\mathcal{O}_{S'}$를 뜻한다(Lemma
\ref{perfect-lemma-quasi-coherence-pullback}을 사용하라).
\end{proof}

\noindent
다음 보조정리는 쌍대화 복합체에 관한 장에서 사용될 것이다.

\begin{lemma}
\label{perfect-lemma-affine-morphism-and-hom-out-of-perfect}
준콤팩트이고 준분리인 스킴들의 데카르트 정사각형
$$
\xymatrix{
X' \ar[r]_{g'} \ar[d]_{f'} & X \ar[d]^f \\
S' \ar[r]^g & S
}
$$
을 생각하자. $g$와 $f$가 Tor 독립이고 $S = \Spec(R)$,
$S' = \Spec(R')$가 아핀이라고 가정하자. $M, K \in D(\mathcal{O}_X)$에 대해
표준 사상
$$
R\Hom_X(M, K) \otimes^\mathbf{L}_R R'
\longrightarrow
R\Hom_{X'}(L(g')^*M, L(g')^*K)
$$
은 다음 두 경우에 $D(R')$에서 동형이다.
\begin{enumerate}
\item $M \in D(\mathcal{O}_X)$는 완전이고 $K \in D_\QCoh(X)$이다.
\item $M \in D(\mathcal{O}_X)$는 유사연접이고,
$K \in D_\QCoh^+(X)$이며, $R'$은 $R$ 위에서 유한 Tor 차원을 갖는다.
\end{enumerate}
\end{lemma}

\begin{proof}
대역 Hom 복합체들의 $D(\Gamma(X, \mathcal{O}_X))$ 안에는 표준 사상
$R\Hom_X(M, K) \to R\Hom_{X'}(L(g')^*M, L(g')^*K)$가 있다. Cohomology,
Section \ref{cohomology-section-global-RHom}을 보라. 스칼라를 제한하면 이를
$D(R)$의 사상으로 볼 수 있다. 그러면 제한과 $- \otimes_R^\mathbf{L} R'$의
수반성을 사용하여 보조정리에 표시된 사상을 얻는다. 사상을 정의했으므로 아벨
군의 유도 범주에서 이것이 동형임을 증명하면 충분하다.

\medskip\noindent
오른쪽 변은 Lemma \ref{perfect-lemma-affine-morphism-pull-push}에 의해
$$
R\Hom_X(M, R(g')_*L(g')^*K) =
R\Hom_X(M, K \otimes_{\mathcal{O}_X}^\mathbf{L} g'_*\mathcal{O}_{X'})
$$
과 같다. 두 경우 모두 Lemma \ref{perfect-lemma-quasi-coherence-internal-hom}에 의해
복합체 $R\SheafHom(M, K)$는 $D_\QCoh(\mathcal{O}_X)$의 대상이다. 자연스러운
사상
$$
R\SheafHom(M, K) \otimes_{\mathcal{O}_X}^\mathbf{L} g'_*\mathcal{O}_{X'}
\longrightarrow
R\SheafHom(M, K \otimes_{\mathcal{O}_X}^\mathbf{L} g'_*\mathcal{O}_{X'})
$$
이 있으며, Lemma \ref{perfect-lemma-internal-hom-evaluate-tensor-isomorphism}에 의해
두 경우 모두 동형이다. 이 보조정리가 (2)의 경우에 적용됨을 보기 위해
$g'_*\mathcal{O}_{X'} = Rg'_*\mathcal{O}_{X'} =
Lf^*g_*\mathcal{O}_{S'}$임에 유의하라. 두 번째 등식은 Lemma
\ref{perfect-lemma-compare-base-change}에 의한다. Lemma
\ref{perfect-lemma-tor-dimension-affine}과 Cohomology, Lemma
\ref{cohomology-lemma-tor-amplitude-pullback}을 사용하면
$g'_*\mathcal{O}_{X'}$가 유한 Tor 차원을 가짐을 얻는다. 따라서 두 경우 모두
$K$를 $R\SheafHom(M, K)$로 바꾸면
$$
R\Gamma(X, K) \otimes^\mathbf{L}_A A' \longrightarrow
R\Gamma(X, K \otimes^\mathbf{L}_{\mathcal{O}_X} g'_*\mathcal{O}_{X'})
$$
이 동형임을 증명하는 것으로 환원된다. Lemma
\ref{perfect-lemma-affine-morphism-pull-push}에 의해 왼쪽 변은
$R\Gamma(X', L(g')^*K)$와 같음에 유의하라. 따라서 결과는 Lemma
\ref{perfect-lemma-compare-base-change}에서 따른다.
\end{proof}

\begin{remark}
\label{perfect-remark-multiplication-map}
Lemma \ref{perfect-lemma-affine-morphism-and-hom-out-of-perfect}에서와 같은 표기법을
사용하자. 그림
$$
\xymatrix{
R\Hom_X(M, Rg'_*L) \otimes_R^\mathbf{L} R' \ar[r] \ar[d]_\mu &
R\Hom_{X'}(L(g')^*M, L(g')^*Rg'_*L) \ar[d]^a \\
R\Hom_X(M, R(g')_*L) \ar@{=}[r] &
R\Hom_{X'}(L(g')^*M, L)
}
$$
은 가환한다. 여기서 위의 수평 화살표는 보조정리의 사상이고, $\mu$는 곱셈
사상이며, $a$는 수반 사상 $L(g')^*Rg'_*L \to L$에서 나온다. 곱셈 사상은
임의의 $K' \in D(R')$에 대한 수반 사상
$K' \otimes_R^\mathbf{L} R' \to K'$이다.
\end{remark}

\begin{lemma}
\label{perfect-lemma-tor-independence-and-tor-amplitude}
스킴들의 데카르트 정사각형
$$
\xymatrix{
X' \ar[r]_{g'} \ar[d]_{f'} & X \ar[d]^f \\
S' \ar[r]^g & S
}
$$
을 생각하자. $g$와 $f$가 Tor 독립이라고 가정하자.
\begin{enumerate}
\item $E \in D(\mathcal{O}_X)$가 $f^{-1}\mathcal{O}_S$-가군의 복합체로서
$[a, b]$에서 Tor 진폭을 가지면, $L(g')^*E$는
$f^{-1}\mathcal{O}_{S'}$-가군의 복합체로서 $[a, b]$에서 Tor 진폭을 갖는다.
\item $\mathcal{G}$가 $S$ 위에서 평탄한 $\mathcal{O}_X$-가군이면
$L(g')^*\mathcal{G} = (g')^*\mathcal{G}$이다.
\end{enumerate}
\end{lemma}

\begin{proof}
줄기에서 Tor 차원을 계산할 수 있다. Cohomology, Lemma
\ref{cohomology-lemma-tor-amplitude-stalk}을 보라. $x' \in X'$의 상이
$x \in X$이면
$$
(L(g')^*E)_{x'} =
E_x \otimes_{\mathcal{O}_{X, x}}^\mathbf{L} \mathcal{O}_{X', x'}
$$
이다. $s' \in S'$와 $s \in S$를 각각 $x'$와 $x$의 상이라 하자. $X$와
$S'$가 $S$ 위에서 Tor 독립이므로 More on Algebra, Lemma
\ref{more-algebra-lemma-base-change-comparison}을 적용하면, 표시된 식의 오른쪽
변이 $D(\mathcal{O}_{S', s'})$에서
$E_x \otimes_{\mathcal{O}_{S, s}}^\mathbf{L} \mathcal{O}_{S', s'}$와 같음을
알 수 있다. 따라서 (1)은 More on Algebra, Lemma
\ref{more-algebra-lemma-pull-tor-amplitude}에서 따른다. (2)를 보기 위해
$\mathcal{G}$의 평탄성은 $\mathcal{G}[0]$이 $[0, 0]$에서 Tor 진폭을 갖는다는
조건과 동치임에 유의하라. (1)을 적용하면 결론을 얻는다.
\end{proof}

\begin{lemma}
\label{perfect-lemma-compare-base-change-closed-immersion}
스킴들의 데카르트 그림
$$
\xymatrix{
Z' \ar[r]_{i'} \ar[d]_g & X' \ar[d]^f \\
Z \ar[r]^i & X
}
$$
을 생각하자. 여기서 $i$는 닫힌 매장이다. $Z$와 $X'$가 $X$ 위에서 Tor
독립이면, 함자 $D(\mathcal{O}_Z) \to D(\mathcal{O}_{X'})$로서
$Ri'_* \circ Lg^* = Lf^* \circ Ri_*$이다.
\end{lemma}

\begin{proof}
보조정리는 모든 $K \in D(\mathcal{O}_Z)$에 대해 성립해야 함에 유의하라.
$i_*$와 $i'_*$는 완전 함자이므로 $Ri_*$와 $Ri'_*$는 임의의 대표에 $i_*$와
$i'_*$를 적용하여 계산된다. 따라서 기저변환 사상
$$
Lf^*(Ri_*(K)) \longrightarrow Ri'_*(Lg^*(K))
$$
은 상이 $z \in Z$인 점 $z' \in Z'$의 줄기에서
$$
K_z \otimes_{\mathcal{O}_{X, z}}^\mathbf{L} \mathcal{O}_{X', z'}
\longrightarrow
K_z \otimes_{\mathcal{O}_{Z, z}}^\mathbf{L} \mathcal{O}_{Z', z'}
$$
로 주어진다. More on Algebra, Lemma
\ref{more-algebra-lemma-base-change-comparison}과 가정한 Tor 독립성에 의해 이
사상은 동형이다.
\end{proof}








\section{K\"unneth 공식 II}
\label{perfect-section-kunneth}

\noindent
기저가 체인 경우는 Varieties, Section \ref{varieties-section-kunneth}을
참조하라. 스킴들의 데카르트 도식을 생각하자.
$$
\xymatrix{
& X \times_S Y \ar[ld]^p \ar[rd]_q \ar[dd]^f \\
X \ar[rd]_a & & Y \ar[ld]^b \\
& S
}
$$
$K \in D(\mathcal{O}_X)$ 및 $M \in D(\mathcal{O}_Y)$라 하자.
표준 사상이 존재한다.
\begin{equation}
\label{perfect-equation-kunneth}
Ra_*K \otimes_{\mathcal{O}_S}^\mathbf{L} Rb_*M
\longrightarrow
Rf_*(Lp^*K \otimes_{\mathcal{O}_{X \times_S Y}}^\mathbf{L} Lq^*M)
\end{equation}
구체적으로 다음 사상들을 사용하고,
$Ra_*K \to Ra_*Rp_* Lp^*K = Rf_*Lp^*K$와
$Rb_*M \to Rb_*Rq_* Lq^*M = Rf_*Lq^*M$를 사용한 뒤
상대 컵곱(Cohomology, Remark \ref{cohomology-remark-cup-product})을
사용하면 된다.

\medskip\noindent
 $A = \Gamma(S, \mathcal{O}_S)$로 놓자. $D(A)$ 안에 대역 K\"unneth 사상이
\begin{equation}
\label{perfect-equation-kunneth-global}
R\Gamma(X, K) \otimes_A^\mathbf{L} R\Gamma(Y, M)
\longrightarrow
R\Gamma(X \times_S Y,
Lp^*K \otimes_{\mathcal{O}_{X \times_S Y}}^\mathbf{L} Lq^*M)
\end{equation}
존재한다. 이 사상은 다음 당김 사상들
$R\Gamma(X, K) \to R\Gamma(X \times_S Y, Lp^*K)$와
$R\Gamma(Y, M) \to R\Gamma(X \times_S Y, Lq^*M)$와
Cohomology, Section \ref{cohomology-section-cup-product}에서
구성한 컵곱을 사용하여 구성된다.

\begin{lemma}
\label{perfect-lemma-kunneth}
위 상황에서 $a$와 $b$가 준콤팩트이고 준분리이며 $X$와 $Y$가
$S$ 위에서 Tor 독립이면, (\ref{perfect-equation-kunneth})는
$K \in D_\QCoh(\mathcal{O}_X)$ 및 $M \in D_\QCoh(\mathcal{O}_Y)$에
대해 동형사상이다. 또한 $S = \Spec(A)$가 아핀이면
(\ref{perfect-equation-kunneth-global})의 사상도 동형사상이다.
\end{lemma}

\begin{proof}[첫 번째 증명]
이는 다음 동형사상 열로부터 따른다.
\begin{align*}
Rf_*(Lp^*K \otimes_{\mathcal{O}_{X \times_S Y}}^\mathbf{L} Lq^*M)
& =
Ra_*Rp_*(Lp^*K \otimes_{\mathcal{O}_{X \times_S Y}}^\mathbf{L} Lq^*M) \\
& =
Ra_*(K \otimes_{\mathcal{O}_X}^\mathbf{L} Rp_*Lq^*M) \\
& =
Ra_*(K \otimes_{\mathcal{O}_X}^\mathbf{L} La^*Rb_*M) \\
& =
Ra_*K \otimes_{\mathcal{O}_S}^\mathbf{L} Rb_*M
\end{align*}
첫 번째 등식은 $f = a \circ p$이기 때문이다. 두 번째 등식은
Lemma \ref{perfect-lemma-cohomology-base-change}, 세 번째 등식은
Lemma \ref{perfect-lemma-compare-base-change}, 네 번째 등식은
Lemma \ref{perfect-lemma-cohomology-base-change}에 따른다.
이 동형사상들의 합성이 (\ref{perfect-equation-kunneth})의 사상과 같다는
확인은 생략한다. $S$가 아핀이면
(\ref{perfect-equation-kunneth-global})의 화살표의 시작 대상과 끝 대상은
(\ref{perfect-equation-kunneth})의 시작 대상과 끝 대상에 $R\Gamma(S, -)$를
적용하여 얻어진다. 따라서 마지막 주장을 얻는다. 자세한 내용은 생략한다.
\end{proof}

\begin{proof}[두 번째 증명]
(\ref{perfect-equation-kunneth})의 화살표 구성은 상대 컵곱의 구성에서
즉시 보이듯이 $S$의 열린 부분 스킴으로 제한하는 것과 양립한다.
따라서 $S$가 아핀일 때 (\ref{perfect-equation-kunneth})가 동형사상임을
증명하면 충분하다.

\medskip\noindent
 $S = \Spec(A)$가 아핀이라고 하자. Leray에 의해
$R\Gamma(S, Rf_*K) = R\Gamma(X, K)$이고 다른 경우도 마찬가지이다.
Cohomology, Lemma \ref{cohomology-lemma-compose-cup-product}에 의해
(\ref{perfect-equation-kunneth})의 사상에 $R\Gamma(S, -)$를 적용하면
(\ref{perfect-equation-kunneth-global})의 사상이 유도된다. 한편
Lemmas \ref{perfect-lemma-quasi-coherence-direct-image}와
\ref{perfect-lemma-quasi-coherence-tensor-product}
에 의해 (\ref{perfect-equation-kunneth})의 사상의 시작 대상과 끝 대상은
$D_\QCoh(\mathcal{O}_S)$에 속한다. 그러므로 Lemma
\ref{perfect-lemma-affine-compare-bounded}에 의해
(\ref{perfect-equation-kunneth-global})가 동형사상임을 증명하면 충분하다.

\medskip\noindent
 $S = \Spec(A)$, $X = \Spec(B)$, $Y = \Spec(C)$가 모두 아핀이라고 하자.
이하에서는 Lemma \ref{perfect-lemma-affine-compare-bounded}를 별도 언급 없이
사용한다. 이 경우 항들이 평탄한 $B$-가군인 K-평탄 복합체
$K^\bullet$를 택하여 $K$가 $\widetilde{K}^\bullet$로 표현되게 할 수 있다.
마찬가지로 항들이 평탄한 $C$-가군인 K-평탄 복합체 $M^\bullet$를 택하여
$M$이 $\widetilde{M}^\bullet$로 표현되게 할 수 있다. More on Algebra,
Lemma \ref{more-algebra-lemma-K-flat-resolution}을 보라.
그러면 $\widetilde{K}^\bullet$는 $\mathcal{O}_X$-가군의 K-평탄 복합체이고
$\widetilde{M}^\bullet$도 마찬가지이다(Lemma \ref{perfect-lemma-affine-K-flat} 참조).
따라서 $La^*K$는 다음으로 표현된다.
$$
a^*\widetilde{K}^\bullet = \widetilde{K^\bullet \otimes_A C}
$$
 $Lb^*M$도 마찬가지이다. 이는 위에서 인용한 보조정리와 More on Algebra,
Lemma \ref{more-algebra-lemma-base-change-K-flat}에 의해
$\mathcal{O}_{X \times_S Y}$-가군의 K-평탄 복합체이다.
따라서 다음에 대응하는 $\mathcal{O}_{X \times_S Y}$-가군 복합체가
$$
\text{Tot}((K^\bullet \otimes_A C) \otimes_{B \otimes_A C}
(B \otimes_A M^\bullet)) =
\text{Tot}(K^\bullet \otimes_A M^\bullet)
$$
은 $X \times_S Y$의 유도 범주에서
$La^*K \otimes_{\mathcal{O}_{X \times_S Y}}^\mathbf{L} Lb^*M$을 표현한다.
대역 단면을 취하면 $\text{Tot}(K^\bullet \otimes_A M^\bullet)$를 얻는데,
이는 물론 $R\Gamma(X, K) \otimes_A^\mathbf{L} R\Gamma(Y, M)$을
표현하는 복합체이기도 하다. 이 동형사상이 컵곱으로 주어진다는 사실은
Cohomology, Section \ref{cohomology-section-cup-product}에서 진정한
컵곱과 순진한 컵곱 사이의 관계로부터 따른다(특히 Cohomology, Lemma
\ref{cohomology-lemma-cup-compatible-with-naive} 및 그 뒤의 논의를 보라).

\medskip\noindent
 $S = \Spec(A)$와 $Y$가 아핀이라고 하자. 다음의 귀납 원리를 사용한다.
Cohomology of Schemes, Lemma \ref{coherent-lemma-induction-principle}
주장을 증명하기 위해 다음을 보이면 충분하다. $X = U \cup V$가
$U$, $V$라는 준콤팩트 열린집합에 의한 열린 덮개이고, 다음 사상
(\ref{perfect-equation-kunneth-global})이
$$
R\Gamma(U, K) \otimes_A^\mathbf{L} R\Gamma(Y, M)
\longrightarrow
R\Gamma(U \times_S Y,
Lp^*K \otimes_{\mathcal{O}_{X \times_S Y}}^\mathbf{L} Lq^*M)
$$
 $S$ 위의 $U$, $Y$에 대해 성립하고, 다음 사상
(\ref{perfect-equation-kunneth-global})이
$$
R\Gamma(V, K) \otimes_A^\mathbf{L} R\Gamma(Y, M)
\longrightarrow
R\Gamma(V \times_S Y,
Lp^*K \otimes_{\mathcal{O}_{X \times_S Y}}^\mathbf{L} Lq^*M)
$$
 $S$ 위의 $V$, $Y$에 대해 성립하며, 다음 사상
(\ref{perfect-equation-kunneth-global})이
$$
R\Gamma(U \cap V, K) \otimes_A^\mathbf{L} R\Gamma(Y, M)
\longrightarrow
R\Gamma((U \cap V) \times_S Y,
Lp^*K \otimes_{\mathcal{O}_{X \times_S Y}}^\mathbf{L} Lq^*M)
$$
 $S$ 위의 $U \cap V$, $Y$에 대해 동형사상이면, $S$ 위의 $X$, $Y$에
대한 사상 (\ref{perfect-equation-kunneth-global})도 동형사상이다.
그러나 Cohomology, Lemma \ref{cohomology-lemma-mayer-vietoris-cup}에
의해 이 사상들은 (\ref{perfect-equation-kunneth-global})을 마지막 변으로 하는
완전 삼각형 사이의 사상을 이룬다. 따라서 Derived Categories, Lemma
\ref{derived-lemma-third-isomorphism-triangle}에 의해 결론을 얻는다.

\medskip\noindent
 $S = \Spec(A)$가 아핀이라고 하자. 증명을 끝내기 위해 다음의 귀납 원리를
Cohomology of Schemes, Lemma \ref{coherent-lemma-induction-principle}
 $Y$에 적용할 수 있다. 즉 위에서 이미 $Y$가 아핀이면 우리 사상이
동형사상임을 알았다. 나머지 논의는 $X$와 $Y$의 역할을 바꾸면
앞 단락과 정확히 같다.
\end{proof}

\begin{lemma}
\label{perfect-lemma-cohomology-de-rham-base-change}
$a : X \to S$를 스킴의 준콤팩트 준분리 사상이라 하자.
$\mathcal{F}^\bullet$를 $a^{-1}\mathcal{O}_S$-가군의 국소 유계
복합체라 하자. 모든 $n \in \mathbf{Z}$에 대해 층 $\mathcal{F}^n$이
평탄한 $a^{-1}\mathcal{O}_S$-가군이고, $\mathcal{F}^n$은 주어진
$a^{-1}\mathcal{O}_S$-가군 구조와 양립하는 준연접
$\mathcal{O}_X$-가군 구조를 가진다고 가정하자(단,
$\mathcal{F}^\bullet$의 미분이 반드시 $\mathcal{O}_X$-선형일 필요는 없다).
그러면 다음이 성립한다.
\begin{enumerate}
\item $Ra_*\mathcal{F}^\bullet$는 국소 유계이다.
\item $Ra_*\mathcal{F}^\bullet$는 $D_\QCoh(\mathcal{O}_S)$에 속한다.
\item $Ra_*\mathcal{F}^\bullet$는 국소적으로 유한 Tor 차원을 갖는다.
\item $\mathcal{G} \otimes_{\mathcal{O}_S}^\mathbf{L} Ra_*\mathcal{F}^\bullet =
Ra_*(a^{-1}\mathcal{G} \otimes_{a^{-1}\mathcal{O}_S} \mathcal{F}^\bullet)$
이는 $\mathcal{G} \in \QCoh(\mathcal{O}_S)$에 대해 성립한다.
\item $K \otimes_{\mathcal{O}_S}^\mathbf{L} Ra_*\mathcal{F}^\bullet =
Ra_*(a^{-1}K \otimes_{a^{-1}\mathcal{O}_S}^\mathbf{L} \mathcal{F}^\bullet)$
이는 $K \in D_\QCoh(\mathcal{O}_S)$에 대해 성립한다.
\end{enumerate}
\end{lemma}

\begin{proof}
(1), (2), (3)은 $S$ 위에서 국소적이므로 $S$가 아핀이라고 가정해도
된다. $a$가 준콤팩트이므로 $X$도 준콤팩트이다. $\mathcal{F}^\bullet$가
국소 유계이므로 $\mathcal{F}^\bullet$는 유계이다.

\medskip\noindent
 (1), (2)에는 다음 첫 번째 스펙트럴 열을 사용할 수 있다.
$R^pa_*\mathcal{F}^q \Rightarrow R^{p + q}a_*\mathcal{F}^\bullet$는
Derived Categories, Lemma \ref{derived-lemma-two-ss-complex-functor}의
것이다. Cohomology of Schemes, Lemma
\ref{coherent-lemma-quasi-coherence-higher-direct-images}와 Homology,
Lemma \ref{homology-lemma-biregular-ss-converges}를 결합하면 결론을 얻는다.

\medskip\noindent
(3)은 Cohomology of Schemes, Lemma
\ref{coherent-lemma-induction-principle}의 귀납 원리로 증명하자.
즉 $X$의 준콤팩트 열린집합 $U$에 대해
$R(a|_U)_*(\mathcal{F}^\bullet|_U)$가 유한 Tor 차원을 갖는다는
조건을 생각한다. $U, V$가 $X$의 준콤팩트 열린집합이면 상대
Mayer--Vietoris 완전 삼각형
$$
R(a|_{U \cup V})_*\mathcal{F}^\bullet|_{U \cup V} \to
R(a|_U)_*\mathcal{F}^\bullet|_U \oplus
R(a|_V)_*\mathcal{F}^\bullet|_V \to
R(a|_{U \cap V})_*\mathcal{F}^\bullet|_{U \cap V} \to
$$
Cohomology, Lemma \ref{cohomology-lemma-unbounded-relative-mayer-vietoris}에
의해 주어진다. 완전 삼각형에서 Tor 진폭의 거동(Cohomology, Lemma
\ref{cohomology-lemma-cone-tor-amplitude} 참조)에 의해 $U$, $V$,
$U \cap V$에 대한 결과를 알면 $U \cup V$에도 결과가 성립한다.
따라서 $X$가 아핀인 경우로 환원된다. 이 경우
$$
Ra_*\mathcal{F}^\bullet = a_*\mathcal{F}^\bullet
$$
Leray 비순환성 보조정리(Derived Categories, Lemma
\ref{derived-lemma-leray-acyclicity})와 아핀 사상 아래 준연접 가군의
고차 직상 소멸(Cohomology of Schemes, Lemma
\ref{coherent-lemma-relative-affine-vanishing})에 의해 성립한다.
가정에 따라 $\mathcal{F}^n$은 $S$-평탄하고 $X$는 아핀이므로,
모든 $n$에 대해 $a_*\mathcal{F}^n$은 평탄하다. 따라서
$a_*\mathcal{F}^\bullet$는 평탄한 $\mathcal{O}_S$-가군의 유계 복합체이고,
유한 Tor 차원을 갖는다.

\medskip\noindent
 (5)를 증명하자. 다음을
$a' : (X, a^{-1}\mathcal{O}_S) \to (S, \mathcal{O}_S)$
환붙은 공간의 자명한 평탄 사상이라 하자. (5)의 주장은 다음과 같다.
$$
K \otimes_{\mathcal{O}_S}^\mathbf{L} Ra'_*\mathcal{F}^\bullet =
Ra'_*(L(a')^*K \otimes_{a^{-1}\mathcal{O}_S}^\mathbf{L}
\mathcal{F}^\bullet)
$$
그러므로 Cohomology, Equation (\ref{cohomology-equation-projection-formula-map})은
왼쪽에서 오른쪽으로 가는 함자적 사상을 주며, 이 사상이 동형사상임을
보이고자 한다. 이 문제는 $S$ 위에서 국소적이므로 $S$가 아핀이라고
가정해도 된다. 나머지 증명은 Lemma \ref{perfect-lemma-cohomology-base-change}의
증명과 {\it 정확히} 같지만, 다음 함자가
$K \mapsto Ra'_*(L(a')^*K \otimes_{a^{-1}\mathcal{O}_S}^\mathbf{L}
\mathcal{F}^\bullet)$가 직합과 교환한다는 것을 보여야 한다.
여기서 $\mathcal{F}^n$이 준연접 $\mathcal{O}_X$-가군 구조를 갖는다는
사실을 사용한다. 즉 다음 함자는
$K \mapsto L(a')^*K
\otimes_{a^{-1}\mathcal{O}_S}^\mathbf{L} \mathcal{F}^\bullet$
임의의 직합과 교환한다. 이제 $\mathcal{F}^\bullet$가 하나의 준연접
$\mathcal{O}_X$-가군으로 이루어져 $\mathcal{F}^\bullet = \mathcal{F}^n[-n]$
이면 다음을 얻는다. $L(a')^*G
\otimes_{a^{-1}\mathcal{O}_S}^\mathbf{L} \mathcal{F}^\bullet =
La^*K \otimes_{\mathcal{O}_X}^\mathbf{L} \mathcal{F}^n[-n]$, see
Cohomology, Lemma \ref{cohomology-lemma-variant-derived-pullback}을 보라.
따라서 이 경우 직합과의 교환은 Lemma
\ref{perfect-lemma-quasi-coherence-pushforward-direct-sums}에서 따른다.
일반적으로 $S$가 아핀(따라서 $X$가 준콤팩트)이고
$\mathcal{F}^\bullet$가 국소 유계이므로 다음 복합체는
$$
\mathcal{F}^\bullet = (\mathcal{F}^a \to \ldots \to \mathcal{F}^b)
$$
유계이다. $b-a$에 대한 귀납법과 다음 완전 삼각형을 사용하면
$$
\mathcal{F}^b[-b] \to (\mathcal{F}^a \to \ldots \to \mathcal{F}^b)
\to (\mathcal{F}^a \to \ldots \to \mathcal{F}^{b - 1}) \to
\mathcal{F}^b[-b + 1]
$$
이 부분의 증명이 끝난다. 일부 세부 사항은 생략한다.

\medskip\noindent
 (4)를 증명하자. $a' : (X, a^{-1}\mathcal{O}_S) \to (S, \mathcal{O}_S)$를
위와 같이 둔다. $\mathcal{F}^\bullet$는 평탄한
$a^{-1}\mathcal{O}_S$-가군의 국소 유계 복합체이므로 다음 복합체는
$a^{-1}\mathcal{G} \otimes_{a^{-1}\mathcal{O}_S} \mathcal{F}^\bullet$
 $L(a')^*\mathcal{G}
\otimes_{a^{-1}\mathcal{O}_S}^\mathbf{L}
\mathcal{F}^\bullet$를 $D(a^{-1}\mathcal{O}_S)$에서 표현한다. 따라서 (4)는
(5)에서 따른다.
\end{proof}

\begin{lemma}
\label{perfect-lemma-K-flat}
 $f : X \to Y$를 $Y = \Spec(A)$가 아핀인 스킴의 사상이라 하자.
$\mathcal{U} : X = \bigcup_{i \in I} U_i$를 유한 아핀 열린 덮개로서,
모든 유한 교집합
$U_{i_0 \ldots i_p} = U_{i_0} \cap \ldots \cap U_{i_p}$
아핀이라고 하자. $\mathcal{F}^\bullet$를
$f^{-1}\mathcal{O}_Y$-가군의 유계 복합체라 하자. 모든 $n \in \mathbf{Z}$에
대해 층 $\mathcal{F}^n$이 평탄한 $f^{-1}\mathcal{O}_Y$-가군이고,
주어진 $p^{-1}\mathcal{O}_Y$-가군 구조와 양립하는 $\mathcal{F}^n$의 준연접
$\mathcal{O}_X$-가군 구조를 갖는다고 가정하자(단, $\mathcal{F}^\bullet$의 미분은
$\mathcal{O}_X$-선형일 필요가 없다). 그러면 다음 복합체
$\text{Tot}(\check{\mathcal{C}}^\bullet(\mathcal{U}, \mathcal{F}^\bullet))$
는 $A$-가군의 복합체로서 K-평탄이다.
\end{lemma}

\begin{proof}
다음과 같이 쓸 수 있다.
$$
\mathcal{F}^\bullet = (\mathcal{F}^a \to \ldots \to \mathcal{F}^b)
$$
 $b-a$에 대한 귀납법으로 다음 완전 삼각형을 생각하고
$$
\mathcal{F}^b[-b] \to (\mathcal{F}^a \to \ldots \to \mathcal{F}^b)
\to (\mathcal{F}^a \to \ldots \to \mathcal{F}^{b - 1}) \to
\mathcal{F}^b[-b + 1]
$$
More on Algebra, Lemma \ref{more-algebra-lemma-K-flat-two-out-of-three}를
사용하면 $\mathcal{F}^\bullet$가 $0$차에 놓인 하나의 준연접
$\mathcal{O}_X$-가군 $\mathcal{F}$만으로 이루어진 경우로 환원된다.
이 경우 $\mathcal{F}$와 $\mathcal{U}$의 {\v C}ech 복합체는 교대
{\v C}ech 복합체와 호모토피 동치이다(Cohomology, Lemma
\ref{cohomology-lemma-alternating-usual} 참조). $U_{i_0 \ldots i_p}$는
항상 아핀이므로 $\mathcal{F}(U_{i_0 \ldots i_p})$는 $A$-평탄하다.
따라서
$\check{\mathcal{C}}_{alt}^\bullet(\mathcal{U}, \mathcal{F})$
는 평탄한 $A$-가군의 유계 복합체이고, More on Algebra, Lemma
\ref{more-algebra-lemma-derived-tor-quasi-isomorphism}에 의해 K-평탄이다.
\end{proof}

\noindent
$X, Y, S, a, b, p, q, f$를 이 절의 도입에서와 같이 두자.
$\mathcal{F}$를 $\mathcal{O}_X$-가군으로, $\mathcal{G}$를
$\mathcal{O}_Y$-가군으로 하자. $A = \Gamma(S, \mathcal{O}_S)$로 놓고
다음 사상을 생각하자.
\begin{equation}
\label{perfect-equation-kunneth-single-sheaves}
R\Gamma(X, \mathcal{F})
\otimes_A^\mathbf{L}
R\Gamma(Y, \mathcal{G})
\longrightarrow
R\Gamma(X \times_S Y,
p^*\mathcal{F}
\otimes_{\mathcal{O}_{X \times_S Y}}
q^*\mathcal{G})
\end{equation}
이는 $D(A)$ 안에 있다. 이 사상은 다음 당김 사상들
$R\Gamma(X, \mathcal{F}) \to R\Gamma(X \times_S Y, p^*\mathcal{F})$ 및
$R\Gamma(Y, \mathcal{G}) \to R\Gamma(X \times_S Y, q^*\mathcal{G})$,
Cohomology, Section \ref{cohomology-section-cup-product}에서 구성한
컵곱, 그리고 다음 표준 사상을 사용하여 구성된다.
$p^*\mathcal{F}
\otimes_{\mathcal{O}_{X \times_S Y}}^\mathbf{L}
q^*\mathcal{G}
\to
p^*\mathcal{F}
\otimes_{\mathcal{O}_{X \times_S Y}}
q^*\mathcal{G}$.

\begin{lemma}
\label{perfect-lemma-kunneth-single-sheaf}
위 상황에서 $S$가 아핀이고 $\mathcal{F}$와 $\mathcal{G}$가 $S$-평탄
준연접이며 $X$, $Y$가 준콤팩트이고 아핀 대각선을 가지면
(\ref{perfect-equation-kunneth-single-sheaves})의 사상은 동형사상이다.
\end{lemma}

\begin{proof}
독자에게 먼저 Varieties, Lemma \ref{varieties-lemma-kunneth}의 증명을
읽을 것을 강하게 권한다. 유한 아핀 열린 덮개
$\mathcal{U} : X = \bigcup_{i \in I} U_i$ 및
$\mathcal{V} : Y = \bigcup_{j \in J} V_j$로 잡자.
이로부터 다음 아핀 열린 덮개가 정해진다.
$\mathcal{W} : X \times_S Y = \bigcup_{(i, j) \in I \times J} U_i \times_S V_j$.
 $\mathcal{W}$는 $\text{pr}_1^{-1}\mathcal{U}$와
$\text{pr}_2^{-1}\mathcal{V}$의 세분임에 유의하라.
따라서 Cohomology, Section
\ref{cohomology-section-cech-cohomology-of-complexes}의 논의에 의해
다음 사상들을 얻는다.
$$
\check{\mathcal{C}}^\bullet(\mathcal{U}, \mathcal{F})
\to
\check{\mathcal{C}}^\bullet(\mathcal{W}, p^*\mathcal{F})
\quad\text{및}\quad
\check{\mathcal{C}}^\bullet(\mathcal{V}, \mathcal{G})
\to
\check{\mathcal{C}}^\bullet(\mathcal{W}, q^*\mathcal{G})
$$
이들은 호모토피까지 잘 정의되고 코호몰로지의 당김 사상과 양립한다.
Cohomology, Equation (\ref{cohomology-equation-needs-signs})에서
복합체의 사상을 구성했는데,
$$
\text{Tot}(
\check{\mathcal{C}}^\bullet(\mathcal{W}, p^*\mathcal{F})
\otimes_A
\check{\mathcal{C}}^\bullet(\mathcal{W}, q^*\mathcal{G}))
\longrightarrow
\check{\mathcal{C}}^\bullet(\mathcal{W},
p^*\mathcal{F} \otimes_{\mathcal{O}_{X \times_S Y}}
q^*\mathcal{G})
$$
이는 Cohomology, Lemma \ref{cohomology-lemma-diagrams-commute}에 의해
코호몰로지의 컵곱과 양립한다. 위의 사상들을 결합하면 다음 복합체의
사상을 얻는다.
\begin{equation}
\label{perfect-equation-kunneth-on-cech}
\text{Tot}(
\check{\mathcal{C}}^\bullet(\mathcal{U}, \mathcal{F})
\otimes_A
\check{\mathcal{C}}^\bullet(\mathcal{V}, \mathcal{G}))
\to
\check{\mathcal{C}}^\bullet(\mathcal{W},
p^*\mathcal{F}
\otimes_{\mathcal{O}_{X \times_S Y}}
q^*\mathcal{G})
\end{equation}
이것이 보조정리의 명제에 나오는 사상이라고 주장한다. 즉 이 화살표의
시작 대상과 끝 대상은 (\ref{perfect-equation-kunneth-single-sheaves})의 것과
같다. 실제로 Cohomology of Schemes, Lemma
\ref{coherent-lemma-quasi-coherent-affine-cohomology-zero} 및
Cohomology, Lemma \ref{cohomology-lemma-cech-complex-complex-computes}에
의해 다음 표준 사상들은
$$
\check{\mathcal{C}}^\bullet(\mathcal{U}, \mathcal{F})
\to
R\Gamma(X, \mathcal{F}),
\quad
\check{\mathcal{C}}^\bullet(\mathcal{V}, \mathcal{G})
\to
R\Gamma(Y, \mathcal{G})
$$
그리고
$$
\check{\mathcal{C}}^\bullet(\mathcal{W},
p^*\mathcal{F} \otimes_{\mathcal{O}_{X \times_S Y}} q^*\mathcal{G})
\to
R\Gamma(X \times_S Y,
p^*\mathcal{F} \otimes_{\mathcal{O}_{X \times_S Y}}
q^*\mathcal{G})
$$
동형사상이다. 한편 다음 복합체는
$\check{\mathcal{C}}^\bullet(\mathcal{U}, \mathcal{F})$
는 Lemma \ref{perfect-lemma-K-flat}에 의해 K-평탄이므로 다음을 얻는다.
$\text{Tot}(
\check{\mathcal{C}}^\bullet(\mathcal{U}, \mathcal{F})
\otimes_A
\check{\mathcal{C}}^\bullet(\mathcal{V}, \mathcal{G}))$
는 유도 텐서곱
$R\Gamma(X, \mathcal{F}) \otimes_A^\mathbf{L} R\Gamma(Y, \mathcal{G})$을
표현한다. 주장한 바와 같다.

\medskip\noindent
이제 (\ref{perfect-equation-kunneth-on-cech})가 준동형사상임을 보여야 한다.
이를 차원 이동으로 보인다. $d(\mathcal{F}) = \max \{d \mid
H^d(X, \mathcal{F}) \not = 0\}$로 놓자. $d(\mathcal{F}) > 0$이라 하고
$U = \coprod\nolimits_{i \in I} U_i$로 놓자. $I$가 유한하므로 이는
아핀 스킴이다. 각 $U_i$에서 포함 $U_i \to X$가 되는 사상
$j : U \to X$를 쓰자. $X$의 대각선이 아핀이므로 $j$도 아핀이다
(Morphisms, Lemma \ref{morphisms-lemma-affine-permanence} 참조).
따라서 $\mathcal{F}' = j_*j^*\mathcal{F}$는 $S$-평탄하다
(Morphisms, Lemma \ref{morphisms-lemma-pushforward-flat-affine} 참조).
또한 Cohomology of Schemes, Lemmas
\ref{coherent-lemma-relative-affine-cohomology}와
\ref{coherent-lemma-quasi-coherent-affine-cohomology-zero}를 결합하면
$d(\mathcal{F}') = 0$이다.
모든 $x \in X$에 대해 $\mathcal{F}_x \to \mathcal{F}'_x$는 직합인자의
포함이다. $x \in U_i$이면 $\mathcal{F}' \to (U_i \to X)_*
\mathcal{F}|_{U_i}$가 분할을 주기 때문이다. 따라서
$\mathcal{F} \to \mathcal{F}'$는 단사이고
$\mathcal{F}'' = \mathcal{F}'/\mathcal{F}$도 $S$-평탄하다. 평탄한
준연접 $\mathcal{O}_X$-가군의 다음 짧은 완전열은
$0 \to \mathcal{F} \to \mathcal{F}' \to \mathcal{F}'' \to 0$
복합체의 짧은 완전열을 만든다.
$$
0 \to
\text{Tot}(
\check{\mathcal{C}}^\bullet(\mathcal{U}, \mathcal{F})
\otimes_A
\check{\mathcal{C}}^\bullet(\mathcal{V}, \mathcal{G})) \to
\text{Tot}(
\check{\mathcal{C}}^\bullet(\mathcal{U}, \mathcal{F}')
\otimes_A
\check{\mathcal{C}}^\bullet(\mathcal{V}, \mathcal{G})) \to
\text{Tot}(
\check{\mathcal{C}}^\bullet(\mathcal{U}, \mathcal{F}'')
\otimes_A
\check{\mathcal{C}}^\bullet(\mathcal{V}, \mathcal{G})) \to 0
$$
그리고 복합체의 짧은 완전열
$$
0 \to
\check{\mathcal{C}}^\bullet(\mathcal{W},
p^*\mathcal{F}
\otimes_{\mathcal{O}_{X \times_S Y}}
q^*\mathcal{G}) \to
\check{\mathcal{C}}^\bullet(\mathcal{W},
p^*\mathcal{F}'
\otimes_{\mathcal{O}_{X \times_S Y}}
q^*\mathcal{G}) \to
\check{\mathcal{C}}^\bullet(\mathcal{W},
p^*\mathcal{F}''
\otimes_{\mathcal{O}_{X \times_S Y}}
q^*\mathcal{G}) \to 0
$$
더구나 이들 사이의 (\ref{perfect-equation-kunneth-on-cech}) 사상은 이 짧은
완전열들과 양립한다. 따라서 $\mathcal{F}'$와 $\mathcal{F}''$에 대해
(\ref{perfect-equation-kunneth-on-cech})가 동형사상임을 보이면 충분하다.
마지막으로 $d(\mathcal{F}'') < d(\mathcal{F})$이다. 이로써
$d(\mathcal{F}) = 0$인 경우로 환원된다.

\medskip\noindent
 $\mathcal{G}$에 대해서도 같은 방식으로 논하면 $\mathcal{F}$와
$\mathcal{G}$ 모두 $0$차에서만 0이 아닌 코호몰로지를 갖는다고 가정할 수
있다. 이는 $\Gamma(X, \mathcal{F})$가 다음 $K$-평탄 복합체와
준동형사상임을 뜻한다.
$\check{\mathcal{C}}^\bullet(\mathcal{U}, \mathcal{F})$
이는 차수 $\geq 0$에 놓인 $A$-가군 복합체이다. 따라서
$\Gamma(X, \mathcal{F})$는 평탄한 $A$-가군이다(이 가군에 대한 고차
Tor를 Čech 복합체와 텐서하여 계산할 수 있기 때문이다).
$V \subset Y$를 아핀 열린집합이라 하자. 다음 아핀 열린 덮개를 생각하자.
$\mathcal{U}_V : X \times_S V = \bigcup_{i \in I} U_i \times_S V$.
다음 등식은 자명하다.
$$
\check{\mathcal{C}}^\bullet(\mathcal{U}, \mathcal{F})
\otimes_A \mathcal{G}(V) =
\check{\mathcal{C}}^\bullet(\mathcal{U}_V,
p^*\mathcal{F} \otimes_{\mathcal{O}_{X \times Y}}
q^*\mathcal{G})
$$
(복합체의 등식). $A$ 위에서 $\mathcal{G}(V)$가 평탄하므로 다음 사상은
$\Gamma(X, \mathcal{F}) \otimes_A \mathcal{G}(V) \to
\check{\mathcal{C}}^\bullet(\mathcal{U}, \mathcal{F})
\otimes_A \mathcal{G}(V)$는 준동형사상이다. 다음 프리셰프의 층화가
$V \mapsto \check{\mathcal{C}}^\bullet(\mathcal{U}_V,
p^*\mathcal{F} \otimes_{\mathcal{O}_{X \times Y}}
q^*\mathcal{G})$ represents
$Rq_*(p^*\mathcal{F} \otimes_{\mathcal{O}_{X \times Y}} q^*\mathcal{G})$를
표현하므로, Cohomology of Schemes, Lemma
\ref{coherent-lemma-separated-case-relative-cech}에 의해 다음을 얻는다.
$$
Rq_*(p^*\mathcal{F} \otimes_{\mathcal{O}_{X \times Y}} q^*\mathcal{G})
\cong
\Gamma(X, \mathcal{F}) \otimes_A \mathcal{G}
$$
 $Y$ 위에서 성립하며, 오른쪽 표기는 다음 가군을 뜻한다.
$$
b^*\widetilde{\Gamma(X, \mathcal{F})} \otimes_{\mathcal{O}_Y} \mathcal{G}
$$
 $q$에 대한 Leray 스펙트럴 열을 사용하면
$$
H^n(X \times_S Y, p^*\mathcal{F} \otimes_{\mathcal{O}_{X \times Y}}
q^*\mathcal{G}) =
H^n(Y,
b^*\widetilde{\Gamma(X, \mathcal{F})} \otimes_{\mathcal{O}_Y} \mathcal{G})
$$
사상 $b : Y \to S = \Spec(A)$에 Lemma
\ref{perfect-lemma-cohomology-base-change}를 적용하고 $\Gamma(X, \mathcal{F})$가
$A$-평탄임을 사용하면 다음을 얻는다.
$H^n(X \times_S Y, p^*\mathcal{F} \otimes_{\mathcal{O}_{X \times Y}}
q^*\mathcal{G})$는 $n > 0$이면 0이고, $n = 0$이면
$H^0(X, \mathcal{F}) \otimes_A H^0(Y, \mathcal{G})$와 동형이다.
여기서 $\mathcal{G}$가 $0$차에서만 코호몰로지를 갖는다는 것도 사용했다.
이로써 증명이 끝난다(단, 이 동형사상이 $0$차에서 실제로 컵곱으로
주어지는지 확인해야 하지만 그 검증은 생략한다).
\end{proof}

\begin{remark}
\label{perfect-remark-annoying-compatibility}
 $S = \Spec(A)$를 아핀 스킴이라 하자. $a : X \to S$와
$b : Y \to S$를 스킴의 사상이라 하자. $\mathcal{F}$, $\mathcal{G}$를
준연접 $\mathcal{O}_X$-가군으로, $\mathcal{E}$를 준연접
$\mathcal{O}_Y$-가군으로 하자. 다음을 생각하자.
$\xi \in H^i(X, \mathcal{G})$와 그 당김
$p^*\xi \in H^i(X \times_S Y, p^*\mathcal{G})$를 잡자.
그러면 다음 도식은 가환한다.
$$
\xymatrix{
R\Gamma(X, \mathcal{F})[-i] \otimes_A^\mathbf{L} R\Gamma(Y, \mathcal{E})
\ar[d] \ar[rr]_-{\xi \otimes \text{id}} & &
R\Gamma(X, \mathcal{G} \otimes_{\mathcal{O}_X} \mathcal{F})
\otimes_A^\mathbf{L} R\Gamma(Y, \mathcal{E}) \ar[d] \\
R\Gamma(X \times_S Y, p^*\mathcal{F} \otimes q^*\mathcal{E})[-i]
\ar[rr]^-{p^*\xi} & &
R\Gamma(X \times_S Y,
p^*(\mathcal{G} \otimes_{\mathcal{O}_X} \mathcal{F}) \otimes q^*\mathcal{E})
}
$$
표시하지 않은 텐서곱은 $\mathcal{O}_{X \times_S Y}$ 위에서 취한다.
수평 화살표는 Cohomology, Remark
\ref{cohomology-remark-cup-with-element-map-total-cohomology}에서 오고,
수직 화살표는 (\ref{perfect-equation-kunneth-global})이므로
$X \times_S Y$에서 당긴 뒤 컵곱을 취하는 사상이다.
이 도식이 가환하는 것은 전역 컵곱($X \times_S Y$에서 층
$p^*\mathcal{G}$, $p^*\mathcal{F}$, $q^*\mathcal{E}$에 대한)이
결합법칙을 만족하기 때문이다(Cohomology, Lemma
\ref{cohomology-lemma-cup-product-associative} 참조).
\end{remark}







\section{K\"unneth 공식 III}
\label{perfect-section-kunneth-complexes}

\noindent
$X, Y, S, a, b, p, q, f$를 Section \ref{perfect-section-kunneth}의 도입에서와
같이 두자. 이 절에서는 $\mathcal{O}_X$-가군 $\mathcal{F}$와
$\mathcal{O}_Y$-가군 $\mathcal{G}$가 주어졌을 때 다음과 같이 놓는다.
$$
\mathcal{F} \boxtimes \mathcal{G} =
p^*\mathcal{F} \otimes_{\mathcal{O}_{X \times_S Y}} q^*\mathcal{G}
$$
뒤의 절에서와 달리 여기서는 유도되지 않은 텐서곱을 취한다는 점에
유의하라.

\medskip\noindent
 $X$ 위에서 $\mathcal{F}^\bullet$를 아벨 군 층의 복합체로 하자.
그 항들은 준연접 $\mathcal{O}_X$-가군이고 미분
$d^i_\mathcal{F} : \mathcal{F}^i \to \mathcal{F}^{i + 1}$
은 $X/S$ 위의 유한 차수 미분 연산자라고 하자(Morphisms, Section
\ref{morphisms-section-differential-operators} 참조). 마찬가지로 $Y$ 위에서
$\mathcal{G}^\bullet$를 아벨 군 층의 복합체로 하되, 그 항들은 준연접
$\mathcal{O}_Y$-가군이고 미분
$d^j_\mathcal{G} : \mathcal{G}^j \to \mathcal{G}^{j + 1}$
은 $Y/S$ 위의 유한 차수 미분 연산자라고 하자. Morphisms, Lemma
\ref{morphisms-lemma-base-change-differential-operators}의 구성을 적용하면
다음 이중 복합체를 얻는다.
$$
\xymatrix{
\ldots &
\ldots &
\ldots &
\ldots \\
\ldots \ar[r] &
\mathcal{F}^i \boxtimes \mathcal{G}^{j + 1}
\ar[r]^{d_1^{i, j + 1}} \ar[u] &
\mathcal{F}^{i + 1} \boxtimes \mathcal{G}^{j + 1} \ar[r] \ar[u] &
\ldots \\
\ldots \ar[r] &
\mathcal{F}^i \boxtimes \mathcal{G}^j
\ar[r]^{d_1^{i, j}} \ar[u]^{d_2^{i, j}} &
\mathcal{F}^{i + 1} \boxtimes \mathcal{G}^j \ar[r] \ar[u]_{d_2^{i + 1, j}} &
\ldots \\
\ldots &
\ldots \ar[u] &
\ldots \ar[u] &
\ldots
}
$$
이는 사상들이 $X \times_S Y / S$ 위의 유한 차수 미분 연산자인
준연접 가군의 이중 복합체이다. Morphisms, Remark
\ref{morphisms-remark-base-change-differential-operators}와 Homology,
Example \ref{homology-example-double-complex-as-tensor-product-of}의
논의를 참조하라. 명시적으로 다음과 같이 놓는다.
$$
d_1^{i, j} = d^i_\mathcal{F} \boxtimes 1
\quad\text{및}\quad
d_2^{i, j} = 1 \boxtimes d^j_\mathcal{G}
$$
아래 논의에서 다음 표기는
$$
\text{Tot}(\mathcal{F}^\bullet \boxtimes \mathcal{G}^\bullet)
$$
이 이중 복합체에 대응하는 전체 복합체를 뜻한다. 이 복합체의 항들은
준연접 $\mathcal{O}_{X \times_S Y}$-가군이고, 미분은
$X \times_S Y / S$ 위의 유한 차수 미분 연산자이다.

\medskip\noindent
위 상황에는 ``상대 컵곱'' 사상이 존재한다.
\begin{equation}
\label{perfect-equation-relative-de-rham-kunneth}
Ra_*(\mathcal{F}^\bullet)
\otimes_{\mathcal{O}_S}^\mathbf{L}
Rb_*(\mathcal{G}^\bullet)
\longrightarrow
Rf_*\left(\text{Tot}(\mathcal{F}^\bullet \boxtimes \mathcal{G}^\bullet)\right)
\end{equation}
즉 다음 사상들을 결합하여 이 사상을 구성할 수 있다.
\begin{enumerate}
\item $Ra_*(\mathcal{F}^\bullet) \to Rf_*(p^{-1}\mathcal{F}^\bullet)$,
\item $Rb_*(\mathcal{G}^\bullet) \to Rf_*(q^{-1}\mathcal{G}^\bullet)$,
\item
$Rf_*(p^{-1}\mathcal{F}^\bullet)
\otimes_{\mathcal{O}_S}^\mathbf{L}
Rf_*(q^{-1}\mathcal{G}^\bullet)
\to
Rf_*(p^{-1}\mathcal{F}^\bullet \otimes_{f^{-1}\mathcal{O}_S}^\mathbf{L}
q^{-1}\mathcal{G}^\bullet)$,
\item
$p^{-1}\mathcal{F}^\bullet \otimes_{f^{-1}\mathcal{O}_S}^\mathbf{L}
q^{-1}\mathcal{G}^\bullet \to
\text{Tot}(p^{-1}\mathcal{F}^\bullet \otimes_{f^{-1}\mathcal{O}_S}
q^{-1}\mathcal{G}^\bullet)$
\item
$\text{Tot}(p^{-1}\mathcal{F}^\bullet \otimes_{f^{-1}\mathcal{O}_S}
q^{-1}\mathcal{G}^\bullet) \to \text{Tot}(\mathcal{F}^\bullet \boxtimes
\mathcal{G}^\bullet)$.
\end{enumerate}
사상 (1), (2)는 당김 사상이고, 사상 (3)은 상대 컵곱이다
(Cohomology, Remark \ref{cohomology-remark-cup-product} 참조).
사상 (4)는 유도 텐서곱과 유도되지 않은 텐서곱을 비교하며,
사상 (5)는 다음 자명한 사상들로 주어진다.
$p^{-1}\mathcal{F}^i \otimes_{f^{-1}\mathcal{O}_S} q^{-1}\mathcal{G}^j
\to \mathcal{F}^i \boxtimes \mathcal{G}^j$로, 바탕 이중 복합체 위에서
정의된다.

\medskip\noindent
 $A = \Gamma(S, \mathcal{O}_S)$로 놓자. ``대역 컵곱'' 사상이 존재한다.
\begin{equation}
\label{perfect-equation-de-rham-kunneth}
R\Gamma(X, \mathcal{F}^\bullet)
\otimes_A^\mathbf{L}
R\Gamma(Y, \mathcal{G}^\bullet)
\longrightarrow
R\Gamma(X \times_S Y,
\text{Tot}(\mathcal{F}^\bullet \boxtimes \mathcal{G}^\bullet))
\end{equation}
이는 $D(A)$에서 성립한다. 위의 상대 컵곱과 유사하게 다음 사상들을 사용하여 구성된다.
\begin{enumerate}
\item $R\Gamma(X, \mathcal{F}^\bullet) \to
R\Gamma(X \times_S Y, p^{-1}\mathcal{F}^\bullet)$
\item $R\Gamma(Y, \mathcal{G}^\bullet) \to
R\Gamma(X \times_S Y, q^{-1}\mathcal{G}^\bullet)$,
\item $R\Gamma(X \times_S Y, p^{-1}\mathcal{F}^\bullet)
\otimes_A^\mathbf{L} R\Gamma(X \times_S Y, q^{-1}\mathcal{G}^\bullet) \to
R\Gamma(X \times_S Y,
p^{-1}\mathcal{F}^\bullet \otimes_{f^{-1}\mathcal{O}_S}^\mathbf{L}
q^{-1}\mathcal{G}^\bullet)$,
\item
$p^{-1}\mathcal{F}^\bullet \otimes_{f^{-1}\mathcal{O}_S}^\mathbf{L}
q^{-1}\mathcal{G}^\bullet \to
\text{Tot}(p^{-1}\mathcal{F}^\bullet \otimes_{f^{-1}\mathcal{O}_S}
q^{-1}\mathcal{G}^\bullet)$
\item
$\text{Tot}(p^{-1}\mathcal{F}^\bullet \otimes_{f^{-1}\mathcal{O}_S}
q^{-1}\mathcal{G}^\bullet) \to \text{Tot}(\mathcal{F}^\bullet \boxtimes
\mathcal{G}^\bullet)$.
\end{enumerate}
여기서 사상 (1), (2)는 당김 사상이고, 사상 (3)은 Cohomology, Section
\ref{cohomology-section-cup-product}에서 구성한 컵곱이다.
사상 (4), (5)는 앞 단락에서 설명한 것과 같다.

\begin{lemma}
\label{perfect-lemma-kunneth-special}
위 상황에서 다음 가정들이 성립하면 컵곱
(\ref{perfect-equation-de-rham-kunneth})은 $D(A)$에서 동형사상이다.
\begin{enumerate}
\item $S = \Spec(A)$는 아핀이다.
\item $X$, $Y$는 준콤팩트이고 아핀 대각선을 갖는다.
\item $\mathcal{F}^\bullet$는 유계이다.
\item $\mathcal{G}^\bullet$는 아래로 유계이다.
\item $\mathcal{F}^n$은 $S$-평탄이다.
\item $\mathcal{G}^m$은 $S$-평탄이다.
\end{enumerate}
\end{lemma}

\begin{proof}
Morphisms, Remark \ref{morphisms-remark-base-change-differential-operators}에서
도입한 표기 $\mathcal{A}_{X/S}$와 $\mathcal{A}_{Y/S}$를 사용한다.
다음과 같은 복합체의 사상들이
$$
\mathcal{F}_1^\bullet \to
\mathcal{F}_2^\bullet \to
\mathcal{F}_3^\bullet \to
\mathcal{F}_1^\bullet[1]
$$
$\mathcal{A}_{X/S}$의 범주에 있다고 하자. 컵곱의 함자성에 의해
다음 가환 도식을 얻는다.
$$
\xymatrix{
R\Gamma(X, \mathcal{F}_1^\bullet)
\otimes_A^\mathbf{L}
R\Gamma(Y, \mathcal{G}^\bullet)
\ar[r] \ar[d] &
R\Gamma(X \times_S Y,
\text{Tot}(\mathcal{F}_1^\bullet \boxtimes \mathcal{G}^\bullet)) \ar[d] \\
R\Gamma(X, \mathcal{F}_2^\bullet)
\otimes_A^\mathbf{L}
R\Gamma(Y, \mathcal{G}^\bullet)
\ar[r] \ar[d] &
R\Gamma(X \times_S Y,
\text{Tot}(\mathcal{F}_2^\bullet \boxtimes \mathcal{G}^\bullet)) \ar[d] \\
R\Gamma(X, \mathcal{F}_3^\bullet)
\otimes_A^\mathbf{L}
R\Gamma(Y, \mathcal{G}^\bullet)
\ar[r] \ar[d] &
R\Gamma(X \times_S Y,
\text{Tot}(\mathcal{F}_3^\bullet \boxtimes \mathcal{G}^\bullet)) \ar[d] \\
R\Gamma(X, \mathcal{F}_1^\bullet[1])
\otimes_A^\mathbf{L}
R\Gamma(Y, \mathcal{G}^\bullet)
\ar[r] &
R\Gamma(X \times_S Y,
\text{Tot}(\mathcal{F}_1^\bullet[1] \boxtimes \mathcal{G}^\bullet))
}
$$
원래의 사상들이 $\mathcal{A}_{X/S}$의 호모토피 범주에서 완전 삼각형을
이루면, 이 도식의 열들은 $D(A)$에서 완전 삼각형을 이룬다.

\medskip\noindent
보조정리의 상황에서 $n < i$이면 $\mathcal{F}^n = 0$이라고 하자.
그러면 항별로 분할되는 다음 복합체의 짧은 완전열을 생각할 수 있다.
$$
0 \to \sigma_{\geq i + 1}\mathcal{F}^\bullet \to
\mathcal{F}^\bullet \to \mathcal{F}^i[-i] \to 0
$$
여기서 절단은 Homology, Section \ref{homology-section-truncations}에서와
같다. 이로부터 다음 완전 삼각형이 생긴다.
$$
\sigma_{\geq i + 1}\mathcal{F}^\bullet \to
\mathcal{F}^\bullet \to
\mathcal{F}^i[-i] \to
(\sigma_{\geq i + 1}\mathcal{F}^\bullet)[1]
$$
이는 $\mathcal{A}_{X/S}$의 호모토피 범주에서 성립하며, 마지막 화살표는
경계 사상 $\mathcal{F}^i \to \mathcal{F}^{i + 1}$로 주어진다.
앞의 논의에 의해 $\mathcal{F}^i[-i]$와
$\sigma_{\geq i + 1}\mathcal{F}^\bullet$에 대해 보조정리를 증명하면
충분하다. $\sigma_{\geq i + 1}\mathcal{F}^\bullet$는 0이 아닌 항이 더
적으므로 귀납법에 의해, $\mathcal{F}^\bullet$가 한 차수에서만 0이
아닌 경우를 증명할 수 있으면 보조정리가 따른다. 따라서
$\mathcal{F}^\bullet$가 한 차수에서만 0이 아니라고 가정해도 된다.

\medskip\noindent
 $\mathcal{F}^\bullet$가 $0$차에 놓인 $S$-평탄 준연접
$\mathcal{O}_X$-가군 $\mathcal{F}$를 가지며 다른 차수에서는 0인 복합체라고
가정하자. 예를 들어 Lemma \ref{perfect-lemma-cohomology-de-rham-base-change}에
의해 $R\Gamma(X, \mathcal{F})$는 유한 Tor 차원을 갖는다. 그 Tor 진폭이
$[i, j]$라고 하자. $N \gg 0$를 택하고 다음 완전 삼각형을 생각하자.
$$
\sigma_{\geq N + 1}\mathcal{G}^\bullet \to
\mathcal{G}^\bullet \to
\sigma_{\leq N}\mathcal{G}^\bullet \to
(\sigma_{\geq N + 1}\mathcal{G}^\bullet)[1]
$$
이는 $\mathcal{A}_{Y/S}$의 호모토피 범주에서 성립한다. 이제 다음 두
복합체가
$$
R\Gamma(X, \mathcal{F})
\otimes_A^\mathbf{L}
R\Gamma(Y, \sigma_{\geq N + 1}\mathcal{G}^\bullet)
\quad\text{및}\quad
R\Gamma(X \times_S Y,
\text{Tot}(\mathcal{F} \boxtimes \sigma_{\geq N + 1}\mathcal{G}^\bullet))
$$
는 차수 $\leq N+i$에서 코호몰로지가 소멸한다. 따라서 위에서 제시한
논의를 사용하면 특정 차수에서 주장을 증명할 때 $\mathcal{G}^\bullet$가
유계라고 가정할 수 있다. 위의 논의를 한 번 더 반복하면
$\mathcal{G}^\bullet$도 한 차수에서만 0이 아니라고 가정할 수 있다.
이 경우는 Lemma \ref{perfect-lemma-kunneth-single-sheaf}가 처리한다.
\end{proof}







\section{K\"unneth 공식 (Ext)}
\label{perfect-section-kunneth-Ext}

\noindent
스킴들의 데카르트 도식을 생각하자.
$$
\xymatrix{
& X \times_S Y \ar[ld]^p \ar[rd]_q \ar[dd]^f \\
X \ar[rd]_a & & Y \ar[ld]^b \\
& S
}
$$
$K \in D(\mathcal{O}_X)$ 및 $M \in D(\mathcal{O}_Y)$에 대해
이 절에서는 다음을 정의하자.
$$
K \boxtimes M =
Lp^*K \otimes_{\mathcal{O}_{X \times_S Y}}^\mathbf{L} Lq^*M
$$
다음과 같은 표준 사상이 존재한다고 하자.
\begin{equation}
\label{perfect-equation-kunneth-ext}
Ra_*R\SheafHom(K, K')
\otimes_{\mathcal{O}_S}^\mathbf{L}
Rb_*R\SheafHom(M, M')
\longrightarrow
Rf_*(R\SheafHom(K \boxtimes M, K' \boxtimes M'))
\end{equation}
이는 $K, K' \in D(\mathcal{O}_X)$ 및 $M, M' \in D(\mathcal{O}_Y)$에
대해 성립한다. 구체적으로 다음 사상의 수반으로 주어진 사상을 취할 수 있다.
$$
\begin{matrix}
Lf^*\left(Ra_*R\SheafHom(K, K')
\otimes_{\mathcal{O}_S}^\mathbf{L}
Rb_* R\SheafHom(M, M')\right) = \\
Lf^* Ra_* R\SheafHom(K, K')
\otimes_{\mathcal{O}_{X \times_S Y}}^\mathbf{L}
Lf^* Rb_* R\SheafHom(M, M') = \\
Lp^* La^* Ra_* R\SheafHom(K, K')
\otimes_{\mathcal{O}_{X \times_S Y}}^\mathbf{L}
Lq^* Lb^* Rb_* R\SheafHom(M, M') \to \\
Lp^* R\SheafHom(K, K')
\otimes_{\mathcal{O}_{X \times_S Y}}^\mathbf{L}
Lq^* R\SheafHom(M, M') \to \\
R\SheafHom(Lp^*K, Lp^*K')
\otimes_{\mathcal{O}_{X \times_S Y}}^\mathbf{L}
R\SheafHom(Lq^*M, Lq^*M') \to \\
R\SheafHom(K \boxtimes M, K' \boxtimes M')
\end{matrix}
$$
여기서 첫 번째 등식은 당김과 텐서곱의 호환성에 따른 것으로,
Cohomology, Lemma \ref{cohomology-lemma-pullback-tensor-product}에
의해 주어진다. 두 번째 등식은 $f = a \circ p = b \circ q$ 및
당김 사상의 합성에 따른 것으로,
Cohomology, Lemma \ref{cohomology-lemma-derived-pullback-composition}에
의해 주어진다. 첫 번째 화살표는 직상과 당김이 수반이라는 사실에 따른
수반 사상 $La^* Ra_* \to \text{id}$ 및
$Lb^* Rb_* \to \text{id}$로 주어지며,
Cohomology, Lemma \ref{cohomology-lemma-adjoint}를 사용한다.
두 번째 화살표는 Cohomology, Remark
\ref{cohomology-remark-prepare-fancy-base-change}에 의해 주어진다.
세 번째이자 마지막 화살표는 Cohomology, Remark
\ref{cohomology-remark-tensor-internal-hom}에 의해 주어진다.
이 사상의 간단한 특수한 경우가 다음 결과이다.

\begin{lemma}
\label{perfect-lemma-kunneth-Ext}
위 상황에서 $a$와 $b$가 준콤팩트이고 준분리이며 $X$와 $Y$가
$S$ 위에서 Tor 독립이라고 하자. $K$가 완전하고
$K' \in D_\QCoh(\mathcal{O}_X)$, $M$이 완전하며
$M' \in D_\QCoh(\mathcal{O}_Y)$이면,
(\ref{perfect-equation-kunneth-ext})는 동형사상이다.
\end{lemma}

\begin{proof}
이 경우
$R\SheafHom(K, K') = K' \otimes^\mathbf{L} K^\vee$,
$R\SheafHom(M, M') = M' \otimes^\mathbf{L} M^\vee$이고,
$$
R\SheafHom(K \boxtimes M, K' \boxtimes M') =
(K' \otimes^\mathbf{L} K^\vee) \boxtimes
(M' \otimes^\mathbf{L} M^\vee)
$$
이는 Cohomology, Lemma \ref{cohomology-lemma-dual-perfect-complex}를
보면 알 수 있으며, 완전성은 당김과 텐서곱에 의해 보존된다는 사실도
사용한다. 따라서 이 경우는 Lemma \ref{perfect-lemma-kunneth}에서 따른다.
이 식별 아래에서 얻는 사상이 원래의 사상과 같다는 확인은 생략한다.
\end{proof}








\section{코호몰로지와 기저변환 V}
\label{perfect-section-cohomology-base-change}

\noindent
Section \ref{perfect-section-cohomology-and-base-change-perfect}에서는 사상들이
Tor 독립일 때 기저변환 정리가 성립함을 보았다. 이러한 조건이 없으면
아핀인 경우에도 기저변환 정리가 성립할 수 없다. More on Algebra,
Section \ref{more-algebra-section-tor-independence}를 보라.
이 절에서는 ‘한 번에 하나의 복합체’에 대해 기저변환 결과를 얻을 수
있는 조건을 분석한다.

\medskip\noindent
이를 위해 다음과 같은 가환 도식이 주어졌다고 하자.
$$
\xymatrix{
X' \ar[r]_{g'} \ar[d]_{f'} &
X \ar[d]^f \\
S' \ar[r]^g &
S
}
$$
이는 스킴들의 도식이며(보통 데카르트라고 가정한다).
$K \in D_\QCoh(\mathcal{O}_X)$를 잡고
$L(g')^*K \to K'$를 $D_\QCoh(\mathcal{O}_{X'})$에서의 사상이라 하자.
점 $x' \in X'$에 대해 $x = g'(x') \in X$,
$s' = f'(x') \in S'$, $s = f(x) = g(s')$로 두자.
그러면 다음 사상들을 생각할 수 있다.
$$
K_x \otimes_{\mathcal{O}_{S, s}}^\mathbf{L} \mathcal{O}_{S', s'} \to
K_x \otimes_{\mathcal{O}_{X, x}}^\mathbf{L} \mathcal{O}_{X', x'} \to
K'_{x'}
$$
첫 번째 화살표는 More on Algebra, Equation
(\ref{more-algebra-equation-comparison-map})이고, 두 번째 화살표는
$(L(g')^*K)_{x'} =
K_x \otimes_{\mathcal{O}_{X, x}}^\mathbf{L} \mathcal{O}_{X', x'}$
및 주어진 사상 $L(g')^*K \to K'$에서 온다.
각 $i \in \mathbf{Z}$에 대해 다음 가군 위에
$\mathcal{O}_{X, x} \otimes_{\mathcal{O}_{S, s}} \mathcal{O}_{S', s'}$-module
다음과 같은 구조를 얻는다.
$H^i(K_x \otimes_{\mathcal{O}_{S, s}}^\mathbf{L} \mathcal{O}_{S', s'})$.
이를 모두 합치면 다음 표준 사상을 얻는다.
\begin{equation}
\label{perfect-equation-bc}
H^i(K_x \otimes_{\mathcal{O}_{S, s}}^\mathbf{L} \mathcal{O}_{S', s'})
\otimes_{(\mathcal{O}_{X, x} \otimes_{\mathcal{O}_{S, s}} \mathcal{O}_{S', s'})}
\mathcal{O}_{X', x'}
\longrightarrow
H^i(K'_{x'})
\end{equation}
이는 $\mathcal{O}_{X', x'}$-가군의 사상이다.

\begin{lemma}
\label{perfect-lemma-single-complex-base-change-condition}
다음과 같이 두자.
$$
\xymatrix{
X' \ar[r]_{g'} \ar[d]_{f'} &
X \ar[d]^f \\
S' \ar[r]^g &
S
}
$$
스킴들의 데카르트 도식이라고 하자. $K \in D_\QCoh(\mathcal{O}_X)$를
잡고 $L(g')^*K \to K'$를 $D_\QCoh(\mathcal{O}_{X'})$에서의 사상이라 하자.
다음 조건들은 서로 동치이다.
\begin{enumerate}
\item 임의의 $x' \in X'$ 및 $i \in \mathbf{Z}$에 대해
(\ref{perfect-equation-bc})의 사상이 동형사상이다.
\item $U \subset X$, $V' \subset S'$가 아핀 열린집합이고 둘 다
아핀 열린집합 $V \subset S$로 사상되며 $U' = V' \times_V U$일 때,
다음 합성
$$
R\Gamma(U, K) \otimes_{\mathcal{O}_S(U)}^\mathbf{L} \mathcal{O}_{S'}(V')
\to
R\Gamma(U, K) \otimes_{\mathcal{O}_X(U)}^\mathbf{L} \mathcal{O}_{X'}(U')
\to
R\Gamma(U', K')
$$
은 $D(\mathcal{O}_{S'}(V'))$에서 동형사상이다.
\item (2)와 같은 네 쌍 $U_i, V_i', V_i, U_i'$의 집합 $I$가 존재하며,
$i \in I$이고 $X' = \bigcup U'_i$이다.
\end{enumerate}
\end{lemma}

\begin{proof}
조건 (2)의 두 번째 화살표는 다음 등식과
$$
R\Gamma(U, K) \otimes_{\mathcal{O}_X(U)}^\mathbf{L} \mathcal{O}_{X'}(U') =
R\Gamma(U', L(g')^*K)
$$
Lemma \ref{perfect-lemma-quasi-coherence-pullback} 및 주어진 화살표
$L(g')^*K \to K'$에서 온다. 조건 (2)의 첫 번째 화살표는
More on Algebra, Equation (\ref{more-algebra-equation-comparison-map})이다.
(2)가 (3)을 함의하는 것은 명백하다. (1)은 $X'$ 위에서 국소적임에
유의하자.
따라서 $X$, $S$, $S'$, $X'$가 모두 아핀일 때 (1)이 다음 조건과
동치임을 보이면 충분하다.
$$
R\Gamma(X, K) \otimes_{\mathcal{O}_S(S)}^\mathbf{L} \mathcal{O}_{S'}(S')
\to
R\Gamma(X, K) \otimes_{\mathcal{O}_X(X)}^\mathbf{L} \mathcal{O}_{X'}(X')
\to
R\Gamma(X', K')
$$
은 $D(\mathcal{O}_{S'}(S'))$에서 동형사상이다. 이제
$S = \Spec(R)$, $X = \Spec(A)$, $S' = \Spec(R')$, $X' = \Spec(A')$라
하고, $K$가 $A$-가군의 복합체 $M^\bullet$에 대응하며
$K'$가 $A'$-가군의 복합체 $N^\bullet$에 대응한다고 하자.
$A' = A \otimes_R R'$임에 유의하자. 위 조건은 다음 합성이
$$
M^\bullet \otimes_R^\mathbf{L} R' \to
M^\bullet \otimes_A^\mathbf{L} A' \to
N^\bullet
$$
은 $D(R')$에서 동형사상이다. 이는 모든 $i \in \mathbf{Z}$에 대해
다음 사상이
$$
H^i(M^\bullet \otimes_R^\mathbf{L} R') \to
H^i(M^\bullet \otimes_A^\mathbf{L} A') \to
H^i(N^\bullet)
$$
동형사상이라는 것과 동치이다. 이는 $A \otimes_R R'$-가군,
즉 $A'$-가군의 사상임에 유의하자. 한편 (1)은 서로 호환되는 소아이디얼
$\mathfrak q' \subset A'$, $\mathfrak q \subset A$,
$\mathfrak p' \subset R'$, $\mathfrak p \subset R$에 대해 다음 합성이
$$
H^i(M^\bullet_\mathfrak q \otimes_{R_\mathfrak p}^\mathbf{L} R'_{\mathfrak p'})
\otimes_{(A_\mathfrak q \otimes_{R_\mathfrak p} R'_{\mathfrak p'})}
A'_{\mathfrak q'} \to
H^i(M^\bullet_{\mathfrak q}
\otimes_{A_\mathfrak q}^\mathbf{L} A'_{\mathfrak q'})
\to H^i(N^\bullet_{\mathfrak q'})
$$
동형사상이라는 조건이다. 다음 등식에 의해
$$
H^i(M^\bullet_\mathfrak q \otimes_{R_\mathfrak p}^\mathbf{L} R'_{\mathfrak p'})
\otimes_{(A_\mathfrak q \otimes_{R_\mathfrak p} R'_{\mathfrak p'})}
A'_{\mathfrak q'} =
H^i(M^\bullet \otimes_R^\mathbf{L} R') \otimes_{A'} A'_{\mathfrak q'}
$$
는 $\mathfrak q'$에서의 국소화이므로, 두 조건은
Algebra, Lemma \ref{algebra-lemma-characterize-zero-local}에 의해
서로 동치이다.
\end{proof}

\begin{lemma}
\label{perfect-lemma-single-complex-base-change}
다음 도식을 생각하자.
$$
\xymatrix{
X' \ar[r]_{g'} \ar[d]_{f'} &
X \ar[d]^f \\
S' \ar[r]^g &
S
}
$$
이는 스킴들의 데카르트 도식이다. $K \in D_\QCoh(\mathcal{O}_X)$를
잡고 $L(g')^*K \to K'$를 $D_\QCoh(\mathcal{O}_{X'})$에서의 사상이라 하자.
다음 조건을 가정하자.
\begin{enumerate}
\item Lemma \ref{perfect-lemma-single-complex-base-change-condition}의
동치인 조건들이 성립한다.
\item $f$가 준콤팩트이고 준분리이다.
\end{enumerate}
그러면 합성 $Lg^*Rf_*K \to Rf'_*L(g')^*K \to Rf'_*K'$는
동형사상이다.
\end{lemma}

\begin{proof}
Lemma \ref{perfect-lemma-compare-base-change}의 증명과 같은 방법으로 이를
증명할 수도 있지만, 여기서는 귀납 원리와 상대 Mayer--Vietoris를
사용하여 증명한다.

\medskip\noindent
이 사상이 동형사상임을 확인하기 위해 $S'$ 위에서 국소적으로 작업할
수 있다. 따라서 $g : S' \to S$가 아핀 스킴들의 사상이라고 가정해도
된다. 특히 $X$는 준콤팩트 준분리 스킴이다.
Cohomology of Schemes, Lemma \ref{coherent-lemma-induction-principle}의
귀납 원리를 사용하여, 임의의 준콤팩트 열린집합 $U \subset X$에 대해
구성한 사상 $Lg^*R(U \to S)_*K|_U \to R(U' \to S')_*K'|_{U'}$가
동형사상임을 보인다. 여기서 $U' = (g')^{-1}(U)$이다.

\medskip\noindent
$U \subset X$가 아핀 열린집합이면 가정에 의해 결과가 성립한다.
이는 Lemma \ref{perfect-lemma-single-complex-base-change-condition}의 (2)와
Lemma \ref{perfect-lemma-affine-compare-bounded} 및
\ref{perfect-lemma-quasi-coherence-pullback}이 제공하는 대수적 해석을
참조하면 된다.

\medskip\noindent
귀납 단계이다. $X = U \cup V$가 열린 덮개이고 $U$, $V$, $U \cap V$가
준콤팩트이며 $U$, $V$, $U \cap V$에 대해 결과가 성립한다고 하자.
$a = f|_U$, $b = f|_V$, $c = f|_{U \cap V}$로 두자.
$a' : U' \to S'$, $b' : V' \to S'$, $c' : U' \cap V' \to S'$를
$a$, $b$, $c$의 기저변환이라 하자.
상대 Mayer--Vietoris에서 얻는 완전 삼각형
(Cohomology, Lemma \ref{cohomology-lemma-unbounded-relative-mayer-vietoris})을
사용하면 다음 가환 도식을 얻는다.
$$
\xymatrix{
Lg^*Rf_*K \ar[r] \ar[d] &
Rf'_* K' \ar[d] \\
Lg^*Ra_* K|_U \oplus
Lg^*Rb_* K|_V \ar[r] \ar[d] &
Ra'_* K'|_{U'} \oplus
Rb'_* K'|_{V'} \ar[d] \\
Lg^*Rc_* K|_{U \cap V} \ar[r] \ar[d] &
Rc'_* K'|_{U' \cap V'} \ar[d] \\
Lg^*Rf_* K[1] \ar[r] &
Rf'_* K'[1]
}
$$
두 번째와 세 번째 수평 화살표가 동형사상이므로 첫 번째 화살표도
동형사상이다(Derived Categories, Lemma
\ref{derived-lemma-third-isomorphism-triangle}). 이로써 보조정리의
증명이 끝난다.
\end{proof}

\begin{lemma}
\label{perfect-lemma-single-complex-base-change-condition-inherited}
다음과 같이 두자.
$$
\xymatrix{
X' \ar[r]_{g'} \ar[d]_{f'} &
X \ar[d]^f \\
S' \ar[r]^g &
S
}
$$
이는 스킴들의 데카르트 도식이다. $K \in D_\QCoh(\mathcal{O}_X)$를
잡고 $L(g')^*K \to K'$를 $D_\QCoh(\mathcal{O}_{X'})$에서의 사상이라 하자.
Lemma \ref{perfect-lemma-single-complex-base-change-condition}의 동치 조건들이
성립한다고 가정하자. 그러면
\begin{enumerate}
\item $E \in D_\QCoh(\mathcal{O}_X)$에 대해
Lemma \ref{perfect-lemma-single-complex-base-change-condition}의 동치 조건들이
$L(g')^*(E \otimes^\mathbf{L} K) \to L(g')^*E \otimes^\mathbf{L} K'$에
대해 성립한다.
\item $E$가 $D(\mathcal{O}_X)$의 대상이고 완전하면
Lemma \ref{perfect-lemma-single-complex-base-change-condition}의 동치 조건들이
$L(g')^*R\SheafHom(E, K) \to R\SheafHom(L(g')^*E, K')$에 대해 성립한다.
\item $K$가 아래로 유계이고 $E$가 $D(\mathcal{O}_X)$의 대상이며
유사연접이면 Lemma \ref{perfect-lemma-single-complex-base-change-condition}의
동치 조건들이 $L(g')^*R\SheafHom(E, K) \to
R\SheafHom(L(g')^*E, K')$에 대해 성립한다.
\end{enumerate}
\end{lemma}

\begin{proof}
이 명제는 관련 복합체들의 코호몰로지 층이 준연접이므로 의미가 있다.
이는 다음 보조정리들에 의해 성립한다.
\ref{perfect-lemma-quasi-coherence-pullback},
\ref{perfect-lemma-quasi-coherence-tensor-product}, 및
\ref{perfect-lemma-quasi-coherence-internal-hom} 및 Cohomology, Lemmas
\ref{cohomology-lemma-pseudo-coherent-pullback}와
\ref{cohomology-lemma-perfect-pullback}를 사용한다.
이제 (1)의 경우에는 텐서곱의 추이적 성질을 사용하여
사상들 (\ref{perfect-equation-bc})가 동형사상인지 확인하고,
(\ref{perfect-equation-bc})의 시작 대상과 끝 대상을 계산하면 그 사상이
동형사상임을 확인할 수 있다. More on Algebra, Lemma
\ref{more-algebra-lemma-triple-tensor-product}를 보라.
(2)와 (3)의 사상은 다음 합성이다.
$$
L(g')^*R\SheafHom(E, K) \to R\SheafHom(L(g')^*E, L(g')^*K)
\to R\SheafHom(L(g')^*E, K')
$$
여기서 첫 번째 화살표는 Cohomology, Remark
\ref{cohomology-remark-prepare-fancy-base-change}이고, 두 번째 화살표는
주어진 사상 $L(g')^*K \to K'$에서 온다.
(\ref{perfect-equation-bc})의 사상들이 동형사상임을 보이기 위해 (2)의 경우
 $E_x$를 유한 생성 사영 $\mathcal{O}_{X. x}$-가군들의 유계 복합체로,
(3)의 경우 유한 자유 가군들의 위로 유계 복합체로 나타내고 화살표의
시작 대상과 끝 대상을 계산한다. 세부 계산은 생략한다.
\end{proof}

\begin{lemma}
\label{perfect-lemma-base-change-tensor}
 $f : X \to S$를 준콤팩트 준분리 스킴 사상이라 하자.
$E \in D_\QCoh(\mathcal{O}_X)$라 하자. $\mathcal{G}^\bullet$를
$S$ 위에서 평탄한 준연접 $\mathcal{O}_X$-가군의 위로 유계 복합체라
하자. 그러면 다음의 구성은
$$
Rf_*(E \otimes^\mathbf{L}_{\mathcal{O}_X} \mathcal{G}^\bullet)
$$
임의의 기저변환과 교환한다(정확한 명제는 증명을 보라).
\end{lemma}

\begin{proof}
이 명제의 뜻은 다음과 같다. $g : S' \to S$를 스킴의 사상이라 하고
다음 기저변환 도식을 생각하자.
$$
\xymatrix{
X' \ar[r]_{g'} \ar[d]_{f'} &
X \ar[d]^f \\
S' \ar[r]^g &
S
}
$$
즉 $X' = S' \times_S X$이다. 이 보조정리는 다음 사상이
$$
Lg^*Rf_*(E \otimes^\mathbf{L}_{\mathcal{O}_X} \mathcal{G}^\bullet)
\longrightarrow
Rf'_*\left(
L(g')^*E \otimes^\mathbf{L}_{\mathcal{O}_{X'}} (g')^*\mathcal{G}^\bullet
\right)
$$
동형사상이라고 주장한다. 우변에서는 $\mathcal{G}^\bullet$에 유도된
당김을 사용하지 않는다는 점에 유의하자. 이를 보이기 위해 Lemmas
\ref{perfect-lemma-single-complex-base-change}와
\ref{perfect-lemma-single-complex-base-change-condition-inherited}를 적용하면,
다음 표준 사상이
$$
L(g')^*\mathcal{G}^\bullet \to (g')^*\mathcal{G}^\bullet
$$
Lemma \ref{perfect-lemma-single-complex-base-change-condition}의 동치 조건들을
만족함을 보이면 충분하다. 이는 줄기에서 조건을 확인하면 따른다.
우리의 $\mathcal{G}^\bullet$에 대한 가정에 의해
$\mathcal{G}^\bullet_x \otimes_{\mathcal{O}_{S, s}} \mathcal{O}_{S', s'}$가
유도 텐서곱을 계산하기 때문이다.
\end{proof}

\begin{lemma}
\label{perfect-lemma-base-change-RHom}
스킴의 사상 $f : X \to S$가 준콤팩트이고 준분리라고 하자.
$E$를 $D(\mathcal{O}_X)$의 대상으로 하자.
$\mathcal{G}^\bullet$를 준연접 $\mathcal{O}_X$-가군들의 복합체라 하자. 다음 중 하나가 성립한다고 하자.
\begin{enumerate}
\item $E$가 완전 대상이고, $\mathcal{G}^\bullet$가 위로 유계이며,
$\mathcal{G}^n$이 $S$ 위에서 평탄하거나,
\item $E$가 유사연접 대상이고, $\mathcal{G}^\bullet$가 유계이며,
$\mathcal{G}^n$이 $S$ 위에서 평탄하다.
\end{enumerate}
그러면
$$
Rf_*R\SheafHom(E, \mathcal{G}^\bullet)
$$
의 구성은 임의의 기저변환과 교환한다(정확한 명제는 증명을 보라).
\end{lemma}

\begin{proof}
이 명제의 뜻은 다음과 같다. 스킴의 사상 $g : S' \to S$를 잡고
다음 기저변환 도식을 생각하자.
$$
\xymatrix{
X' \ar[r]_{g'} \ar[d]_{f'} &
X \ar[d]^f \\
S' \ar[r]^g &
S
}
$$
즉 $X' = S' \times_S X$이다. 이 보조정리는 다음 사상이
$$
Lg^*Rf_*R\SheafHom(E, \mathcal{G}^\bullet)
\longrightarrow
R(f')_*R\SheafHom(L(g')^*E, (g')^*\mathcal{G}^\bullet)
$$
동형사상이라고 주장한다. 우변에서는 $\mathcal{G}^\bullet$에 유도된
당김을 사용하지 않는다는 점에 유의하자. 이를 보이기 위해 Lemmas
\ref{perfect-lemma-single-complex-base-change}와
\ref{perfect-lemma-single-complex-base-change-condition-inherited}를 적용하면
다음 표준 사상이
$$
L(g')^*\mathcal{G}^\bullet \to (g')^*\mathcal{G}^\bullet
$$
Lemma \ref{perfect-lemma-single-complex-base-change-condition}의 동치 조건들을
만족함을 보이면 충분하다.
이는 Lemma \ref{perfect-lemma-base-change-tensor}의 증명에서 보였다.
\end{proof}





\section{완전 복합체의 구성}
\label{perfect-section-producing-perfect}

\noindent
다음 보조정리는 완전 복합체를 구성하는 우리의 주요 기술 도구이다.
이 결과의 이후 버전들은 Noetherian 근사를 통해 이 보조정리로 환원된다.
\ref{perfect-section-cohomology-and-base-change-final}절을 보라.

\begin{lemma}
\label{perfect-lemma-perfect-direct-image}
Noether 스킴 $S$를 잡자. 스킴의 사상 $f : X \to S$가 국소적으로 유한형이라고
하자. $E \in D(\mathcal{O}_X)$라 하자. 다음 조건을 만족한다고 하자.
\begin{enumerate}
\item $E \in D^b_{\textit{Coh}}(\mathcal{O}_X)$,
\item 모든 $i$에 대해 $H^i(E)$의 지지가 $S$ 위에서 고유하고,
\item $E$는 $D(f^{-1}\mathcal{O}_S)$의 대상으로서 유한 Tor 차원을 갖는다.
\end{enumerate}
그러면 $Rf_*E$는 $D(\mathcal{O}_S)$의 완전 대상이다.
\end{lemma}

\begin{proof}
Lemma \ref{perfect-lemma-direct-image-coherent}에 의해 $Rf_*E$는
$D^b_{\textit{Coh}}(\mathcal{O}_S)$의 대상임을 알 수 있다. 따라서 $Rf_*E$는
유사연접 대상이다(Lemma \ref{perfect-lemma-identify-pseudo-coherent-noetherian}).
그러므로 $Rf_*E$가 유한 Tor 차원을 갖는다는 것을 보이면 충분하다. 이는
Cohomology, Lemma \ref{cohomology-lemma-perfect}을 보라.
Lemma \ref{perfect-lemma-tor-qc-qs}에 의해 다음을 확인하면 충분하다.
$Rf_*(E) \otimes_{\mathcal{O}_S}^\mathbf{L} \mathcal{F}$
가 $S$ 위의 모든 준연접 층 $\mathcal{F}$에 대해 유계 코호몰로지를 갖는다는 것을
확인하면 된다. 위로 유계라는 것은 $Rf_*(E)$가 위로 유계이므로 자명하다.
모든 $i$에 대한 $H^i(E)$의 지지들의 합집합을 $T \subset X$라 하자.
그러면 가정 (1), (2)에 의해 $T$는 $S$ 위에서 고유하다. Cohomology of
Schemes, Lemma \ref{coherent-lemma-union-closed-proper-over-base}를 보라.
특히 $T$를 포함하는 준콤팩트 열린집합 $X' \subset X$가 존재한다.
$f' = f|_{X'}$로 두면 $E$가 $X \setminus T$에서 0으로 제한되므로
$Rf_*(E) = Rf'_*(E|_{X'})$이다. 따라서 $X$를 $X'$로 바꾸어 $f$가
준콤팩트라고 가정할 수 있다. 또한 Morphisms, Lemma
\ref{morphisms-lemma-finite-type-Noetherian-quasi-separated}에 의해 $f$는
준분리이다. 이제
$$
Rf_*(E) \otimes_{\mathcal{O}_S}^\mathbf{L} \mathcal{F} =
Rf_*\left(E \otimes_{\mathcal{O}_X}^\mathbf{L} Lf^*\mathcal{F}\right) =
Rf_*\left(E \otimes_{f^{-1}\mathcal{O}_S}^\mathbf{L} f^{-1}\mathcal{F}\right)
$$
이는
Lemma \ref{perfect-lemma-cohomology-base-change}
및
Cohomology, Lemma \ref{cohomology-lemma-variant-derived-pullback}.
가정 (3)에 의해 복합체
$E \otimes_{f^{-1}\mathcal{O}_S}^\mathbf{L} f^{-1}\mathcal{F}$
의 코호몰로지 층들은 어떤 유한한 범위, 가령 $[a, b]$에 놓인다. 따라서
그 복합체의 $Rf_*$는 $[a, \infty)$ 범위에 코호몰로지를 가지며 결론을 얻는다.
\end{proof}

\begin{lemma}
\label{perfect-lemma-tensor-perfect}
$S$를 뇌터 스킴이라 하고 $f : X \to S$를 국소 유한형인 스킴 사상이라 하자.
$E \in D(\mathcal{O}_X)$를 완전이라 하자. $\mathcal{G}^\bullet$을 $S$ 위에서
평탄하고 지지가 $S$ 위에서 고유인 연접 $\mathcal{O}_X$-가군들의 유계
복합체라 하자. 그러면
$K = Rf_*(E \otimes_{\mathcal{O}_X}^\mathbf{L} \mathcal{G}^\bullet)$는
$D(\mathcal{O}_S)$의 완전 대상이다.
\end{lemma}

\begin{proof}
Lemma \ref{perfect-lemma-perfect-direct-image}에 의해 대상 $K$는 완전이다. 이
보조정리가 적용됨을 확인하자. 국소적으로 $E$는 유한 자유
$\mathcal{O}_X$-가군들의 유한 복합체와 동형이다. 따라서 국소적으로
$E \otimes^\mathbf{L}_{\mathcal{O}_X} \mathcal{G}^\bullet$은 각 항이 다음
꼴인 유한 복합체와 동형이다.
$$
\bigoplus\nolimits_{i = a, \ldots, b} (\mathcal{G}^i)^{\oplus r_i}
$$
여기서 $a, b, r_a, \ldots, r_b$는 어떤 정수들이다. 이로부터 코호몰로지 층
$H^i(E \otimes^\mathbf{L}_{\mathcal{O}_X} \mathcal{G})$가 연접임이 즉시
따른다. 또한 $\mathcal{G}^i$가 $f^{-1}\mathcal{O}_S$ 위에서 평탄이므로
Tor 차원에 관한 가정도 따른다.
\end{proof}

\begin{lemma}
\label{perfect-lemma-ext-perfect}
$S$를 뇌터 스킴이라 하고 $f : X \to S$를 국소 유한형인 스킴 사상이라 하자.
$E \in D(\mathcal{O}_X)$를 완전이라 하자. $\mathcal{G}^\bullet$을 $S$ 위에서
평탄하고 지지가 $S$ 위에서 고유인 연접 $\mathcal{O}_X$-가군들의 유계
복합체라 하자. 그러면 $K = Rf_*R\SheafHom(E, \mathcal{G}^\bullet)$는
$D(\mathcal{O}_S)$의 완전 대상이다.
\end{lemma}

\begin{proof}
$E$가 완전 복합체이므로 쌍대 완전 복합체 $E^\vee$가 존재한다. Cohomology,
Lemma \ref{cohomology-lemma-dual-perfect-complex}를 보라. 다음에 유의하라.
$R\SheafHom(E, \mathcal{G}^\bullet) =
E^\vee \otimes^\mathbf{L}_{\mathcal{O}_X} \mathcal{G}^\bullet$.
따라서 $K$가 완전임은 Lemma \ref{perfect-lemma-tensor-perfect}에서 따른다.
\end{proof}

\noindent
다음 보조정리는 Lemma \ref{perfect-lemma-flat-proper-perfect-direct-image-general}에서
일반 기저 위의 평탄 고유 사상으로, More on Morphisms, Lemma
\ref{more-morphisms-lemma-perfect-proper-perfect-direct-image}에서 완전 고유
사상으로 일반화하겠다.

\begin{lemma}
\label{perfect-lemma-flat-proper-perfect-direct-image}
$S$를 뇌터 스킴이라 하고 $f : X \to S$를 평탄 고유 스킴 사상이라 하자.
$E \in D(\mathcal{O}_X)$가 완전이면 $Rf_*E$는
$D(\mathcal{O}_S)$의 완전 대상이다.
\end{lemma}

\begin{proof}
Lemma \ref{perfect-lemma-perfect-direct-image}가 적용된다고 주장한다. 조건 (1)과
(2)는 즉시 성립한다. 조건 (3)은 $X$ 위에서 국소적이다. 따라서 $X$와
$S$가 아핀이고 $E$가 $\mathcal{O}_X$-가군들의 엄밀 완전 복합체로
표현된다고 가정할 수 있다. $\mathcal{O}_X$는
$f^{-1}\mathcal{O}_S$-가군의 층으로서 평탄이므로 조건 (3)이 성립한다.
\end{proof}






\section{Ext에 대한 사영 공식}
\label{perfect-section-ext}

\noindent
Lemma \ref{perfect-lemma-compute-ext}(또는 문헌의 유사한 결과)은 Artin의 조건들 중
하나를 확인하는 데 때때로 사용된다. 예를 들어 Quot 함자, Hilbert 스킴 및
다른 모듈라이 문제에서 그렇다. $f : X \to S$가 고유하고 평탄하며 유한
표현 가능한 스킴의 사상이고 $E \in D(\mathcal{O}_X)$가 완전하다고 하자.
이때 보조정리는 다음을 말한다.
$$
\Ext^i_X(E, f^*\mathcal{F}) =
\Ext^i_S((Rf_*E^\vee)^\vee, \mathcal{F})
$$
$\mathcal{F}$가 $S$ 위의 준연접 가군일 때 성립한다.
이렇게 쓰면 Ext에 대한 사영 공식처럼 보이며, 실제로 이 결과는
Lemma \ref{perfect-lemma-cohomology-base-change}에서 상당히 쉽게 따른다.

\begin{lemma}
\label{perfect-lemma-compute-tensor-perfect}
Lemma \ref{perfect-lemma-tensor-perfect}와 같은 가정과 표기를 사용하자.
그러면 다음 함자적 동형사상들이 존재한다.
$$
H^i(S, K \otimes^\mathbf{L}_{\mathcal{O}_S} \mathcal{F})
\longrightarrow
H^i(X, E \otimes_{\mathcal{O}_X}^\mathbf{L}
(\mathcal{G}^\bullet \otimes_{\mathcal{O}_X} f^*\mathcal{F}))
$$
$\mathcal{F}$가 $S$ 위의 준연접 가군일 때 성립하며,
경계 사상과 양립한다(증명을 보라).
\end{lemma}

\begin{proof}
다음 등식이 성립한다.
$$
\mathcal{G}^\bullet \otimes_{\mathcal{O}_X}^\mathbf{L} Lf^*\mathcal{F} =
\mathcal{G}^\bullet \otimes_{f^{-1}\mathcal{O}_S}^\mathbf{L} f^{-1}\mathcal{F} =
\mathcal{G}^\bullet \otimes_{f^{-1}\mathcal{O}_S} f^{-1}\mathcal{F} =
\mathcal{G}^\bullet \otimes_{\mathcal{O}_X} f^*\mathcal{F}
$$
첫 번째 등식은
Cohomology, Lemma \ref{cohomology-lemma-variant-derived-pullback},
두 번째는 $\mathcal{G}^n$이 평탄한 $f^{-1}\mathcal{O}_S$-가군이라는 사실에,
세 번째는 당김의 정의에 따른다. 따라서 다음을 얻는다.
\begin{align*}
H^i(X, E \otimes^\mathbf{L}_{\mathcal{O}_X}
(\mathcal{G}^\bullet \otimes_{\mathcal{O}_X} f^*\mathcal{F}))
& =
H^i(X, E \otimes^\mathbf{L}_{\mathcal{O}_X} \mathcal{G}^\bullet
\otimes_{\mathcal{O}_X}^\mathbf{L} Lf^*\mathcal{F}) \\
& =
H^i(S,
Rf_*(E \otimes^\mathbf{L}_{\mathcal{O}_X} \mathcal{G}^\bullet
\otimes^\mathbf{L}_{\mathcal{O}_X} Lf^*\mathcal{F})) \\
& =
H^i(S, Rf_*(E \otimes^\mathbf{L}_{\mathcal{O}_X} \mathcal{G}^\bullet)
\otimes^\mathbf{L}_{\mathcal{O}_S} \mathcal{F}) \\
& =
H^i(S, K \otimes^\mathbf{L}_{\mathcal{O}_S} \mathcal{F}) 
\end{align*}
첫 번째 등식은 위의 결과에, 두 번째는 Leray 정리(Cohomology, Lemma
\ref{cohomology-lemma-before-Leray})에, 세 번째 등식은
Lemma \ref{perfect-lemma-cohomology-base-change}에 따른다.
경계 사상과의 양립은 다음을 뜻한다. 준연접 $\mathcal{O}_S$-가군들의 짧은
완전열 $0 \to \mathcal{F}_1 \to \mathcal{F}_2 \to \mathcal{F}_3 \to 0$이
주어지면, 이 동형사상들은 다음 가환 도식에 들어간다.
$$
\xymatrix{
H^i(S, K \otimes^\mathbf{L}_{\mathcal{O}_S} \mathcal{F}_3)
\ar[r] \ar[d]_\delta &
H^i(X, E \otimes^\mathbf{L}_{\mathcal{O}_X}
(\mathcal{G}^\bullet \otimes_{\mathcal{O}_X} f^*\mathcal{F}_3))
\ar[d]^\delta \\
H^{i + 1}(S, K \otimes^\mathbf{L}_{\mathcal{O}_S} \mathcal{F}_1)
\ar[r] &
H^{i + 1}(X, E \otimes^\mathbf{L}_{\mathcal{O}_X}
(\mathcal{G}^\bullet \otimes_{\mathcal{O}_X} f^*\mathcal{F}_1))
}
$$
여기서 경계 사상은 다음 구별 삼각형에서 온다.
$$
K \otimes^\mathbf{L}_{\mathcal{O}_S} \mathcal{F}_1 \to
K \otimes^\mathbf{L}_{\mathcal{O}_S} \mathcal{F}_2 \to
K \otimes^\mathbf{L}_{\mathcal{O}_S} \mathcal{F}_3 \to
K \otimes^\mathbf{L}_{\mathcal{O}_S} \mathcal{F}_1[1]
$$
그리고 다음 짧은 완전열에 대응하는 $D(\mathcal{O}_X)$ 안의 구별 삼각형에서도
온다.
$$
0 \to
\mathcal{G}^\bullet \otimes_{\mathcal{O}_X} f^*\mathcal{F}_1 \to
\mathcal{G}^\bullet \otimes_{\mathcal{O}_X} f^*\mathcal{F}_2 \to
\mathcal{G}^\bullet \otimes_{\mathcal{O}_X} f^*\mathcal{F}_3 \to 0
$$
이는 $\mathcal{O}_X$-가군 복합체들의 열이다. 이 열은
$\mathcal{G}^n$이 $S$ 위에서 평탄하므로 완전하다.
표시한 도식의 가환성 확인은 생략한다.
\end{proof}

\begin{lemma}
\label{perfect-lemma-compute-ext-perfect}
Lemma \ref{perfect-lemma-ext-perfect}와 같은 가정과 표기를 사용하자.
그러면 다음 함자적 동형사상들이 존재한다.
$$
H^i(S, K \otimes^\mathbf{L}_{\mathcal{O}_S} \mathcal{F})
\longrightarrow
\Ext^i_{\mathcal{O}_X}(E,
\mathcal{G}^\bullet \otimes_{\mathcal{O}_X} f^*\mathcal{F})
$$
$\mathcal{F}$가 $S$ 위의 준연접 가군일 때 성립하며,
경계 사상과 양립한다(증명을 보라).
\end{lemma}

\begin{proof}
Lemma \ref{perfect-lemma-ext-perfect}의 증명에서와 같이 쌍대 완전 복합체 $E^\vee$를
잡고 다음을 기억하자.
$K = Rf_*(E^\vee \otimes_{\mathcal{O}_X}^\mathbf{L} \mathcal{G}^\bullet)$.
또한 다음이 성립한다.
$$
\Ext^i_{\mathcal{O}_X}(E,
\mathcal{G}^\bullet \otimes_{\mathcal{O}_X} f^*\mathcal{F})
=
H^i(X, E^\vee \otimes^\mathbf{L}_{\mathcal{O}_X}
(\mathcal{G}^\bullet \otimes_{\mathcal{O}_X} f^*\mathcal{F}))
$$
$E^\vee$의 구성에 의해, 이 동형사상들의 존재는 $E^\vee$와
$\mathcal{G}^\bullet$에 Lemma \ref{perfect-lemma-compute-tensor-perfect}를 적용하면 따른다.
경계 사상과의 양립은 다음을 뜻한다. 짧은 완전열
$0 \to \mathcal{F}_1 \to \mathcal{F}_2 \to \mathcal{F}_3 \to 0$이
주어지면 이 동형사상들은 다음 가환 도식에 들어간다.
$$
\xymatrix{
H^i(S, K \otimes^\mathbf{L}_{\mathcal{O}_S} \mathcal{F}_3)
\ar[r] \ar[d]_\delta &
\Ext^i_{\mathcal{O}_X}(E,
\mathcal{G}^\bullet \otimes_{\mathcal{O}_X} f^*\mathcal{F}_3) \ar[d]^\delta \\
H^{i + 1}(S, K \otimes^\mathbf{L}_{\mathcal{O}_S} \mathcal{F}_1)
\ar[r] &
\Ext^{i + 1}_{\mathcal{O}_X}(E,
\mathcal{G}^\bullet \otimes_{\mathcal{O}_X} f^*\mathcal{F}_1)
}
$$
여기서 경계 사상은 다음 구별 삼각형에서 온다.
$$
K \otimes^\mathbf{L}_{\mathcal{O}_S} \mathcal{F}_1 \to
K \otimes^\mathbf{L}_{\mathcal{O}_S} \mathcal{F}_2 \to
K \otimes^\mathbf{L}_{\mathcal{O}_S} \mathcal{F}_3 \to
K \otimes^\mathbf{L}_{\mathcal{O}_S} \mathcal{F}_1[1]
$$
그리고 다음 짧은 완전열에 대응하는 $D(\mathcal{O}_X)$ 안의 구별 삼각형에서도
온다.
$$
0 \to
\mathcal{G}^\bullet \otimes_{\mathcal{O}_X} f^*\mathcal{F}_1 \to
\mathcal{G}^\bullet \otimes_{\mathcal{O}_X} f^*\mathcal{F}_2 \to
\mathcal{G}^\bullet \otimes_{\mathcal{O}_X} f^*\mathcal{F}_3 \to 0
$$
이는 복합체들의 열이다. 이 열은 $\mathcal{G}$가 $S$ 위에서 평탄하므로
완전하다. 표시한 도식의 가환성 확인은 생략한다.
\end{proof}

\begin{lemma}
\label{perfect-lemma-compute-ext}
스킴의 사상 $f : X \to S$, 대상 $E \in D(\mathcal{O}_X)$ 및
$\mathcal{G}^\bullet$인 $\mathcal{O}_X$-가군 복합체를 잡자.
다음을 가정하자.
\begin{enumerate}
\item $S$는 Noether 스킴이다.
\item $f$는 국소적으로 유한형이다.
\item $E \in D^-_{\textit{Coh}}(\mathcal{O}_X)$이다.
\item $\mathcal{G}^\bullet$는 $S$ 위에서 평탄하고 지지가 $S$ 위에서 고유한
준연접 $\mathcal{O}_X$-가군들의 유계 복합체이다.
\end{enumerate}
그러면 다음 두 명제가 성립한다.
\begin{enumerate}
\item[(A)] 모든 $m \in \mathbf{Z}$에 대해 $D(\mathcal{O}_S)$의 완전 대상 $K$와
함자적 사상
$$
\alpha^i_\mathcal{F} :
\Ext^i_{\mathcal{O}_X}(E,
\mathcal{G}^\bullet \otimes_{\mathcal{O}_X} f^*\mathcal{F})
\longrightarrow
H^i(S, K \otimes^\mathbf{L}_{\mathcal{O}_S} \mathcal{F})
$$
가 존재한다. 이 사상은 $\mathcal{F}$가 $S$ 위의 준연접 가군일 때 경계
사상과 양립하며(증명을 보라), $i \leq m$이면
$\alpha^i_\mathcal{F}$는 동형사상이다.
\item[(B)] 유사연접 대상 $L \in D(\mathcal{O}_S)$와 함자적 동형사상
$$
\Ext^i_{\mathcal{O}_S}(L, \mathcal{F}) \longrightarrow
\Ext^i_{\mathcal{O}_X}(E,
\mathcal{G}^\bullet \otimes_{\mathcal{O}_X} f^*\mathcal{F})
$$
$\mathcal{F}$가 $S$ 위의 준연접 가군일 때 경계 사상과 양립한다.
\end{enumerate}
\end{lemma}

\begin{proof}
 (A)의 증명.
$\mathcal{G}^i$가 $i \in [a, b]$에서만 0이 아니라고 하자.
$X$를 $\mathcal{G}^i$의 지지들의 합집합을 포함하는 준콤팩트 열린
근방으로 바꾸어도 된다. 따라서 $X$가 Noether 스킴이라고 가정할 수 있다.
이 경우 $X$와 $f$는 준콤팩트이고 준분리이다.
$(X, E, -m - 1 + a)$에 대해 완전 복합체 $P$에 의한 근사
$P \to E$를 택하자.
(이는 Theorem \ref{perfect-theorem-approximation}에 의해 가능하다).
그러면 유도된 사상은
$$
\Ext^i_{\mathcal{O}_X}(E,
\mathcal{G}^\bullet \otimes_{\mathcal{O}_X} f^*\mathcal{F})
\longrightarrow
\Ext^i_{\mathcal{O}_X}(P,
\mathcal{G}^\bullet \otimes_{\mathcal{O}_X} f^*\mathcal{F})
$$
이는 $i \leq m$일 때 동형사상이다. 구체적으로 이 사상의 핵과 여핵은 각각
다음 가군의 몫과 부분가군이다.
$$
\Ext^i_{\mathcal{O}_X}(C,
\mathcal{G}^\bullet \otimes_{\mathcal{O}_X} f^*\mathcal{F})
\quad\text{각각}\quad
\Ext^{i + 1}_{\mathcal{O}_X}(C,
\mathcal{G}^\bullet \otimes_{\mathcal{O}_X} f^*\mathcal{F})
$$
여기서 $C$는 $P \to E$의 사상원뿔이다. $C$의 코호몰로지 층은
차수 $\geq -m - 1 + a$에서 소멸하므로, 다음 $\Ext$-군들은
$i \leq m + 1$일 때
Derived Categories, Lemma \ref{derived-lemma-negative-exts}.
에 의해 0이다. 따라서 문제는 $E$가 완전 복합체인 경우로 환원되며,
이는 Lemma \ref{perfect-lemma-compute-ext-perfect}이다.
경계에 관한 명제는 Lemma \ref{perfect-lemma-compute-ext-perfect}의 증명에서
설명했다.

\medskip\noindent
 (B)의 증명. (A)의 증명에서처럼 $X$가 Noether 스킴이라고 가정할 수 있다.
$E$는 다음에 의해 유사연접이다.
Lemma \ref{perfect-lemma-identify-pseudo-coherent-noetherian}.
Lemma \ref{perfect-lemma-pseudo-coherent-hocolim}에 의해
$E = \text{hocolim} E_n$로 쓸 수 있다. 여기서 $E_n$은 완전하고
$E_n \to E$는 절단 $\tau_{\geq -n}$에서 동형을 유도한다. 쌍대 완전 복합체
$E_n^\vee$를 잡자.
(Cohomology, Lemma \ref{cohomology-lemma-dual-perfect-complex}).
그러면 완전 대상들의 역계 $\ldots \to E_3^\vee \to E_2^\vee \to E_1^\vee$를
얻는다. 이로부터 다음 역계가 생긴다.
$$
\ldots \to K_3 \to K_2 \to K_1\quad\text{여기서}\quad
K_n = Rf_*(E_n^\vee \otimes_{\mathcal{O}_X}^\mathbf{L} \mathcal{G}^\bullet)
$$
$S$ 위에서 완전하다. Lemma \ref{perfect-lemma-tensor-perfect}를 보라.
Lemma \ref{perfect-lemma-compute-ext-perfect}와 그 증명, 그리고 앞 단락의 논거
($P = E_n$)에 의해 $S$ 위의 임의의 준연접 $\mathcal{F}$에 대해 다음
함자적 표준 사상들이 존재한다.
$$
\xymatrix{
& \Ext^i_{\mathcal{O}_X}(E,
\mathcal{G}^\bullet \otimes_{\mathcal{O}_X} f^*\mathcal{F})
\ar[ld] \ar[rd] \\
H^i(S, K_{n + 1} \otimes_{\mathcal{O}_S}^\mathbf{L} \mathcal{F})
\ar[rr] & &
H^i(S, K_n \otimes_{\mathcal{O}_S}^\mathbf{L} \mathcal{F})
}
$$
이 사상들은 $i \leq n + a$일 때 동형사상이다.
$L_n = K_n^\vee$를 쌍대 완전 복합체라 하자.
그러면 $L_1 \to L_2 \to L_3 \to \ldots$는 $D(\mathcal{O}_S)$ 안의 완전
대상들로 이루어진 계이며, $S$ 위의 임의의 준연접 $\mathcal{F}$에 대해
다음 사상들이
$$
\Ext^i_{\mathcal{O}_S}(L_{n + 1}, \mathcal{F})
\longrightarrow
\Ext^i_{\mathcal{O}_S}(L_n, \mathcal{F})
$$
 $i \leq n + a - 1$일 때 동형사상이다. 따라서
$L_n \to L_{n + 1}$는 절단 $\tau_{\geq -n - a + 2}$에서 동형을 유도한다
(힌트: $L_n \to L_{n + 1}$의 사상원뿔을 취하고 마지막으로 소멸하지 않는
코호몰로지 층을 보라).
그러므로 $L = \text{hocolim} L_n$은 유사연접이다. Lemma
\ref{perfect-lemma-pseudo-coherent-hocolim}을 보라. 호모토피 귀납극한의 사상
성질에 의해 다음이 성립한다.
$\Ext^i_{\mathcal{O}_S}(L, \mathcal{F}) =
\Ext^i_{\mathcal{O}_S}(L_n, \mathcal{F})$
$i \leq n + a - 3$일 때 성립하며, 이로써 증명이 끝난다.
\end{proof}

\begin{remark}
\label{perfect-remark-base-change-of-L}
$L$인 유사연접 복합체는 Lemma \ref{perfect-lemma-compute-ext}의 (B) 부분의
상황에 표준적으로 대응한다. 예를 들어 (B)에서 $L$을 구성하는 것은
기저변환과 양립한다. 즉, 다음 데카르트 도식이 주어지면
$$
\xymatrix{
X' \ar[r]_{g'} \ar[d]_{f'} &
X \ar[d]^f \\
S' \ar[r]^g &
S
}
$$
스킴들에 대해 다음 표준 함자적 동형사상들을 갖는다.
$$
\Ext^i_{\mathcal{O}_{S'}}(Lg^*L, \mathcal{F}') \longrightarrow
\Ext^i_{\mathcal{O}_X}(L(g')^*E,
(g')^*\mathcal{G}^\bullet \otimes_{\mathcal{O}_{X'}} (f')^*\mathcal{F}')
$$
$\mathcal{F}'$가 $S'$ 위의 준연접 가군일 때 성립한다. 우변에서는
$\mathcal{G}^\bullet$에 유도된 당김을 사용하지 않는다는 점에 유의하자.
언젠가 이것이 필요해지면 여기에서 정확한 명제를 정식화하고 자세한 증명을
제시할 것이다.
\end{remark}






\section{극한과 유도 범주}
\label{perfect-section-limits}

\noindent
이 절에서는 스킴의 역계의 극한으로 주어진 스킴의 유도 범주에 관한
몇 가지 결과를 모은다. 더 정확히는 다음 상황에서 논한다.

\begin{situation}
\label{perfect-situation-descent}
$S = \lim_{i \in I} S_i$를 아핀 전이 사상
$f_{i'i} : S_{i'} \to S_i$를 갖는 스킴의 유향계의 극한이라 하자.
모든 $i \in I$에 대해 $S_i$가 준콤팩트이고 준분리라고 가정한다.
사영을 $f_i : S \to S_i$로 표기한다. 또한 원소 $0 \in I$를 하나 고정한다.
\end{situation}

\begin{lemma}
\label{perfect-lemma-descend-homomorphisms}
Situation \ref{perfect-situation-descent}에서 $E_0$와 $K_0$를
$D(\mathcal{O}_{S_0})$의 대상이라 하자.
$i \geq 0$에 대해 $E_i = Lf_{i0}^*E_0$와 $K_i = Lf_{i0}^*K_0$로 놓고,
$E = Lf_0^*E_0$와 $K = Lf_0^*K_0$로 놓자. 그러면 다음 중 하나가
성립할 때 다음 사상은
$$
\colim_{i \geq 0} \Hom_{D(\mathcal{O}_{S_i})}(E_i, K_i)
\longrightarrow
\Hom_{D(\mathcal{O}_S)}(E, K)
$$
동형사상이다.
\begin{enumerate}
\item $E_0$가 완전이고 $K_0 \in D_\QCoh(\mathcal{O}_{S_0})$이다.
\item $E_0$가 유사연접이고
$K_0 \in D_\QCoh(\mathcal{O}_{S_0})$가 유한 Tor 차원을 갖는다.
\end{enumerate}
\end{lemma}

\begin{proof}
각 열린집합 $U_0 \subset S_0$에 대해 다음 표준 사상이 동형사상이라는
조건 $P$를 생각하자.
$$
\colim_{i \geq 0} \Hom_{D(\mathcal{O}_{U_i})}(E_i|_{U_i}, K_i|_{U_i})
\longrightarrow
\Hom_{D(\mathcal{O}_U)}(E|_U, K|_U)
$$
여기서 $U = f_0^{-1}(U_0)$이고 $U_i = f_{i0}^{-1}(U_0)$이다.
Cohomology of Schemes, Lemma \ref{coherent-lemma-induction-principle}의
귀납원리를 사용하여 모든 준콤팩트 열린집합 $U_0$에 대해 $P$가 성립함을
보이겠다. 이 보조정리의 조건 (2)는 유도 범주에서의 Hom에 대한
Mayer--Vietoris로부터 즉시 따른다. Cohomology, Lemma
\ref{cohomology-lemma-mayer-vietoris-hom}를 보라. 따라서 $S_0$가 아핀인
경우에 보조정리를 증명하면 충분하다.

\medskip\noindent
$S_0$가 아핀이라고 가정하자. $S_0 = \Spec(A_0)$,
$S_i = \Spec(A_i)$, $S = \Spec(A)$라 하자. 이하에서는 별도의 언급 없이
Lemma \ref{perfect-lemma-affine-compare-bounded}를 사용한다.

\medskip\noindent
경우 (1)에서 대상 $E_0^\bullet$은 유한 사영 $A_0$-가군의 유한 복합체에
대응한다. Lemma \ref{perfect-lemma-perfect-affine}를 보라. 대상 $K_0$는
$A_0$-가군의 K-평탄 복합체 $K_0^\bullet$로 나타낼 수 있다. 이 상황에서
다음 사상이 동형사상임을 증명하려는 것이다.
$$
\colim_{i \geq 0} \Hom_{D(A_i)}(E_0^\bullet \otimes_{A_0} A_i,
K_0^\bullet \otimes_{A_0} A_i)
\longrightarrow
\Hom_{D(A)}(E_0^\bullet \otimes_{A_0} A, K_0^\bullet \otimes_{A_0} A)
$$
$E_0^\bullet$이 사영 가군의 위로 유계인 복합체이므로 이를 다음과 같이
다시 쓸 수 있다.
$$
\colim_{i \geq 0} \Hom_{K(A_0)}(E_0^\bullet,
K_0^\bullet \otimes_{A_0} A_i)
\longrightarrow
\Hom_{K(A_0)}(E_0^\bullet, K_0^\bullet \otimes_{A_0} A)
$$
0이 아닌 가군 $E_0^n$은 유한 개뿐이고 이들은 모두 유한 표시 가군이므로,
이 사상은 동형사상이다.

\medskip\noindent
경우 (2)에서 대상 $E_0$는 유한 자유 $A_0$-가군의 위로 유계인 복합체
$E_0^\bullet$에 대응한다. Lemma \ref{perfect-lemma-pseudo-coherent-affine}를 보라.
$K_0$는 평탄 $A_0$-가군의 유한 복합체 $K_0^\bullet$로 나타낼 수 있다.
Lemma \ref{perfect-lemma-tor-dimension-affine}와 More on Algebra, Lemma
\ref{more-algebra-lemma-tor-amplitude}를 보라. 특히 $K_0^\bullet$은
K-평탄이며, 앞에서와 같이 논하면 다음 사상에 이른다.
$$
\colim_{i \geq 0} \Hom_{K(A_0)}(E_0^\bullet,
K_0^\bullet \otimes_{A_0} A_i)
\longrightarrow
\Hom_{K(A_0)}(E_0^\bullet, K_0^\bullet \otimes_{A_0} A)
$$
이 사상이 동형사상임은 명백하다($K_0^\bullet$이 유계이므로 유한 개의
항만 관여한다).
\end{proof}

\begin{lemma}
\label{perfect-lemma-descend-perfect}
Situation \ref{perfect-situation-descent}에서 $D(\mathcal{O}_S)$의 완전 대상들의
범주는 $D(\mathcal{O}_{S_i})$의 완전 대상들의 범주들의 귀납극한이다.
\end{lemma}

\begin{proof}
각 열린집합 $U_0 \subset S_0$에 대해 다음 함자가 동치라는 조건 $P$를
생각하자.
$$
\colim_{i \geq 0} D_{perf}(\mathcal{O}_{U_i})
\longrightarrow
D_{perf}(\mathcal{O}_U)
$$
여기서 ${}_{perf}$는 완전 대상들로 이루어진 충만 부분범주를 나타내며,
$U = f_0^{-1}(U_0)$이고 $U_i = f_{i0}^{-1}(U_0)$이다. Cohomology of Schemes,
Lemma \ref{coherent-lemma-induction-principle}의 귀납원리를 사용하여 모든
준콤팩트 열린집합 $U_0$에 대해 $P$가 성립함을 보이겠다. 우선 Lemma
\ref{perfect-lemma-descend-homomorphisms}에 의해 이 함자가 충실충만함은 이미
알고 있다. 따라서 본질적 전사성을 증명하면 충분하다.

\medskip\noindent
먼저 귀납원리의 조건 (2)를 확인하자. $S_0 = U_0 \cup V_0$이고
$U_0$, $V_0$, $U_0 \cap V_0$에 대해 $P$가 성립한다고 하자.
$E$를 $D(\mathcal{O}_S)$의 완전 대상이라 하자. 어떤 $i \geq 0$와
$U_i$ 위의 완전 대상 $E_{U, i}$ 및 $V_i$ 위의 완전 대상 $E_{V, i}$를
찾을 수 있으며, 이들을 $U$와 $V$로 역상한 것은 각각 $E|_U$와
$E|_V$에 동형이다. 다음 사상들을
$$
a : E_{U, i} \to (Rf_{i, *}E)|_{U_i}
\quad\text{및}\quad
b : E_{V, i} \to (Rf_{i, *}E)|_{V_i}
$$
$Lf_i^*E_{U, i} \to E|_U$와 $Lf_i^*E_{V, i} \to E|_V$에 수반하는
사상이라 하자. 충실충만성에 의해 $i$를 크게 잡은 뒤에는 다음 동형사상을
찾을 수 있다.
$c : E_{U, i}|_{U_i \cap V_i} \to E_{V, i}|_{U_i \cap V_i}$
이 사상의 역상은 다음 식별이 된다.
$$
Lf_i^*E_{U, i}|_{U \cap V} \to E|_{U \cap V} \to Lf_i^*E_{V, i}|_{U \cap V}.
$$
Cohomology, Lemma \ref{cohomology-lemma-glue}를 적용하면 $S_i$ 위의 대상
$E_i$와 사상 $d : E_i \to Rf_{i, *}E$를 얻으며, 이 사상은 $U_i$와 $V_i$
위에서 각각 사상 $a$와 $b$로 제한된다. 그러면 $E_i$가 완전이고
$d$가 동형사상 $Lf_i^*E_i \to E$에 수반함은 명백하다.

\medskip\noindent
끝으로 귀납원리의 조건 (1)을 확인한다. 다시 말해 $S_0$가 아핀일 때
보조정리가 성립함을 확인한다. $S_0 = \Spec(A_0)$,
$S_i = \Spec(A_i)$, $S = \Spec(A)$라 하자. Lemmas
\ref{perfect-lemma-affine-compare-bounded}와 \ref{perfect-lemma-perfect-affine}를 사용하면
다음을 보여야 함을 알 수 있다.
$$
D_{perf}(A) = \colim D_{perf}(A_i)
$$
이는 환 위의 완전 복합체가 유한 사영(따라서 유한 표시) 가군의 유한
복합체로 주어진다는 사실로부터 명백하다. 자세한 내용은 More on Algebra,
Lemma \ref{more-algebra-lemma-colimit-perfect-complexes}를 보라.
\end{proof}





\section{코호몰로지와 기저변환, VI}
\label{perfect-section-cohomology-and-base-change-final}

\noindent
Sections
\ref{perfect-section-cohomology-and-base-change-perfect},
\ref{perfect-section-cohomology-base-change}, 그리고
\ref{perfect-section-producing-perfect}의 논의를 이어 코호몰로지와 기저변환을
다루는 마지막 절이다. 쉽게 이해할 수 있는 특수한 경우는 Remark
\ref{perfect-remark-explain-perfect-direct-image}에 제시한다.

\begin{lemma}
\label{perfect-lemma-base-change-tensor-perfect}
$f : X \to S$를 유한 표시 사상이라 하자.
$E \in D(\mathcal{O}_X)$를 완전 대상이라 하자. $\mathcal{G}^\bullet$을
유한 표시 $\mathcal{O}_X$-가군의 유계 복합체로서 $S$ 위에서 평탄하고
그 지지가 $S$ 위에서 고유인 것이라 하자. 그러면
$$
K = Rf_*(E \otimes_{\mathcal{O}_X}^\mathbf{L} \mathcal{G}^\bullet)
$$
는 $D(\mathcal{O}_S)$의 완전 대상이며, 그 구성은 임의의 기저변환과
가환한다.
\end{lemma}

\begin{proof}
기저변환에 관한 명제는 Lemma \ref{perfect-lemma-base-change-tensor}이다.
따라서 $K$가 완전 대상임을 보이면 충분하다. $S$가 뇌터이면 이는
Lemma \ref{perfect-lemma-tensor-perfect}에서 따른다. 뇌터 근사를 사용하여
이 경우로 환원하겠다. 독자는 이 증명의 나머지를 건너뛰어도 좋다.

\medskip\noindent
이 문제는 $S$ 위에서 국소적이므로 $S$가 아핀이라고 가정해도 된다.
$S = \Spec(R)$라 하자. 뇌터 환 $R_i$들의 여과 귀납극한으로
$R = \colim R_i$라 쓴다. Limits, Lemma
\ref{limits-lemma-descend-finite-presentation}에 의해 어떤 $i$와
$R_i$ 위의 유한 표시 스킴 $X_i$가 존재하여 이를 $R$로 기저변환하면
$X$가 된다. Limits, Lemma
\ref{limits-lemma-descend-modules-finite-presentation}에 의해 $i$를 크게
잡은 뒤에는 유한 표시 $\mathcal{O}_{X_i}$-가군의 유계 복합체
$\mathcal{G}_i^\bullet$이 존재하고, 이를 $X$로 역상하면
$\mathcal{G}^\bullet$이 된다고 가정할 수 있다. $i$를 크게 잡아
$\mathcal{G}_i^n$이 $R_i$ 위에서 평탄하다고 가정할 수 있다. Limits,
Lemma \ref{limits-lemma-descend-module-flat-finite-presentation}를 보라.
다시 $i$를 크게 잡아 $\mathcal{G}_i^n$의 지지가 $R_i$ 위에서 고유라고
가정할 수 있다. Limits, Lemma
\ref{limits-lemma-eventually-proper-support}와 Cohomology of Schemes,
Lemma \ref{coherent-lemma-module-support-proper-over-base}를 보라. 끝으로
Lemma \ref{perfect-lemma-descend-perfect}에 의해 $i$를 크게 잡은 뒤에는
$D(\mathcal{O}_{X_i})$의 완전 대상 $E_i$가 존재하여 이를 $X$로
역상하면 $E$가 된다고 가정할 수 있다. $X_i \to \Spec(R_i)$,
$E_i$, $\mathcal{G}_i^\bullet$에 Lemma \ref{perfect-lemma-tensor-perfect}를
적용하고 이미 보인 기저변환 성질을 사용하면 결과를 얻는다.
\end{proof}

\begin{remark}
\label{perfect-remark-explain-perfect-direct-image}
$R$을 환이라 하자. $X$를 $R$ 위의 유한 표시 스킴이라 하자.
$\mathcal{G}$를 유한 표시 $\mathcal{O}_X$-가군으로서 $R$ 위에서 평탄하고
그 지지가 $R$ 위에서 고유인 것이라 하자. Lemma
\ref{perfect-lemma-base-change-tensor-perfect}에 의해 유한 사영 $R$-가군의 유한
복합체 $M^\bullet$이 존재하여
$$
R\Gamma(X_{R'}, \mathcal{G}_{R'}) = M^\bullet \otimes_R R'
$$
가 모든 $R$-대수 $R'$에 대해 함자적으로 성립한다.
\end{remark}

\begin{lemma}
\label{perfect-lemma-base-change-tensor-pseudo-coherent}
$f : X \to S$를 유한 표시 사상이라 하자.
$E \in D(\mathcal{O}_X)$를 유사연접 대상이라 하자.
$\mathcal{G}^\bullet$을 유한 표시 $\mathcal{O}_X$-가군의 위로 유계인
복합체로서 $S$ 위에서 평탄하고 그 지지가 $S$ 위에서 고유인 것이라 하자.
그러면
$$
K = Rf_*(E \otimes_{\mathcal{O}_X}^\mathbf{L} \mathcal{G}^\bullet)
$$
는 $D(\mathcal{O}_S)$의 유사연접 대상이며, 그 구성은 임의의 기저변환과
가환한다.
\end{lemma}

\begin{proof}
기저변환에 관한 명제는 Lemma \ref{perfect-lemma-base-change-tensor}이다.
따라서 $K$가 유사연접 대상임을 보이면 충분하다. 이는 완전 복합체에 의한
근사를 통해 Lemma \ref{perfect-lemma-base-change-tensor-perfect}에서 따를 것이다.
독자는 이 증명의 나머지를 건너뛰어도 좋다.

\medskip\noindent
이 문제는 $S$ 위에서 국소적이므로 $S$가 아핀이라고 가정해도 된다.
그러면 $X$는 준콤팩트이고 준분리이다. 더욱이 Lemma
\ref{perfect-lemma-quasi-coherence-direct-image}에서 설명한 바와 같이 전체 직접상
$Rf_* : D_\QCoh(\mathcal{O}_X) \to D_\QCoh(\mathcal{O}_S)$의
코호몰로지 차원이 $N$이 되는 정수 $N$이 존재한다. $i > b$에 대해
$\mathcal{G}^i = 0$이 되도록 정수 $b$를 택한다. 모든 $m$에 대해 $K$가
$m$-유사연접임을 보이면 충분하다. $(X, E, m - N - 1 - b)$에 대해 완전
복합체 $P$에 의한 근사 $P \to E$를 택한다. 이는 Theorem
\ref{perfect-theorem-approximation}에 의해 가능하다. 다음 구별 삼각형을
$$
P \to E \to C \to P[1]
$$
$D_\QCoh(\mathcal{O}_X)$에서 택하자. $C$의 코호몰로지 층들은 차수
$\geq m - N - 1 - b$에서 0이다. 따라서
$C \otimes^\mathbf{L} \mathcal{G}^\bullet$의 코호몰로지 층들은 차수
$\geq m - N - 1$에서 0이다. 그러므로
$Rf_*(C \otimes^\mathbf{L} \mathcal{G}^\bullet)$의 코호몰로지 층들은
차수 $\geq m - 1$에서 0이다. 따라서
$$
Rf_*(P \otimes^\mathbf{L} \mathcal{G}^\bullet) \to
Rf_*(E \otimes^\mathbf{L} \mathcal{G}^\bullet)
$$
는 차수 $\geq m$의 코호몰로지 층들 위에서 동형사상이다. 다음으로
$i > a$에 대해 $H^i(P) = 0$이라고 하자. 그러면
$
P \otimes^\mathbf{L} \sigma_{\geq m - N - 1 - a}\mathcal{G}^\bullet
\longrightarrow
P \otimes^\mathbf{L} \mathcal{G}^\bullet
$
는 차수 $\geq m - N - 1$의 코호몰로지 층들 위에서 동형사상이다.
따라서 다시 다음 사상
$$
Rf_*(P \otimes^\mathbf{L} \sigma_{\geq m - N - 1 - a}\mathcal{G}^\bullet) \to
Rf_*(P \otimes^\mathbf{L} \mathcal{G}^\bullet)
$$
이 차수 $\geq m$의 코호몰로지 층들 위에서 동형사상임을 얻는다. Lemma
\ref{perfect-lemma-base-change-tensor-perfect}에 의해 그 원천은 완전 복합체이다.
따라서 원한 대로 $K$가 $m$-유사연접임을 결론짓는다.
\end{proof}

\begin{lemma}
\label{perfect-lemma-flat-proper-perfect-direct-image-general}
$S$를 스킴이라 하자. $f : X \to S$를 유한 표시 고유 사상이라 하자.
\begin{enumerate}
\item $E \in D(\mathcal{O}_X)$가 완전이고 $f$가 평탄하다고 하자. 그러면
$Rf_*E$는 $D(\mathcal{O}_S)$의 완전 대상이며, 그 구성은 임의의
기저변환과 가환한다.
\item $\mathcal{G}$를 유한 표시 $\mathcal{O}_X$-가군으로서 $S$ 위에서
평탄한 것이라 하자. 그러면 $Rf_*\mathcal{G}$는 $D(\mathcal{O}_S)$의
완전 대상이며, 그 구성은 임의의 기저변환과 가환한다.
\end{enumerate}
\end{lemma}

\begin{proof}
Lemma \ref{perfect-lemma-base-change-tensor-perfect}에서 (1)은
$\mathcal{G}^\bullet$을 차수 $0$에 놓인 $\mathcal{O}_X$와 같게 하고,
(2)는 $E = \mathcal{O}_X$로 두며 $\mathcal{G}^\bullet$을 차수 $0$에
놓인 $\mathcal{G}$로 이루어지게 한 특수한 경우이다.
\end{proof}

\begin{lemma}
\label{perfect-lemma-flat-proper-pseudo-coherent-direct-image-general}
$S$를 스킴이라 하자. $f : X \to S$를 유한 표시 평탄 고유 사상이라 하자.
$E \in D(\mathcal{O}_X)$가 유사연접이라고 하자. 그러면 $Rf_*E$는
$D(\mathcal{O}_S)$의 유사연접 대상이며, 그 구성은 임의의 기저변환과
가환한다.
\end{lemma}

\noindent
더 일반적으로 $f : X \to S$가 고유이고 $X$ 위의 $E$가 $S$에 대해
상대적으로 유사연접이면(More on Morphisms, Definition
\ref{more-morphisms-definition-relative-pseudo-coherence}), $Rf_*E$는
유사연접이다(그러나 이 일반성에서는 그 구성이 기저변환과 가환하지 않는다).
\cite{Kiehl}를 보라.

\begin{proof}
Lemma \ref{perfect-lemma-base-change-tensor-pseudo-coherent}에서
$\mathcal{G}^\bullet$을 차수 $0$에 놓인 $\mathcal{O}_X$와 같게 한
특수한 경우이다.
\end{proof}

\begin{lemma}
\label{perfect-lemma-pullback-and-limits}
$R$을 환이라 하자. $X$를 스킴이라 하고
$f : X \to \Spec(R)$를 고유이고 평탄하며 유한 표시인 사상이라 하자.
$(M_n)$을 전이 사상이 전사인 $R$-가군의 역계라 하자. 그러면 표준 사상
$$
\mathcal{O}_X \otimes_R (\lim M_n)
\longrightarrow
\lim \mathcal{O}_X \otimes_R M_n
$$
은 그 원천에서 그 대상에 $DQ_X$를 적용한 것으로 가는 동형사상을 유도한다.
\end{lemma}

\begin{proof}
이 명제는 $D_\QCoh(\mathcal{O}_X)$의 임의의 대상 $E$에 대해 유도된 사상
$$
\Hom(E, \mathcal{O}_X \otimes_R (\lim M_n))
\longrightarrow
\Hom(E, \lim \mathcal{O}_X \otimes_R M_n)
$$
이 동형사상이라는 뜻이다. $D_\QCoh(\mathcal{O}_X)$는 완전 생성자를
가지므로(Theorem \ref{perfect-theorem-bondal-van-den-Bergh}), 완전인 $E$에 대해
이를 확인하면 충분하다. Lemma \ref{perfect-lemma-Rlim-quasi-coherent}에 의해
$\lim \mathcal{O}_X \otimes_R M_n = R\lim \mathcal{O}_X \otimes_R M_n$이다.
Cohomology, Section \ref{cohomology-section-global-RHom}의 완전 함자
$R\Hom_X(E, -) : D_\QCoh(\mathcal{O}_X) \to D(R)$는 곱과 가환하고,
따라서 유도 극한과도 가환한다. 그러므로
$$
R\Hom_X(E, \lim \mathcal{O}_X \otimes_R M_n) =
R\lim R\Hom_X(E, \mathcal{O}_X \otimes_R M_n)
$$
$E^\vee$를 쌍대 완전 복합체라 하자. Cohomology, Lemma
\ref{cohomology-lemma-dual-perfect-complex}를 보라. 다음을 얻는다.
$$
R\Hom_X(E, \mathcal{O}_X \otimes_R M_n) =
R\Gamma(X, E^\vee \otimes_{\mathcal{O}_X}^\mathbf{L} Lf^*M_n) =
R\Gamma(X, E^\vee) \otimes_R^\mathbf{L} M_n
$$
이는 Lemma \ref{perfect-lemma-cohomology-base-change}에 의한다. Lemma
\ref{perfect-lemma-flat-proper-perfect-direct-image-general}에서
$R\Gamma(X, E^\vee)$가 $R$-가군의 완전 복합체임을 알 수 있다.
특히 이는 유사연접 복합체이므로 More on Algebra, Lemma
\ref{more-algebra-lemma-pseudo-coherent-tensor-limit}에 의해
$$
R\lim R\Gamma(X, E^\vee) \otimes_R^\mathbf{L} M_n =
R\Gamma(X, E^\vee) \otimes_R^\mathbf{L} \lim M_n
$$
를 얻으며, 이것이 원하는 바이다.
\end{proof}

\begin{lemma}
\label{perfect-lemma-base-change-RHom-perfect}
$f : X \to S$를 유한 표시 사상이라 하자.
$E \in D(\mathcal{O}_X)$를 완전 대상이라 하자. $\mathcal{G}^\bullet$을
유한 표시 $\mathcal{O}_X$-가군의 유계 복합체로서 $S$ 위에서 평탄하고
그 지지가 $S$ 위에서 고유인 것이라 하자. 그러면
$$
K = Rf_*R\SheafHom(E, \mathcal{G}^\bullet)
$$
는 $D(\mathcal{O}_S)$의 완전 대상이며, 그 구성은 임의의 기저변환과
가환한다.
\end{lemma}

\begin{proof}
기저변환에 관한 명제는 Lemma \ref{perfect-lemma-base-change-RHom}이다.
따라서 $K$가 완전 대상임을 보이면 충분하다. $S$가 뇌터이면 이는
Lemma \ref{perfect-lemma-ext-perfect}에서 따른다. 뇌터 근사를 사용하여 이 경우로
환원하겠다. 독자는 이 증명의 나머지를 건너뛰어도 좋다.

\medskip\noindent
이 문제는 $S$ 위에서 국소적이므로 $S$가 아핀이라고 가정해도 된다.
$S = \Spec(R)$라 하자. 뇌터 환 $R_i$들의 여과 귀납극한으로
$R = \colim R_i$라 쓴다. Limits, Lemma
\ref{limits-lemma-descend-finite-presentation}에 의해 어떤 $i$와
$R_i$ 위의 유한 표시 스킴 $X_i$가 존재하여 이를 $R$로 기저변환하면
$X$가 된다. Limits, Lemma
\ref{limits-lemma-descend-modules-finite-presentation}에 의해 $i$를 크게
잡은 뒤에는 유한 표시 $\mathcal{O}_{X_i}$-가군의 유계 복합체
$\mathcal{G}_i^\bullet$이 존재하고, 이를 $X$로 역상하면
$\mathcal{G}^\bullet$이 된다고 가정할 수 있다. $i$를 크게 잡아
$\mathcal{G}_i^n$이 $R_i$ 위에서 평탄하다고 가정할 수 있다. Limits,
Lemma \ref{limits-lemma-descend-module-flat-finite-presentation}를 보라.
다시 $i$를 크게 잡아 $\mathcal{G}_i^n$의 지지가 $R_i$ 위에서 고유라고
가정할 수 있다. Limits, Lemma
\ref{limits-lemma-eventually-proper-support}와 Cohomology of Schemes,
Lemma \ref{coherent-lemma-module-support-proper-over-base}를 보라. 끝으로
Lemma \ref{perfect-lemma-descend-perfect}에 의해 $i$를 크게 잡은 뒤에는
$D(\mathcal{O}_{X_i})$의 완전 대상 $E_i$가 존재하여 이를 $X$로
역상하면 $E$가 된다고 가정할 수 있다. $X_i \to \Spec(R_i)$,
$E_i$, $\mathcal{G}_i^\bullet$에 Lemma \ref{perfect-lemma-ext-perfect}를 적용하고
이미 보인 기저변환 성질을 사용하면 결과를 얻는다.
\end{proof}









\section{완전 복합체}
\label{perfect-section-perfect-complexes}

\noindent
먼저 완전 복합체의 베티 수가 점프하는 자취를 논한다. 스킴 $X$ 위의 복합체
$E$와 $X$의 점 $x$가 주어졌을 때, 표준 사상 $i_x : x \to X$에 대해 더
정확한 표기 $Li_x^*E$ 대신 흔히
$E \otimes_{\mathcal{O}_X}^\mathbf{L} \kappa(x)$라고 쓴다.

\begin{lemma}
\label{perfect-lemma-jump-loci}
$X$를 스킴이라 하자. $E \in D(\mathcal{O}_X)$가 유사연접이라고 하자
(예를 들어 완전이라고 하자). 임의의 $i \in \mathbf{Z}$에 대해 다음 함수를
생각하자.
$$
\beta_i : X \longrightarrow \{0, 1, 2, \ldots\},\quad
x \longmapsto
\dim_{\kappa(x)}
H^i(E \otimes_{\mathcal{O}_X}^\mathbf{L} \kappa(x))
$$
그러면 다음이 성립한다.
\begin{enumerate}
\item $\beta_i$의 구성은 임의의 기저변환과 가환한다.
\item 함수 $\beta_i$는 상반연속이다.
\item $\beta_i$의 수준집합들은 $X$에서 국소적으로 구성 가능하다.
\end{enumerate}
\end{lemma}

\begin{proof}
스킴의 사상 $f : Y \to X$와 점 $y \in Y$를 생각하자. $x$를 $y$의
상이라 하고 다음 가환 그림을 생각하자.
$$
\xymatrix{
y \ar[r]_j \ar[d]_g & Y \ar[d]^f \\
x \ar[r]^i & X
}
$$
그러면 $Lg^* \circ Li^* = Lj^* \circ Lf^*$임을 알 수 있다. 따라서
유사연접 복합체 $Lf^*E$에 대응하는 함수 $\beta'_i$는 함수 $\beta_i$의
역상이다. 식으로 쓰면 $\beta'_i = \beta_i \circ f$이다. 이것이 (1)의
의미이다.

\medskip\noindent
$i$를 고정하고 $x \in X$라 하자. $x$의 어떤 열린 근방에서 (2)와 (3)이
성립함을 보이면 충분하므로 $X$가 아핀이라고 가정해도 된다. 그러면 $E$를
유한 자유 가군의 위로 유계인 복합체 $\mathcal{F}^\bullet$로 나타낼 수
있다(Lemma \ref{perfect-lemma-lift-pseudo-coherent}). 이때
$P = \sigma_{\geq i - 1}\mathcal{F}^\bullet$은 완전 대상이고, $P \to E$는
모든 $x' \in X$에 대해 동형사상
$$
H^i(P \otimes_{\mathcal{O}_X}^\mathbf{L} \kappa(x')) \to
H^i(E \otimes_{\mathcal{O}_X}^\mathbf{L} \kappa(x'))
$$
을 유도한다. 따라서 $E$가 완전이라고 가정해도 된다. 이 경우 More on
Algebra, Lemma \ref{more-algebra-lemma-lift-perfect-from-residue-field}에
의해 $x$의 아핀 열린 근방 $U$와 $a \leq b$가 존재하여 $E|_U$는 다음
복합체로 나타난다.
$$
\ldots \to 0 \to \mathcal{O}_U^{\oplus \beta_a(x)}
\to \mathcal{O}_U^{\oplus \beta_{a + 1}(x)} \to
\ldots \to
\mathcal{O}_U^{\oplus \beta_{b - 1}(x)} \to
\mathcal{O}_U^{\oplus \beta_b(x)} \to 0
\to \ldots
$$
(여기서는 문제를 대수로 바꾸는 앞의 결과들, 예를 들어 Lemmas
\ref{perfect-lemma-affine-compare-bounded}와 \ref{perfect-lemma-perfect-affine}도 사용한다.)
모든 $x' \in U$에 대해 $\beta_i(x') \leq \beta_i(x)$임이 즉시 따른다.
따라서 $\beta_i$는 상반연속이다.

\medskip\noindent
(3)을 증명하기 위해 $X$가 아핀이고 $E$가 유한 자유
$\mathcal{O}_X$-가군의 복합체로 주어진다고 가정해도 된다(예를 들어 앞
문단처럼 논하거나 Cohomology, Lemma
\ref{cohomology-lemma-perfect-on-locally-ringed}를 사용한다). 따라서 다음
복합체가 주어졌을 때
$$
\mathcal{O}_X^{\oplus a} \to
\mathcal{O}_X^{\oplus b} \to
\mathcal{O}_X^{\oplus c}
$$
점 $x \in X$에 복합체
$\kappa_x^{\oplus a} \to \kappa_x^{\oplus b} \to \kappa_x^{\oplus c}$의
가운데 코호몰로지 차원을 대응시키는 함수의 수준집합들이 구성 가능함을
보여야 한다. $A \in \text{Mat}(a \times b, \Gamma(X, \mathcal{O}_X))$를
첫째 화살표의 행렬이라 하자. $\text{Mat}(a \times b, \kappa(x))$에서
$A$의 상의 계수가 $r$인 것은 $A$의 모든 $(r + 1) \times (r + 1)$
소행렬식이 $x$에서 사라지고, $x$에서 사라지지 않는 $A$의 어떤
$r \times r$ 소행렬식이 있는 것과 동치이다. 따라서 계수가 $r$인 점들의
집합은 구성 가능한 국소 닫힌집합이다. 둘째 화살표에 대해서도 마찬가지로
논하고 이들을 합치면 원하는 결과를 얻는다.
\end{proof}

\begin{lemma}
\label{perfect-lemma-chi-locally-constant}
$X$를 스킴이라 하자. $E \in D(\mathcal{O}_X)$가 완전이라고 하자. 함수
$$
\chi_E : X \longrightarrow \mathbf{Z},\quad
x \longmapsto \sum (-1)^i
\dim_{\kappa(x)} H^i(E \otimes_{\mathcal{O}_X}^\mathbf{L} \kappa(x))
$$
는 $X$에서 국소 상수이다.
\end{lemma}

\begin{proof}
Cohomology, Lemma \ref{cohomology-lemma-perfect-on-locally-ringed}에 의해
$X$ 위에서 국소적으로 $E$를 유한 자유 $\mathcal{O}_X$-가군의 유한
복합체 $\mathcal{E}^\bullet$로 나타낼 수 있다. 그러한 열린집합에서 함수
$\chi_E$는 값 $\sum (-1)^i \text{rank}(\mathcal{E}^i)$를 갖는 상수이다.
\end{proof}

\begin{lemma}
\label{perfect-lemma-open-where-cohomology-in-degree-i-rank-r}
$X$를 스킴이라 하자. $E \in D(\mathcal{O}_X)$가 완전이라고 하자.
$i, r \in \mathbf{Z}$가 주어지면 다음으로 특징지어지는 열린 부분스킴
$U \subset X$가 존재한다.
\begin{enumerate}
\item $E|_U \cong H^i(E|_U)[-i]$이고 $H^i(E|_U)$는 계수가 $r$인
국소 자유 $\mathcal{O}_U$-가군이다.
\item 사상 $f : Y \to X$가 $U$를 통해 분해될 필요충분조건은 $Lf^*E$가
차수 $i$에 놓인 계수 $r$의 국소 자유 가군과 동형인 것이다.
\end{enumerate}
\end{lemma}

\begin{proof}
$j \in \mathbf{Z}$에 대해 $\beta_j : X \to \{0, 1, 2, \ldots\}$를
Lemma \ref{perfect-lemma-jump-loci}의 함수라 하자. 그러면 집합
$$
W = \{x \in X \mid \beta_j(x) \leq 0\text{ 모든 }j \not = i\text{에 대하여}\}
$$
은 $X$에서 열려 있고, 그 구성은 $X$ 위의 임의의 $Y$로의 역상과
가환한다. 이는 임의의 점의 근방에서 0이 아닌 $\beta_j$가 선험적으로
유한 개뿐이라는 사실과 그 보조정리에서 따른다. 따라서 $X$를 $W$로
바꾸어 모든 $x \in X$와 모든 $j \not = i$에 대해 $\beta_j(x) = 0$이라고
가정해도 된다. 이 경우 $H^i(E)$는 유한 국소 자유 가군이고
$E \cong H^i(E)[-i]$이다. 예를 들어 More on Algebra, Lemma
\ref{more-algebra-lemma-lift-perfect-from-residue-field}를 보라. 따라서
$X$는 $H^i(E)$의 계수가 고정된 열린 부분스킴들의 서로소 합집합이며,
결론을 얻는다.
\end{proof}

\begin{lemma}
\label{perfect-lemma-locally-closed-where-H0-locally-free}
$X$를 스킴이라 하자. $E \in D(\mathcal{O}_X)$가 완전이고 어떤
$a, b \in \mathbf{Z}$에 대해 Tor-진폭이 $[a, b]$에 들어 있다고 하자.
$r \geq 0$이라 하자. 그러면 다음으로 특징지어지는 국소 닫힌 부분스킴
$j : Z \to X$가 존재한다.
\begin{enumerate}
\item $H^a(Lj^*E)$는 계수가 $r$인 국소 자유 $\mathcal{O}_Z$-가군이다.
\item 사상 $f : Y \to X$가 $Z$를 통해 분해될 필요충분조건은 모든 사상
$g : Y' \to Y$에 대해 $\mathcal{O}_{Y'}$-가군
$H^a(L(f \circ g)^*E)$가 계수 $r$의 국소 자유 가군인 것이다.
\end{enumerate}
더욱이 $j : Z \to X$는 유한 표시이고 다음이 성립한다.
\begin{enumerate}
\item[(3)] $f : Y \to X$가 $Y \xrightarrow{g} Z \to X$로 분해되면
$H^a(Lf^*E) = g^*H^a(Lj^*E)$이다.
\item[(4)] 모든 $x \in X$에 대해 $\beta_a(x) \leq r$이면 $j$는 닫힌
몰입이고, 주어진 $f : Y \to X$에 대해 다음은 동치이다.
\begin{enumerate}
\item $f : Y \to X$는 $Z$를 통해 분해된다.
\item $H^a(Lf^*E)$는 계수가 $r$인 국소 자유 $\mathcal{O}_Y$-가군이다.
\end{enumerate}
$r = 1$이면 이들은 또한 다음과 동치이다.
\begin{enumerate}
\item[(c)] $\mathcal{O}_Y \to \SheafHom_{\mathcal{O}_Y}(H^a(Lf^*E), H^a(Lf^*E))$
는 단사이다.
\end{enumerate}
\end{enumerate}
\end{lemma}

\begin{proof}
먼저 $U \subset X$를 Lemma \ref{perfect-lemma-jump-loci}의 함수 $\beta_a$가
$\leq r$인 값을 갖는 국소적으로 구성 가능한 열린 부분스킴이라 하자.
(2)에서와 같은 $f : Y \to X$를 택하자. 그러면 임의의 $y \in Y$에 대해
$\beta_a(Lf^*E) = r$이므로 Lemma \ref{perfect-lemma-jump-loci}에 의해 $y$는
$U$로 사상된다. 따라서 (2)에서와 같은 $f$는 $U$를 통해 분해된다.
그러므로 $X$를 $U$로 바꾸어 모든 $x \in X$에 대해
$\beta_a(x) \in \{0, 1, \ldots, r\}$이라고 가정해도 된다. 이 경우
$r = 1$일 때 (4)(a), (b), (4)(c)의 동치로 특징지어지고 (3)을 만족하는,
유한형 준연접 아이디얼로 잘라낸 닫힌 부분스킴 $Z \subset X$가 있음을
보이겠다. 그러면 특히 (2)에서와 같은 사상들이 $Z$를 통해 분해됨을
보이므로 증명이 끝난다.

\medskip\noindent
$x \in X$이고 $\beta_a(x) < r$이면, $\beta_a < r$인 $x$의 열린 근방이
존재한다(Lemma \ref{perfect-lemma-jump-loci}). 이로부터 적어도 집합론적으로
$Z$가 닫힌 부분집합임을 알 수 있다.

\medskip\noindent
스킴론적 구조를 얻기 위해 $\beta_a(x) = r$인 점 $x \in X$를 생각하자.
$\beta = \beta_{a + 1}(x)$로 놓는다. More on Algebra, Lemma
\ref{more-algebra-lemma-lift-perfect-from-residue-field}에 의해 $x$의 아핀
열린 근방 $U$가 존재하여 $K|_U$는 다음 복합체로 나타난다.
$$
\ldots \to 0 \to \mathcal{O}_U^{\oplus r}
\xrightarrow{(f_{ij})} \mathcal{O}_U^{\oplus \beta} \to
\ldots \to
\mathcal{O}_U^{\oplus \beta_{b - 1}(x)} \to
\mathcal{O}_U^{\oplus \beta_b(x)} \to 0
\to \ldots
$$
(여기서는 문제를 대수로 바꾸는 앞의 결과들, 예를 들어 Lemmas
\ref{perfect-lemma-affine-compare-bounded}와 \ref{perfect-lemma-perfect-affine}도 사용한다.)
이제 $g : Y \to U$가 스킴의 임의의 사상이고 어떤 쌍 $i, j$에 대해
$g^\sharp(f_{ij})$가 0이 아니면, $H^a(Lg^*E)$는 계수가 $r$인 국소 자유
$\mathcal{O}_Y$-가군이 아니다. More on Algebra, Lemma
\ref{more-algebra-lemma-coker-injective-free}를 보라. 모든 $i, j$에 대해
$g^\sharp(f_{ij}) = 0$이면 $H^a(Lg^*E)$가 국소 자유
$\mathcal{O}_Y$-가군임은 자명하다. 따라서 $U$ 위에서는 모든 $f_{ij}$로
잘라낸 닫힌 부분스킴이 (3)을 만족하고, (4)(a)와 (b)가 동치임을 알 수 있다.
$Z$의 특징짓기에 의해 국소적으로 구성한 조각들이 접합된다(세부사항은
생략한다). 끝으로 $r = 1$이면, 이 경우 국소적으로
$H^a(Lg^*E) \subset \mathcal{O}_Y$가 원소 $g^\sharp(f_{ij})$들이 생성하는
아이디얼의 소멸자이므로 (4)(c)는 (4)(b)와 동치이다.
\end{proof}





\section{응용}
\label{perfect-section-applications}

\noindent
대부분 코호몰로지와 기저변환의 응용이다. 추후 이 결과들을 Lemma
\ref{perfect-lemma-base-change-tensor-perfect}에서 논한 상황으로 일반화할 수 있다.

\begin{lemma}
\label{perfect-lemma-jump-loci-geometric}
$f : X \to S$를 유한 표시 고유 사상이라 하자. $\mathcal{F}$를 유한 표시
$\mathcal{O}_X$-가군으로서 $S$ 위에서 평탄한 것이라 하자. 고정된
$i \in \mathbf{Z}$에 대해 다음 함수를 생각하자.
$$
\beta_i : S \to \{0, 1, 2, \ldots\},\quad
s \longmapsto \dim_{\kappa(s)} H^i(X_s, \mathcal{F}_s)
$$
그러면 다음이 성립한다.
\begin{enumerate}
\item $\beta_i$의 구성은 임의의 기저변환과 가환한다.
\item 함수 $\beta_i$는 상반연속이다.
\item $\beta_i$의 수준집합들은 $S$에서 국소적으로 구성 가능하다.
\end{enumerate}
\end{lemma}

\begin{proof}
코호몰로지와 기저변환에 의해(더 정확히는 Lemma
\ref{perfect-lemma-flat-proper-perfect-direct-image-general}에 의해), 대상
$K = Rf_*\mathcal{F}$는 $S$의 유도 범주의 완전 대상이고 그 구성은
임의의 기저변환과 가환한다. 특히
$$
H^i(X_s, \mathcal{F}_s) = H^i(K \otimes_{\mathcal{O}_S}^\mathbf{L} \kappa(s))
$$
를 얻는다. 따라서 보조정리는 Lemma \ref{perfect-lemma-jump-loci}에서 따른다.
\end{proof}

\begin{lemma}
\label{perfect-lemma-chi-locally-constant-geometric}
$f : X \to S$를 유한 표시 고유 사상이라 하자. $\mathcal{F}$를 유한 표시
$\mathcal{O}_X$-가군으로서 $S$ 위에서 평탄한 것이라 하자. 함수
$$
s \longmapsto \chi(X_s, \mathcal{F}_s)
$$
는 $S$에서 국소 상수이다. 이 함수의 구성은 기저변환과 가환한다.
\end{lemma}

\begin{proof}
코호몰로지와 기저변환에 의해(더 정확히는 Lemma
\ref{perfect-lemma-flat-proper-perfect-direct-image-general}에 의해), 대상
$K = Rf_*\mathcal{F}$는 $S$의 유도 범주의 완전 대상이고 그 구성은
임의의 기저변환과 가환한다. 따라서 다음 사상이
$$
s \longmapsto \sum (-1)^i \dim_{\kappa(s)}
H^i(K \otimes^\mathbf{L}_{\mathcal{O}_S} \kappa(s))
$$
$S$에서 국소 상수임을 보여야 한다. 이는 Lemma
\ref{perfect-lemma-chi-locally-constant}이다.
\end{proof}

\begin{lemma}
\label{perfect-lemma-open-where-cohomology-in-degree-i-rank-r-geometric}
$f : X \to S$를 유한 표시 고유 사상이라 하자. $\mathcal{F}$를 유한 표시
$\mathcal{O}_X$-가군으로서 $S$ 위에서 평탄한 것이라 하자.
$i, r \in \mathbf{Z}$를 고정하자. 그러면 다음 성질을 갖는 열린 부분스킴
$U \subset S$가 존재한다. 사상 $T \to S$가 $U$를 통해 분해될
필요충분조건은 $Rf_{T, *}\mathcal{F}_T$가 차수 $i$에 놓인 계수 $r$의
유한 국소 자유 가군과 동형인 것이다.
\end{lemma}

\begin{proof}
코호몰로지와 기저변환에 의해(더 정확히는 Lemma
\ref{perfect-lemma-flat-proper-perfect-direct-image-general}에 의해), 대상
$K = Rf_*\mathcal{F}$는 $S$의 유도 범주의 완전 대상이고 그 구성은
임의의 기저변환과 가환한다. 따라서 이 보조정리는 Lemma
\ref{perfect-lemma-open-where-cohomology-in-degree-i-rank-r}에서 즉시 따른다.
\end{proof}

\begin{lemma}
\label{perfect-lemma-vanishing-implies-locally-free}
$f : X \to S$를 유한 표시 사상이라 하자. $\mathcal{F}$를 유한 표시
$\mathcal{O}_X$-가군으로서 $S$ 위에서 평탄하고 그 지지가 $S$ 위에서
고유인 것이라 하자. $i > 0$에 대해 $R^if_*\mathcal{F} = 0$이면,
$f_*\mathcal{F}$는 국소 자유이고 그 구성은 임의의 기저변환과 가환한다
(설명은 증명을 보라).
\end{lemma}

\begin{proof}
Lemma \ref{perfect-lemma-base-change-tensor-perfect}에 의해 $D(\mathcal{O}_S)$의
대상 $E = Rf_*\mathcal{F}$는 완전이고 그 구성은 임의의 기저변환과
가환한다. 즉 임의의 데카르트 그림
$$
\xymatrix{
X' \ar[r]_{g'} \ar[d]_{f'} &
X \ar[d]^f \\
S' \ar[r]^g &
S
}
$$
에 대해 $Rf'_*(g')^*\mathcal{F} = Lg^*E$이다. 차수 $< 0$에서는
코호몰로지가 전혀 없으므로, 어떤 $b$에 대해 $E$의 Tor-진폭이 국소적으로
$[0, b]$에 들어 있음을 알 수 있다. $i > 0$에 대해
$H^i(E) = R^if_*\mathcal{F} = 0$이면 $E$의 Tor-진폭은 $[0, 0]$에 들어
있다. 따라서 $E = H^0(E)[0]$이다. More on Algebra, Lemma
\ref{more-algebra-lemma-perfect}에 의해(또한 Algebra, Lemma
\ref{algebra-lemma-finite-projective}의 유한 사영 가군의 특징짓기에 의해)
$H^0(E) = f_*\mathcal{F}$가 유한 국소 자유임을 결론짓는다. 기저변환과의
가환은 위와 같은 그림에 대해 $g^*f_*\mathcal{F} = f'_*(g')^*\mathcal{F}$라는
뜻이며, 이는 이미 확립한 $E$의 기저변환과의 가환성에서 따른다.
\end{proof}

\begin{lemma}
\label{perfect-lemma-proper-flat-h0}
$f : X \to S$를 스킴의 사상이라 하자. 다음을 가정하자.
\begin{enumerate}
\item $f$는 고유이고 평탄이며 유한 표시이다.
\item 모든 $s \in S$에 대해 $\kappa(s) = H^0(X_s, \mathcal{O}_{X_s})$이다.
\end{enumerate}
그러면 다음이 성립한다.
\begin{enumerate}
\item[(a)] $f_*\mathcal{O}_X = \mathcal{O}_S$이고 이는 임의의 기저변환
뒤에도 성립한다.
\item[(b)] $S$ 위에서 국소적으로
$$
Rf_*\mathcal{O}_X = \mathcal{O}_S \oplus P
$$
가 $D(\mathcal{O}_S)$에서 성립하며, 여기서 $P$는 완전이고 Tor-진폭이
$[1, \infty)$에 들어 있다.
\end{enumerate}
\end{lemma}

\begin{proof}
코호몰로지와 기저변환에 의해(Lemma
\ref{perfect-lemma-flat-proper-perfect-direct-image-general}), 복합체
$E = Rf_*\mathcal{O}_X$는 완전이고 그 구성은 임의의 기저변환과
가환한다. 첫째, 이로부터 $E$의 Tor-진폭이 $[0, \infty)$에 들어 있음을
얻는다. 둘째, $s \in S$에 대해 다음을 얻는다.
$H^0(E \otimes^\mathbf{L} \kappa(s)) =
H^0(X_s, \mathcal{O}_{X_s}) = \kappa(s)$.
따라서 사상 $\mathcal{O}_S \to Rf_*\mathcal{O}_X = E$는
동형사상
$\mathcal{O}_S \otimes \kappa(s) \to H^0(E \otimes^\mathbf{L} \kappa(s))$를
유도한다. 그러므로
$H^0(E) \otimes \kappa(s) \to H^0(E \otimes^\mathbf{L} \kappa(s))$는
전사이고, More on Algebra, Lemma
\ref{more-algebra-lemma-better-cut-complex-in-two}를 적용하면 $S$를 $s$의
아핀 열린 근방으로 바꾼 뒤 분해 $E = H^0(E) \oplus \tau_{\geq 1}E$를
얻는다. 여기서 $\tau_{\geq 1}E$는 완전이고 Tor-진폭이
$[1, \infty)$에 들어 있다. $E$의 Tor-진폭이 $[0, \infty)$에 들어
있으므로 $H^0(E)$는 평탄 $\mathcal{O}_S$-가군이다. 따라서 $H^0(E)$는
평탄 완전 $\mathcal{O}_S$-가군이고, 이에 따라 유한 국소 자유이다. More on
Algebra, Lemma \ref{more-algebra-lemma-perfect}를 보라(또한 Algebra, Lemma
\ref{algebra-lemma-finite-projective}에 의해 유한 사영 가군은 유한 국소
자유이다). 사상 $\mathcal{O}_S \to H^0(E)$가 동형사상인지 여부는 잉여체와
텐서한 뒤 확인할 수 있으므로 이 사상은 동형사상이다(Algebra, Lemma
\ref{algebra-lemma-cokernel-flat}).
\end{proof}

\begin{lemma}
\label{perfect-lemma-proper-flat-geom-red-connected}
$f : X \to S$를 스킴의 사상이라 하자. 다음을 가정하자.
\begin{enumerate}
\item $f$는 고유이고 평탄이며 유한 표시이다.
\item $f$의 기하적 올들은 축약이며 연결이다.
\end{enumerate}
그러면 $f_*\mathcal{O}_X = \mathcal{O}_S$이고 이는 임의의 기저변환
뒤에도 성립한다.
\end{lemma}

\begin{proof}
Lemma \ref{perfect-lemma-proper-flat-h0}에 의해 모든 $s \in S$에 대해
$\kappa(s) = H^0(X_s, \mathcal{O}_{X_s})$임을 보이면 충분하다. 이는
Varieties, Lemma
\ref{varieties-lemma-proper-geometrically-reduced-global-sections}와
$X_s$가 기하적으로 연결이고 기하적으로 축약이라는 사실에서 따른다.
\end{proof}

\begin{lemma}
\label{perfect-lemma-proper-idempotent-on-fibre}
$f : X \to S$를 스킴의 고유 사상이라 하자. $s \in S$라 하고
$e \in H^0(X_s, \mathcal{O}_{X_s})$를 멱등원이라 하자. 그러면 $e$는 사상
$(f_*\mathcal{O}_X)_s \to H^0(X_s, \mathcal{O}_{X_s})$의 상에 속한다.
\end{lemma}

\begin{proof}
$X_s = T_1 \amalg T_2$를 $e$에 대응하는 서로소 합집합 분해라 하자.
여기서 $T_1$과 $T_2$는 공집합이 아니며 $X_s$에서 열리고 닫혀 있고,
$e$는 $T_1$에서 항등적으로 $1$, $T_2$에서 항등적으로 $0$이다.

\medskip\noindent
$S$가 뇌터라고 가정하자. 형식 함수 정리를 Cohomology of Schemes,
Lemma \ref{coherent-lemma-formal-functions-stalk}의 형태로 사용하겠다.
이에 따르면
$$
(f_*\mathcal{O}_X)_s^\wedge = \lim_n H^0(X_n, \mathcal{O}_{X_n})
$$
이며, 여기서 $X_n$은 $X_s$의 $n$번째 무한소 근방이다. $X_n$의 밑
위상공간은 $X_s$의 밑 위상공간과 같으므로, 모든 $n $에 대해 스킴의
서로소 합집합 분해 $X_n = T_{1, n} \amalg T_{2, n}$을 얻는다. 여기서
$i = 1, 2$에 대해 $T_{i, n}$의 밑 위상공간은 $T_i$이다. 따라서
$H^0(X_n, \mathcal{O}_{X_n})$은 자명하지 않은 멱등원 $e_n$을 포함한다.
이는 $T_{1, n}$에서 항등적으로 $1$이고 $T_{2, n}$에서 항등적으로 $0$인
함수이다. $e_{n + 1}$을 $X_n$으로 제한하면 $e_n$이 됨은 명백하다.
따라서 $e_\infty = \lim e_n$은 그 극한의 자명하지 않은 멱등원이다.
그러므로 $e_\infty$는 $(f_*\mathcal{O}_X)_s$의 완비화의 원소이고,
$H^0(X_s, \mathcal{O}_{X_s})$의 $e$로 사상된다. 사상
$(f_*\mathcal{O}_X)_s^\wedge \to H^0(X_s, \mathcal{O}_{X_s})$는
$(f_*\mathcal{O}_X)^\wedge_s / \mathfrak m_s (f_*\mathcal{O}_X)_s^\wedge =
(f_*\mathcal{O}_X)_s / \mathfrak m_s (f_*\mathcal{O}_X)_s$
를 통해 분해되므로(Algebra, Lemma
\ref{algebra-lemma-hathat-finitely-generated}), 원하는 대로 $e$가 사상
$(f_*\mathcal{O}_X)_s \to H^0(X_s, \mathcal{O}_{X_s})$의 상에 속한다고
결론짓는다.

\medskip\noindent
일반적인 경우에는 극한 논법으로 뇌터인 경우에 환원한다. 독자는 이 증명을
건너뛰어도 좋다. $S$를 $s$의 아핀 열린 근방으로 바꾸어도 된다. 따라서
$S$가 아핀이라고 가정한다. Limits, Lemma
\ref{limits-lemma-proper-limit-of-proper-finite-presentation-noetherian}에 의해
$(f : X \to S) = \lim (f_i : X_i \to S_i)$로 쓸 수 있다. 여기서 $f_i$는
고유이고 $S_i$는 뇌터이다. $s_i \in S_i$를 $s$의 상이라 하자. 그러면
$s = \lim s_i$이다. Limits, Lemma
\ref{limits-lemma-inverse-limit-irreducibles}를 보라. 극한들은 극한들과
가환하므로(Categories, Lemma \ref{categories-lemma-colimits-commute}),
$X_s = X \times_S s = \lim X_i \times_{S_i} s_i = \lim X_{i, s_i}$이다.
따라서 Limits, Lemma \ref{limits-lemma-descend-section}에 의해 $e$는 어떤
멱등원 $e_i \in H^0(X_{i, s_i}, \mathcal{O}_{X_{i, s_i}})$의 상이다.
뇌터인 경우에 의해 줄기 $(f_{i, *}\mathcal{O}_{X_i})_{s_i}$에는 $e_i$로
사상되는 원소 $\tilde e_i$가 있다. $\tilde e_i$의 역상을 취하면
$(f_*\mathcal{O}_X)_s$의 원소 $\tilde e$를 얻고 이는 $e$로 사상된다.
이로써 증명이 끝난다.
\end{proof}

\begin{lemma}
\label{perfect-lemma-proper-flat-geom-red}
$f : X \to S$를 스킴의 사상이라 하고 $s \in S$라 하자. 다음을 가정하자.
\begin{enumerate}
\item $f$는 고유이고 평탄이며 유한 표시이다.
\item 올 $X_s$는 기하적으로 축약이다.
\end{enumerate}
그러면 $S$를 $s$의 열린 근방으로 바꾼 뒤 직접합 분해
$Rf_*\mathcal{O}_X = f_*\mathcal{O}_X \oplus P$
가 $D(\mathcal{O}_S)$에서 존재하며, 여기서 $f_*\mathcal{O}_X$는 유한 에탈
$\mathcal{O}_S$-대수이고 $P$는 완전이며 Tor-진폭이 $[1, \infty)$에
들어 있다.
\end{lemma}

\begin{proof}
이 보조정리의 증명은 Lemma \ref{perfect-lemma-proper-flat-h0}의 증명과 비슷하다.
독자는 그 증명을 먼저 읽기를 권한다. 코호몰로지와 기저변환에 의해(Lemma
\ref{perfect-lemma-flat-proper-perfect-direct-image-general}), 복합체
$E = Rf_*\mathcal{O}_X$는 완전이고 그 구성은 임의의 기저변환과
가환한다. 이로부터 먼저 $E$의 Tor-진폭이 $[0, \infty)$에 들어 있음을
얻는다.

\medskip\noindent
$S$를 $s$의 열린 근방으로 바꾼 뒤 $D(\mathcal{O}_S)$에서 직접합 분해
$E = H^0(E) \oplus \tau_{\geq 1}E$를 찾을 수 있고, 여기서
$\tau_{\geq 1}E$의 Tor-진폭은 $[1, \infty)$에 들어 있다고 주장한다.
우선 이 주장이 참이고 이 분해를 갖도록 기저를 바꾸었다고 가정하자.
$E$의 Tor-진폭이 $[0, \infty)$에 들어 있으므로 $H^0(E)$는 평탄
$\mathcal{O}_S$-가군이다. 따라서 $H^0(E)$는 평탄 완전
$\mathcal{O}_S$-가군이고, 이에 따라 유한 국소 자유이다. More on Algebra,
Lemma \ref{more-algebra-lemma-perfect}를 보라(또한 Algebra, Lemma
\ref{algebra-lemma-finite-projective}에 의해 유한 사영 가군은 유한 국소
자유이다). 물론 $H^0(E) = f_*\mathcal{O}_X$는 $\mathcal{O}_S$-대수이다.
코호몰로지와 기저변환에 의해
$H^0(E) \otimes \kappa(s) = H^0(X_s, \mathcal{O}_{X_s})$를 얻는다.
Varieties, Lemma
\ref{varieties-lemma-proper-geometrically-reduced-global-sections}와
$X_s$가 기하적으로 축약이라는 가정에 의해
$\kappa(s) \to H^0(E) \otimes \kappa(s)$는 유한 에탈이다. 유한 국소
자유 사상 $\underline{\Spec}_S(H^0(E)) \to S$에 Morphisms, Lemma
\ref{morphisms-lemma-set-points-where-fibres-etale}를 적용하면, $S$를
축소한 뒤 $\mathcal{O}_S$-대수 $H^0(E)$가 유한 에탈이라고 결론짓는다.

\medskip\noindent
주장을 증명하는 일만 남았다. 이를 위해서는 사상
$$
(f_*\mathcal{O}_X)_s
\longrightarrow
H^0(X_s, \mathcal{O}_{X_s}) = H^0(E \otimes^\mathbf{L} \kappa(s))
$$
이 전사임을 증명하면 충분하다. More on Algebra, Lemma
\ref{more-algebra-lemma-better-cut-complex-in-two}를 보라. $A$의 잉여체
$k$가 대수적으로 닫혀 있도록 평탄 국소 환 준동형
$\mathcal{O}_{S, s} \to A$를 택하자. Algebra, Lemma
\ref{algebra-lemma-flat-local-given-residue-field}를 보라. 평탄 기저변환에
의해(Cohomology of Schemes, Lemma
\ref{coherent-lemma-flat-base-change-cohomology}),
$H^0(X_A, \mathcal{O}_{X_A}) =
(f_*\mathcal{O}_X)_s \otimes_{\mathcal{O}_{S, s}} A$ 및
$H^0(X_k, \mathcal{O}_{X_k}) =
H^0(X_s, \mathcal{O}_{X_s}) \otimes_{\kappa(s)} k$를 얻는다. 따라서
$H^0(X_A, \mathcal{O}_{X_A}) \to H^0(X_k, \mathcal{O}_{X_k})$가 전사임을
증명하면 충분하다. $X_k$는 $k$ 위의 축약 고유 스킴이고 $k$는 대수적으로
닫혀 있으므로, 앞에서 사용한 Varieties, Lemma
\ref{varieties-lemma-proper-geometrically-reduced-global-sections}에 의해
$H^0(X_k, \mathcal{O}_{X_k})$는 $k$의 유한 개 복사본의 곱이다. Lemma
\ref{perfect-lemma-proper-idempotent-on-fibre}에 의해 이 $k$-대수의 멱등원들은
$H^0(X_A, \mathcal{O}_{X_A}) \to H^0(X_k, \mathcal{O}_{X_k})$의 상에
속하므로 결론을 얻는다.
\end{proof}





\section{그 밖의 응용}
\label{perfect-section-other-applications}

\noindent
이 절에서는 앞에서 전개한 이론으로부터 도출할 수 있는 몇 가지 결과를
서술하고 증명한다.

\begin{lemma}
\label{perfect-lemma-cohomology-over-coherent-ring}
$R$을 연접 환이라 하자. $X$를 $R$ 위의 유한 표시 스킴이라 하자.
$\mathcal{G}$를 유한 표시 $\mathcal{O}_X$-가군으로서 $R$ 위에서 평탄하고
그 지지가 $R$ 위에서 고유인 것이라 하자. 그러면
$H^i(X, \mathcal{G})$는 연접 $R$-가군이다.
\end{lemma}

\begin{proof}
Lemma \ref{perfect-lemma-base-change-tensor-perfect}와 More on Algebra, Lemmas
\ref{more-algebra-lemma-coherent-pseudo-coherent} 및
\ref{more-algebra-lemma-perfect}를 결합한다.
\end{proof}

\begin{lemma}
\label{perfect-lemma-countable-cohomology}
$X$를 준콤팩트이고 준분리인 스킴이라 하자. $K$를
$D_\QCoh(\mathcal{O}_X)$의 대상으로서 코호몰로지 층 $H^i(K)$의 아핀
열린집합 위 절단들의 집합이 가산인 것이라 하자. 그러면 임의의 준콤팩트
열린집합 $U \subset X$와 $D(\mathcal{O}_X)$의 임의의 완전 대상 $E$에 대해
집합들
$$
H^i(U, K \otimes^\mathbf{L} E),\quad \Ext^i(E|_U, K|_U)
$$
은 가산이다.
\end{lemma}

\begin{proof}
Cohomology, Lemma \ref{cohomology-lemma-dual-perfect-complex}를 사용하면
군 $H^i(U, K \otimes^\mathbf{L} E)$에 대해 결과를 증명하면 충분함을 알 수
있다. 귀납원리를 사용하여 보조정리를 증명하겠다. Cohomology of Schemes,
Lemma \ref{coherent-lemma-induction-principle}를 보라.

\medskip\noindent
먼저 $U = \Spec(A)$가 아핀일 때 성립함을 보이자. 즉 $K$를 $A$-가군의
복합체 $K^\bullet$로, $E$를 유한 사영 $A$-가군의 유한 복합체
$P^\bullet$로 나타낼 수 있다. Lemmas \ref{perfect-lemma-affine-compare-bounded}와
\ref{perfect-lemma-perfect-affine}, 그리고 $A$-가군의 완전 복합체에 대한 정의
(More on Algebra, Definition \ref{more-algebra-definition-perfect})를 보라.
그러면 $(E \otimes^\mathbf{L} K)|_U$는 이중복합체
$P^\bullet \otimes_A K^\bullet$에 대응하는 전체 복합체로 나타난다
(Lemma \ref{perfect-lemma-quasi-coherence-tensor-product}). 복합체 $P^\bullet$의
길이에 대한 귀납법을 사용하면(또는 적절한 스펙트럼 열을 사용하면), 각
$a$에 대해 $H^i(P^a \otimes_A K^\bullet)$가 가산임을 보이면 충분함을 알
수 있다. 어떤 $n$에 대해 $P^a$는 $A^{\oplus n}$의 직합인자이므로, 이는
코호몰로지 군 $H^i(K^\bullet)$가 가산이라는 가정에서 따른다.

\medskip\noindent
증명을 끝내기 위해 다음을 보이면 충분하다. $U = V \cup W$이고 $V$,
$W$, $V \cap W$에 대해 결과가 성립하면 $U$에 대해서도 결과가 성립한다.
이는 Mayer--Vietoris 열의 즉각적인 귀결이다. Cohomology, Lemma
\ref{cohomology-lemma-unbounded-mayer-vietoris}를 보라.
\end{proof}

\begin{lemma}
\label{perfect-lemma-countable}
$X$를 준콤팩트이고 준분리인 스킴으로서 $\mathcal{O}_X$의 아핀 열린집합
위 절단들의 집합이 가산인 것이라 하자. $K$를
$D_\QCoh(\mathcal{O}_X)$의 대상이라 하자. 다음은 동치이다.
\begin{enumerate}
\item $K = \text{hocolim} E_n$이고 $E_n$은 $D(\mathcal{O}_X)$의 완전
대상이다.
\item 코호몰로지 층 $H^i(K)$의 아핀 열린집합 위 절단들의 집합은 가산이다.
\end{enumerate}
\end{lemma}

\begin{proof}
(1)이 참이면 (2)도 참이다. 실제로 호모토피 귀납극한은 코호몰로지 층을
취하는 것과 가환하고(Derived Categories, Lemma
\ref{derived-lemma-cohomology-of-hocolim}), 완전 복합체는 국소적으로 유한
자유 $\mathcal{O}_X$-가군의 유한 복합체와 동형이므로 $X$에 대한 가정에
의해 (2)를 만족한다.

\medskip\noindent
(2)를 가정하자. $K$를 나타내는 K-단사 복합체 $\mathcal{K}^\bullet$을
택한다. $D_\QCoh(\mathcal{O}_X)$의 완전 생성자 $E$를 택하고 이를 K-단사
복합체 $\mathcal{I}^\bullet$로 나타내자. Theorem
\ref{perfect-theorem-DQCoh-is-Ddga}와 그 증명에 따르면 삼각 범주의 동치
$F : D_\QCoh(\mathcal{O}_X) \to D(A, \text{d})$가 존재하며, 여기서
$(A, \text{d})$는 미분 등급 대수
$$
(A, \text{d}) =
\Hom_{\text{Comp}^{dg}(\mathcal{O}_X)}
(\mathcal{I}^\bullet, \mathcal{I}^\bullet)
$$
이고 이 함자는 $K$를 미분 등급 가군
$$
M = \Hom_{\text{Comp}^{dg}(\mathcal{O}_X)}
(\mathcal{I}^\bullet, \mathcal{K}^\bullet)
$$
로 보낸다. $H^i(A) = \Ext^i(E, E)$이고 $H^i(M) = \Ext^i(E, K)$임을
주목하자. 더욱이 $F$가 동치이므로 이 함자와 그 준역함자는 호모토피
귀납극한과 가환한다. 따라서 $M$을 $D(A, \text{d})$의 콤팩트 대상들의
호모토피 귀납극한으로 쓰면 충분하다. Differential Graded Algebra, Lemma
\ref{dga-lemma-countable}에 의해 각 $i$에 대해 $\Ext^i(E, E)$와
$\Ext^i(E, K)$가 가산임을 보이면 충분하다. 이는 Lemma
\ref{perfect-lemma-countable-cohomology}에서 따른다.
\end{proof}

\begin{lemma}
\label{perfect-lemma-computing-sections-as-colim}
$A$를 환이라 하자. $X$를 $A$ 위의 유한 표시 스킴이라 하자.
$f : U \to X$를 유한 표시 평탄 사상이라 하자. 그러면 다음이 성립한다.
\begin{enumerate}
\item $D(\mathcal{O}_X)$의 완전 대상 $L_n$들의 역계가 존재하여
$$
R\Gamma(U, Lf^*K) = \text{hocolim}\ R\Hom_X(L_n, K)
$$
가 $D_\QCoh(\mathcal{O}_X)$의 $K$에 대해 함자적으로 $D(A)$에서
성립한다.
\item $D(\mathcal{O}_X)$의 완전 대상 $E_n$들의 계가 존재하여
$$
R\Gamma(U, Lf^*K) = \text{hocolim}\ R\Gamma(X, E_n \otimes^\mathbf{L} K)
$$
가 $D_\QCoh(\mathcal{O}_X)$의 $K$에 대해 함자적으로 $D(A)$에서
성립한다.
\end{enumerate}
\end{lemma}

\begin{proof}
Lemma \ref{perfect-lemma-cohomology-base-change}에 의해
$$
R\Gamma(U, Lf^*K) = R\Gamma(X, Rf_*\mathcal{O}_U \otimes^\mathbf{L} K)
$$
가 $K$에 대해 함자적으로 성립한다. Lemma
\ref{perfect-lemma-quasi-coherence-pushforward-direct-sums}에 의해
$R\Gamma(X, -)$는 직합과 가환하므로 호모토피 귀납극한과도 가환한다.
마찬가지로 $- \otimes^\mathbf{L} K$는 유도 귀납극한과 가환한다. 실제로
$- \otimes^\mathbf{L} K$는 직합과 가환한다
($D(\mathcal{O}_X)$의 직합은 나타내는 복합체들의 직합으로
주어지기 때문이다). 따라서 (2)를 증명하려면 $D(\mathcal{O}_X)$의 완전
대상 $E_n$들의 어떤 계에 대해
$Rf_*\mathcal{O}_U = \text{hocolim} E_n$이라고 쓰면 충분하다. 이것을 한
뒤에는 $L_n = E_n^\vee$로 놓아 (1)을 얻는다. Cohomology, Lemma
\ref{cohomology-lemma-dual-perfect-complex}를 보라.

\medskip\noindent
$A = \colim A_i$로 쓰되 $A_i$가 $\mathbf{Z}$ 위에서 유한형이 되게 하자.
Limits, Lemma \ref{limits-lemma-descend-finite-presentation}에 의해 어떤
$i$와 유한 표시 사상 $U_i \to X_i \to \Spec(A_i)$를 찾을 수 있으며,
이를 $\Spec(A)$로 기저변환하면 $U \to X \to \Spec(A)$를 복원한다.
$i$를 크게 잡아 $f_i : U_i \to X_i$가 평탄이라고 가정할 수 있다.
Limits, Lemma \ref{limits-lemma-descend-flat-finite-presentation}를 보라.
Lemma \ref{perfect-lemma-compare-base-change}에 의해 $g : X \to X_i$에 따른
$Rf_{i, *}\mathcal{O}_{U_i}$의 유도 역상은 $Rf_*\mathcal{O}_U$와 같다.
$Lg^*$는 유도 귀납극한과 가환하므로 $f_i$에 대해 원하는 것을 증명하면
충분하다. 따라서 $U$와 $X$가 $\mathbf{Z}$ 위에서 유한형이라고 가정해도
된다.

\medskip\noindent
$f : U \to X$가 $\mathbf{Z}$ 위의 유한형 스킴들의 사상이라고 가정하자.
증명을 끝내기 위해 $Rf_*\mathcal{O}_U$가 완전 복합체들의 호모토피
귀납극한임을 보이겠다. 이를 위해 Lemma \ref{perfect-lemma-countable}을 적용한다.
따라서 $R^if_*\mathcal{O}_U$의 아핀 열린집합 위 절단들의 집합이 가산임을
보이면 충분하다. 이는 구조층에 Lemma
\ref{perfect-lemma-countable-cohomology}를 적용하면 따른다.
\end{proof}




\section{유사연접 복합체의 특징짓기, II}
\label{perfect-section-pseudo-coherent}

\noindent
이 절은 Section \ref{perfect-section-pseudo-coherent-hocolim}의 연속이다. 여기서는
코호몰로지를 이용한 유사연접 복합체의 특징짓기를 논한다. 이러한 종류의
결과는 More on Morphisms, Section
\ref{more-morphisms-section-characterize-pseudo-coherent}에서 더 찾을 수 있다.

\begin{lemma}
\label{perfect-lemma-pseudo-coherent-over-algebra}
$A$를 환이라 하자. $R$을 (비가환일 수도 있는) $A$-대수로서 $A$-가군으로
유한 자유인 것이라 하자. 그러면 $D(A)$에서 유사연접인 $D(R)$의 임의의
대상 $M$은 유한 자유 (오른쪽) $R$-가군의 위로 유계인 복합체로 나타낼 수
있다.
\end{lemma}

\begin{proof}
$M$을 나타내는 오른쪽 $R$-가군의 복합체 $M^\bullet$을 택한다. $M$이
유사연접이므로 충분히 큰 $i$에 대해 $H^i(M) = 0$이다. $H^m(M)$이 0이
아닌 가장 작은 지표를 $m$이라 하자. More on Algebra, Lemma
\ref{more-algebra-lemma-finite-cohomology}에 의해 $H^m(M)$은 유한
$A$-가군이다. 따라서 유한 자유 $R$-가군 $F^m$과 사상 $F^m \to M^m$을
택하여 $F^m \to M^m \to M^{m + 1}$이 0이고 $F^m \to H^m(M)$이 전사가
되게 할 수 있다. 그림은 다음과 같다.
$$
\xymatrix{
& F^m \ar[d]^\alpha \ar[r] & 0 \ar[d] \ar[r] & \ldots \\
M^{m - 1} \ar[r] & M^m \ar[r] & M^{m + 1} \ar[r] & \ldots
}
$$
$n \leq m$에 대한 하강 귀납법으로 $i \geq n$에 대한 유한 자유
$R$-가군 $F^i$, $i \geq n$에 대한 미분 $d^i : F^i \to F^{i + 1}$, 그리고
미분과 양립하는 사상 $\alpha : F^i \to K^i$를 다음 조건을 만족하도록
구성하겠다. (1) $H^i(\alpha)$는 $i > n$에 대해 동형사상이고 $i = n$에
대해 전사이다. (2) $i > m$에 대해 $F^i = 0$이다. 그림은 다음과 같다.
$$
\xymatrix{
& F^n \ar[r] \ar[d]^\alpha & F^{n + 1} \ar[d]^\alpha \ar[r] & \ldots \ar[r]
& F^i \ar[d]^\alpha \ar[r] & 0 \ar[d] \ar[r] & \ldots \\
M^{n - 1} \ar[r] & M^n \ar[r] & M^{n + 1} \ar[r] & \ldots \ar[r] &
M^i \ar[r] & M^{i + 1} \ar[r] & \ldots
}
$$
기초 단계는 위에서 다룬 $n = m$이다. 귀납 단계를 보자.
$C^\bullet$을 $\alpha$의 원뿔이라 하자(Derived Categories, Definition
\ref{derived-definition-cone}). 코호몰로지의 긴 완전열에 의해 $i \geq n$에
대해 $H^i(C^\bullet) = 0$이다. $R$은 $A$-가군으로 유한 자유이므로
$F^\bullet$은 $A$-가군의 복합체로서 유사연접임을 주목하자. 따라서 More on
Algebra, Lemma \ref{more-algebra-lemma-cone-pseudo-coherent}에 의해
$C^\bullet$은 $A$-가군의 복합체로서 $(n - 1)$-유사연접이다. More on
Algebra, Lemma \ref{more-algebra-lemma-finite-cohomology}에 의해
$H^{n - 1}(C^\bullet)$은 유한 $A$-가군이다. 유한 자유 $R$-가군
$F^{n - 1}$과 사상 $\beta : F^{n - 1} \to C^{n - 1}$을 택하여 합성
$F^{n - 1} \to C^{n - 1} \to C^n$이 0이고 $F^{n - 1}$에서
$H^{n - 1}(C^\bullet)$로의 사상이 전사가 되게 하자.
$C^{n - 1} = M^{n - 1} \oplus F^n$이므로
$\beta = (\alpha^{n - 1}, -d^{n - 1})$로 쓸 수 있다. 합성
$F^{n - 1} \to C^{n - 1} \to C^n$의 소멸은 이 사상들이 다음 복합체의
사상을 이룸을 뜻한다.
$$
\xymatrix{
& F^{n - 1} \ar[d]^{\alpha^{n - 1}} \ar[r]_{d^{n - 1}} &
F^n \ar[r] \ar[d]^\alpha &
F^{n + 1} \ar[d]^\alpha \ar[r] & \ldots \\
\ldots \ar[r] &
M^{n - 1} \ar[r] & M^n \ar[r] & M^{n + 1} \ar[r] & \ldots
}
$$
더욱이 이 사상들은 다음 구별 삼각형의 사상을 정의한다.
$$
\xymatrix{
(F^n \to \ldots) \ar[r] \ar[d] &
(F^{n - 1} \to \ldots) \ar[r] \ar[d] &
F^{n - 1} \ar[r] \ar[d]_\beta &
(F^n \to \ldots)[1] \ar[d] \\
(F^n \to \ldots) \ar[r] &
M^\bullet \ar[r] &
C^\bullet \ar[r] &
(F^n \to \ldots)[1]
}
$$
따라서 $\beta$의 선택에 의해 복합체의 사상
$(F^{n - 1} \to \ldots) \to M^\bullet$은 차수 $\geq n$의 코호몰로지에서
동형사상을 유도하고 차수 $n - 1$에서 전사를 유도한다. 이로써 보조정리의
증명이 끝난다.
\end{proof}

\begin{lemma}
\label{perfect-lemma-pseudo-coherent-on-projective-space}
$A$를 환이라 하고 $n \geq 0$이라 하자.
$K \in D_\QCoh(\mathcal{O}_{\mathbf{P}^n_A})$라 하자. 다음은 동치이다.
\begin{enumerate}
\item $K$는 유사연접이다.
\item $D(\mathcal{O}_{\mathbf{P}^n_A})$의 각 유사연접 대상 $E$에 대해
$R\Gamma(\mathbf{P}^n_A, E \otimes^\mathbf{L} K)$는 $D(A)$의 유사연접
대상이다.
\item $D(\mathcal{O}_{\mathbf{P}^n_A})$의 각 완전 대상 $E$에 대해
$R\Gamma(\mathbf{P}^n_A, E \otimes^\mathbf{L} K)$는 $D(A)$의 유사연접
대상이다.
\item $D(\mathcal{O}_{\mathbf{P}^n_A})$의 각 완전 대상 $E$에 대해
$R\Hom_{\mathbf{P}^n_A}(E, K)$는 $D(A)$의 유사연접 대상이다.
\item $R\Gamma(\mathbf{P}^n_A,
K \otimes^\mathbf{L} \mathcal{O}_{\mathbf{P}^n_A}(d))$는
$d = 0, 1, \ldots, n$에 대해 $D(A)$의 유사연접 대상이다.
\end{enumerate}
\end{lemma}

\begin{proof}
정의에 의해
$$
R\Hom_{\mathbf{P}^n_A}(E, K) =
R\Gamma(\mathbf{P}^n_A, R\SheafHom_{\mathcal{O}_{\mathbf{P}^n_A}}(E, K))
$$
임을 상기하자. Cohomology, Section \ref{cohomology-section-global-RHom}을
보라. 따라서 Cohomology, Lemma
\ref{cohomology-lemma-dual-perfect-complex}에 의해 (4)와 (3)은 동치이다.

\medskip\noindent
모든 완전 복합체는 유사연접이므로 (2)가 (3)을 함의함은 명백하다.

\medskip\noindent
(1)이 성립한다고 가정하자. 그러면 모든 유사연접 $E$에 대해
$E \otimes^\mathbf{L} K$는 유사연접이다. Cohomology, Lemma
\ref{cohomology-lemma-tensor-pseudo-coherent}을 보라. Lemma
\ref{perfect-lemma-flat-proper-pseudo-coherent-direct-image-general}에 의해 그러한
유사연접 복합체의 직접상은 유사연접이므로 (2)가 참이다.

\medskip\noindent
$d = 0, 1, \ldots, n$에 대해 $E = \mathcal{O}_{\mathbf{P}^n_A}(d)$로
둘 수 있으므로 (3)은 (5)를 함의한다.

\medskip\noindent
증명을 끝내려면 (5)가 (1)을 함의함을 보여야 한다. $P$를
(\ref{perfect-equation-generator-Pn})에서와 같이, $R$을
(\ref{perfect-equation-algebra-for-Pn})에서와 같이 택하자. Lemma
\ref{perfect-lemma-Pn-module-category}에 의해 다음 동치가 있다.
$$
- \otimes^\mathbf{L}_R P :
D(R) \longrightarrow D_\QCoh(\mathcal{O}_{\mathbf{P}^n_A})
$$
$M \otimes^\mathbf{L} P = K$가 되는 대상 $M \in D(R)$을 택하자.
Differential Graded Algebra, Lemma
\ref{dga-lemma-upgrade-tensor-with-complex-derived}에 의해 동형사상
$$
R\Hom(R, M) = R\Hom_{\mathbf{P}^n_A}(P, K)
$$
이 $D(A)$에서 성립한다. 위와 같이 논하면 다음을 얻는다.
$$
R\Hom_{\mathbf{P}^n_A}(P, K) = R\Gamma(\mathbf{P}^n_A,
R\SheafHom_{\mathcal{O}_{\mathbf{P}^n_A}}(E, K)) =
R\Gamma(\mathbf{P}^n_A,
P^\vee \otimes^\mathbf{L}_{\mathcal{O}_{\mathbf{P}^n_A}} K).
$$
$P^\vee$가 $d = 0, 1, \ldots, n$에 대한
$\mathcal{O}_{\mathbf{P}^n_A}(d)$들의 직합이라는 사실과 (5)를 사용하면,
$R\Hom(R, M)$이 $A$-가군의 복합체로서 유사연접이라고 결론짓는다. 물론
$D(A)$에서 $M = R\Hom(R, M)$이다. 따라서 $M$은 $A$-가군의 복합체로서
유사연접이다. Lemma \ref{perfect-lemma-pseudo-coherent-over-algebra}에 의해 $M$을
유한 자유 $R$-가군의 위로 유계인 복합체 $F^\bullet$로 나타낼 수 있다.
그러면
$F^\bullet = \bigcup_{p \geq 0} \sigma_{\geq p}F^\bullet$
은 $F^\bullet$이 성질 (P)를 갖는 미분 등급 $R$-가군임을 보여 주는
여과이다. Differential Graded Algebra, Section
\ref{dga-section-P-resolutions}를 보라. 따라서
$K = M \otimes^\mathbf{L}_R P$는 $F^\bullet \otimes_R P$로 나타난다
(이는 유도 텐서 함자의 구성에서 따르며, 예를 들어 Differential Graded
Algebra, Lemma \ref{dga-lemma-tensor-with-complex-derived}의 증명을 보라).
$F^\bullet \otimes_R P$는 각 항이 $P$의 복사본들의 직합인 위로 유계인
복합체이므로 보조정리가 참이라고 결론짓는다.
\end{proof}

\begin{lemma}
\label{perfect-lemma-perfect-enough}
$A$를 환이라 하자. $X$를 준콤팩트이고 준분리인 $A$ 위의 스킴이라 하자.
$K \in D^-_\QCoh(\mathcal{O}_X)$라 하자. $D(\mathcal{O}_X)$의 모든 완전
$E$에 대해 $R\Gamma(X, E \otimes^\mathbf{L} K)$가 $D(A)$에서
유사연접이면, $D(\mathcal{O}_X)$의 모든 유사연접 $E$에 대해서도
$R\Gamma(X, E \otimes^\mathbf{L} K)$는 $D(A)$에서 유사연접이다.
\end{lemma}

\begin{proof}
Lemma \ref{perfect-lemma-quasi-coherence-direct-image}에서 설명한 바와 같이
$R\Gamma(X, -) : D_\QCoh(\mathcal{O}_X) \to D(A)$의 코호몰로지 차원이
$N$이 되는 정수 $N$이 존재한다. $i > b$에 대해 $H^i(K) = 0$이 되도록
$b \in \mathbf{Z}$를 택하자. $E$가 $X$ 위에서 유사연접이라고 하자.
모든 $m$에 대해 $R\Gamma(X, E \otimes^\mathbf{L} K)$가
$m$-유사연접임을 보이면 충분하다. $(X, E, m - N - 1 - b)$에 대해 완전
복합체 $P$에 의한 근사 $P \to E$를 택한다. 이는 Theorem
\ref{perfect-theorem-approximation}에 의해 가능하다. 다음 구별 삼각형을
$$
P \to E \to C \to P[1]
$$
$D_\QCoh(\mathcal{O}_X)$에서 택하자. $C$의 코호몰로지 층들은 차수
$\geq m - N - 1 - b$에서 0이다. 따라서 $C \otimes^\mathbf{L} K$의
코호몰로지 층들은 차수 $\geq m - N - 1$에서 0이다. 그러므로
$R\Gamma(X, C \otimes^\mathbf{L} K)$의 코호몰로지는 차수
$\geq m - 1$에서 0이다. 따라서
$$
R\Gamma(X, P \otimes^\mathbf{L} K) \to R\Gamma(X, E \otimes^\mathbf{L} K)
$$
는 차수 $\geq m$의 코호몰로지에서 동형사상이다. 가정에 의해 그 원천은
유사연접이다. 따라서 원하는 대로 $R\Gamma(X, E \otimes^\mathbf{L} K)$가
$m$-유사연접이라고 결론짓는다.
\end{proof}









\section{상대적 완전 대상}
\label{perfect-section-relatively-perfect}

\noindent
이 절에서는 \cite{lieblich-complexes}의 개념을 도입한다.

\begin{definition}
\label{perfect-definition-relatively-perfect}
$f : X \to S$를 평탄하고 국소적으로 유한 표시인 스킴의 사상이라 하자.
$D(\mathcal{O}_X)$의 대상 $E$가 유사연접이고
(Cohomology, Definition \ref{cohomology-definition-pseudo-coherent}),
또한 $E$가 국소적으로 $D(f^{-1}\mathcal{O}_S)$의 대상으로서
유한 Tor 차원을 가지면
(Cohomology, Definition \ref{cohomology-definition-tor-amplitude}),
$E$를 {\it $S$에 관해 완전} 또는 {\it $S$-완전}이라 한다.
\end{definition}

\noindent
이에 관한 논의는 Remark \ref{perfect-remark-discuss-rel-perfect}를 보라.

\begin{example}
\label{perfect-example-relatively-perfect-field}
$k$를 체라 하고, $X$를 $k$ 위의 유한 표시 스킴이라 하자
(특히 $X$는 준콤팩트이다). 그러면 $D(\mathcal{O}_X)$의 대상 $E$가
$k$-완전일 필요충분조건은 (정의에 의해) 유계이고 유사연접인 것이며,
이는 (Lemma \ref{perfect-lemma-identify-pseudo-coherent-noetherian}에 의해)
$D^b_{\textit{Coh}}(X)$에 속하는 것과 동치이다.
따라서 상대적으로 완전하다는 것은 ``각 올 위에서 완전하다''는 뜻이
{\bf 아니다}.
\end{example}

\noindent
이에 대응하는 대수적 개념은 More on Algebra, Section
\ref{more-algebra-section-relatively-perfect}에서 다룬다.
스킴에 대한 개념과 대수적 개념은 다음과 같이 연결된다.

\begin{lemma}
\label{perfect-lemma-affine-locally-rel-perfect}
$f : X \to S$를 평탄하고 국소적으로 유한 표시인 스킴의 사상이라 하고,
$E$를 $D_\QCoh(\mathcal{O}_X)$의 대상이라 하자. 다음은 서로 동치이다.
\begin{enumerate}
\item $E$는 $S$-완전이다.
\item 아핀 열린집합 $V \subset S$ 안으로 사상되는 임의의 아핀 열린집합
$U \subset X$에 대해 복합체 $R\Gamma(U, E)$는 $\mathcal{O}_S(V)$-완전이다.
\item 아핀 열린 덮개 $S = \bigcup V_i$가 존재하고, 각 $i$에 대해
아핀 열린 덮개 $f^{-1}(V_i) = \bigcup U_{ij}$가 존재하여 복합체
$R\Gamma(U_{ij}, E)$가 $\mathcal{O}_S(V_i)$-완전이다.
\end{enumerate}
\end{lemma}

\begin{proof}
유사연접성은 국소적 성질이고, ``국소적으로 유한 Tor 차원을 갖는 것''도
국소적 성질이다. 따라서 이 보조정리는 곧 다음 명제로 환원된다.
$X$와 $S$가 아핀이면, $E$가 $S$-완전일 필요충분조건은
$K = R\Gamma(X, E)$가 $\mathcal{O}_S(S)$-완전인 것이다.
$X = \Spec(A)$, $S = \Spec(R)$라 하고, $E$가 $K \in D(A)$에 대응한다고
하자. 즉, $K = R\Gamma(X, E)$이다. Lemma
\ref{perfect-lemma-affine-compare-bounded}를 보라.

\medskip\noindent
$K$가 $R$-완전일 필요충분조건은 $K$가 유사연접이고 $R$-가군의
복합체로서 유한 Tor 차원을 갖는 것임을 관찰하자(More on Algebra,
Definition \ref{more-algebra-definition-relatively-perfect}).
Lemma \ref{perfect-lemma-pseudo-coherent-affine}에 의해 $E$가 유사연접일
필요충분조건은 $K$가 유사연접인 것이다.
Lemma \ref{perfect-lemma-tor-dimension-rel-affine}에 의해 $E$가
$f^{-1}\mathcal{O}_S$ 위에서 유한 Tor 차원을 가질 필요충분조건은
$K$가 $R$-가군의 복합체로서 유한 Tor 차원을 갖는 것이다.
\end{proof}

\begin{lemma}
\label{perfect-lemma-triangulated}
$f : X \to S$를 평탄하고 국소적으로 유한 표시인 스킴의 사상이라 하자.
$D(\mathcal{O}_X)$에서 $S$-완전 대상들로 이루어진 충만 부분범주는
포화된\footnote{Derived Categories, Definition
\ref{derived-definition-saturated}.} 삼각 부분범주이다.
\end{lemma}

\begin{proof}
이는 Cohomology, Lemmas
\ref{cohomology-lemma-cone-pseudo-coherent},
\ref{cohomology-lemma-summands-pseudo-coherent},
\ref{cohomology-lemma-cone-tor-amplitude}, 및
\ref{cohomology-lemma-summands-tor-amplitude}에서 따른다.
\end{proof}

\begin{lemma}
\label{perfect-lemma-perfect-relatively-perfect}
$f : X \to S$를 평탄하고 국소적으로 유한 표시인 스킴의 사상이라 하자.
$D(\mathcal{O}_X)$의 완전 대상은 $S$-완전이다.
$K, M \in D(\mathcal{O}_X)$이고 $K$가 완전이며 $M$이 $S$-완전이면,
$K \otimes_{\mathcal{O}_X}^\mathbf{L} M$은 $S$-완전이다.
\end{lemma}

\begin{proof}
첫째 증명: Lemma \ref{perfect-lemma-affine-locally-rel-perfect}를 사용하여
아핀인 경우에 환원한 뒤 More on Algebra, Lemma
\ref{more-algebra-lemma-perfect-relatively-perfect}를 적용한다.
\end{proof}

\begin{lemma}
\label{perfect-lemma-base-change-relatively-perfect}
$f : X \to S$를 평탄하고 국소적으로 유한 표시인 스킴의 사상이라 하자.
$g : S' \to S$를 스킴의 사상이라 하자. $X' = S' \times_S X$로 놓고,
사영을 $g' : X' \to X$로 나타내자.
$K \in D(\mathcal{O}_X)$가 $S$-완전이면 $L(g')^*K$는 $S'$-완전이다.
\end{lemma}

\begin{proof}
첫째 증명: Lemma \ref{perfect-lemma-affine-locally-rel-perfect}를 사용하여
아핀인 경우에 환원한 뒤 More on Algebra, Lemma
\ref{more-algebra-lemma-base-change-relatively-perfect}를 적용한다.

\medskip\noindent
둘째 증명: Cohomology, Lemma
\ref{cohomology-lemma-pseudo-coherent-pullback}에 의해 $L(g')^*K$는
유사연접이고, 유계 Tor 차원 성질은 Lemma
\ref{perfect-lemma-tor-independence-and-tor-amplitude}에서 따른다.
\end{proof}

\begin{situation}
\label{perfect-situation-relative-descent}
$S = \lim_{i \in I} S_i$를 아핀 전이 사상
$g_{i'i} : S_{i'} \to S_i$를 갖는 스킴의 유향계의 극한이라 하자.
모든 $i \in I$에 대해 $S_i$가 준콤팩트이고 준분리라고 가정한다.
사영을 $g_i : S \to S_i$로 나타내자. 원소 $0 \in I$와 유한 표시인
평탄 사상 $X_0 \to S_0$를 고정한다.
$X_i = S_i \times_{S_0} X_0$ 및 $X = S \times_{S_0} X_0$로 놓고,
전이 사상을 $f_{i'i} : X_{i'} \to X_i$, 사영을
$f_i : X \to X_i$로 나타내자.
\end{situation}

\begin{lemma}
\label{perfect-lemma-relative-descend-homomorphisms}
Situation \ref{perfect-situation-relative-descent}에서 생각하자.
$K_0$와 $L_0$를 $D(\mathcal{O}_{X_0})$의 대상이라 하자.
$i \geq 0$에 대해 $K_i = Lf_{i0}^*K_0$ 및 $L_i = Lf_{i0}^*L_0$로
놓고, $K = Lf_0^*K_0$ 및 $L = Lf_0^*L_0$로 놓자. 그러면 사상
$$
\colim_{i \geq 0} \Hom_{D(\mathcal{O}_{X_i})}(K_i, L_i)
\longrightarrow
\Hom_{D(\mathcal{O}_X)}(K, L)
$$
$K_0$가 유사연접이고 $L_0 \in D_\QCoh(\mathcal{O}_{X_0})$가
$D((X_0 \to S_0)^{-1}\mathcal{O}_{S_0})$의 대상으로서 (국소적으로)
유한 Tor 차원을 가지면 이 사상은 동형이다.
\end{lemma}

\begin{proof}
각 준콤팩트 열린집합 $U_0 \subset X_0$에 대해 다음이 성립한다는
조건 $P$를 생각하자.
$$
\colim_{i \geq 0} \Hom_{D(\mathcal{O}_{U_i})}(K_i|_{U_i}, L_i|_{U_i})
\longrightarrow
\Hom_{D(\mathcal{O}_U)}(K|_U, L|_U)
$$
여기서 $U = f_0^{-1}(U_0)$이고 $U_i = f_{i0}^{-1}(U_0)$이다.
$U_0$, $V_0$, $U_0 \cap V_0$에 대해 $P$가 성립하면, 유도 범주에서의
hom에 대한 Mayer--Vietoris에 의해 $U_0 \cup V_0$에 대해서도 성립한다.
Cohomology, Lemma \ref{cohomology-lemma-mayer-vietoris-hom}을 보라.

\medskip\noindent
주어진 사상을 $\pi_0 : X_0 \to S_0$로 나타내자.
먼저 $W_0 \subset S_0$가 준콤팩트 열린집합인
$U_0 = \pi_0^{-1}(W_0)$를 생각할 수 있다. $S_0$의 준콤팩트
열린집합들에 Cohomology of Schemes, Lemma
\ref{coherent-lemma-induction-principle}의 귀납 원리와 위의 관찰을
적용하면, $W_0$가 아핀인 $U_0 = \pi_0^{-1}(W_0)$에 대해 $P$를
증명하면 충분함을 알 수 있다. 다시 말해 $S_0$가 아핀인 경우로 환원했다.
다음으로 이번에는 $X_0$의 모든 준콤팩트 준분리 열린집합에 귀납 원리를
다시 적용하여 $X_0$도 아핀인 경우로 환원한다.

\medskip\noindent
$X_0$와 $S_0$가 아핀이면 결과는 More on Algebra, Lemma
\ref{more-algebra-lemma-colimit-relatively-perfect}에서 따른다.
실제로 Lemmas \ref{perfect-lemma-pseudo-coherent}와
\ref{perfect-lemma-affine-compare-bounded}에 의해 명제는 가군의 유도 범주에서
hom을 계산하는 문제로 옮겨진다. 이어서 Lemma
\ref{perfect-lemma-pseudo-coherent-affine}는 $K_0$에 대응하는 가군 복합체가
유사연접임을 보인다. 또한 Lemma \ref{perfect-lemma-tor-dimension-rel-affine}는
$L_0$에 대응하는 가군 복합체가 $\mathcal{O}_{S_0}(S_0)$ 위에서
유한 Tor 차원을 가짐을 보인다. 따라서 More on Algebra, Lemma
\ref{more-algebra-lemma-colimit-relatively-perfect}의 가정들이 충족되어
결론을 얻는다.
\end{proof}

\begin{lemma}
\label{perfect-lemma-descend-relatively-perfect}
Situation \ref{perfect-situation-relative-descent}에서 $D(\mathcal{O}_X)$의
$S$-완전 대상들의 범주는 $D(\mathcal{O}_{X_i})$의 $S_i$-완전
대상들의 범주들의 귀납극한이다.
\end{lemma}

\begin{proof}
각 준콤팩트 열린집합 $U_0 \subset X_0$에 대해 함자
$$
\colim_{i \geq 0} D_{S_i\text{-perfect}}(\mathcal{O}_{U_i})
\longrightarrow
D_{S\text{-perfect}}(\mathcal{O}_U)
$$
가 동치라는 조건 $P$를 생각하자. 여기서 $U = f_0^{-1}(U_0)$이고
$U_i = f_{i0}^{-1}(U_0)$이다. Lemma
\ref{perfect-lemma-relative-descend-homomorphisms}에 의해 이 함자가 이미
충실충만임을 알고 있다. 따라서 본질적 전사성을 증명하면 충분하다.

\medskip\noindent
$X_0$의 준콤팩트 열린집합 $U_0$, $V_0$에 대해 $P$가 성립한다고 하자.
$U_0 \cup V_0$에 대해서도 $P$가 성립함을 보이겠다. 다음 표기를 쓴다.
$U_i = f_{i0}^{-1}U_0$, $U = f_0^{-1}U_0$, $V_i = f_{i0}^{-1}V_0$,
$V = f_0^{-1}V_0$. 또한 사상 $U \to U_i$, $V \to V_i$,
$U \cap V \to U_i \cap V_i$, $U \cup V \to U_i \cup V_i$를 모두
다소 남용하여 $f_i$로 나타낸다.
$E$를 $D(\mathcal{O}_{U \cup V})$의 $S$-완전 대상이라 하자.
목표는 $E$가 이 함자의 본질적 상에 속함을 보이는 것이다.
가정에 의해 $i \geq 0$, $U_i$ 위의 $S_i$-완전 대상 $E_{U, i}$,
$V_i$ 위의 $S_i$-완전 대상 $E_{V, i}$ 및 동형
$Lf_i^*E_{U, i} \to E|_U$, $Lf_i^*E_{V, i} \to E|_V$를 찾을 수 있다.
이 동형들에 수반되는 사상을
$$
a : E_{U, i} \to (Rf_{i, *}E)|_{U_i}
\quad\text{및}\quad
b : E_{V, i} \to (Rf_{i, *}E)|_{V_i}
$$
로 놓자. 이들은 각각 동형 $Lf_i^*E_{U, i} \to E|_U$ 및
$Lf_i^*E_{V, i} \to E|_V$에 수반된다.
충실충만성에 의해 $i$를 증가시킨 뒤, 역상이 다음 식별이 되는
동형 $c : E_{U, i}|_{U_i \cap V_i} \to E_{V, i}|_{U_i \cap V_i}$를
찾을 수 있다.
$$
Lf_i^*E_{U, i}|_{U \cap V} \to E|_{U \cap V} \to Lf_i^*E_{V, i}|_{U \cap V}.
$$
Cohomology, Lemma \ref{cohomology-lemma-glue}를 적용하여
$U_i \cup V_i$ 위의 대상 $E_i$와, $U_i$ 및 $V_i$ 위에서 각각
$a$와 $b$로 제한되는 사상 $d : E_i \to Rf_{i, *}E$를 얻는다.
그러면 $E_i$가 $S_i$-완전이고(상대적 완전성은 국소적 성질이므로),
$d$가 동형 $Lf_i^*E_i \to E$에 수반됨은 분명하다.

\medskip\noindent
귀납 원리(Cohomology of Schemes, Lemma
\ref{coherent-lemma-induction-principle})를 사용한 Lemma
\ref{perfect-lemma-relative-descend-homomorphisms}의 증명과 정확히 같은 논증으로
$X_0$와 $S_0$가 모두 아핀인 경우로 환원한다. (먼저 $S_0$의
열린집합들을 다루어 $S_0$가 아핀인 경우로 환원하고, 이어서 $X_0$의
열린집합들을 다루어 $X_0$가 아핀인 경우로 환원한다.) 아핀인 경우 결과는
More on Algebra, Lemma \ref{more-algebra-lemma-colimit-relatively-perfect}에서
따른다. 대수로의 번역은 Lemma
\ref{perfect-lemma-affine-locally-rel-perfect}에 의해 이루어진다.
\end{proof}

\begin{lemma}
\label{perfect-lemma-derived-pushforward-rel-perfect}
$f : X \to S$를 평탄하고 고유이며 유한 표시인 스킴의 사상이라 하고,
$E \in D(\mathcal{O}_X)$를 $S$-완전이라 하자. 그러면 $Rf_*E$는
$D(\mathcal{O}_S)$의 완전 대상이고, 그 형성은 임의의 기저변환과 가환한다.
\end{lemma}

\begin{proof}
기저변환에 관한 명제는 Lemma \ref{perfect-lemma-compare-base-change}이다.
따라서 $Rf_*E$가 완전 대상임을 보이면 충분하다. 극한 논증으로
$S$가 뇌터 아핀인 경우에 환원하겠다.

\medskip\noindent
문제는 $S$ 위에서 국소적이므로 $S$가 아핀이라고 가정할 수 있다.
$S = \Spec(R)$라 하자. 이를 뇌터 환 $R_i$들의 여과 귀납극한
$R = \colim R_i$로 쓴다. Limits, Lemma
\ref{limits-lemma-descend-finite-presentation}에 의해 어떤 $i$와
$R_i$ 위의 유한 표시 스킴 $X_i$가 존재하여, 이를 $R$로 기저변환한 것이
$X$이다. Limits, Lemmas \ref{limits-lemma-eventually-proper}와
\ref{limits-lemma-descend-flat-finite-presentation}에 의해 $X_i$가
$R_i$ 위에서 고유하고 평탄하다고 가정할 수 있다. Lemma
\ref{perfect-lemma-descend-relatively-perfect}에 의해 $D(\mathcal{O}_{X_i})$의
$R_i$-완전 대상 $E_i$가 존재하여 그 $X$로의 역상이 $E$라고 가정할 수 있다.
$X_i \to \Spec(R_i)$와 $E_i$에 Lemma
\ref{perfect-lemma-perfect-direct-image}를 적용하고 이미 보인 기저변환 성질을
사용하면 결과를 얻는다.
\end{proof}

\begin{lemma}
\label{perfect-lemma-compute-ext-rel-perfect}
$f : X \to S$를 스킴의 사상이라 하고, $E, K \in D(\mathcal{O}_X)$라 하자.
다음을 가정하자.
\begin{enumerate}
\item $S$는 준콤팩트이고 준분리이다.
\item $f$는 고유하고 평탄하며 유한 표시이다.
\item $E$는 $S$-완전이다.
\item $K$는 유사연접이다.
\end{enumerate}
그러면 유사연접인 $L \in D(\mathcal{O}_S)$가 존재하여
$$
Rf_*R\SheafHom(K, E) = R\SheafHom(L, \mathcal{O}_S)
$$
이고, 임의의 기저변환 뒤에도 마찬가지이다. 즉, 다음 그림이 주어지면
$$
\vcenter{
\xymatrix{
X' \ar[r]_{g'} \ar[d]_{f'} &
X \ar[d]^f \\
S' \ar[r]^g &
S
}
}
\quad\quad
\begin{matrix}
\text{가 데카르트이면 다음이 성립한다. } \\
Rf'_*R\SheafHom(L(g')^*K, L(g')^*E) \\
= R\SheafHom(Lg^*L, \mathcal{O}_{S'})
\end{matrix}
$$
\end{lemma}

\begin{proof}
$S$가 준콤팩트이고 준분리이므로 $X$도 그러하다. Lemma
\ref{perfect-lemma-pseudo-coherent-hocolim}에 의해 $K_n$이 완전이고
$K_n \to K$가 절단 $\tau_{\geq -n}$ 위에서 동형을 유도하도록
$K = \text{hocolim} K_n$으로 쓸 수 있다. $K_n^\vee$를 쌍대 완전
복합체라 하자(Cohomology, Lemma
\ref{cohomology-lemma-dual-perfect-complex}). 그러면 완전 대상들의 역계
$\ldots \to K_3^\vee \to K_2^\vee \to K_1^\vee$를 얻는다.
Lemma \ref{perfect-lemma-perfect-relatively-perfect}에 의해
$K_n^\vee \otimes_{\mathcal{O}_X} E$는 $S$-완전이다. 따라서
$K_n^\vee \otimes_{\mathcal{O}_X} E$에 Lemma
\ref{perfect-lemma-derived-pushforward-rel-perfect}를 적용하여 역계
$$
\ldots \to M_3 \to M_2 \to M_1
$$
를 얻는다. 이는 $S$ 위의 완전 복합체들로 이루어지며
$$
M_n = Rf_*(K_n^\vee \otimes_{\mathcal{O}_X}^\mathbf{L} E) =
Rf_*R\SheafHom(K_n, E)
$$
이다. 또한 이 복합체들의 형성은 임의의 기저변환과 가환한다. 즉,
$Lg^*M_n =
Rf'_*((L(g')^*K_n)^\vee \otimes_{\mathcal{O}_{X'}}^\mathbf{L} L(g')^*E) =
Rf'_*R\SheafHom(L(g')^*K_n, L(g')^*E)$.

\medskip\noindent
$K_n \to K$가 $\tau_{\geq -n}$ 위에서 동형을 유도하므로,
$K_n \to K_{n + 1}$도 $\tau_{\geq -n}$ 위에서 동형을 유도한다.
$K_n^\vee = R\SheafHom(K_n, \mathcal{O}_X)$이므로
$K_{n + 1}^\vee \to K_n^\vee$는 $\tau_{\leq n}$ 위에서 동형을 유도한다.
$E$가 $f^{-1}\mathcal{O}_Y$-가군의 복합체로서 $[a, b]$ 안의
Tor-진폭을 갖는다고 하자. 그러면 임의의 기저변환 뒤에도 마찬가지이다.
Lemma \ref{perfect-lemma-tor-independence-and-tor-amplitude}를 보라. 따라서 사상
$K_{n + 1}^\vee \otimes_{\mathcal{O}_X} E \to
K_n^\vee \otimes_{\mathcal{O}_X} E$
은 $\tau_{\leq n + a}$ 위에서 동형을 유도하고, 임의의 기저변환 뒤에도
마찬가지이다. 오른쪽 유도 함자 $Rf_*$를 적용하면 사상
$M_{n + 1} \to M_n$이 $\tau_{\leq n + a}$ 위에서 동형을 유도하며,
임의의 기저변환 뒤에도 마찬가지임을 알 수 있다. 완전 삼각형
$$
M_{n + 1} \to M_n \to C_n \to M_{n + 1}[1]
$$
을 택하자. $S'$를 점 $s \in S$에서의 잔여체의 스펙트럼으로 놓고
역상하면 $C_n \otimes_{\mathcal{O}_S}^\mathbf{L} \kappa(s)$가
$\geq n + a$인 차수에서만 영이 아닌 코호몰로지를 가짐을 알 수 있다.
More on Algebra, Lemma
\ref{more-algebra-lemma-lift-perfect-from-residue-field}에 의해 어떤 정수
$m_n$에 대해 완전 복합체 $C_n$은 $[n + a, m_n]$ 안의 Tor-진폭을 갖는다.
특히 쌍대 완전 복합체 $C_n^\vee$는 $[-m_n, -n - a]$ 안의
Tor-진폭을 갖는다.

\medskip\noindent
$L_n = M_n^\vee$를 쌍대 완전 복합체라 하자. 앞 문단의 논의에서
$L_n \to L_{n + 1}$이 $\tau_{\geq -n - a}$ 위에서 동형을 유도한다는
결론을 얻는다. 따라서 $L = \text{hocolim} L_n$은 유사연접이다.
Lemma \ref{perfect-lemma-pseudo-coherent-hocolim}을 보라. 이제
$$
R\SheafHom(K, E) = R\SheafHom(\text{hocolim} K_n, E) =
R\lim R\SheafHom(K_n, E) = R\lim K_n^\vee \otimes_{\mathcal{O}_X} E
$$
이고(Cohomology, Lemma \ref{cohomology-lemma-colim-and-lim-of-duals}),
$R\lim$은 $Rf_*$와 가환하므로
$$
Rf_*R\SheafHom(K, E) = R\lim M_n = R\lim R\SheafHom(L_n, \mathcal{O}_S) =
R\SheafHom(L, \mathcal{O}_S)
$$
을 얻는다. 이로써 $S$ 위의 공식을 증명했다. $M_n$의 구성은
기저변환과 양립하므로 임의의 기저변환 뒤에도 이 공식이 계속 성립한다.
\end{proof}

\begin{remark}
\label{perfect-remark-compare-L}
독자는 Lemma \ref{perfect-lemma-compute-ext-rel-perfect}와 Lemma
\ref{perfect-lemma-compute-ext} 사이의 유사성을 알아차렸을 것이다.
실제로 Lemma \ref{perfect-lemma-compute-ext-rel-perfect}의 유사연접 복합체 $L$은
다음 함자적 동형들이 존재하는 $S$ 위의 유일한 유사연접 복합체로
특징지을 수 있다.
$$
\Ext^i_{\mathcal{O}_S}(L, \mathcal{F}) \longrightarrow
\Ext^i_{\mathcal{O}_X}(K,
E \otimes_{\mathcal{O}_X}^\mathbf{L} Lf^*\mathcal{F})
$$
여기서 $\mathcal{F}$는 $\QCoh(\mathcal{O}_S)$를 달리며, 이 동형들은
경계 사상과 양립한다. 이를 사용할 필요가 생기면 여기에 정확한 결과를
정식화하고 상세한 증명을 제시할 것이다.
\end{remark}

\begin{lemma}
\label{perfect-lemma-bounded-on-fibres}
$f : X \to S$를 평탄하고 국소적으로 유한 표시인 스킴의 사상이라 하고,
$E$를 $D(\mathcal{O}_X)$의 유사연접 대상이라 하자. 다음은 서로 동치이다.
\begin{enumerate}
\item $E$는 $S$-완전이다.
\item $E$는 국소적으로 아래로 유계이고, 모든 점 $s \in S$에 대해
$D(\mathcal{O}_{X_s})$의 대상 $L(X_s \to X)^*E$는 국소적으로
아래로 유계이다.
\end{enumerate}
\end{lemma}

\begin{proof}
모든 것이 국소적이므로 $X$와 $S$가 아핀인 경우로 즉시 환원한다.
Lemma \ref{perfect-lemma-affine-locally-rel-perfect}를 보라.
$X \to S$가 $\Spec(A) \to \Spec(R)$에 대응하고, $E$가 $D(A)$의
$K$에 대응한다고 하자. $s$가 소 아이디얼 $\mathfrak p \subset R$에
대응하면, $R \to A$가 평탄이므로 $L(X_s \to X)^*E$는
$K \otimes_R^\mathbf{L} \kappa(\mathfrak p)$에 대응한다. 예를 들어
Lemma \ref{perfect-lemma-compare-base-change}를 보라. 따라서 이 보조정리는
대응하는 대수적 결과에서 따른다. More on Algebra, Lemma
\ref{more-algebra-lemma-bounded-on-fibres-relatively-perfect}를 보라.
\end{proof}

\begin{remark}
\label{perfect-remark-discuss-rel-perfect}
상대적 완전 복합체에 관한 Definition
\ref{perfect-definition-relatively-perfect}은 우리의 정의를 적용할 수 있는 경우
\cite{lieblich-complexes}에 주어진 정의와 동치이다\footnote{이를 보려면
Lemma \ref{perfect-lemma-affine-locally-rel-perfect}와 More on Algebra, Lemma
\ref{more-algebra-lemma-structure-relatively-perfect}를 사용하라.}.
다음으로 $f : X \to S$가 단지 국소적으로 유한형이라고 가정하자
(평탄이거나 국소적으로 유한 표시일 필요는 없다). 위에서 인용한 논문의
정의에 따르면 $E \in D(\mathcal{O}_X)$가 상대적으로 완전이라는 것은
다음을 뜻한다.
\begin{enumerate}
\item[(A)] $X$ 위에서 국소적으로 대상 $E$는 $S$-평탄이고 유한 표시인
$\mathcal{O}_X$-가군들로 이루어진 유한 복합체와 준동형이다.
\end{enumerate}
한편 Definition \ref{perfect-definition-relatively-perfect}의 자연스러운
일반화는 다음과 같다.
\begin{enumerate}
\item[(B)] $E$는 $S$에 관해 유사연접이고(More on Morphisms, Definition
\ref{more-morphisms-definition-relative-pseudo-coherence}),
$E$는 국소적으로 $D(f^{-1}\mathcal{O}_S)$의 대상으로서 유한 Tor 차원을
갖는다(Cohomology, Definition
\ref{cohomology-definition-tor-amplitude}).
\end{enumerate}
조건 (B)의 장점은 그것이 분명히 $D(\mathcal{O}_X)$의 삼각 부분범주를
정의한다는 점인 반면, 조건 (A)에 대해서는 그렇지 않으리라 예상한다.
조건 (A)의 장점은 특히 극한과 관련하여 다루기 더 쉽다는 점이다.
\end{remark}









\section{분해 성질}
\label{perfect-section-resolution-property}

\noindent
이 개념은 논문 \cite{totaro_resolution}에서 논의되며, 그 논의는
\cite{Gross-thesis}, \cite{Gross-surface}, \cite{Gross-stack}에서 계속된다.
체 위의 고유 스킴이 항상 분해 성질을 갖는지, 아니면 이것이 거짓인지는
현재 알려져 있지 않다. 이 질문의 답을 알고 있다면
\href{mailto:stacks.project@gmail.com}{stacks.project@gmail.com}.
로 이메일을 보내 주기 바란다.

\medskip\noindent
일반적인 스킴에 대해 이를 생각하는 것은 거의 의미가 없지만,
다음 정의를 내릴 수 있다.

\begin{definition}
\label{perfect-definition-resolution-property}
$X$를 스킴이라 하자. 유한형인 모든 준연접 $\mathcal{O}_X$-가군이
유한 국소 자유 $\mathcal{O}_X$-가군의 몫이면, $X$가
{\it 분해 성질}을 갖는다고 한다.
\end{definition}

\noindent
$X$가 준콤팩트 준분리 스킴이면, 유한 표시인 모든 $\mathcal{O}_X$-가군
(자동으로 준연접이다)이 유한 국소 자유 $\mathcal{O}_X$-가군의 몫임을
확인하는 것으로 충분하다. Properties, Lemma
\ref{properties-lemma-finite-directed-colimit-surjective-maps}.
를 보라. $X$가 뇌터 스킴이면 유한형 준연접 가군은 정확히 연접
$\mathcal{O}_X$-가군이다. Cohomology of Schemes, Lemma
\ref{coherent-lemma-coherent-Noetherian}을 보라.

\begin{lemma}
\label{perfect-lemma-resolution-property-ample}
$X$를 스킴이라 하자. $X$가 풍부한 가역 $\mathcal{O}_X$-가군을 가지면,
$X$는 분해 성질을 갖는다.
\end{lemma}

\begin{proof}
Properties, Proposition \ref{properties-proposition-characterize-ample}의
즉각적인 결과이다.
\end{proof}

\begin{lemma}
\label{perfect-lemma-resolution-property-ample-relative}
$f : X \to Y$를 스킴의 사상이라 하자. 다음을 가정하자.
\begin{enumerate}
\item $Y$는 준콤팩트이고 준분리이며 분해 성질을 갖는다.
\item $X$ 위에 $f$-풍부 가역 가군이 존재한다.
\end{enumerate}
그러면 $X$는 분해 성질을 갖는다.
\end{lemma}

\begin{proof}
$\mathcal{F}$를 유한형 준연접 $\mathcal{O}_X$-가군이라 하고,
$\mathcal{L}$을 $f$-풍부 가역 가군이라 하자. 아핀 열린 덮개
$Y = V_1 \cup \ldots \cup V_m$을 택하고 $U_j = f^{-1}(V_j)$로 놓자.
Properties, Proposition \ref{properties-proposition-characterize-ample}에
의해 각 $j$에 대해 함께 전사인 유한 개의 사상
$s_{j, i} : \mathcal{L}^{\otimes n_{j, i}}|_{U_j} \to \mathcal{F}|_{U_j}$
가 존재한다. 다음 준연접 $\mathcal{O}_Y$-가군을 생각하자.
$$
\mathcal{H}_n =
f_*(\mathcal{F} \otimes_{\mathcal{O}_X} \mathcal{L}^{\otimes n})
$$
$s_{j, i}$를 층 $\mathcal{H}_{-n_{j, i}}$의 $V_j$ 위의 절단으로
생각할 수 있다. 유한 국소 자유 $\mathcal{O}_Y$-가군
$\mathcal{E}_{i, j}$와 사상
$\mathcal{E}_{i, j} \to \mathcal{H}_{-n_{j, i}}$를 찾아
$s_{j, i}$가 상에 속하게 할 수 있다고 하자. 그러면 대응하는 사상
$$
f^*\mathcal{E}_{i, j}
\otimes_{\mathcal{O}_X}
\mathcal{L}^{\otimes n_{i, j}} \longrightarrow \mathcal{F}
$$
들은 함께 전사이므로 보조정리가 증명된다. Properties, Lemma
\ref{properties-lemma-quasi-coherent-colimit-finite-type}에 의해 각
$i, j$에 대해 $V_j$ 위의 절단 $s_{i, j}$를 포함하는 유한형 준연접
부분가군 $\mathcal{H}'_{i, j} \subset  \mathcal{H}_{-n_{j, i}}$를
찾을 수 있다. 따라서 $Y$의 분해 성질로부터 전사
$\mathcal{E}_{i, j} \to \mathcal{H}'_{j, i}$를 얻어 결론이 따른다.
\end{proof}

\begin{lemma}
\label{perfect-lemma-resolution-property-goes-up-affine}
$f : X \to Y$를 스킴의 아핀 또는 준아핀 사상이라 하고,
$Y$가 준콤팩트이고 준분리라고 하자. $Y$가 분해 성질을 가지면
$X$도 분해 성질을 갖는다.
\end{lemma}

\begin{proof}
Morphisms, Lemma \ref{morphisms-lemma-quasi-affine-O-ample}에 의해 이는
Lemma \ref{perfect-lemma-resolution-property-ample-relative}의 특수한 경우이다.
\end{proof}

\noindent
다음은 분해 성질이 아래로 내려감을 증명할 수 있는 경우이다.

\begin{lemma}
\label{perfect-lemma-resolution-property-goes-down-finite-flat}
$f : X \to Y$를 전사인 유한 국소 자유 스킴 사상이라 하자.
$X$가 분해 성질을 가지면 $Y$도 분해 성질을 갖는다.
\end{lemma}

\begin{proof}
이 조건은 $f$가 아핀이고 $f_*\mathcal{O}_X$가 양의 계수를 갖는 유한
국소 자유 $\mathcal{O}_Y$-가군임을 뜻한다. $\mathcal{G}$를 유한형
준연접 $\mathcal{O}_Y$-가군이라 하자. 가정에 의해 어떤 유한 국소 자유
$\mathcal{O}_X$-가군 $\mathcal{E}$와 전사
$\mathcal{E} \to f^*\mathcal{G}$가 존재한다. $f_*$는 준연접 가군들
위에서 완전하므로(Cohomology of Schemes, Lemma
\ref{coherent-lemma-relative-affine-vanishing}) 전사
$$
f_*\mathcal{E}
\longrightarrow
f_*f^*\mathcal{G} = \mathcal{G} \otimes_{\mathcal{O}_Y} f_*\mathcal{O}_X
$$
를 얻는다. 쌍대를 취하면 전사
$$
f_*\mathcal{E}
\otimes_{\mathcal{O}_Y}
\SheafHom_{\mathcal{O}_Y}(f_*\mathcal{O}_X, \mathcal{O}_Y)
\longrightarrow
\mathcal{G}
$$
를 얻는다. $f_*\mathcal{E}$는 유한 국소 자유이므로\footnote{실제로
$A \to B$가 유한 국소 자유 환 사상이고 $N$이 유한 국소 자유
$B$-가군이면, $N$은 유한 국소 자유 $A$-가군이다. 이를 보려면 먼저
$N$이 $B$ 위에서 유한 국소 자유라는 사실로부터 $N$이 평탄하고
유한 표시인 $B$-가군임을 관찰하라. Algebra, Lemma
\ref{algebra-lemma-finite-projective}를 보라. 이어서 Algebra, Lemma
\ref{algebra-lemma-finite-finitely-presented-extension}에 의해 $N$은
유한 표시 $A$-가군이고, Algebra, Lemma
\ref{algebra-lemma-composition-flat}에 의해 평탄 $A$-가군이다. 이제
$A$ 위에서 Algebra, Lemma \ref{algebra-lemma-finite-projective}를
사용하여 결론을 얻는다.}, 결론이 따른다.
\end{proof}

\begin{lemma}
\label{perfect-lemma-resolution-property-finite-number}
$X$를 스킴이라 하자. 다음이 주어졌다고 하자.
\begin{enumerate}
\item 유한 아핀 열린 덮개 $X = U_1 \cup \ldots \cup U_m$,
\item $V(\mathcal{I}_j) = X \setminus U_j$를 만족하는 유한형 준연접
아이디얼 $\mathcal{I}_j$.
\end{enumerate}
그러면 $X$가 분해 성질을 가질 필요충분조건은 $j = 1, \ldots, m$에
대해 $\mathcal{I}_j$가 유한 국소 자유 $\mathcal{O}_X$-가군의 몫인 것이다.
\end{lemma}

\begin{proof}
한 방향은 자명하다. 반대 방향을 위해 $\mathcal{E}_j$가 유한 국소
자유이고 $\mathcal{E}_j \to \mathcal{I}_j$가 전사라고 하자.
다음 문단에서는 뇌터인 경우로 환원한다. 독자는 이 문단을 건너뛰어도 좋다.

\medskip\noindent
먼저 $U_j \cap U_{j'}$는 아핀 스킴 $U_j$ 위의 유한형 준연접 아이디얼
$\mathcal{I}_{j'}|{U_j}$의 영점 스킴의 여집합이므로 준콤팩트이다.
Properties, Lemma
\ref{properties-lemma-quasi-coherent-finite-type-ideals}를 보라.
따라서 $X$는 준콤팩트이고 준분리이다. Schemes, Lemma
\ref{schemes-lemma-characterize-quasi-separated}를 보라. Limits,
Proposition \ref{limits-proposition-approximate}에 의해 $X = \lim X_i$를
아핀 전이 사상을 갖는 뇌터 스킴들의 유향계의 극한으로 쓸 수 있다.
각 $j$에 대해 어떤 $i$와 유한 국소 자유 $\mathcal{O}_{X_i}$-가군
$\mathcal{E}_{i, j}$를 찾아 그 역상이 $\mathcal{E}_j$가 되게 할 수 있다.
Limits, Lemma \ref{limits-lemma-descend-invertible-modules}를 보라.
$i$를 증가시킨 뒤 합성 $\mathcal{E}_j \to \mathcal{I}_j \to
\mathcal{O}_X$가 어떤 사상 $\mathcal{E}_{i, j} \to \mathcal{O}_{X_i}$의
역상이라고 가정할 수 있다. Limits, Lemma
\ref{limits-lemma-descend-modules-finite-presentation}을 보라.
이 사상의 상을 $\mathcal{I}_{i, j} \subset \mathcal{O}_{X_i}$로 나타내자.
이는 뇌터 스킴 $X_i$ 위의 준연접 아이디얼 층이며, 그 $X$로의 역상은
$\mathcal{I}_j$이다. 여집합인 열린집합을 $U_{i, j} \subset X_i$로
나타내면, 모든 $i, j$에 대해 이들이 아핀이라고 가정할 수 있다.
Limits, Lemma \ref{limits-lemma-limit-affine}을 보라.
$X_i$ 위의 열린집합 $U_{i, j}$와 아이디얼 층 $\mathcal{I}_{i, j}$에 대해
보조정리를 증명할 수 있다면, $X$는 $X_i$ 위에서 아핀이므로 Lemma
\ref{perfect-lemma-resolution-property-goes-up-affine}에 의해 분해 성질을 갖는다.
이로써 뇌터 스킴인 경우로 환원한다.

\medskip\noindent
$X$가 뇌터라고 가정하자. 모든 연접 가군 $\mathcal{F}$에 대해,
$\mathcal{F}$의 $U_j$로의 제한을 생성하는 유한한 절단 목록
$s_{jk} \in \mathcal{F}(U_j)$, $k = 1, \ldots, e_j$를 택할 수 있다.
Cohomology of Schemes, Lemma \ref{coherent-lemma-homs-over-open}에 의해
어떤 $n_{jk} \geq 1$에 대해 $s_{jk}$를 사상
$s'_{jk} : \mathcal{I}_i^{n_{jk}} \to \mathcal{F}$로 연장할 수 있다.
이제 합성
$$
\mathcal{E}_j^{\otimes n_{jk}} \to \mathcal{I}_j^{n_{jk}} \to \mathcal{F}
$$
들을 생각하여 결론을 얻는다.
\end{proof}

\begin{lemma}
\label{perfect-lemma-resolution-property-ample-family}
$X$를 스킴이라 하자. $X$가 가역 가군들의 풍부한 족을 가지면
(Morphisms, Definition
\ref{morphisms-definition-family-ample-invertible-modules}),
$X$는 분해 성질을 갖는다.
\end{lemma}

\begin{proof}
$X$가 준콤팩트이므로 어떤 $n$과 쌍 $(\mathcal{L}_i, s_i)$,
$i = 1, \ldots, n$이 존재하여, $\mathcal{L}_i$는 가역
$\mathcal{O}_X$-가군이고 $s_i \in \Gamma(X, \mathcal{L}_i)$는 절단이며,
$s_i$가 영이 아닌 점들의 집합 $U_i \subset X$는 아핀이고
$X = U_1 \cup \ldots \cup U_n$이다. $\mathcal{I}_i \subset
\mathcal{O}_X$를 $s_i : \mathcal{L}_i^{\otimes -1} \to \mathcal{O}_X$의
상이라 하자. Lemma \ref{perfect-lemma-resolution-property-finite-number}를
적용하면 $X$가 분해 성질을 가짐을 알 수 있다.
\end{proof}

\begin{lemma}
\label{perfect-lemma-regular-resolution-property}
$X$를 대각 사상이 아핀인 준콤팩트 정칙 스킴이라 하자.
그러면 $X$는 분해 성질을 갖는다.
\end{lemma}

\begin{proof}
Divisors, Lemma \ref{divisors-lemma-regular-ample-family}와 위의 Lemma
\ref{perfect-lemma-resolution-property-ample-family}를 결합한다.
\end{proof}

\begin{lemma}
\label{perfect-lemma-resolution-property-descends}
$X = \lim X_i$를 아핀 전이 사상을 갖는 준콤팩트 준분리 스킴들의
유향계의 극한이라 하자. 그러면 $X$가 분해 성질을 가질 필요충분조건은
어떤 $i$에 대해 $X_i$가 분해 성질을 갖는 것이다.
\end{lemma}

\begin{proof}
$X_i$가 분해 성질을 가지면 Lemma
\ref{perfect-lemma-resolution-property-goes-up-affine}에 의해 $X$도 그러하다.
$X$가 분해 성질을 갖는다고 가정하자. $i \in I$와 유한 아핀 열린 덮개
$X_i = U_{i, 1} \cup \ldots \cup U_{i, m}$을 택한다. 각 $j$에 대해
$X_i \setminus V(\mathcal{I}_{i, j}) = U_{i, j}$가 되도록 유한형
준연접 아이디얼 층 $\mathcal{I}_{i, j} \subset \mathcal{O}_{X_i}$를
택한다. Properties, Lemma
\ref{properties-lemma-quasi-coherent-finite-type-ideals}를 보라.
$U_{i, j}$의 역상을 $U_j \subset X$로 나타내고,
$\mathcal{I}_{i, j}$의 역상을 $\mathcal{I}_j \subset \mathcal{O}_X$로
나타내자. $X$가 분해 성질을 가지므로 유한 국소 자유
$\mathcal{O}_X$-가군 $\mathcal{E}_j$와 전사
$\mathcal{E}_j \to \mathcal{I}_j$를 택할 수 있다.
Limits, Lemmas \ref{limits-lemma-descend-invertible-modules}와
\ref{limits-lemma-descend-modules-finite-presentation}에 의해 $i$를
증가시킨 뒤, 유한 국소 자유 $\mathcal{O}_{X_i}$-가군
$\mathcal{E}_{i, j}$와 사상 $\mathcal{E}_{i, j} \to \mathcal{O}_{X_i}$를
찾아, 이를 $X$로 기저변환하면 합성
$\mathcal{E}_j \to \mathcal{I}_j \to \mathcal{O}_X$,
$j = 1, \ldots, m$을 되찾게 할 수 있다. 유한 표시
$\mathcal{O}_{X_i}$-가군
$\Coker(\mathcal{E}_{i, j} \to \mathcal{O}_{X_i})$와
$\mathcal{O}_{X_i}/\mathcal{I}_{i, j}$의 $X$로의 역상은
$\mathcal{O}_X$의 몫으로서 일치한다. 따라서 Limits, Lemma
\ref{limits-lemma-descend-modules-finite-presentation}에 의해 이들이
일치한다고 가정할 수 있다. 다시 말해
$\mathcal{E}_{i, j} \to \mathcal{O}_{X_i}$의 상이
$\mathcal{I}_{i, j}$와 같다고 가정할 수 있다. 그러면 Lemma
\ref{perfect-lemma-resolution-property-finite-number}에 의해 $X_i$가 분해
성질을 가짐을 얻는다.
\end{proof}

\begin{lemma}
\label{perfect-lemma-resolution-property-affine-diagonal}
\begin{reference}
\cite[Proposition 1.3]{totaro_resolution}의 특수한 경우이다.
\end{reference}
$X$를 분해 성질을 갖는 준콤팩트 준분리 스킴이라 하자.
그러면 $X$의 대각 사상은 아핀이다.
\end{lemma}

\begin{proof}
Limits, Proposition \ref{limits-proposition-approximate}와 Lemma
\ref{perfect-lemma-resolution-property-descends}를 결합하면 $X$가 뇌터인 경우로
환원된다(작은 세부사항은 생략한다). $X$가 뇌터라고 가정하자.
$X \times X$는 $X$의 아핀 열린집합 $U$, $V$에 대한 아핀 열린집합
$U \times V$들로 덮임을 상기하자. Schemes, Section
\ref{schemes-section-fibre-products}를 보라. 따라서 대각 사상
$\Delta : X \to X \times X$가 아핀임을 보이려면, $X$의 모든 아핀
열린집합 $U$, $V$에 대해 $U \cap V = \Delta^{-1}(U \times V)$가
아핀임을 보이면 충분하다. Morphisms, Lemma
\ref{morphisms-lemma-characterize-affine}을 보라. 특히 $U$가 $X$의
아핀 열린집합일 때 포함 사상 $j : U \to X$가 아핀임을 보이면 충분하다.
Cohomology of Schemes, Lemma
\ref{coherent-lemma-criterion-affine-morphism}에 의해 임의의 준연접
$\mathcal{O}_U$-가군 $\mathcal{G}$에 대해 $R^1j_*\mathcal{G} = 0$임을
보이면 충분하다. Proposition \ref{perfect-proposition-Noetherian}에 의해
(여기서 뇌터인 경우로 환원했다는 사실을 사용한다), $Rj_*\mathcal{G}$를
준연접 $\mathcal{O}_X$-가군들의 복합체 $\mathcal{H}^\bullet$로
나타낼 수 있다. 모순을 얻기 위해
$$
H^1(\mathcal{H}^\bullet) =
\Ker(\mathcal{H}^1 \to \mathcal{H}^2)/\Im(\mathcal{H}^0 \to \mathcal{H}^1)
$$
이 영이 아니라고 가정하자. 그러면 연접 $\mathcal{O}_X$-가군
$\mathcal{F}$와 사상
$$
\mathcal{F} \longrightarrow
\Ker(\mathcal{H}^1 \to \mathcal{H}^2)
$$
를 찾아, 이를 $H^1(\mathcal{H}^\bullet)$로의 사영과 합성한 것이 영이
아니게 할 수 있다. 실제로 Properties, Lemma
\ref{properties-lemma-quasi-coherent-colimit-finite-type}에 의해
$\Ker(\mathcal{H}^1 \to \mathcal{H}^2)$를 그 연접 부분가군들의 여과
합집합으로 쓸 수 있으므로, 그 가운데 하나가 원하는 역할을 한다.
다음으로 $X$의 분해 성질을 사용하여 유한 국소 자유
$\mathcal{O}_X$-가군 $\mathcal{E}$와 전사
$\mathcal{E} \to \mathcal{F}$를 택한다. 이로부터 유도 범주에서 사상
$$
\mathcal{E}[-1] \longrightarrow Rj_*\mathcal{G}
$$
을 얻는데, 이는 코호몰로지 층 위에서 영이 아니므로
$D(\mathcal{O}_X)$에서도 영이 아니다. 수반에 의해 이는 사상
$$
j^*\mathcal{E}[-1] \to \mathcal{G}
$$
이 $D(\mathcal{O}_U)$에서 영이 아닌 것과 같다. $\mathcal{E}$가
유한 국소 자유이므로 이는
$$
H^1(U, j^*\mathcal{E}^\vee \otimes_{\mathcal{O}_U} \mathcal{G})
$$
의 영이 아닌 원소와 같다. 여기서
$\mathcal{E}^\vee = \SheafHom_{\mathcal{O}_X}(\mathcal{E}, \mathcal{O}_X)$
는 쌍대 유한 국소 자유 가군이다. 그러나 Cohomology of Schemes, Lemma
\ref{coherent-lemma-quasi-coherent-affine-cohomology-zero}에 의해 이 군은
영이며, 이는 원하는 모순이다. (쌍대를 사용하는 단계가 의심스럽다면 더
일반적인 Cohomology, Lemma
\ref{cohomology-lemma-dual-perfect-complex}를 보라.)
\end{proof}






\section{분해 성질과 완전 복합체}
\label{perfect-section-resolution-property-perfect}

\noindent
이 절에서는 분해 성질을 갖는 스킴 위의 완전 복합체와 엄밀 완전
복합체 사이의 관계를 논의한다.

\begin{lemma}
\label{perfect-lemma-construct-strictly-perfect}
$X$를 분해 성질을 갖는 준콤팩트 준분리 스킴이라 하자.
$\mathcal{F}^\bullet$을 $D(\mathcal{O}_X)$의 완전 대상을 나타내는,
준연접 $\mathcal{O}_X$-가군들의 아래로 유계인 복합체라 하자.
그러면 유한 국소 자유 $\mathcal{O}_X$-가군들의 유계 복합체
$\mathcal{E}^\bullet$과 준동형
$\mathcal{E}^\bullet \to \mathcal{F}^\bullet$이 존재한다.
\end{lemma}

\begin{proof}
$\mathcal{F}^\bullet$이 $[a, b]$ 안의 Tor-진폭을 갖고 $n < a$에 대해
$\mathcal{F}^n = 0$이 되게 하는 정수 $a, b \in \mathbf{Z}$를 택하자.
그러한 정수쌍의 존재는 Cohomology, Lemma
\ref{cohomology-lemma-perfect}와 $X$의 준콤팩트성에서 따른다.
$b < a$이면 $\mathcal{F}^\bullet$은 유도 범주에서 영이므로 보조정리가
성립한다. $b - a \geq 0$에 관한 귀납법으로, $\mathcal{E}^i$가
유한 국소 자유인 복합체 $\mathcal{E}^a \to \ldots \to \mathcal{E}^b$와
준동형 $\mathcal{E}^\bullet \to \mathcal{F}^\bullet$이 존재함을
증명하겠다.

\medskip\noindent
기초 단계는 $b - a = 0$인 경우이다. 이 경우
$H^b(\mathcal{F}^\bullet) = H^a(\mathcal{F}^\bullet) =
\Ker(\mathcal{F}^a \to \mathcal{F}^{a + 1})$
는 유한 국소 자유이다. 실제로 이는 Tor 차원 $0$인 유한 표시
$\mathcal{O}_X$-가군이므로 유한 국소 자유이다. Cohomology, Lemmas
\ref{cohomology-lemma-perfect}, \ref{cohomology-lemma-finite-cohomology}와
Properties, Lemma \ref{properties-lemma-finite-locally-free}를 보라.
따라서 $\mathcal{E}^\bullet$을 차수 $b$에 놓인
$H^b(\mathcal{F}^\bullet)$로 택할 수 있다. 증명의 나머지는 귀납 단계이다.

\medskip\noindent
$b > a$라고 가정하자. 다음 가군을 관찰하자.
$$
H^b(\mathcal{F}^\bullet) =
\Ker(\mathcal{F}^b \to \mathcal{F}^{b + 1})/
\Im(\mathcal{F}^{b - 1} \to \mathcal{F}^b)
$$
이는 유한형 준연접 $\mathcal{O}_X$-가군이다. Cohomology, Lemmas
\ref{cohomology-lemma-perfect}와
\ref{cohomology-lemma-finite-cohomology}를 보라. 그러면 유한형 준연접
$\mathcal{O}_X$-가군 $\mathcal{F}$와 사상
$$
\mathcal{F} \longrightarrow
\Ker(\mathcal{F}^b \to \mathcal{F}^{b + 1})
$$
를 찾아, 이를 $H^b(\mathcal{F}^\bullet)$로의 사영과 합성한 것이
전사가 되게 할 수 있다. 실제로 Properties, Lemma
\ref{properties-lemma-quasi-coherent-colimit-finite-type}에 의해
$\Ker(\mathcal{F}^b \to \mathcal{F}^{b + 1})$를 그 유한형 준연접
부분가군들의 여과 합집합으로 쓸 수 있으므로, 그 가운데 하나가 원하는
역할을 한다. 다음으로 $X$의 분해 성질을 사용하여 유한 국소 자유
$\mathcal{O}_X$-가군 $\mathcal{E}^b$와 전사
$\mathcal{E}^b \to \mathcal{F}$를 택한다. 복합체의 사상
$$
\alpha : \mathcal{E}^b[-b] \to \mathcal{F}^\bullet
$$
와 그 원뿔 $C(\alpha)^\bullet$을 생각하자. Derived Categories,
Definition \ref{derived-definition-cone}을 보라. $C(\alpha)^\bullet$은
$\geq a$인 차수에서만 영이 아니고, Cohomology, Lemma
\ref{cohomology-lemma-cone-tor-amplitude}에 의해 $[a, b]$ 안의
Tor-진폭을 가지며, 구성에 의해 $H^b(C(\alpha)^\bullet) = 0$이다.
따라서 실제로 $C(\alpha)^\bullet$은 $[a, b - 1]$ 안의 Tor-진폭을
갖는다. 그러므로 $C(\alpha)^\bullet$에 귀납 가설을 적용하여 복합체의 사상
$$
(\mathcal{E}^a \to \ldots \to \mathcal{E}^{b - 1})
\longrightarrow
C(\alpha)^\bullet
$$
을 얻는데, 이는 귀납 가설에 진술된 성질들을 갖는다. 원뿔의 정의를
풀어 쓰면 가환 그림
$$
\xymatrix{
\ldots \ar[r] &
\mathcal{E}^{b - 2} \ar[r] \ar[d] &
\mathcal{E}^{b - 1} \ar[r] \ar[d] &
0 \ar[r] \ar[d] &
\ldots \\
\ldots \ar[r] &
\mathcal{F}^{b - 2} \ar[r] &
\mathcal{F}^{b - 1} \oplus \mathcal{E}^b \ar[r] &
\mathcal{F}^b \ar[r] &
\ldots
}
$$
을 얻는다. 복합체의 사상
$(\mathcal{E}^a \to \ldots \to \mathcal{E}^b) \to \mathcal{F}^\bullet$을
얻음은 분명하다. 이 사상이 준동형이라는 확인은 생략한다.
\end{proof}

\begin{lemma}
\label{perfect-lemma-resolution-property-perfect-complex}
$X$를 분해 성질을 갖는 준콤팩트 준분리 스킴이라 하자. 그러면
$D(\mathcal{O}_X)$의 모든 완전 대상은 유한 국소 자유
$\mathcal{O}_X$-가군들의 유계 복합체로 나타낼 수 있다.
\end{lemma}

\begin{proof}
$E$를 $D(\mathcal{O}_X)$의 완전 대상이라 하자. Lemma
\ref{perfect-lemma-resolution-property-affine-diagonal}에 의해 $X$의 대각 사상은
아핀이다. 따라서 Proposition
\ref{perfect-proposition-quasi-compact-affine-diagonal}에 의해 $E$를 준연접
$\mathcal{O}_X$-가군들의 복합체 $\mathcal{F}^\bullet$로 나타낼 수 있다.
$X$가 준콤팩트이므로 $E$는 $D^b(\mathcal{O}_X)$에 속함을 관찰하자.
따라서 $n \ll 0$이면 $\tau_{\geq n}\mathcal{F}^\bullet$은 $E$를
나타내는 준연접 $\mathcal{O}_X$-가군들의 아래로 유계인 복합체이다.
이제 Lemma \ref{perfect-lemma-construct-strictly-perfect}를 적용하여 결론을 얻는다.
\end{proof}

\begin{lemma}
\label{perfect-lemma-resolution-property-map-perfect-complex}
$X$를 분해 성질을 갖는 준콤팩트 준분리 스킴이라 하자.
$\mathcal{E}^\bullet$과 $\mathcal{F}^\bullet$을 유한 국소 자유
$\mathcal{O}_X$-가군들의 유한 복합체라 하자. 그러면 임의의
$\alpha \in \Hom_{D(\mathcal{O}_X)}(\mathcal{E}^\bullet, \mathcal{F}^\bullet)$
는 그림
$$
\mathcal{E}^\bullet \leftarrow \mathcal{G}^\bullet \to \mathcal{F}^\bullet
$$
으로 나타낼 수 있다. 여기서 $\mathcal{G}^\bullet$은 유한 국소 자유
$\mathcal{O}_X$-가군들의 유계 복합체이고,
$\mathcal{G}^\bullet \to \mathcal{E}^\bullet$은 준동형이다.
\end{lemma}

\begin{proof}
Lemma \ref{perfect-lemma-resolution-property-affine-diagonal}에 의해 $X$의
대각 사상은 아핀이다. 따라서 Proposition
\ref{perfect-proposition-quasi-compact-affine-diagonal}에 의해 $\alpha$를 그림
$$
\mathcal{E}^\bullet \leftarrow \mathcal{H}^\bullet \to \mathcal{F}^\bullet
$$
으로 나타낼 수 있다. 여기서 $\mathcal{H}^\bullet$은 준연접
$\mathcal{O}_X$-가군들의 복합체이고,
$\mathcal{H}^\bullet \to \mathcal{E}^\bullet$은 준동형이다.
$\mathcal{E}^\bullet$과 $\mathcal{F}^\bullet$이 유계 복합체이기 때문에,
$n \ll 0$이면 사상 $\mathcal{H}^\bullet \to \mathcal{E}^\bullet$과
$\mathcal{H}^\bullet \to \mathcal{F}^\bullet$은 준동형
$\mathcal{H}^\bullet \to \tau_{\geq n}\mathcal{H}^\bullet$을 통해
분해된다. 따라서 $\mathcal{H}^\bullet$을
$\tau_{\geq n}\mathcal{H}^\bullet$으로 바꾸어
$\mathcal{H}^\bullet$이 아래로 유계라고 가정할 수 있다. 이제 Lemma
\ref{perfect-lemma-construct-strictly-perfect}를 적용하여 결론을 얻는다.
\end{proof}

\begin{lemma}
\label{perfect-lemma-resolution-property-homotopy-map-perfect-complex}
$X$를 분해 성질을 갖는 준콤팩트 준분리 스킴이라 하자.
$\mathcal{E}^\bullet$과 $\mathcal{F}^\bullet$을 유한 국소 자유
$\mathcal{O}_X$-가군들의 유한 복합체라 하자.
$\alpha^\bullet, \beta^\bullet :\mathcal{E}^\bullet \to \mathcal{F}^\bullet$
을 $D(\mathcal{O}_X)$에서 같은 사상을 정의하는 두 복합체 사상이라 하자.
그러면 준동형 $\gamma^\bullet : \mathcal{G}^\bullet \to
\mathcal{E}^\bullet$이 존재하여, $\mathcal{G}^\bullet$은 유한 국소 자유
$\mathcal{O}_X$-가군들의 유계 복합체이고,
$\alpha^\bullet \circ \gamma^\bullet$과
$\beta^\bullet \circ \gamma^\bullet$은 서로 호모토픽인 복합체 사상이다.
\end{lemma}

\begin{proof}
Lemma \ref{perfect-lemma-resolution-property-affine-diagonal}에 의해 $X$의 대각
사상은 아핀이다. 따라서 Proposition
\ref{perfect-proposition-quasi-compact-affine-diagonal}와 유도 범주의 정의에 의해
준동형 $\gamma^\bullet : \mathcal{G}^\bullet \to
\mathcal{E}^\bullet$이 존재하여, $\mathcal{G}^\bullet$은 준연접
$\mathcal{O}_X$-가군들의 복합체이고,
$\alpha^\bullet \circ \gamma^\bullet$과
$\beta^\bullet \circ \gamma^\bullet$은 서로 호모토픽인 복합체 사상이다.
이를 나타내는 호모토피 $h^i : \mathcal{G}^i \to
\mathcal{F}^{i - 1}$를 택하고 $n \ll 0$을 택하자.
$\mathcal{E}^\bullet$이 아래로 유계이므로 사상 $\gamma^\bullet$은
몫사상 $\mathcal{G}^\bullet \to \tau_{\geq n}\mathcal{G}^\bullet$을
통해 표준적으로 분해된다. 정확히 같은 이유로 사상 $h^i$들은 전사
$\mathcal{G}^i \to (\tau_{\geq n}\mathcal{G})^i$를 통해 분해된다.
따라서 $\mathcal{G}^\bullet$을 $\tau_{\geq n}\mathcal{G}^\bullet$으로
바꿀 수 있다. 이제 Lemma \ref{perfect-lemma-construct-strictly-perfect}를
적용하여 결론을 얻는다.
\end{proof}

\begin{proposition}
\label{perfect-proposition-perfect-resolution-property}
$X$를 분해 성질을 갖는 준콤팩트 준분리 스킴이라 하자. 다음과 같이
나타내자.
\begin{enumerate}
\item $\mathcal{A}$는 유한 국소 자유 $\mathcal{O}_X$-가군들의 가법 범주이다.
\item $K^b(\mathcal{A})$는 $\mathcal{A}$ 안의 유계 복합체들의 호모토피
범주이다. Derived Categories, Section \ref{derived-section-homotopy}를 보라.
\item $D_{perf}(\mathcal{O}_X)$는 $D(\mathcal{O}_X)$에서 완전 대상들로
이루어진 엄밀 충만 포화 삼각 부분범주이다.
\end{enumerate}
이 표기 아래 자명한 함자
$$
K^b(\mathcal{A}) \longrightarrow D_{perf}(\mathcal{O}_X)
$$
는 삼각 범주의 완전 함자이며, 삼각 범주의 동치
$S^{-1}K^b(\mathcal{A}) \to D_{perf}(\mathcal{O}_X)$를 통해 분해된다.
여기서 $S$는 $K^b(\mathcal{A})$ 안의 준동형들로 이루어진 포화 곱셈계이다.
\end{proposition}

\begin{proof}
명제의 진술을 이해할 수 있다면 이 첫 문단은 건너뛰어도 좋다.
사용된 정의들 가운데 일부는
Derived Categories, Definition
\ref{derived-definition-triangulated-subcategory}
(삼각 부분범주),
Derived Categories, Definition \ref{derived-definition-saturated}
(포화 삼각 부분범주),
Derived Categories, Definition \ref{derived-definition-localization}
(삼각 구조와 양립하는 곱셈계), 그리고 Categories,
Definition \ref{categories-definition-saturated-multiplicative-system}
(포화 곱셈계)를 보라. Cohomology, Lemmas
\ref{cohomology-lemma-two-out-of-three-perfect}와
\ref{cohomology-lemma-summands-perfect}에 의해
$D_{perf}(\mathcal{O}_X)$는 $D(\mathcal{O}_X)$의 포화 삼각 부분범주이다.
또한 $K^b(\mathcal{A})$는 삼각 범주이다. Derived Categories, Lemma
\ref{derived-lemma-bounded-triangulated-subcategories}를 보라.

\medskip\noindent
이 함자가 완전 삼각형을 완전 삼각형으로 보내므로 완전함은 분명하다.
그러면 Derived Categories, Lemma
\ref{derived-lemma-triangle-functor-localize}에 의해 $S$는
$K^b(\mathcal{A})$의 삼각 구조와 양립하는 포화 곱셈계이다. 따라서
Derived Categories, Proposition
\ref{derived-proposition-construct-localization}에 의해 국소화
$S^{-1}K^b(\mathcal{A})$가 존재하며 삼각 범주이다. Derived Categories,
Lemma \ref{derived-lemma-universal-property-localization}에 의해 완전한 분해
$S^{-1}K^b(\mathcal{A}) \to D_{perf}(\mathcal{O}_X)$를 얻는다. Lemmas
\ref{perfect-lemma-resolution-property-perfect-complex},
\ref{perfect-lemma-resolution-property-map-perfect-complex},
\ref{perfect-lemma-resolution-property-homotopy-map-perfect-complex}에 의해 이 함자는
동치이다. 마지막으로 Derived Categories, Lemma
\ref{derived-lemma-exact-equivalence}에 의해 함자
$S^{-1}K^b(\mathcal{A}) \to D_{perf}(\mathcal{O}_X)$는 삼각 범주의
동치이다(완전 삼각형들이 서로 대응한다는 의미에서).
\end{proof}






\section{K-군}
\label{perfect-section-K-groups}

\noindent
스킴과 연관된 여러 범주의 $K_0$에 관해 조금만 다룬다. 앞선 내용은
Algebra, Section \ref{algebra-section-K-groups},
Homology, Section \ref{homology-section-K-groups},
Derived Categories, Section \ref{derived-section-K-groups}, 그리고
More on Algebra, Lemma
\ref{more-algebra-lemma-perfect-to-K-group-universal}에서 찾을 수 있다.

\medskip\noindent
Algebra, Section \ref{algebra-section-K-groups}와 유사하게 임의의 뇌터
스킴 $X$에 대해 두 $K$-군 $K'_0(X)$와 $K_0(X)$를 정의하겠다.
첫째 것은 연접 $\mathcal{O}_X$-가군을 사용하고, 둘째 것은 유한 국소 자유
$\mathcal{O}_X$-가군을 사용한다.

\begin{lemma}
\label{perfect-lemma-Noetherian-Kprime}
$X$를 뇌터 스킴이라 하자. 그러면
$$
K_0(\textit{Coh}(\mathcal{O}_X)) =
K_0(D^b(\textit{Coh}(\mathcal{O}_X)) =
K_0(D^b_{\textit{Coh}}(\mathcal{O}_X))
$$
\end{lemma}

\begin{proof}
첫째 등식은 Derived Categories, Lemma
\ref{derived-lemma-K-bounded-derived}이다. Derived Categories, Lemma
\ref{derived-lemma-DBA-map-K}에 의해
$K_0(\textit{Coh}(\mathcal{O}_X)) =
K_0(D^b_{\textit{Coh}}(\mathcal{O}_X))$이다. 이로써 보조정리가 증명된다.
(둘째 등식을 보기 위해 Proposition \ref{perfect-proposition-DCoh}에 의한
$D^b(\textit{Coh}(\mathcal{O}_X)) = D^b_{\textit{Coh}}(\mathcal{O}_X)$도
사용할 수 있다.)
\end{proof}

\noindent
정의는 다음과 같다.

\begin{definition}
\label{perfect-definition-K-group}
$X$를 스킴이라 하자.
\begin{enumerate}
\item $K_0(X)$를 {\it $X$의 그로텐디크 군}이라 한다. 이는
$D(\mathcal{O}_X)$에서 완전 대상들로 이루어진 엄밀 충만 포화 삼각
부분범주 $D_{perf}(\mathcal{O}_X)$의 제0 K-군이다. 식으로 쓰면
$$
K_0(X) = K_0(D_{perf}(\mathcal{O}_X))
$$
\item $X$가 국소 뇌터이면 $K'_0(X)$를
{\it $X$ 위의 연접층들의 그로텐디크 군}이라 한다. 이는 연접
$\mathcal{O}_X$-가군들의 아벨 범주의 제0 $K$-군이다. 식으로 쓰면
$$
K'_0(X) = K_0(\textit{Coh}(\mathcal{O}_X))
$$
\end{enumerate}
\end{definition}

\noindent
$X$가 분해 성질을 가지면(예를 들어 $X$가 풍부한 가역 가군을 가지면),
$K_0(X)$에 대한 우리의 정의가 유한 국소 자유 가군을 사용하는 흔한
정의와 일치함을 보이겠다. Lemma \ref{perfect-lemma-K-is-old-K}를 보라.

\begin{lemma}
\label{perfect-lemma-K-agrees-affine}
$X = \Spec(R)$를 아핀 스킴이라 하자. 그러면 $K_0(X) = K_0(R)$이고,
$R$이 뇌터이면 $K'_0(X) = K'_0(R)$이다.
\end{lemma}

\begin{proof}
$K'_0(R)$과 $K_0(R)$이 Algebra, Section
\ref{algebra-section-K-groups}에서 정의되었음을 상기하자.

\medskip\noindent
More on Algebra, Lemma
\ref{more-algebra-lemma-perfect-to-K-group-universal}에 의해
$K_0(R) = K_0(D_{perf}(R))$이다. Lemmas \ref{perfect-lemma-perfect-affine}와
\ref{perfect-lemma-affine-compare-bounded}에 의해
$D_{perf}(R) = D_{perf}(\mathcal{O}_X)$이다. 이로써 등식
$K_0(R) = K_0(X)$가 증명된다.

\medskip\noindent
Cohomology of Schemes, Lemma \ref{coherent-lemma-coherent-Noetherian}에
의해 $\textit{Coh}(\mathcal{O}_X)$는 유한 $R$-가군들의 범주와
동치이므로 $K'_0(R) = K'_0(X)$이다. 또한 정의로부터 $K'_0(R)$이
유한 $R$-가군들의 범주의 제0 K-군임은 분명하다.
\end{proof}

\noindent
$X$를 뇌터 스킴이라 하자. 그러면 $K'_0(X)$와 $K_0(X)$가 모두
정의된다. 이 경우 표준 사상
$$
K_0(X) = K_0(D_{perf}(\mathcal{O}_X))
\longrightarrow
K_0(D^b_{\textit{Coh}}(\mathcal{O}_X)) = K'_0(X)
$$
이 존재한다. 실제로 완전 복합체는
$D^b_{\textit{Coh}}(\mathcal{O}_X)$에 속하고(Lemma
\ref{perfect-lemma-identify-pseudo-coherent-noetherian}에 의해), 포함 함자
$D_{perf}(\mathcal{O}_X) \to D^b_{\textit{Coh}}(\mathcal{O}_X)$는
제0 $K$-군 위의 사상을 유도하며(Derived Categories, Lemma
\ref{derived-lemma-map-K}), 오른쪽 등식은 Lemma
\ref{perfect-lemma-Noetherian-Kprime}에 의해 성립한다.

\begin{lemma}
\label{perfect-lemma-Kprime-K}
$X$를 뇌터 정칙 스킴이라 하자. 그러면 사상
$K_0(X) \to K'_0(X)$는 동형이다.
\end{lemma}

\begin{proof}
Lemma \ref{perfect-lemma-perfect-on-regular}와 위에서 구성한 사상
$K_0(X) \to K'_0(X)$로부터 즉시 따른다.
\end{proof}

\noindent
$X$를 스킴이라 하자. 유한 국소 자유 $\mathcal{O}_X$-가군들의 범주를
$\textit{Vect}(X)$로 나타내자. 일반적으로 $\textit{Vect}(X)$는 아벨
범주가 아니지만, $\textit{Vect}(X)$의 짧은 완전열이 무엇인지는 분명하다.
다음 성질들을 갖는 유일한 아벨 군을 $K_0(\textit{Vect}(X))$로
나타내자\footnote{여기서 올바른 일반성은 임의의 완전 범주에 대해
$K_0$을 정의하는 것이다. Injectives, Remark
\ref{injectives-remark-embed-exact-category}를 보라.}.
\begin{enumerate}
\item 모든 유한 국소 자유 $\mathcal{O}_X$-가군 $\mathcal{E}$에 대해
$K_0(\textit{Vect}(X))$의 원소 $[\mathcal{E}]$가 주어진다.
\item 유한 국소 자유 $\mathcal{O}_X$-가군들의 모든 짧은 완전열
$0 \to \mathcal{E}' \to \mathcal{E} \to \mathcal{E}'' \to 0$
에 대해 $K_0(\textit{Vect}(X))$에서 관계
$[\mathcal{E}] = [\mathcal{E}'] + [\mathcal{E}'']$가 성립한다.
\item 군 $K_0(\textit{Vect}(X))$는 원소 $[\mathcal{E}]$들로 생성된다.
\item $K_0(\textit{Vect}(X))$에서 생성원 $[\mathcal{E}]$들 사이의 모든
관계는 위와 같은 완전열에서 오는 관계들의 $\mathbf{Z}$-선형결합이다.
\end{enumerate}
$K_0(\textit{Vect}(X))$의 상세한 구성은 생략한다. 자연스러운 사상
$$
K_0(\textit{Vect}(X)) \longrightarrow K_0(X)
$$
이 존재한다. 실제로 유한 국소 자유 $\mathcal{O}_X$-가군
$\mathcal{E}$가 주어지면, 차수 $0$에는 $\mathcal{E}$가 놓이고 다른
차수에서는 영인 $X$ 위의 완전 복합체를 $\mathcal{E}[0]$로 나타내자.
유한 국소 자유 $\mathcal{O}_X$-가군들의 짧은 완전열
$0 \to \mathcal{E} \to \mathcal{E}' \to \mathcal{E}'' \to 0$
이 주어지면 완전 삼각형
$\mathcal{E}[0] \to \mathcal{E}'[0] \to \mathcal{E}''[0] \to \mathcal{E}[1]$,
을 얻는다. Derived Categories, Section
\ref{derived-section-canonical-delta-functor}를 보라. 따라서
$K_0(\textit{Vect}(X)) \to K_0(D_{perf}(\mathcal{O}_X)) = K_0(X)$
를 얻는데, 이는 $[\mathcal{E}]$를 $[\mathcal{E}[0]]$로 보낸다.
끔찍한 표기는 양해를 구한다.

\begin{lemma}
\label{perfect-lemma-K-is-old-K}
$X$를 분해 성질을 갖는 준콤팩트 준분리 스킴이라 하자. 그러면 사상
$K_0(\textit{Vect}(X)) \to K_0(X)$는 동형이다.
\end{lemma}

\begin{proof}
이 보조정리는 Lemmas \ref{perfect-lemma-resolution-property-perfect-complex},
\ref{perfect-lemma-resolution-property-map-perfect-complex}, 및
\ref{perfect-lemma-resolution-property-homotopy-map-perfect-complex}
에서 직접 따른다. 이 결과들은 더 언급하지 않고 사용하겠다.
역사상
$$
c : K_0(X) = K_0(D_{perf}(\mathcal{O}_X)) \longrightarrow
K_0(\textit{Vect}(X))
$$
을 구성하자. 실제로 $D_{perf}(\mathcal{O}_X)$의 임의의 대상은 유한 국소
자유 $\mathcal{O}_X$-가군들의 유계 복합체 $\mathcal{E}^\bullet$로
나타낼 수 있다. 다음과 같이 놓는다.
$$
c([\mathcal{E}^\bullet]) = \sum (-1)^i[\mathcal{E}^i]
$$
물론 이것이 잘 정의됨을 보여야 한다. 우선 $c$를 유한 국소 자유
$\mathcal{O}_X$-가군들의 유계 복합체들 위에 정의된 사상으로 생각한다.

\medskip\noindent
유한 국소 자유 $\mathcal{O}_X$-가군들의 유계 복합체 사이에 전사 사상
$\mathcal{E}^\bullet \to \mathcal{F}^\bullet$이 주어졌다고 하자.
$\mathcal{K}^\bullet$을 그 핵이라 하자. 그러면 $\mathcal{O}_X$-가군들의
짧은 완전열
$$
0 \to \mathcal{K}^n \to \mathcal{E}^n \to \mathcal{F}^n \to 0
$$
을 얻는다. $\mathcal{F}^n$이 유한 국소 자유이므로 이 열들은 국소적으로
분할된다. 따라서 $\mathcal{K}^\bullet$도 유한 국소 자유
$\mathcal{O}_X$-가군들의 유계 복합체이며, $K_0(\textit{Vect}(X))$에서
$c(\mathcal{E}^\bullet) = c(\mathcal{K}^\bullet) + c(\mathcal{F}^\bullet)$
가 성립한다.

\medskip\noindent
유한 국소 자유 $\mathcal{O}_X$-가군들의 비순환 유계 복합체
$\mathcal{E}^\bullet$이 주어졌다고 하자. $n \not \in [a, b]$일 때
$\mathcal{E}^n = 0$이라고 하자. 그러면 $\mathcal{E}^\bullet$을 다음
짧은 완전열들로 분해할 수 있다.
$$
\begin{matrix}
0 \to \mathcal{E}^a \to \mathcal{E}^{a + 1} \to \mathcal{F}^{a + 1} \to 0,\\
0 \to \mathcal{F}^{a + 1} \to \mathcal{E}^{a + 2} \to
\mathcal{F}^{a + 3} \to 0, \\
\ldots \\
0 \to \mathcal{F}^{b - 3} \to \mathcal{E}^{b - 2} \to
\mathcal{F}^{b - 2} \to 0, \\
0 \to \mathcal{F}^{b - 2} \to \mathcal{E}^{b - 1} \to \mathcal{E}^b \to 0
\end{matrix}
$$
하강 귀납법을 적용하면
$\mathcal{F}^{b - 2}, \ldots, \mathcal{F}^{a + 1}$
은 유한 국소 자유 $\mathcal{O}_X$-가군들이며,
$$
c(\mathcal{E}^\bullet) = \sum (-1)[\mathcal{E}^n] =
\sum (-1)^n([\mathcal{F}^{n - 1}] + [\mathcal{F}^n]) = 0
$$
임을 알 수 있다. 따라서 우리의 구성은 비순환 복합체에서 영을 준다.

\medskip\noindent
앞의 두 문단의 결과로부터 $c$가 잘 정의됨이 따른다. 실제로 유한 국소
자유 $\mathcal{O}_X$-가군들의 유계 복합체 $\mathcal{E}^\bullet$과
$\mathcal{F}^\bullet$이 $D(\mathcal{O}_X)$의 같은 대상을 나타낸다고 하자.
그러면 준동형사상
$a : \mathcal{G}^\bullet \to \mathcal{E}^\bullet$와
$b : \mathcal{G}^\bullet \to \mathcal{F}^\bullet$
을 찾을 수 있으며, 여기서 $\mathcal{G}^\bullet$은 유한 국소 자유
$\mathcal{O}_X$-가군들의 유계 복합체이다. 복합체들의 짧은 완전열
$$
0 \to \mathcal{E}^\bullet \to C(a)^\bullet \to \mathcal{G}^\bullet[1] \to 0
$$
을 얻는다. Derived Categories, Definition \ref{derived-definition-cone}을
보라. $a$가 준동형사상이므로 원뿔 $C(a)^\bullet$은 비순환이다(예를 들어
Derived Categories, Section \ref{derived-section-canonical-delta-functor}의
논의로부터 따른다). 따라서
$$
0 = c(C(f)^\bullet) = c(\mathcal{E}^\bullet) + c(\mathcal{G}^\bullet[1]) =
c(\mathcal{E}^\bullet) - c(\mathcal{G}^\bullet)
$$
를 얻는다. $b$를 사용한 같은 논증은
$0 = c(\mathcal{F}^\bullet) - c(\mathcal{G}^\bullet)$임을 보인다.
따라서 $c(\mathcal{E}^\bullet) = c(\mathcal{F}^\bullet)$이고 $c$는 잘 정의된다.

\medskip\noindent
유한 국소 자유 $\mathcal{O}_X$-가군들의 유계 복합체 사이의 사상
$\mathcal{E}^\bullet \to \mathcal{F}^\bullet$의 원뿔을 사용하는 비슷한
논증으로, $X \to Y \to Z$가 $D_{perf}(\mathcal{O}_X)$의 완전 삼각형이면
$c(Y) = c(X) + c(Z)$임을 알 수 있다. 자세한 내용은 생략한다. 따라서
원하는 아벨 군 준동형 $c : K_0(X) \to K_0(\textit{Vect}(X))$를 얻는다.

\medskip\noindent
합성
$K_0(\textit{Vect}(X)) \to K_0(X) \to K_0(\textit{Vect}(X))$
이 항등사상임은 분명하다. 한편 $\mathcal{E}^\bullet$을 유한 국소 자유
$\mathcal{O}_X$-가군들의 유계 복합체라 하자. ``단순 절단''의 완전 삼각형
(Homology, Section \ref{homology-section-truncations}를 보라)
$$
\sigma_{\geq n}\mathcal{E}^\bullet \to
\sigma_{\geq n - 1}\mathcal{E}^\bullet \to
\mathcal{E}^{n - 1}[-n + 1] \to
(\sigma_{\geq n}\mathcal{E}^\bullet)[1]
$$
이 존재한다는 사실과 귀납법으로
$$
[\mathcal{E}^\bullet] = \sum (-1)^i[\mathcal{E}^i[0]]
$$
임을 알 수 있다. 이 등식은 $K_0(X) = K_0(D_{perf}(\mathcal{O}_X))$에서
성립하며, 표기는 양해를 구한다. 따라서 사상
$K_0(\textit{Vect}(X)) \to K_0(D_{perf}(\mathcal{O}_X)) = K_0(X)$는
전사이고, 이로써 증명이 끝난다.
\end{proof}

\begin{remark}
\label{perfect-remark-K-ring}
$X$를 스킴이라 하자. K-군 $K_0(X)$는 표준적으로 가환환이다. 실제로
유도 텐서곱
$$
\otimes = \otimes^\mathbf{L}_{\mathcal{O}_X} :
D_{perf}(\mathcal{O}_X) \times D_{perf}(\mathcal{O}_X)
\longrightarrow
D_{perf}(\mathcal{O}_X)
$$
과 Derived Categories, Lemma \ref{derived-lemma-bilinear-map-K}를 사용하면
쌍선형 곱셈을 얻는다. $K \otimes L \cong L \otimes K$이므로 이 곱은
가환이다. 또한
$(K \otimes L) \otimes M = K \otimes (L \otimes M)$
이므로 이 곱은 결합적이다. 마지막으로 $K_0(X)$의 단위원은 원소
$1 = [\mathcal{O}_X]$이다.

\medskip\noindent
$\textit{Vect}(X)$와 $K_0(\textit{Vect}(X))$가 위와 같다면,
$K_0(\textit{Vect}(X))$도 환 구조를 갖는다. 즉 $\mathcal{E}$와
$\mathcal{F}$가 유한 국소 자유 $\mathcal{O}_X$-가군이면 다음과 같이 놓는다.
$$
[\mathcal{E}] \cdot [\mathcal{F}] =
[\mathcal{E} \otimes_{\mathcal{O}_X} \mathcal{F}]
$$
이것이 실제로 쌍선형이고 가환이며 결합적인 곱을 정의함을 독자는 쉽게
확인할 수 있다. 자세한 내용은 생략한다. 위에서 구성한 사상
$$
K_0(\textit{Vect}(X)) \longrightarrow K_0(X)
$$
은 이 정의들에 관해 환 준동형이다.

\medskip\noindent
이제 $X$가 뇌터라고 가정하자. 유도 텐서곱은 또한 사상
$$
\otimes = \otimes^\mathbf{L}_{\mathcal{O}_X} :
D_{perf}(\mathcal{O}_X) \times D^b_{\textit{Coh}}(\mathcal{O}_X)
\longrightarrow
D^b_{\textit{Coh}}(\mathcal{O}_X)
$$
을 준다. Lemma \ref{perfect-lemma-Noetherian-Kprime}에 의해
$K'_0(X) = K_0(D^b_{\textit{Coh}}(\mathcal{O}_X))$이므로, 다시
Derived Categories, Lemma \ref{derived-lemma-bilinear-map-K}를 사용하면
쌍선형 곱셈 $K_0(X) \times K'_0(X) \to K'_0(X)$을 얻는다. 이로써
$K'_0(X)$에 환 $K_0(X)$ 위의 가군 구조가 주어짐을 독자는 쉽게 보일 수 있다.
\end{remark}

\begin{remark}
\label{perfect-remark-pushforward-K}
$f : X \to Y$를 국소 뇌터 스킴들의 고유 사상이라 하자. 사상
$$
f_* : K'_0(X) \longrightarrow K'_0(Y)
$$
이 존재하며, 이는 $[\mathcal{F}]$를
$$
[\bigoplus\nolimits_{i \geq 0} R^{2i}f_*\mathcal{F}] -
[\bigoplus\nolimits_{i \geq 0} R^{2i + 1}f_*\mathcal{F}]
$$
로 보낸다. 층 $R^if_*\mathcal{F}$들이 연접이고(Cohomology of Schemes,
Proposition \ref{coherent-proposition-proper-pushforward-coherent}), 국소적으로
유한 개만 영이 아니며, $X$ 위의 연접층들의 짧은 완전열이 $Y$ 위에서
$R^if_*$의 긴 완전열을 주므로 이 사상은 잘 정의된다. $Y$가 준콤팩트이면
(실제로 가장 자주 사용하는 유일한 경우), 위 식을 다음과 같이 다시 쓸 수 있다.
$$
f_*[\mathcal{F}] = \sum (-1)^i[R^if_*\mathcal{F}] = [Rf_*\mathcal{F}]
$$
여기서는 Lemma \ref{perfect-lemma-Noetherian-Kprime}의 등식
$K'_0(Y) = K_0(D^b_{\textit{Coh}}(Y))$을 사용했다.
\end{remark}

\begin{lemma}
\label{perfect-lemma-projection-formula}
$f : X \to Y$를 국소 뇌터 스킴들의 고유 사상이라 하자. 그러면
$\alpha \in K'_0(X)$와 $\beta \in K_0(Y)$에 대해
$f_*(\alpha \cdot f^*\beta) = f_*\alpha \cdot \beta$이다.
\end{lemma}

\begin{proof}
Lemma \ref{perfect-lemma-cohomology-base-change}, Remark \ref{perfect-remark-pushforward-K}의
논의, 그리고 Remark \ref{perfect-remark-K-ring}의 곱
$K'_0(X) \times K_0(X) \to K'_0(X)$의 정의로부터 따른다.
\end{proof}

\begin{remark}
\label{perfect-remark-perf-Z}
$X$를 스킴이라 하고 $Z \subset X$를 닫힌 부분스킴이라 하자. 엄밀 충만하고
포화된 삼각 부분범주
$$
D_{Z, perf}(\mathcal{O}_X) \subset D(\mathcal{O}_X)
$$
를 생각하자. 이는 코호몰로지 층들이 집합론적으로 $Z$ 위에 지지되는
$\mathcal{O}_X$-가군들의 완전 복합체들로 이루어진다. 이 삼각 범주의
0차 $K$-군 $K_0(D_{Z, perf}(\mathcal{O}_X))$는 때때로 $K_Z(X)$ 또는
$K_{0, Z}(X)$로 표기한다. Remark \ref{perfect-remark-K-ring}에서와 정확히 같이
유도 텐서곱을 사용하면 $K_0(D_{Z, perf}(\mathcal{O}_X))$가 결합적이고
가환인 곱셈을 가짐을 알 수 있지만, 일반적으로
$K_0(D_{Z, perf}(\mathcal{O}_X))$에는 단위원이 없다.
\end{remark}













\section{복합체의 행렬식}
\label{perfect-section-det-two-terms}

\noindent
이 절은 More on Algebra, Section
\ref{more-algebra-section-determinants-complexes}의 연속이다. 임의의 환 달린
공간 $(X, \mathcal{O}_X)$에 대해 함자
$$
\det :
\left\{
\begin{matrix}
\text{완전 복합체들의 범주} \\
\text{사상은 동형사상}
\end{matrix}
\right\}
\longrightarrow
\left\{
\begin{matrix}
\text{가역 가군들의 범주} \\
\text{사상은 동형사상}
\end{matrix}
\right\}
$$
가 존재한다. 더욱이 여과가 유한하고 등급 부분들이 완전 복합체인 여과 유도
범주 $DF(\mathcal{O}_X)$의 대상 $(L, F)$가 주어지면 표준 동형사상
$\det(\text{gr}L) \to \det(L)$이 존재한다. 최초의 해설은
\cite{determinant}를 보라. 이 내용은 나중에 추가하겠다(향후 참고문헌 삽입).

\medskip\noindent
당장은 $X$가 스킴이고 $D(\mathcal{O}_X)$에서 Tor-진폭이 $[-1, 0]$에
들어 있는 완전 대상 $L$을 생각하는 경우에 대한 특별한 구성을 제시하겠다.

\begin{lemma}
\label{perfect-lemma-determinant-two-term-complexes}
$X$를 스킴이라 하자. 함자
$$
\det :
\left\{
\begin{matrix}
\text{완전 복합체들의 범주} \\
\text{Tor-진폭이 }[-1, 0]\text{에 들어 있음} \\
\text{사상은 동형사상}
\end{matrix}
\right\}
\longrightarrow
\left\{
\begin{matrix}
\text{가역 가군들의 범주} \\
\text{사상은 동형사상}
\end{matrix}
\right\}
$$
가 존재한다. 또한 $D(\mathcal{O}_X)$에서 계수가 $0$이고 Tor-진폭이
$[-1, 0]$에 들어 있는 완전 대상 $L$이 주어지면 표준 원소
$\delta(L) \in \Gamma(X, \det(L))$이 존재하여, $D(\mathcal{O}_X)$의 임의의
동형사상 $a : L \to K$에 대해 $\det(a)(\delta(L)) = \delta(K)$가 성립한다.
더욱이 이 구성은 아핀 국소적으로 More on Algebra, Section
\ref{more-algebra-section-determinants-complexes}의 구성으로 주어진다.
\end{lemma}

\begin{proof}
$L$을 왼쪽 범주의 대상이라 하자. $\Spec(A) = U \subset X$가 아핀
열린집합이면 $L|_U$는 Tor-진폭이 $[-1, 0]$에 들어 있는 $A$-가군들의
완전 복합체 $L^\bullet$에 대응한다. Lemmas
\ref{perfect-lemma-affine-compare-bounded},
\ref{perfect-lemma-tor-dimension-affine}, 및
\ref{perfect-lemma-perfect-affine}를 보라. 그러면 More on Algebra, Lemma
\ref{more-algebra-lemma-determinant-two-term-complexes}에서 구성한 가역
$A$-가군 $\det(L^\bullet)$을 생각할 수 있다. $\Spec(B) = V \subset U$가
$U$에 포함된 또 다른 아핀 열린집합이면
$\det(L^\bullet) \otimes_A B = \det(L^\bullet \otimes_A B)$이다. 따라서
이 구성은 제한 사상과 양립한다(Lemma \ref{perfect-lemma-quasi-coherence-pullback}를
보고 $A \to B$가 평탄임에 유의하라). 그러므로 이 가역 가군들을 붙여서
$X$ 위의 가역 가군 $\det(L)$을 얻을 수 있다. 함자성과 표준 단면도 정확히
같은 방식으로 구성한다. 자세한 내용은 생략한다.
\end{proof}

\begin{remark}
\label{perfect-remark-functorial-det}
Lemma \ref{perfect-lemma-determinant-two-term-complexes}의 구성은 당김과 양립한다.
더 정확히 말해 스킴들의 사상 $f : X \to Y$와 $D(\mathcal{O}_Y)$에서
Tor-진폭이 $[-1, 0]$에 들어 있는 완전 대상 $K$가 주어지면, $Lf^*K$는
$D(\mathcal{O}_X)$에서 Tor-진폭이 $[-1, 0]$에 들어 있는 완전 대상 $K$이고
표준 식별
$$
f^*\det(K) \longrightarrow \det(Lf^*K)
$$
을 갖는다. 더욱이 $K$의 계수가 $0$이면 이 사상을 통해 $\delta(K)$는
$\delta(Lf^*K)$로 당겨진다. 이는 행렬식의 아핀 국소 구성에서 분명하다.
\end{remark}






\section{유계성 판별}
\label{perfect-section-detecting-boundedness}

\noindent
이 절에서는 Section \ref{perfect-section-generating}에서 구성한 준콤팩트 준분리
스킴의 $D_\QCoh$의 콤팩트 생성자가 특별한 성질을 가짐을 보인다. 그 절은
이 절과 매우 비슷하므로 먼저 읽기를 권한다.

\begin{lemma}
\label{perfect-lemma-orthogonal-koszul-first-variant}
Situation \ref{perfect-situation-complex}에서 열린 매장을 $j : U \to X$로 나타내고,
$K$를 $A$ 위의 $f_1, \ldots, f_r$에 대한 코줄 복합체에 대응하는
$D(\mathcal{O}_X)$의 완전 대상이라 하자. $E \in D_\QCoh(\mathcal{O}_X)$와
$a \in \mathbf{Z}$에 대해 다음 조건들을 생각하자.
\begin{enumerate}
\item 표준 사상 $\tau_{\geq a}E \to \tau_{\geq a} Rj_*(E|_U)$는
동형사상이다.
\item 모든 $n \geq a$에 대해
$\Hom_{D(\mathcal{O}_X)}(K[-n], E) = 0$이다.
\end{enumerate}
그러면 (2)는 (1)을 함의하고, (1)은 $a$를 $a + 1$로 바꾼 (2)를 함의한다.
\end{lemma}

\begin{proof}
완전 삼각형 $N \to E \to Rj_*(E|_U) \to N[1]$을 선택하자. 그러면 (1)은
$\tau_{\geq a + 1} N = 0$을 함의하고, $\tau_{\geq a}N = 0$은 (1)을
함의한다. 다음에 유의하라.
$$
\Hom_{D(\mathcal{O}_X)}(K[-n], Rj_*(E|_U)) =
\Hom_{D(\mathcal{O}_U)}(K|_U[-n], E) = 0
$$
$K|_U = 0$이므로 이는 모든 $n$에 대해 성립한다. 따라서 (2)는 모든
$n \geq a$에 대해 $\Hom_{D(\mathcal{O}_X)}(K[-n], N) = 0$인 것과
동치이다. 코줄 복합체들의 다음 완전 삼각형들이 존재함에 유의하라.
$$
K^\bullet(f_1^{e_1}, \ldots, f_i^{e'_i}, \ldots, f_r^{e_r}) \to
K^\bullet(f_1^{e_1}, \ldots, f_i^{e'_i + e''_i}, \ldots, f_r^{e_r}) \to
K^\bullet(f_1^{e_1}, \ldots, f_i^{e''_i}, \ldots, f_r^{e_r}) \to \ldots
$$
More on Algebra, Lemma \ref{more-algebra-lemma-koszul-mult}를 보라. 따라서
모든 $n \geq a$에 대해
$\Hom_{D(\mathcal{O}_X)}(K[-n], N) = 0$인 것은 Lemma
\ref{perfect-lemma-represent-cohomology-class-on-closed}에서와 같은 $K_e$에 대해,
모든 $n \geq a$와 모든 $e \geq 1$에 대해
$\Hom_{D(\mathcal{O}_X)}(K_e[-n], N) = 0$인 것과 동치이다. $N|_U = 0$이므로
그 보조정리에 의해 이는 다시 $n \geq a$에 대해 $H^n(X, N) = 0$인 것과
동치이다. $N$은 $A$-가군 복합체 $R\Gamma(X, N)$에 의해 결정되므로, (2)는
$\tau_{\geq a}N = 0$과 동치이다. Lemma
\ref{perfect-lemma-affine-compare-bounded}를 보라. 따라서 보조정리가 성립한다.
\end{proof}

\begin{lemma}
\label{perfect-lemma-orthogonal-koszul-second-variant}
Situation \ref{perfect-situation-complex}에서 열린 매장을 $j : U \to X$로 나타내고,
$K$를 $A$ 위의 $f_1, \ldots, f_r$에 대한 코줄 복합체에 대응하는
$D(\mathcal{O}_X)$의 완전 대상이라 하자. $E \in D_\QCoh(\mathcal{O}_X)$와
$a \in \mathbf{Z}$에 대해 다음 조건들을 생각하자.
\begin{enumerate}
\item 표준 사상 $\tau_{\leq a}E \to \tau_{\leq a} Rj_*(E|_U)$는
동형사상이다.
\item 모든 $n \leq a$에 대해
$\Hom_{D(\mathcal{O}_X)}(K[-n], E) = 0$이다.
\end{enumerate}
그러면 (2)는 (1)을 함의하고, (1)은 $a$를 $a - 1$로 바꾼 (2)를 함의한다.
\end{lemma}

\begin{proof}
완전 삼각형 $E \to Rj_*(E|_U) \to N \to E[1]$을 선택하자. 그러면 (1)은
$\tau_{\leq a - 1}N = 0$을 함의하고, $\tau_{\leq a}N = 0$은 (1)을
함의한다. 다음에 유의하라.
$$
\Hom_{D(\mathcal{O}_X)}(K[-n], Rj_*(E|_U)) =
\Hom_{D(\mathcal{O}_U)}(K|_U[-n], E) = 0
$$
$K|_U = 0$이므로 이는 모든 $n$에 대해 성립한다. 따라서 (2)는 모든
$n \leq a$에 대해 $\Hom_{D(\mathcal{O}_X)}(K[-n], N) = 0$인 것과
동치이다. 코줄 복합체들의 다음 완전 삼각형들이 존재함에 유의하라.
$$
K^\bullet(f_1^{e_1}, \ldots, f_i^{e'_i}, \ldots, f_r^{e_r}) \to
K^\bullet(f_1^{e_1}, \ldots, f_i^{e'_i + e''_i}, \ldots, f_r^{e_r}) \to
K^\bullet(f_1^{e_1}, \ldots, f_i^{e''_i}, \ldots, f_r^{e_r}) \to \ldots
$$
More on Algebra, Lemma \ref{more-algebra-lemma-koszul-mult}를 보라. 따라서
모든 $n \leq a$에 대해
$\Hom_{D(\mathcal{O}_X)}(K[-n], N) = 0$인 것은 Lemma
\ref{perfect-lemma-represent-cohomology-class-on-closed}에서와 같은 $K_e$에 대해,
모든 $n \leq a$와 모든 $e \geq 1$에 대해
$\Hom_{D(\mathcal{O}_X)}(K_e[-n], N) = 0$인 것과 동치이다. $N|_U = 0$이므로
그 보조정리에 의해 이는 다시 $n \leq a$에 대해 $H^n(X, N) = 0$인 것과
동치이다. $N$은 $A$-가군 복합체 $R\Gamma(X, N)$에 의해 결정되므로, (2)는
$\tau_{\leq a}N = 0$과 동치이다. Lemma
\ref{perfect-lemma-affine-compare-bounded}를 보라. 따라서 보조정리가 성립한다.
\end{proof}

\begin{lemma}
\label{perfect-lemma-bounded-truncation}
$X$를 준콤팩트이고 준분리인 스킴이라 하자.
$P \in D_{perf}(\mathcal{O}_X)$와 $E \in D_{\QCoh}(\mathcal{O}_X)$라 하고,
$a \in \mathbf{Z}$라 하자. 다음은 동치이다.
\begin{enumerate}
\item $\Hom_{D(\mathcal{O}_X)}(P[-i], E) = 0$이 $i \gg 0$에 대해 성립한다.
\item $\Hom_{D(\mathcal{O}_X)}(P[-i], \tau_{\geq a} E) = 0$이
$i \gg 0$에 대해 성립한다.
\end{enumerate}
\end{lemma}

\begin{proof}
삼각형 $ \tau_{< a} E \to E \to \tau_{\geq a} E \to$를 사용하면, $i \gg 0$에
대해
$$
\Hom_{D(\mathcal{O}_X)}(P[-i], \tau_{< a} E) =
\Hom_{D(\mathcal{O}_X)}(P, (\tau_{< a} E)[i]) = 0 
$$
임을 보이면 동치가 따름을 알 수 있다. $P$가 완전이므로 이는 Lemma
\ref{perfect-lemma-ext-from-perfect-into-bounded-QCoh}에 의해 참이다.
\end{proof}

\begin{lemma}
\label{perfect-lemma-bounded-below-truncation}
$X$를 준콤팩트이고 준분리인 스킴이라 하자.
$P \in D_{perf}(\mathcal{O}_X)$와 $E \in D_{\QCoh}(\mathcal{O}_X)$라 하고,
$a \in \mathbf{Z}$라 하자. 다음은 동치이다.
\begin{enumerate}
\item $\Hom_{D(\mathcal{O}_X)}(P[-i], E) = 0$이 $i \ll 0$에 대해 성립한다.
\item $\Hom_{D(\mathcal{O}_X)}(P[-i], \tau_{\leq a} E) = 0$이
$i \ll 0$에 대해 성립한다.
\end{enumerate}
\end{lemma}

\begin{proof}
삼각형 $ \tau_{\leq a} E \to E \to \tau_{> a} E \to$를 사용하면, $i \ll 0$에
대해
$$
\Hom_{D(\mathcal{O}_X)}(P[-i], \tau_{> a} E) =
\Hom_{D(\mathcal{O}_X)}(P, (\tau_{> a} E)[i]) = 0
$$
임을 보이면 동치가 따름을 알 수 있다. $P$가 완전이므로 이는 Lemma
\ref{perfect-lemma-ext-from-perfect-into-bounded-QCoh}에 의해 참이다.
\end{proof}

\begin{proposition}
\label{perfect-proposition-detecting-bounded-above}
$X$를 준콤팩트이고 준분리인 스킴이라 하자.
$G \in D_{perf}(\mathcal{O}_X)$를 $D_\QCoh (\mathcal{O}_X)$를 생성하는 완전
복합체라 하고, $E \in D_\QCoh (\mathcal{O}_X)$라 하자. 다음은 동치이다.
\begin{enumerate}
\item $E \in D^-_\QCoh (\mathcal{O}_X)$이다.
\item $\Hom_{D(\mathcal{O}_X)}(G[-i], E) = 0$이 $i \gg 0$에 대해 성립한다.
\item $\Ext^i_X(G, E) = 0$이 $i \gg 0$에 대해 성립한다.
\item $R\Hom_X(G, E)$는 $D^-(\mathbf{Z})$에 속한다.
\item $H^i(X, G^\vee \otimes_{\mathcal{O}_X}^\mathbf{L} E) = 0$
$i \gg 0$에 대해 성립한다.
\item $R\Gamma(X, G^\vee \otimes_{\mathcal{O}_X}^\mathbf{L} E)$
는 $D^-(\mathbf{Z})$에 속한다.
\item $D(\mathcal{O}_X)$의 모든 완전 대상 $P$에 대해 다음이 성립한다.
\begin{enumerate}
\item $G$를 $P$로 바꾼 명제 (2), (3), (4)가 성립한다.
\item $H^i(X, P \otimes_{\mathcal{O}_X}^\mathbf{L} E) = 0$이
$i \gg 0$에 대해 성립한다.
\item $R\Gamma(X, P \otimes_{\mathcal{O}_X}^\mathbf{L} E)$
는 $D^-(\mathbf{Z})$에 속한다.
\end{enumerate}
\end{enumerate}
\end{proposition}

\begin{proof}
(1)을 가정하자. 등식
$\Hom_{D(\mathcal{O}_X)}(G[-i], E) = \Hom_{D(\mathcal{O}_X)}(G, E[i])$
과 Lemma \ref{perfect-lemma-ext-from-perfect-into-bounded-QCoh}에 의해 이는
$i \gg 0$에 대해 영이다. 따라서 (1)은 (2)를 함의한다.

\medskip\noindent
Cohomology, Section \ref{cohomology-section-global-RHom}의 논의에 의해
(2), (3), (4)는 동치이다. 정의에 의해
$H^i(X, -) = H^i(R\Gamma(X, -))$이므로 (5)와 (6)은 동치이다. Cohomology,
Lemma \ref{cohomology-lemma-dual-perfect-complex}와
$(G[-i])^\vee = G^\vee[i]$라는 사실에 의해, 동치인 조건 (2), (3), (4)는
동치인 조건 (5), (6)과 동치이다.

\medskip\noindent
(7)이 (2)를 함의함은 분명하다. 반대로 동치인 조건 (2)--(6)이 (7)을
함의함을 증명하자. Remark \ref{perfect-remark-classical-generator}에 의해 $G$는
$D_{perf}(\mathcal{O}_X)$의 고전적 생성자임을 상기하라.
$P \in D_{perf}(\mathcal{O}_X)$에 대해 $T(P)$를
$R\Hom_X(P, E)$가 $D^-(\mathbf{Z})$에 속한다는 명제라 하자. 분명히 $T$는
직합과 직합인자에 의해 유전되고, 완전 삼각형에 대해 셋 중 둘 성질을
만족하며, 이동에 의해 보존된다. 따라서 Derived Categories, Remark
\ref{derived-remark-check-on-generator}에 의해 (4)는 모든
$D_{perf}(\mathcal{O}_X)$에서 $T$가 성립함을 함의한다. 다른 모든 성질에도
같은 논증이 적용되지만, (7)(b)와 (7)(c)에는 $P \mapsto P^\vee$가
$D_{perf}(\mathcal{O}_X)$의 자기 동치라는 사실도 사용한다. 작은 세부사항은
생략한다.

\medskip\noindent
Cohomology of Schemes, Lemma \ref{coherent-lemma-induction-principle}의
귀납 원리를 사용하여 동치인 조건 (2)--(7)이 (1)을 함의함을 증명하겠다.

\medskip\noindent
먼저 $X$가 아핀일 때 (2)--(7) $\Rightarrow$ (1)을 증명하자.
$P = \mathcal{O}_X[0]$로 놓자. (7)에서 $i \gg 0$에 대해
$H^i (X, E) = 0$을 얻는다. $E$는 $R\Gamma (X, E)$에 의해 결정되므로 (1)이
따른다. Lemma \ref{perfect-lemma-affine-compare-bounded}를 보라.

\medskip\noindent
이제 $X = U \cup V$이고, $U$는 $X$의 준콤팩트 열린집합이며, $V$는 아핀
열린집합이라고 하자. 또한 스킴 $U$, $V$, $U \cap V$에 대해 함의
(2)--(7) $\Rightarrow$ (1)이 알려졌다고 하자.
$E \in D_\QCoh(\mathcal{O}_X)$가 (2)--(7)을 만족한다고 하자. Lemma
\ref{perfect-lemma-direct-summand-of-a-restriction}과 Theorem
\ref{perfect-theorem-bondal-van-den-Bergh}에 의해 $Q|_U$가
$D_\QCoh (\mathcal{O}_U)$를 생성하는 $X$ 위의 완전 복합체 $Q$가 존재한다.
$V$의 부분집합으로서 $V \setminus U = V(f_1, \dots , f_r)$가 되도록
$f_1, \dots , f_r \in \Gamma (V, \mathcal{O}_V)$를 선택하자.
$K \in D_{perf}(\mathcal{O}_V)$를 $f_1, \dots , f_r$에 대한 코줄 복합체에
대응하는 대상이라 하자. $K' \in D_{perf}(\mathcal{O}_X)$를
\begin{equation}
\label{perfect-equation-detecting-bounded-above}
K' = R (V \to X)_* K = R (V \to X)_! K,
\end{equation}
로 놓자. Cohomology, Lemmas \ref{cohomology-lemma-pushforward-restriction}와
\ref{cohomology-lemma-pushforward-perfect}를 보라. 이는 닫힌집합
$X \setminus U \subset V$에 지지되고 $V$ 위에서 $K$와 동형인 $X$ 위의
완전 복합체이다. 가정에 의해 $R\Hom_{\mathcal{O}_X}(Q, E)$와
$R\Hom_{\mathcal{O}_X}(K', E)$는 위로 유계이다.

\medskip\noindent
(\ref{perfect-equation-detecting-bounded-above})의 $K'$에 대한 두 번째 표현에 의해
$$
\Hom_{D(\mathcal{O}_V)}(K[-i], E|_V) = \Hom_{D(\mathcal{O}_X)}(K'[-i], E) = 0
$$
이 $i \gg 0$에 대해 성립한다. 따라서 Lemma
\ref{perfect-lemma-orthogonal-koszul-first-variant}를 $E|_V$에 적용하여
다음을 만족하는 정수 $a$를 얻을 수 있다.
$\tau_{\geq a}(E|_V) = \tau_{\geq a} R (U \cap V \to V)_* (E|_{U \cap V})$.
그러면 $\tau_{\geq a} E = \tau_{\geq a} R (U \to X)_* (E |_U)$이다(표준
사상이 $U$와 $V$로 제한한 뒤 동형사상인지 확인하라). 따라서 Lemma
\ref{perfect-lemma-bounded-truncation}을 두 번 사용하면
$$
\Hom_{D(\mathcal{O}_U)}(Q|_U [-i], E|_U) =
\Hom_{D(\mathcal{O}_X)}(Q[-i], R (U \to X)_* (E|_U)) = 0
$$
이 $i \gg 0$에 대해 성립함을 알 수 있다. 이 명제가 $U$와 생성자 $Q|_U$에
대해 성립하므로 $E|_U \in D^-_\QCoh(\mathcal{O}_U)$이다. 이제 함자
$R (U \to X)_*$가 $D^-_\QCoh$를 보존하므로(Lemma
\ref{perfect-lemma-quasi-coherence-direct-image}에 의해),
$\tau_{\geq a}E \in D^-_\QCoh(\mathcal{O}_X)$를 얻는다. 따라서
$E \in D^-_\QCoh (\mathcal{O}_X)$이다.
\end{proof}

\begin{proposition}
\label{perfect-proposition-detecting-bounded-below}
$X$를 준콤팩트이고 준분리인 스킴이라 하자.
$G \in D_{perf}(\mathcal{O}_X)$를 $D_\QCoh (\mathcal{O}_X)$를 생성하는 완전
복합체라 하고, $E \in D_\QCoh (\mathcal{O}_X)$라 하자. 다음은 동치이다.
\begin{enumerate}
\item $E \in D^+_\QCoh (\mathcal{O}_X)$이다.
\item $\Hom_{D(\mathcal{O}_X)}(G[-i], E) = 0$이 $i \ll 0$에 대해 성립한다.
\item $\Ext^i_X(G, E) = 0$이 $i \ll 0$에 대해 성립한다.
\item $R\Hom_X(G, E)$는 $D^+(\mathbf{Z})$에 속한다.
\item $H^i(X, G^\vee \otimes_{\mathcal{O}_X}^\mathbf{L} E) = 0$
$i \ll 0$에 대해 성립한다.
\item $R\Gamma(X, G^\vee \otimes_{\mathcal{O}_X}^\mathbf{L} E)$
는 $D^+(\mathbf{Z})$에 속한다.
\item $D(\mathcal{O}_X)$의 모든 완전 대상 $P$에 대해 다음이 성립한다.
\begin{enumerate}
\item $G$를 $P$로 바꾼 명제 (2), (3), (4)가 성립한다.
\item $H^i(X, P \otimes_{\mathcal{O}_X}^\mathbf{L} E) = 0$이
$i \ll 0$에 대해 성립한다.
\item $R\Gamma(X, P \otimes_{\mathcal{O}_X}^\mathbf{L} E)$
는 $D^+(\mathbf{Z})$에 속한다.
\end{enumerate}
\end{enumerate}
\end{proposition}

\begin{proof}
(1)을 가정하자. 등식
$\Hom_{D(\mathcal{O}_X)}(G[-i], E) = \Hom_{D(\mathcal{O}_X)}(G, E[i])$
과 Lemma \ref{perfect-lemma-ext-from-perfect-into-bounded-QCoh}에 의해 이는
$i \ll 0$에 대해 영이다. 따라서 (1)은 (2)를 함의한다.

\medskip\noindent
Cohomology, Section \ref{cohomology-section-global-RHom}의 논의에 의해
(2), (3), (4)는 동치이다. 정의에 의해
$H^i(X, -) = H^i(R\Gamma(X, -))$이므로 (5)와 (6)은 동치이다. Cohomology,
Lemma \ref{cohomology-lemma-dual-perfect-complex}와
$(G[-i])^\vee = G^\vee[i]$라는 사실에 의해, 동치인 조건 (2), (3), (4)는
동치인 조건 (5), (6)과 동치이다.

\medskip\noindent
(7)이 (2)를 함의함은 분명하다. 반대로 동치인 조건 (2)--(6)이 (7)을
함의함을 증명하자. Remark \ref{perfect-remark-classical-generator}에 의해 $G$는
$D_{perf}(\mathcal{O}_X)$의 고전적 생성자임을 상기하라.
$P \in D_{perf}(\mathcal{O}_X)$에 대해 $T(P)$를
$R\Hom_X(P, E)$가 $D^+(\mathbf{Z})$에 속한다는 명제라 하자. 분명히 $T$는
직합과 직합인자에 의해 유전되고, 완전 삼각형에 대해 셋 중 둘 성질을
만족하며, 이동에 의해 보존된다. 따라서 Derived Categories, Remark
\ref{derived-remark-check-on-generator}에 의해 (4)는 모든
$D_{perf}(\mathcal{O}_X)$에서 $T$가 성립함을 함의한다. 다른 모든 성질에도
같은 논증이 적용되지만, (7)(b)와 (7)(c)에는 $P \mapsto P^\vee$가
$D_{perf}(\mathcal{O}_X)$의 자기 동치라는 사실도 사용한다. 작은 세부사항은
생략한다.

\medskip\noindent
Cohomology of Schemes, Lemma \ref{coherent-lemma-induction-principle}의
귀납 원리를 사용하여 동치인 조건 (2)--(7)이 (1)을 함의함을 증명하겠다.

\medskip\noindent
먼저 $X$가 아핀일 때 (2)--(7) $\Rightarrow$ (1)을 증명하자.
$P = \mathcal{O}_X[0]$로 놓자. (7)에서 $i \ll 0$에 대해
$H^i (X, E) = 0$을 얻는다. $E$는 $R\Gamma (X, E)$에 의해 결정되므로 (1)이
따른다. Lemma \ref{perfect-lemma-affine-compare-bounded}를 보라.

\medskip\noindent
이제 $X = U \cup V$이고, $U$는 $X$의 준콤팩트 열린집합이며, $V$는 아핀
열린집합이라고 하자. 또한 스킴 $U$, $V$, $U \cap V$에 대해 함의
(2)--(7) $\Rightarrow$ (1)이 알려졌다고 하자.
$E \in D_\QCoh(\mathcal{O}_X)$가 (2)--(7)을 만족한다고 하자. Lemma
\ref{perfect-lemma-direct-summand-of-a-restriction}과 Theorem
\ref{perfect-theorem-bondal-van-den-Bergh}에 의해 $Q|_U$가
$D_\QCoh (\mathcal{O}_U)$를 생성하는 $X$ 위의 완전 복합체 $Q$가 존재한다.
$V$의 부분집합으로서 $V \setminus U = V(f_1, \dots , f_r)$가 되도록
$f_1, \dots , f_r \in \Gamma (V, \mathcal{O}_V)$를 선택하자.
$K \in D_{perf}(\mathcal{O}_V)$를 $f_1, \dots , f_r$에 대한 코줄 복합체에
대응하는 대상이라 하자. $K' \in D_{perf}(\mathcal{O}_X)$를
\begin{equation}
\label{perfect-equation-detecting-bounded-below}
K' = R (V \to X)_* K = R (V \to X)_! K,
\end{equation}
로 놓자. Cohomology, Lemmas \ref{cohomology-lemma-pushforward-restriction}와
\ref{cohomology-lemma-pushforward-perfect}를 보라. 이는 닫힌집합
$X \setminus U \subset V$에 지지되고 $V$ 위에서 $K$와 동형인 $X$ 위의
완전 복합체이다. 가정에 의해 $R\Hom_{\mathcal{O}_X}(Q, E)$와
$R\Hom_{\mathcal{O}_X}(K', E)$는 아래로 유계이다.

\medskip\noindent
(\ref{perfect-equation-detecting-bounded-below})의 $K'$에 대한 두 번째 표현에 의해
$$
\Hom_{D(\mathcal{O}_V)}(K[-i], E|_V) = \Hom_{D(\mathcal{O}_X)}(K'[-i], E) = 0
$$
이 $i \ll 0$에 대해 성립한다. 따라서 Lemma
\ref{perfect-lemma-orthogonal-koszul-second-variant}를 $E|_V$에 적용하여
다음을 만족하는 정수 $a$를 얻을 수 있다.
$\tau_{\leq a}(E|_V) = \tau_{\leq a} R(U \cap V \to V)_*(E|_{U \cap V})$.
그러면 $\tau_{\leq a} E = \tau_{\leq a} R(U \to X)_*(E|_U)$이다(표준
사상이 $U$와 $V$로 제한한 뒤 동형사상인지 확인하라). 따라서 Lemma
\ref{perfect-lemma-bounded-below-truncation}을 두 번 사용하면
$$
\Hom_{D(\mathcal{O}_U)}(Q|_U [-i], E|_U) =
\Hom_{D(\mathcal{O}_X)}(Q[-i], R (U \to X)_* (E|_U)) = 0
$$
이 $i \ll 0$에 대해 성립함을 알 수 있다. 이 명제가 $U$와 생성자 $Q|_U$에
대해 성립하므로 $E|_U \in D^+_\QCoh (\mathcal{O}_U)$이다. 이제 함자
$R(U \to X)_*$가 아래로 유계인 대상들을 보존하므로(Cohomology, Section
\ref{cohomology-section-derived-functors}를 보라),
$\tau_{\leq a} E \in D^+_\QCoh(\mathcal{O}_X)$를 얻는다. 따라서
$E \in D^+_\QCoh (\mathcal{O}_X)$이다.
\end{proof}










\section{유도 범주의 준연접 대상}
\label{perfect-section-QC}

\noindent
$X$를 스킴이라 하자. $X_{affine, Zar}$가 표준 자리스키 덮개로 주어지는
위상을 갖춘 $X$의 아핀 열린집합들의 범주를 나타냄을 상기하라. Topologies,
Definition \ref{topologies-definition-big-small-Zariski}를 보라.
$X_{affine, Zar}$의 토포스는 $X$의 작은 자리스키 토포스이다. Topologies,
Lemma \ref{topologies-lemma-alternative-zariski}를 보라. 사이트
$X_{affine, Zar}$에는 구조층 $\mathcal{O}$가 주어져 있고, 환 달린 토포스의
동치
$$
(\Sh(X_{affine, Zar}), \mathcal{O})
\longrightarrow
(\Sh(X_{Zar}), \mathcal{O})
$$
가 존재한다. Descent, Equation
(\ref{descent-equation-alternative-small-ringed})과 그 주변의 논의를 보라.
Descent, Section \ref{descent-section-alternative-quasi-coherent}에서는 약간
다른 표기법을 사용한다.

\medskip\noindent
이 절에서는 $X_{affine, Zar}$의 바탕 범주에 층과 전층이 일치하는 혼돈
위상을 준 것을 $X_{affine}$로 나타낸다. 특히 구조층 $\mathcal{O}$는
$X_{affine}$ 위에서도 층이 된다. 환 달린 사이트의 사상
$$
\epsilon :
(X_{affine, Zar}, \mathcal{O})
\longrightarrow
(X_{affine}, \mathcal{O})
$$
을 얻는다. Cohomology on Sites, Section
\ref{sites-cohomology-section-compare}를 보라. 이 절에서는
$D_\QCoh(\mathcal{O}_X)$를 Cohomology on Sites, Section
\ref{sites-cohomology-section-modules-cohomology}에서 도입한 범주
$\mathit{QC}(X_{affine}, \mathcal{O})$와 동일시하겠다.

\begin{lemma}
\label{perfect-lemma-DQCoh-alternative-small}
위의 상황에서 다음 삼각 범주들 사이에 표준 완전 동치들이 존재한다.
\begin{enumerate}
\item $D_\QCoh(\mathcal{O}_X)$,
\item $D_\QCoh(X_{Zar}, \mathcal{O})$,
\item $D_\QCoh(X_{affine, Zar}, \mathcal{O})$,
\item $D_\QCoh(X_{affine}, \mathcal{O}_X)$,
\item $\mathit{QC}(X_{affine}, \mathcal{O})$.
\end{enumerate}
\end{lemma}

\begin{proof}
$U \subset V \subset X$가 아핀 열린집합들이면 환 준동형
$\mathcal{O}(V) \to \mathcal{O}(U)$는 평탄하다. 따라서 (4)와 (5) 사이의
동치는 Cohomology on Sites, Lemma
\ref{sites-cohomology-lemma-cartesian-flat-quasi-coherent}의 특수한 경우이다
(그 증명은 명제도 명확히 해 준다).

\medskip\noindent
Descent, Remark \ref{descent-remark-Zariski-site-space}에 의해 환 달린 사이트
$(X_{Zar}, \mathcal{O})$와 환 달린 공간 $(X, \mathcal{O}_X)$의 가군 범주들은
같다. Descent, Proposition
\ref{descent-proposition-equivalence-quasi-coherent}에 의해 이 동치 아래에서
준연접 가군들이 서로 대응한다. 따라서 (1)과 (2)의 삼각 범주 사이에 표준
완전 동치를 얻는다.

\medskip\noindent
보조정리 앞의 논의는 환 달린 토포스의 동치
$(\Sh(X_{affine, Zar}), \mathcal{O}) \to (\Sh(X_{Zar}), \mathcal{O})$
와 그에 따른 가군들의 아벨 범주 사이의 동치를 준다. 준연접 가군의 개념은
내재적이므로(Modules on Sites, Lemma
\ref{sites-modules-lemma-special-locally-free}), 이 동치는 준연접 가군의
부분범주들을 보존한다. 따라서 (2)와 (3)의 삼각 범주 사이에 표준 완전
동치를 얻는다.

\medskip\noindent
(3)과 (4)의 삼각 범주 사이의 완전 동치를 얻기 위해 Cohomology on Sites,
Lemma \ref{sites-cohomology-lemma-compare-topologies-derived-adequate-modules}를
위의 사상 $\epsilon : (X_{affine, Zar}, \mathcal{O}) \to
(X_{affine}, \mathcal{O})$에 적용하겠다. $\mathcal{B} = \Ob(X_{affine})$로
놓고, $\mathcal{A} \subset \textit{PMod}(X_{affine}, \mathcal{O})$를 사상
$\mathcal{F}(V) \otimes_{\mathcal{O}(V)} \mathcal{O}(U)
\to \mathcal{F}(U)$이 동형사상인 전층 $\mathcal{F}$들로 이루어진 충만
부분범주라 하자. Descent, Lemma
\ref{descent-lemma-quasi-coherent-alternative-small}에 의해
$\mathcal{A}$의 대상들은 자리스키 위상에서 $(X_{affine, Zar}, \mathcal{O})$
위의 준연접 가군인 층들과 정확히 같다. 한편 Modules on Sites, Lemma
\ref{sites-modules-lemma-chaotic-quasi-coherent}에 의해 $\mathcal{A}$의
대상들은 $(X_{affine}, \mathcal{O})$ 위, 즉 혼돈 위상의 준연접 가군들과
정확히 같다. 따라서 Cohomology on Sites, Lemma
\ref{sites-cohomology-lemma-compare-topologies-derived-adequate-modules}가
적용됨을 보이면, $\epsilon^*$와 $R\epsilon_*$를 사용하여 실제로 (3)과
(4)의 범주 사이의 표준 동치를 얻는다.

\medskip\noindent
네 조건을 확인해야 한다.
\begin{enumerate}
\item $\mathcal{A}$의 모든 대상은 자리스키 위상의 층이다. 이는 위에서 보았다.
\item $\mathcal{A}$는 $\textit{Mod}(X_{affine, Zar}, \mathcal{O})$의 약한 세르
부분범주이다. 위에서 $\mathcal{A} = \QCoh(X_{affine, Zar}, \mathcal{O})$임을
보았고, 동치 $\textit{Mod}(X_{affine, Zar}, \mathcal{O}) =
\textit{Mod}(X, \mathcal{O}_X)$ 아래에서 이들이 $X$ 위의 준연접 가군들에
대응함도 보았다. 따라서 결과는 Schemes, Section
\ref{schemes-section-quasi-coherent}의 논의에서 따른다.
\item $X_{affine}$의 모든 대상은 혼돈 위상에서 구성원들이
$\mathcal{B}$의 원소인 덮개를 갖는다. $\mathcal{B}$가 모든 대상을
포함하므로 이는 성립한다.
\item $X_{affine}$의 모든 대상 $U$와 $\mathcal{A}$의 모든
$\mathcal{F}$에 대해 $p > 0$이면 $H^p_{Zar}(U, \mathcal{F}) = 0$이다.
이는 아핀 위 준연접 가군의 코호몰로지가 소멸한다는 사실에서 따른다.
Cohomology of Schemes, Lemma
\ref{coherent-lemma-quasi-coherent-affine-cohomology-zero}를 보라.
\end{enumerate}
이로써 증명이 끝난다.
\end{proof}

\begin{remark}
\label{perfect-remark-QC-compare}
$S$를 스킴이라 하자. 나중에
$\mathit{QC}((\textit{Aff}/S), \mathcal{O})$
도 $D_\QCoh(\mathcal{O}_S)$와 표준적으로 동치임을 보이겠다. Sheaves on
Stacks, Proposition \ref{stacks-sheaves-proposition-QC-compare}를 보라.
\end{remark}
